\documentclass[utf8]{article}
\usepackage{graphicx} % Required for inserting images
\usepackage{ctex}
\usepackage{amsmath}
\usepackage{float}
\usepackage{caption}
\usepackage{amsmath, amssymb}
\usepackage{float}
\usepackage{array}
\usepackage{geometry}
\geometry{left=3cm,right=3cm,top=2cm,bottom=2cm}
\usepackage{indentfirst}
\usepackage{amsfonts}
\usepackage{ctex}
\usepackage[T1]{fontenc}


% 1-32页
\setlength{\parindent}{0em}
\begin{document}
\begin{center}
\Large\bfseries{第二章\ 确定信号分析 }
\end{center}
---------------------------------------

1．设$s_{1}\left(t\right)$,$s_{2}\left(t\right)$是任意的复信号，$S_{1}\left(f\right)$,$S_{2}\left(f\right)$分别是$s_{1}\left(t\right)$,$s_{2}\left(t\right)$的傅氏变换。证明$\textstyle{\int_{-\infty}^{\infty}s_{1}^{*}\left(t\right)s_{2}\left(t\right)d t}=\textstyle{\int_{-\infty}^{\infty}S_{1}^{*}\left(f\right)S_{2}\left(f\right)d f}$。

++++++++++++

证：
$$
\int_{-\infty}^{\infty}s_{1}^{*}\left(t\right)s_{2}\left(t\right)d t = \left.\int_{-\infty}^{\infty}\left[\right.\int_{-\infty}^{\infty}S_{1}\left(f\right)e^{j2\pi f t}d f\right]^{\ast}\left[\int_{-\infty}^{\infty}S_{2}\left(u\right)e^{j2\pi u t}d u\right]dt
$$

$$
=\int_{-\infty}^{\infty}\int_{-\infty}^{\infty} S_{1}^{*}\left(f\right)S_{2}\left(u\right)\biggl[\int_{-\infty}^{\infty}e^{-j2\pi f t+j2\pi u t}d t\biggr]d u d f
$$

$$=\int_{-\infty}^{\infty}\int_{-\infty}^{\infty}S_{1}^{*}\left(f\right)S_{2}\left(u\right)\delta\left(f-u\right)d u d f$$

$$=\int_{-\infty}^{\infty}S_{1}^{*}\left(f\right)S_{2}\left(f\right)d f$$

注：Paserval定理的一般形式是$$\textstyle{\int_{-\infty}^{\infty}s_{1}^{*}\left(t\right)s_{2}\left(t\right)d t}=\textstyle{\int_{-\infty}^{\infty}S_{1}^{*}\left(f\right)S_{2}\left(f\right)d f},$$
其中$s_{1}\left(t\right)$,$s_{2}\left(t\right)$是任意的复信号，$S_{1}\left(f\right)$,$S_{2}\left(f\right)$分别是$s_{1}\left(t\right)$,$s_{2}\left(t\right)$的傅氏变换。它的一个特例是$s_{1}\left(t\right)=s_{2}\left(t\right)=s\left(t\right)$，此时$$\textstyle{\int_{-\infty}^{\infty}}\left|s\left(t\right)\right|^{2}\,d t=\int_{-\infty}^{\infty}\left|S\left(f\right)\right|^{2}\,d f$$
其中$S\left(f\right)$是$s\left(t\right)$的傅氏变换。

---------------------------------------

2．已知信号
$$
g\left(t\right)=\left\{\begin{array}{c c}{{\cos\frac{\pi t}{T_{s}}}}&{{-\frac{T_{s}}{2}\leq t< \frac{T_{s}}{2}}}\\ {{0}}&{{e l s e}}\end{array}\right.
$$

(1)求其傅氏变换$G(f)$； 

(2)$\int_{-\infty}^{\infty}\left|G\left(f\right)\right|^{2}d f=$？

++++++++++++

解：(1)令
$$
g^{\prime}(t)=\left\{\begin{array}{c c}{{1}}&{{-\frac{T_{s}}{2}\leq t< \frac{T_{s}}{2}}}\\{{0}}&{{e l s e}}\end{array}\right.
$$
则
$$g\left(t\right)=g^{\prime}\left(t\right)\cos\frac{\pi t}{T_{s}}=\frac{g'(t)}{2} \left(e^{j\frac{\pi t}{T_{s}}}+e^{-j\frac{\pi t}{T_{s}}} \right)
$$
$g′(t)$的傅氏变换为 $$G^{\prime}\left(f\right)=T_{s}\sin c (f T_{s})$$
因此 $g(t)$ 的傅氏变换为 
$$G(t)=\frac{1}{2}\Bigg[G^{\prime}\Bigg(f-\frac{1}{2T_{s}}\Bigg)+G^{\prime}\Bigg(f+\frac{1}{2T_{s}}\Bigg)\Bigg]\ = \frac{T_{s}}{2} \left[ \frac{ \sin \pi \left( fT_{s}+ \frac{1}{2} \right)}{ \pi \left( fT_{s}+ \frac{1}{2} \right)}+ \frac{ \sin \pi \left( fT_{s}- \frac{1}{2} \right)}{ \pi \left( fT_{s}- \frac{1}{2} \right)} \right]$$
$$
=\frac{T_{{}_{ s}}}{2} \left[\frac{\cos\pi f T_{{}_{ s}}}{\pi\left(f T_{{}_{ s}} + \frac{1}{2}\right)} - \frac{\cos\pi f T_{ {}_{ s}}}{\pi\left(f T_{{}_{ s}} - \frac{1}{2}\right)} \right] = - \frac{T_{{}_{ s}}\cos\pi f T_{{}_{ s}}}{2\pi} \times \frac{1}{\left(f T_{ {}_{ s}}\right)^{ 2}- \frac{1}{4}}\\ 
$$
$$
= \frac{2T_{{}_{ s}}\cos\pi f T_{{}_{ s}}}{\pi\left( 1 - 4f^{ 2}T_{{}_{ s}}^{ 2}\right)}
$$
(2)由 Paserval 定理
$$\int_{- \infty}^{ \infty} \left| G \left( f \right) \right|^{2}df= \int_{- \infty}^{ \infty}g^{2} \left( t \right) dt= \frac{T_{s}}{2}$$

---------------------------------------

3．已知周期信号$s \left( t \right)= \sum_{n=- \infty}^{ \infty}g \left( t-nT \right)$，其中$$g \left( t \right)= \left \{ \begin{matrix} 2 \big/T&-T \big/4 \leq t<T \big/4 \\ 0&else \end{matrix} \right.$$
求 s(t)的功率谱密度。

++++++++++++

解：s(t)可看成是冲激序列$x \left( t \right)= \sum_{n=- \infty}^{ \infty} \delta \left( t-nT \right)$通过一个冲激响应为g(t)的线性系统的输出。将x(t)展成傅氏级数，得
$$x \left( t \right)= \sum_{n=- \infty}^{ \infty} \delta \left( t-nT \right)= \frac{1}{T} \sum_{m=- \infty}^{ \infty}e^{j2 \pi \frac{m}{T}t}$$
所以x(t)的功率谱密度为
$$P_{x} \left( f \right)= \frac{1}{T^{2}} \sum_{m=- \infty}^{ \infty} \delta \left( f- \frac{m}{T} \right)$$
g(t)的傅氏变换是$G \left( f \right)= \sin{c} \left( \frac{fT}{2} \right)$，因此
$$P_{s} \left( f \right)= \frac{1}{T^{2}} \sum_{m=- \infty}^{ \infty} sinc^{2} \left( \frac{m}{2} \right) \delta \left( f- \frac{m}{T} \right)$$
$$= \frac{1}{T^{2}} \delta \left( f \right)+ \frac{1}{T^{2}} \sum_{k=- \infty}^{ \infty} \frac{4}{ \left( 2k-1 \right)^{2} \pi^{2}} \delta \left( f- \frac{2k-1}{T} \right)$$

---------------------------------------

4．对任意实信号$g_{1} \left( t \right),g_{2} \left( t \right),\left(- \infty<t< \infty \right)$，令$u \left( t \right)=g_{1} \left( t \right)+g_{2} \left( t \right)$、$v \left( t \right)=g_{1} \left( t \right)-g_{2} \left( t \right)$。求u(t)与v(t)正交的条件。

++++++++++++

解：
$$\begin{array}{c}\int_{-\infty}^{\infty} u \left( t \right) v \left( t \right) dt= \int_{-\infty}^{ \infty} \left[ g_{1} \left( t \right)+g_{2} \left( t \right) \right] \left[ g_{1} \left( t \right)-g_{2} \left( t \right) \right] dt \\ \qquad \qquad= \int_{-\infty}^{ \infty} \left[ g_{1}^{2} \left( t \right)-g_{2}^{2} \left( t \right) \right] dt=E_{1}-E_{2}\end{array}$$
因此，u(t)与v(t)正交的条件就是$g_{1}\left(t\right)$和$g_{2}\left(t\right)$等能量。 

---------------------------------------

5.$\int_{-\infty}^{\infty}sinc(t-y)sinc(y)dy =$?

++++++++++++

解：令s(t)=sinc(t)，则其傅氏变换是
$$S \left( f \right)= \left \{ \begin{matrix} 1& \left| f \right| \leq \frac{1}{2} \\ 0&else \end{matrix} \right.$$
$u \left( t \right)= \int_{- \infty}^{ \infty} sinc \left( t-y \right) sinc \left( y \right) dy$是sinc(t)和sinc(t)的卷积。时域卷积对应频域乘积，所以u(t)的傅氏变换是
$$U \left( f \right)=S \left( f \right) S \left( f \right)=S \left( f \right)$$
因此u(t)=sinc(t)

---------------------------------------

6.证明$\sum_{k=- \infty}^{ \infty} sinc \left( k+ \epsilon \right)=1$，其中0≤ε<1。

++++++++++++

证：对直流信号m(t)以$f_{s}$=1Hz的采样速率在$t= \frac{k}{f_{s}}+ \epsilon=kT_{s}+ \epsilon$时刻进行理想采样，得到的采样信号是
$$s \left( t \right)= \sum_{k=- \infty}^{ \infty}m \left( kT_{s}+ \epsilon \right) \delta \left( t-kT_{s}- \epsilon \right)= \sum_{k=- \infty}^{ \infty} \delta \left( t-kT_{s}- \epsilon \right)$$
再将s(t)通过一个带宽为$\frac{f_{s}}{2}=0.5Hz$的理想低通滤波器，则输出还是直流m(t)=1。

另外，此理想低通的冲激响应是sinc(t)，因此滤波器输出是$\sum_{k=- \infty}^{ \infty} sinc \left( t-k- \epsilon \right)$，令t=0，并考虑到sinc(x)是偶函数，得$\sum_{k=- \infty}^{ \infty} sinc \left( k+ \epsilon \right)=1$。

---------------------------------------

7．设基带信号m(t)的频谱范围是$\left[ f_{L},f_{H} \right]$，其Hilbert变换是$\hat{m} \left( t \right)$。求下列信号的Hilbert变换（$f_{c}>>f_{H}$）。
$$\left( 1 \right) x_{1} \left( t \right)=m \left( t \right) \cos 2 \pi f_{c}t；$$
$$\left( 2 \right) x_{2} \left( t \right)=m \left( t \right)+j \hat{m} \left( t \right)；$$
$$\left(3 \right) x_{3} \left( t \right)=m \left( t \right) \cos 2 \pi f_{c}t- \hat{m} \left( t \right) \sin 2 \pi f_{c}t；$$
$$\left( 4 \right) x_{4} \left( t \right)=Re \left \{ Ae^{j2 \pi \left( f_{c}t+K \int_{- \infty}^{t}m \left( \tau \right) d \tau \right)} \right \}。$$
++++++++++++

解：设m(t)的傅氏变换为M(f)

(1) $x_{1}(t)$的频谱为 
$$X_{1} \left( f \right)= \frac{1}{2} \left[ M \left( f-f_{c} \right)+M \left( f+f_{c} \right) \right]= \left \{ \begin{matrix} \frac{M \left( f+f_{c} \right)}{2}&f<0 \\ \frac{M \left( f-f_{c} \right)}{2}&f>0 \end{matrix} \right.$$
因此其Hilbert变换$\hat{x}_{1} \left( t \right)$的频谱为
$$\hat{X}_{1} \left( f \right)= \left \{ \begin{matrix} j \frac{M \left( f+f_{c} \right)}{2}&f<0 \\ -j \frac{M \left( f-f_{c} \right)}{2}&f>0 \end{matrix} \right.$$
也即
$$\hat{X}_{1} \left( f \right)=j \frac{M \left( f+f_{c} \right)}{2}-j \frac{M \left( f-f_{c} \right)}{2}$$
做傅氏反变换得
$$\widehat{x}_{{}_{1}} \left( t \right)= \frac{j}{2}m \left( t \right) e^{-j2 \pi f_{c}t}- \frac{j}{2}m \left( t \right) e^{j2 \pi f_{c}t}=m \left( t \right) \sin 2 \pi f_{c}t$$
(2)$\hat{x}_{{}_{2}} \left( t \right)= \hat{m} \left( t \right)+j \hat{ \hat{m}} \left( t \right)= \hat{m} \left( t \right)-jm \left( t \right)=-jx_{2}^{*} \left( t \right)$

(3)$\hat{x}_{{}_{3}} \left( t \right)=m \left( t \right) \sin 2 \pi f_{c}t+ \hat{m} \left( t \right) \cos 2 \pi f_{c}t$

(4)
$$x_{4} \left(t \right)=Re \left \{ Ae^{j2 \pi \left(f_{c}t+K \int_{- \infty}^{t}m \left( \tau \right) d \tau \right)} \right \}$$
$$=A \cos \left[ 2 \pi f_{c}t+2 \pi K \int_{- \infty}^{t}m \left( \tau \right) d \tau \right]$$
$$\hat{x}_{4} \left(t \right)=A \cos \left[ 2 \pi f_{c}t+2 \pi K \int_{- \infty}^{t}m \left( \tau \right) d \tau \right]$$
$$= \mathrm{Re}\left\{-jAe^{j\left(2\pi f_{c}t+2\pi K\int_{-\infty}^ {t}m\left(\tau\right)d\tau\right)}\right\}$$
注1：窄带信号的Hilbert变换就是将载波相位后移$90^{ \circ}$

注2：对窄带信号进行Hilbert变换就是对复包络乘−j 

---------------------------------------

8．已知某带通信号为
$$s\left(t \right)=m_{{}_{1}} \left(t \right) \cos \left(2 \pi f_{c}t- \varphi \right)-m_{2} \left(t \right) \sin \left(2 \pi f_{c}t- \varphi \right)=Re \left \{ s_{L} \left(t \right) e^{j2 \pi f_{c}t} \right \}$$

其中基带信号$m_{1}(t)$和$m_{2}(t)$的带宽都是W且$W<<f_{ c}$，$s_{L} \left( t \right)$是$s_{L} \left( t \right)$的复包络。

(1)求出$s_{L}(t)$的实部$m_{o}(t) =Re \{s_{L}(t)\}$

(2)当$m_{2}(t)$和$m_{1}(t)$之间满足何种关系时，s(t)的频谱在$\left| f \right|<f_{c}$范围内是0？ 

++++++++++++

解：(1)
$$s \left( t \right)=m_{1} \left( t \right) \cos \left( 2 \pi f_{c}t- \varphi \right)-m_{2} \left( t \right) \sin \left( 2 \pi f_{c}t- \varphi \right)$$
$$= \mathrm{Re}\left\{\left[\left.m_{1}\left(t\right)+jm_{{}_{2} }\left(t\right)\right.\right]e^{j\left(2\pi f_{c}t-\varphi\right)}\right\}= \mathrm{Re}\left\{\left[\left.m_{{}_{1}}\left(t\right)+jm_{2}\left(t \right)\right.\right]e^{-j\varphi}e^{j2\pi f_{c}t}\right\}$$
因此，复包络是
$$s_{L} \left( t \right)= \left[ m_{1} \left( t \right)+jm_{2} \left( t \right) \right] e^{-j \varphi}$$
其实部为
$$m_{O}\left(t \right)=Re \left \{ s_{L}\left(t \right) \right \}=m_{1} \left( t \right) \cos \varphi+m_{2} \left( t \right) \sin \varphi$$
(2)s(t)的频谱在$\left| f \right|<f_{c}$范围内是0，说明$s_{L}$只具有正频率分量，也即$m_{1} \left( t \right)+jm_{2} \left( t \right)$只具有正频率分量。满足这种情况的情形是：$m_{1} \left( t \right)$是$m_{2} \left( t \right)$的Hilbert变换，此时由于$m_{1} \left( t \right)+jm_{2} \left( t \right)$是解析信号，所以只有正频率分量。 

---------------------------------------

9．若带通信号s(t)的复包络$s_{L}$的傅氏变换是$S_{L}$，求 s(t)的傅氏变换S(f)。

++++++++++++

解：$s\left(t\right) = \mathrm{Re} \left\{s_{L}\left(t\right)e^{j2\pi f_{c} t}\right\} = \frac{s_{L}\left(t\right)}{2}e^{j2\pi f_{c} t}+ \frac{s_{L}^{*}\left(t\right)}{2}e^{-j2\pi f_{c} t}$

由于$s_{L}^{*}\left(t\right)$的傅氏变换是$S_{L}^{*}\left(-f\right)$，所以
$$S \left( f \right)= \frac{1}{2}S_{L} \left( f-f_{c} \right)+ \frac{1}{2}S_{L}^{*} \left(-f-f_{c} \right)$$

---------------------------------------

10．求以下带通信号傅氏变换，已知基带信号m(t)是实信号，其傅氏变换是M(f)，带宽为W，$f_{c} \rightarrow W$，$\hat{m} \left( t \right)$是m(t)的希尔伯特变换。

(1)$s \left( t \right)=m \left( t \right) \cos 2 \pi f_{c}t;$

(2)$s \left( t \right)=m \left( t \right) \cos 2 \pi f_{c}t- \hat{m} \left( t \right) \sin 2 \pi f_{c}t。$

++++++++++++

解：(1)$s_{L} \left( t \right)=m \left( t \right)$，所以
$$S \left( f \right)= \frac{M \left( f-f_{c} \right)+M^{*} \left(-f-f_{c} \right)}{2}$$
又由于$M^{*} \left(-f \right)=M \left( f \right)$，所以
$$S \left( f \right)= \frac{M \left( f-f_{c} \right)+M \left( f+f_{c} \right)}{2}$$
也可以这样做：
$$s \left( t \right)= \frac{m \left( t \right)}{2}e^{j2 \pi f_{c}t}+ \frac{m \left( t \right)}{2}e^{-j2 \pi f_{c}t}$$
利用傅氏变换的频移特性得到
$$S \left( f \right)= \frac{M \left( f-f_{c} \right)+M \left( f+f_{c} \right)}{2}$$
(2)$s_{L} \left( t \right)=m \left( t \right)+j \hat{m} \left( t \right)$是解析信号，其傅氏变换为
$$S_{L} \left( f \right)= \left \{ \begin{matrix} 2M \left( f \right)&f>0 \\ 0&f<0 \end{matrix} \right.$$
$$S_{L}^{*} \left(f \right)= \left \{ \begin{matrix} 2M^{*} \left(f \right)&f>0 \\ 0&f<0 \end{matrix} \right.= \left \{ \begin{matrix} 2M \left(-f \right)&f>0 \\ 0&f<0 \end{matrix} \right.$$
$$S_{L}^{*} \left(-f \right)=\left \{ \begin{matrix} 2M \left( f \right)&f<0 \\ 0&f>0 \end{matrix} \right.$$
$$S \left( f \right)= \frac{S_{L} \left( f+f_{c} \right)+S_{L}^{*} \left(-f-f_{c} \right)}{2}= \left \{ \begin{matrix} M \left( f-f_{c} \right)&f>f_{c} \\ M \left( f+f_{c} \right)&f<-f_{c} \\ 0&else \end{matrix}\right.$$

---------------------------------------

11．在下图所示的线性系统中，$H \left( f \right)= \left \{ \begin{matrix} T_{b}& \left| f \right| \leq \frac{1}{2T_{b}} \\ 0&else \end{matrix} \right.$。若在图中A点加上一个激励δ(t)，求B点的波形y(t)及其频谱Y(f)，并请画出输出信号的振幅谱图。 
\begin{figure}[H]
  \centering
  \includegraphics{images/image1.png}
  \caption*{题11图}
\end{figure}

++++++++++++

解：
$$Y\left(\begin{smallmatrix}f\\ \end{smallmatrix}\right)=H\left(\begin{smallmatrix}f\\ \end{smallmatrix}\right)\left(1+e^{-j2\pi fT_{b}}\right)=H\left(\begin{smallmatrix} f\\ \end{smallmatrix}\right)\left(e^{j\pi fT_{b}}+e^{-j\pi fT_{b}}\right)e^{-j\pi fT_{b}}$$
$$=2H \left( f \right) e^{-j \pi fT_{b}} \cos \pi fT_{b}$$
$$= \left\{\begin{aligned} 2T_{b}& e^{-j\pi f T_{b}} \cos\pi f T_{b}&\big|f\big| \leq \frac{1}{2 T_{b}}\\ & 0&\mathit{else}\end{aligned}\right.$$
其振幅频谱图如下
\begin{figure}[H]
  \centering
  \includegraphics{images/image2.png}
\end{figure}

H(f)对应的冲激响应是
$$h \left( t \right)= \sin c \left( \frac{t}{T_{b}} \right)$$
所以
$$y \left( t \right)=h \left( t \right)+h \left( t-T_{b} \right) $$
$$ = \sin c \left( \frac{t}{T_{b}} \right)+ \sin c \left( \frac{t-T_{b}}{T_{b}} \right)$$
$$= \frac{T_{b}}{ \pi t} \sin \frac{ \pi t}{T_{b}}+ \frac{T_{b}}{ \pi \left( t-T_{b} \right)} \sin \frac{ \pi \left( t-T_{b} \right)}{T_{b}}$$
$$= \frac{T_{b}}{ \pi t} \sin \frac{ \pi t}{T_{b}}- \frac{T_{b}}{ \pi \left( t-T_{b} \right)} \sin \frac{ \pi t}{T_{b}}$$
$$= \frac{T_{b}}{ \pi} \sin \frac{ \pi t}{T_{b}} \left( \frac{1}{t}- \frac{1}{t-T_{b}} \right)$$
$$= \frac{T_{b}^{2} \sin \frac{ \pi t}{T_{b}}}{ \pi t \left( T_{b}-t \right)}$$
---------------------------------------

12．设有周期信号$u \left( t \right)= \sum_{n=- \infty}^{ \infty}x^{2} \left( t-nT \right)$，其中x(t)是带宽为W 的任意实信号。请问

(1)W满足何种条件时，u(t)具有$a+b \cos \left( \frac{2 \pi{t}}{T}+ \theta \right)$的形式？

(2)W满足何种条件时，u(t)是直流？

++++++++++++

解：u(t)是$\sum_{n=- \infty}^{ \infty} \delta \left( t-nT \right)$通过冲激响应为$h \left( t \right)=x^{2} \left( t \right)$后的输出。$\sum_{n=- \infty}^{ \infty} \delta \left( t-nT \right)$的傅氏变换是$\frac{1}{T_{s}} \sum_{m=- \infty}^{ \infty} \delta \left( f- \frac{m}{T} \right)$。设x(t)的频谱是X(f)，则h(t)的频谱H(f)是X(f)和X(f)的卷积，因此H(f)的带宽是2W。 

(1)若H(f)的带宽小于$\frac{2}{T_{S}}$，即若$W< \frac{1}{T_{s}}$，则$\sum_{n=- \infty}^{ \infty} \delta \left( t-nT \right)$经过H(f)后的输出不包含频率为$\frac{2}{T}$或更高的谐波。此时u(t)必然具有$a+b \cos \left( \frac{2 \pi t}{T}+ \theta \right)$的形式。

(2)若H(f)的带宽小于$\frac{1}{T_{S}}$，即若$W< \frac{1}{2T_{s}}$，则$\sum_{n=- \infty}^{ \infty} \delta \left( t-nT \right)$经过H(f)后的不包任何谐波分量，此时u(t)是直流。

---------------------------------------

13．设
$$s_{1} \left( t \right)= \left \{ \begin{matrix} A \cos 2 \pi f_{1}t&0 \leq t<T \\ 0&else \end{matrix} \right.$$
$$s_{2} \left( t \right)= \left \{ \begin{matrix} A \cos 2 \pi f_{2}t&0 \leq t<T \\ 0&else \end{matrix} \right.$$
其中$f_{1} > f_{2}$，$f_{1} + f_{2}$是$\frac{1}{T}$的整倍数，求

(1)$s_{1}\left( t \right)$，$s_{2}\left( t \right)$的自相关函数$R_{1} \left( \tau \right)$、$R_{2} \left( \tau \right)$；

(2)$s_{1}\left( t \right)$，$s_{2}\left( t \right)$的归一化互相关函数$\rho_{12} \left( \tau \right)$；

(3)$f_{1} - f_{2}$满足何种关系时，$\rho_{12} \left( 0 \right)= 0$？

++++++++++++

解：(1)
$$R_{1} \left( \tau \right)= \int_{0}^{T}A \cos 2 \pi f_{1}t \times A \cos \left( 2 \pi f_{1}t+2 \pi f_{1} \tau \right) dt$$
$$= \frac{A^{2}}{2} \left[ \int_{0}^{T} \cos 2 \pi f_{1} \tau dt+ \int_{0}^{T} \cos \left( 4 \pi f_{1}t+2 \pi f \tau \right) dt \right]$$
$$= \frac{A^{2}T}{2} \cos 2 \pi f_{1} \tau$$
同理
$$R_{2} \left( \tau \right)= \frac{A^{2}T}{2} \cos 2 \pi f_{2} \tau 。$$
(2)
$$\rho_{12} \left( \tau \right)= \frac{ \int_{0}^{T}A \cos 2 \pi f_{1}t \times A \cos \left( 2 \pi f_{2}t+2 \pi f_{2} \tau \right) dt}{ \sqrt{R_{1} \left( 0 \right) R_{2} \left( 0 \right)}}$$
$$= \frac{1}{T} \left[ \int_{0}^{T} \cos \left[ 2 \pi \left( f_{1}-f_{2} \right) t-2 \pi f_{2} \tau \right] dt+ \int_{0}^{T} \cos \left[ 2 \pi \left( f_{1}+f_{2} \right) t+2 \pi f_{2} \tau \right] dt \right]$$
$$= \frac{1}{T} \left[ \cos 2 \pi f_{2} \tau \int_{0}^{T} \cos 2 \pi \left( f_{1}-f_{2} \right) tdt+ \sin 2 \pi f_{2} \tau \int_{0}^{T} \sin 2 \pi \left( f_{1}-f_{2} \right) tdt \right]$$
$$= \frac{1}{T} \left[\cos 2\pi f_{2} \tau \times \frac{ \sin 2\pi \big( f_{1} - f_{2} \big)T}{2\pi \big( f_{1} - f_{2} \big)} + \frac{\sin 2\pi f_{2} \tau}{2\pi \big( f_{1} - f_{2} \big)} \times \left( 1 - \cos 2\pi \big( f_{1} - f_{2} \big)T\right)\right]$$
$$= \frac{ \cos 2 \pi f_{2} \tau \sin \left[ 2 \pi \left( f_{1}-f_{2} \right) T \right]+2 \sin 2 \pi f_{2} \tau \sin^2 \left[ \pi \left( f_{1}-f_{2} \right) T \right]}{2 \pi \left( f_{1}-f_{2} \right) T}$$
(3)
$$\rho_{12} \left( 0 \right)= \frac{ \sin \left[ 2 \pi \left( f_{1}-f_{2} \right) T \right]}{2 \pi \left( f_{1}-f_{2} \right) T}= \sin c \left[ 2 \left( f_{1}-f_{2} \right) T \right]$$

当$2 \left( f_{1}-f_{2} \right) T$为正整数时，$\rho_{12} \left( 0 \right)=0$，即要求$f_{1}-f_{2} $是$\frac{1}{2T}$的整数倍。

---------------------------------------

14．带通信号$s \left( t \right)= \left \{ \begin{matrix} A \cos 2 \pi f_{c}t&0 \leq t<T \\ 0&else \end{matrix} \right.$通过一个冲激响应为h(t)的线性系统，输出为 y (t) 。若$f_{c}= \frac{8}{T}$，$h \left( t \right)= \left \{ \begin{matrix} 2 \cos 2 \pi f_{c}t&0 \leq t<T \\ 0&else \end{matrix} \right.$，试求

(1)s(t)的复包络$s_{L} \left( t \right)$；

(2)y(t)的复包络$y_{L} \left( t \right)$；

(3)画出y(t)的波形。 

++++++++++++

解：(1)$s_{L} \left( t \right)= \left \{ \begin{matrix} A&0 \leq t<T_{s} \\ 0&else \end{matrix} \right.$

(2)h(t)的等效低通响应是

$$h_{eq} \left( t \right)= \frac{1}{2}h_{L} \left( t \right)= \left \{ \begin{matrix} 1&0 \leq t<T_{s} \\ 0&else \end{matrix} \right.$$
因此
$$y_{L}\left(t \right)= \int_{-\infty}^{ \infty}s_{L} \left(t- \tau \right) h_{eq} \left( \tau \right) d \tau= \int_{0}^{T}s_{L} \left(t- \tau \right) d \tau$$
$$= \left \{ \begin{matrix} At&0 \leq t<T \\ AT \left( 2- \frac{t}{T} \right)&T \leq t<2T \\ 0&else \end{matrix} \right.$$
(3)
$$y \left( t \right)=Re \left \{ y_{L} \left( t \right) e^{j2 \pi f_{c}t} \right \} = \left \{ \begin{matrix} At\cos 2\pi f_{c}t& 0 \leq t<T \\ AT\left(2-\frac{t}{T}\right)\cos 2\pi f_{c}t& T \leq t<2T \\ 0&else\end{matrix}\right.$$
其图形如下
\begin{figure}[H]
  \centering
  \includegraphics{images/image3.png}
\end{figure}

---------------------------------------

15．带通信号$s \left( t \right)= \left \{ \begin{matrix} A \sin 2 \pi f_{c}t&0 \leq t<T \\ 0&else \end{matrix} \right.$通过一个冲激响应为 h(t) 的线性系统，输出为y(t) 。若$ f_{c} = \frac{8}{T}$，$ h\left(t\right) = \left\{ \begin{aligned} 2\cos 2\pi f_{{}_{c}}& t& 0 \leq t < T\\ 0&& else\end{aligned}\right.$，试求

(1)s(t)的复包络$s_{L} \left( t \right)$；

(2)y(t)的复包络$y_{L} \left( t \right)$；

(3)画出y(t)的波形。 

(4)写出T时刻y(t)及其包络$\left| y_{L} \left( t \right) \right|$的瞬时值y(T)及$\left| y_{L} \left( T \right) \right|$。

++++++++++++

解：(1) 
$$s_{L} \left( t \right)= \left \{ \begin{matrix}-jA&0 \leq t<T_{s} \\ 0&else \end{matrix} \right.$$
(2) h(t) 的等效低通响应是
$$h_{eq} \left( t \right)= \frac{1}{2}h_{L} \left( t \right)= \left \{ \begin{matrix} 1&0 \leq t<T_{s} \\ 0&else \end{matrix} \right.$$
因此 
$$y_{L}\left(t \right)= \int_{-\infty}^{ \infty}s_{L} \left( t- \tau \right) h_{eq} \left( \tau \right) d \tau$$
$$= \int_{0}^{T}s_{L} \left( t- \tau \right) d \tau$$
$$= \left\{\begin{matrix} -jAt& 0\leq t < T\\ -j AT\left(2 - \frac{t}{T}\right)& T\leq t < 2T\\ 0& else \end{matrix} \right.$$
(3) 
$$y \left( t \right)=Re \left \{ y_L \left( t \right) e^{j2 \pi f_c t} \right \}$$
$$= \left\{\begin{matrix} At\sin 2\pi f_{c}t& 0 \leq t<T\\ AT\left(2-\frac{t}{T}\right) \sin 2\pi f_{c}t & T \leq t<2T\\ 0&else \end{matrix}\right.$$
其图形如下
\begin{figure}[H]
  \centering
  \includegraphics{images/image4.png}
\end{figure}

(4)y(T)=0，$\left|y_{L} \left( T \right) \right|=AT$

---------------------------------------

16．请证明，$\sum_{n=- \infty}^{ \infty} \delta \left( f- \frac{m}{T_{s}} \right)=T_{s} \sum_{n=- \infty}^{ \infty}e^{j2 \pi mfT_{s}}$。

++++++++++++

解：因为
$\sum_{n=- \infty}^{ \infty} \delta \left( f- \frac{m}{T_{s}} \right)$
是频率的周期函数，周期为1/$T_{s}$。可将它展开为傅氏级数： 
$$\sum_{n=- \infty}^{ \infty} \delta \left( f- \frac{m}{T_{s}} \right)= \sum_{m=- \infty}^{ \infty} F_{m}e^{j2 \pi mfT_{s}}$$
其中
$$F_{m}= \frac{1}{1/T_{s}} \int_{- \frac{1}{2T_{s}}}^{ \frac{1}{2T_{s}}} \left[ \sum_{n=- \infty}^{ \infty} \delta \left( f- \frac{n}{T_{s}} \right) \right] e^{-j2 \pi mfT_{s}}df$$
$$=T_{s} \int_{- \frac{1}{2T_{s}}}^{ \frac{1}{2T_{s}}} \delta \left( f \right) e^{-j2 \pi mfT_{s}}df$$
$$=T_{s}$$
因此
$$\sum_{n=- \infty}^{ \infty} \delta \left( f- \frac{m}{T_{s}} \right)=T_{s} \sum_{m=- \infty}^{ \infty}e^{j2 \pi mfT_{s}}$$

---------------------------------------

17．已知实信号x(t)的傅氏变换是X(f)，分别定义其等效矩形时宽α和等效矩形带宽β为$\alpha\triangleq\frac{\int\limits_{-\infty}^{+\infty}\left|x\left(t\right) \right|dt}{\left|x\left(t\right)\right|_{\max}}$，$\beta\triangleq\frac{\int\limits_{-\infty}^{+\infty}\left|X\left(f\right) \right|df}{\left|X\left(f\right)\right|_{\max}}$。

(1)证明$\alpha \beta \geq 1$； 

(2)指出等号成立的条件，举出满足条件的x(t)。

++++++++++++

解：设$t^{*}$是x(t)达到绝对值最大的时间，即$\left| x \left( t^{*} \right) \right|= \left| x \left( t \right) \right|_{ \max}$。设$f^{*}$是 X ( f ) 达到绝对值最大的频率，即$\left|X \left( f^{*} \right) \right|= \left|X \left( f \right) \right|_{ \max}$。

(1)因为
$$\begin{array}{l}\int\limits_{-\infty}^{+\infty}\left|x\left(t\right)\right| dt=\int\limits_{-\infty}^{+\infty}\left|x\left(t\right)e^{-j2\pi ft}\right|dt \geq\int\limits_{-\infty}^{+\infty}x\left(t\right)e^{-j2\pi ft}dt=X\left(f \right)\end{array}$$
$$\begin{array}{l}\int \limits_{-\infty}^{+ \infty} \left| x \left( t \right) \right| dt= \int \limits_{-\infty}^{+ \infty} \left|-x \left( t \right) e^{-j2 \pi ft} \right| dt \geq- \int \limits_{- \infty}^{+ \infty}x \left( t \right) e^{-j2 \pi ft}dt=-X \left( f \right) \end{array}$$
故$\int \limits_{-\infty}^{+ \infty} \big|x \big( t \big) \big| dt \geq \left| X \left( f \right) \right|$，同理可得$\int \limits_{-\infty}^{+ \infty} \left| X \left( f \right) \right| df \geq \left| x \left( t \right) \right|$。于是 
$$\alpha \beta= \frac{ \int_{- \infty}^{ \infty} \left| x \left( t \right) \right| dt}{ \left| x \left( t \right) \right|_{ \max}} \times \frac{ \int_{- \infty}^{ \infty} \left| X \left( f \right) \right| df}{ \left| X \left( f \right) \right|_{ \max}} \geq \frac{ \left| X \left( f \right) \right|}{ \left| x \left( t \right) \right|_{ \max}} \times \frac{ \left| x \left( t \right) \right|}{ \left| X \left( f \right) \right|_{ \max}}$$
这个不等式对所有$f \in(- \infty, \infty)$及所有$t \in(- \infty, \infty)$都成立，自然也对$t=t^{*}$及$f^{*}$成立，由此得$\alpha \beta \geq 1$。 

(2)由刚才的推导知 
$$\alpha \beta= \frac{ \int_{- \infty}^{ \infty} \left| X \left( f \right) \right| df}{ \left| x \left( t \right) \right|_{ \max}} \times \frac{ \int_{- \infty}^{ \infty} \left| x \left( t \right) \right| dt}{ \left| X \left( f \right) \right|_{ \max}}$$
且有
$$\frac{ \int_{- \infty}^{ \infty} \left| X \left( f \right) \right| df}{ \left| x \left( t \right) \right|_{ \max}} \geq 1, \frac{ \int_{- \infty}^{ \infty} \left| x \left( t \right) \right| dt}{ \left| X \left( f \right) \right|_{ \max}} \geq 1$$
故欲$\alpha \beta$=1，必须同时满足 
$$\int_{- \infty}^{ \infty} \left| X \left( f \right) \right| df= \left| x \left( t \right) \right|_{ \max} ,
 \int \limits_{-\infty}^{+ \infty} \left| x \left( t \right) \right| dt= \left| X \left( f \right) \right|_{ \max}$$
依前面对 $f^{*}$ 的定义，$\left|X \left( f^{*} \right) \right|= \left|X \left( f \right) \right|_{ \max}$，即$X \left( f^{*} \right)= \left|X \left( f \right) \right|_{ \max}e^{j \theta}$，其中$\theta$是 $X\left(f^{*}\right)$的相位。由$\int_{- \infty}^{ \infty} \left| x \left( t \right) \right| dt= \left| X \left( f \right) \right|_{ \max}$知
$$\int_{-\infty}^{ \infty} \left| x \left( t \right) \right| dt=X \left( f^{*} \right) e^{-j \theta}=e^{-j \theta} \int_{- \infty}^{ \infty}x \left( t \right) e^{-j2 \pi f^{*}t}dt$$
即
$$\int_{-\infty}^{ \infty} \left| x \left( t \right) e^{-j \left( 2 \pi f^{*}t+ \theta \right)} \right| dt= \int_{- \infty}^{ \infty}x \left( t \right) e^{-j \left( 2 \pi f^{*}t+ \theta \right)}dt$$
这只能是$\left|x \left( t \right) e^{-j \left( 2 \pi f^{*}t+ \theta \right)} \right|=x \left( t \right) e^{-j \left( 2 \pi f^{*}t+ \theta \right)}$，说明$x \left( t \right) e^{-j \left( 2 \pi f^{*}t+ \theta \right)}$必然是非负的实函数，这样必然有$f^{*} =0$，且x(t)非负（$ \theta=0$）或者非正（$\theta= \pi $）。

假设$x \left( t \right) \geq 0$，则$x \left( t^{*} \right)= \left| x \left( t \right) \right|_{ \max}$，由$\int_{- \infty}^{ \infty} \left| X \left( f \right) \right| df= \left| x \left( t \right) \right|_{ \max}$知： 
$$\int_{-\infty}^{ \infty} \left| X \left( f \right) \right| df=x \left( t^{*} \right)= \int_{- \infty}^{ \infty}X \left( f \right) e^{j2 \pi ft^{*}}df$$
即
$$\int_{- \infty}^{ \infty} \left| X \left( f \right) e^{j2 \pi ft^{*}} \right| df= \int_{- \infty}^{ \infty}X \left( f \right) e^{j2 \pi ft^{*}}df$$
即
$$\left|X \left( f \right) e^{j2 \pi ft^{*}} \right|=X \left( f \right) e^{j2 \pi ft^{*}}$$
即$X \left( f \right) e^{j2 \pi ft^{*}}$是非负的实函数。令$X \left( f \right)=A \left( f \right) e^{j \varphi \left( f \right)}$，其中 $A \left( f \right)= \left| X \left( f \right) \right|$，则$A\left( f \right) e^{j \left( 2 \pi ft^{*}+ \varphi \left( f \right) \right)}$是非负的实函数，，这必然有$t^{*}=0$,$X \left( f \right)= \left| X \left( f \right) \right|$。X ( f ) 是实函数表明 x (t) 是偶函数。

若假设$x \left( t \right) \leq 0$，同理也可得到$t^{*}=0$， X ( f ) 是非正的实函数， x (t) 是偶函数。

总结起来就是要求：x(t)是非负或者非正的偶函数，且在t =0处达到绝对值最大；同时X(f)也必须是非负或非正的实函数。符合这种情况的实例如三角信号： 
\begin{figure}[H]
  \centering
  \includegraphics{images/image5.png}
\end{figure}

由$\left| x \left( t \right) \right|_{ \max}= A$、$\int_{- \infty}^{ \infty} \left| x \left( t \right) \right| dt= A \tau$，可求得$ \alpha= \frac{ A \tau}{ A}= \tau$。由$X \left( f \right)=A \tau \mathrm{sinc}^{2} \left( f \tau \right)$、$\int_{-\infty}^{ \infty} \left| X \left( f \right) \right| df= \int_{- \infty}^{ \infty}X \left( f \right) df=x \left( 0 \right)=A$、$ \left| X \left( f \right) \right|_{ \max}= A\tau $，可求得$\beta= \frac{A}{A \tau}= \frac{1}{ \tau}$。因此$\alpha \beta = 1$。
---------------------------------------
\begin{center}
\Large\bfseries{第三章\ 随机过程 }
\end{center}
---------------------------------------
1．设$Y \left( t \right)=X \left( t \right) \cos \left( 2 \pi f_{c}t+ \theta \right)$，其中 X (t)与$\theta$统计独立， X (t)为 0 均值的平稳随机过程，自相关函数与功率谱密度分别为$R_{{}_{X}}\left(\tau\right)$，$P_{{}_{X}}\left(f\right)$。

(1)若$\theta$在$\left(0,2 \pi \right)$均匀分布，求Y (t) 的均值、自相关函数和功率谱密度； 

(2)若$\theta$为常数，求Y (t) 的均值、自相关函数和功率谱密度。

++++++++++++

解：无论是(1)还是(2)，都有
$$E \left[ Y \left( t \right) \right]=E \left[ X \left( t \right) \right] E \left[ \cos \left( 2 \pi f_{c}t+ \theta \right) \right]=0$$

$$\begin{aligned}
R_{Y}(\tau) & =E[Y(t) Y(t+\tau)] \\
& =E\left[X(t) \cos \left(2 \pi f_{c} t+\theta\right) X(t+\tau) \cos \left(2 \pi f_{c} t+\theta+2 \pi f_{c} \tau\right)\right] \\
& =E[X(t) X(t+\tau)] E\left[\cos \left(2 \pi f_{c} t+\theta\right) \cos \left(2 \pi f_{c} t+\theta+2 \pi f_{c} \tau\right)\right] \\
& =\frac{1}{2} R_{X}(\tau) E\left[\cos 2 \pi f_{c} \tau+\cos \left(4 \pi f_{c} t+2 \theta+2 \pi f_{c} \tau\right)\right] \\
& =\frac{1}{2} R_{X}(\tau) \cos 2 \pi f_{c} \tau+\frac{1}{2} R_{X}(\tau) E\left[\cos \left(4 \pi f_{c} t+2 \theta+2 \pi f_{c} \tau\right)\right]
\end{aligned}$$
在(1)的条件下，$\theta$的概率密度函数为
$$p(\theta)=\begin{cases}\dfrac{1}{2\pi}&\theta\in[0,2\pi)\\0&\text{else}\end{cases}$$
于是
$$E\Big[\cos\Big(4\pi f_ct+2\theta+2\pi f_c\tau\Big)\Big]=\frac1{2\pi}\int_0^{2\pi}\cos\Big(4\pi f_ct+2\theta+2\pi f_c\tau\Big)d\theta=0$$
因此
$$R_Y\left(\tau\right)=\frac12R_X\left(\tau\right)\cos2\pi f_c\tau $$
$$\begin{aligned}P_{Y}\left(f\right)&=\int_{-\infty}^\infty R_Y\left(\tau\right)e^{-j2\pi f\tau}d\tau\\&=\int_{-\infty}^\infty\frac{R_X\left(\tau\right)\cos2\pi f_c\tau}2e^{-j2\pi f\tau}d\tau\\&=\frac{P_X\left(f-f_c\right)+P_X\left(f+f_c\right)}4\end{aligned}$$
在(2)的条件下
$$R_Y\left(\tau\right)=\frac12R_X\left(\tau\right)\cos2\pi f_c\tau+\frac12R_X\left(\tau\right)\cos\left(4\pi f_ct+2\theta+2\pi f_c\tau\right)$$
表明$Y(t)$是循环平稳过程。对时间 $t$ 平均，由于$\overline{\cos\left(4\pi f_ct+2\theta+2\pi f_c\tau\right)}=0$,所以$Y(t)$的平均自相关函数是
$$\overline{R}_Y\left(\tau\right)=\frac12R_X\left(\tau\right)\cos2\pi f_c\tau $$
因此平均功率谱密度是
$$P_Y\left(f\right)=\frac{P_X\left(f-f_c\right)+P_X\left(f+f_c\right)}4$$

---------------------------------------

2. 双边功率谱密度为$\frac{N_0}2$ 的白噪声经过传递函数为$H(f)$ 的滤波器后成为$X(t)$,若
$$H(f)=\begin{cases}\dfrac{T_s}{2}\Big(1+\cos\pi fT_s\Big)&\Big|f\Big|\le\dfrac{1}{T_s}\\\\0&else\end{cases}$$
求 $X(t)$的功率谱密度及功率。

++++++++++++

解：$X( t)$的功率谱密度为
$$P_X\left(f\right)=\frac{N_0}{2}\Big|H\left(f\right)\Big|^2=\begin{cases}\dfrac{N_0T_s^2}{8}\Big(1+\cos\pi fT_s\Big)^2&\Big|f\Big|\leq\dfrac{1}{T_s}\\0&\textit{else}\end{cases}$$
$X(t)$的功率为
$$P=\int_{-\infty}^\infty P_X\left(f\right)df=\int_{-\frac1{T_s}}^{\frac1{T_s}}\frac{N_0T_s^2}8\left(1+\cos\pi T_sf\right)^2df=\frac{3N_0T_s}8$$

---------------------------------------

3．Y(t)是平稳白噪声n(t)通过图 a 所示电路的输出，图a中滤波器的传递函数H(f)如图b示，求Y (t) 的同相分量及正交分量的功率谱密度，并画出图形。 
\begin{figure}[H]
  \centering
  \includegraphics[scale=0.7]{images/image6.png}
  \caption*{图(a)}
\end{figure}

\begin{figure}[H]
  \centering
  \includegraphics{images/image7.png}
  \caption*{图(b)}
\end{figure}
++++++++++++

解：首先平稳过程通过线性系统还是平稳过程。所以$Y(t)$是窄带平稳过程，其带宽为 $B$。若$Y(t)$的功率谱密度为$P_Y(f)$,则根据窄带平稳过程的性质可知$Y(t)$的同相分量$Y_c(t)$、 正交分量$Y_s\left(t\right)$的功率谱密度$P_{Y_{c}}\left(f\right)$及$P_{Y_s}\left(f\right)$满足
$$P_{Y_c}(f)=P_{Y_s}(f)=\begin{cases}P_Y\left(f+f_c\right)+P_Y\left(f-f_c\right)&\left|f\right|\leq\dfrac{B}{2}\\0&\textit{else}\end{cases}$$
其中
$$P_Y\left(f\right)=\frac{N_0}2\left|H(f)\right|^2\left|j2\pi f\right|^2=\begin{cases}2N_0\left(\pi f\right)^2&\left|f\pm f_c\right|\leq\frac B2\\0&else \end{cases}$$
因此
$$\begin{aligned}P_{Y_c}(f)&=P_{Y_s}(f)\\&=\begin{cases}2N_0\pi^2\left(f+f_c\right)^2+2N_0\pi^2\left(f-f_c\right)^2&\left|f\right|\leq\dfrac{B}{2}\\0&else\end{cases}\\&=\begin{cases}4N_0\pi^2\left(f^2+f_c^2\right)&\left|f\right|\leq\dfrac{B}{2}\\0&else\end{cases}\end{aligned}$$
其图形如下
\begin{figure}[H]
  \centering
  \includegraphics[scale=0.7]{images/image8.png}
\end{figure}
---------------------------------------

4.设$\xi_{1}=\int_{0}^{T}n\left(t\right)\varphi_{1}\left(t\right)dt $,$\xi_{2}=\int_{0}^{T}n\left(t\right)\varphi_{1}\left(t\right)dt $,其中$n\left(t\right)$是双边功率谱密度为$\frac{N_{0}}{2} $ 白高斯噪声,$\varphi_{1}\left(t\right)$和 $\varphi_{2}\left(t\right)$为确定函数,求 $\xi_{1}$ 和$\xi_{2}$统计独立的条件。

++++++++++++

解:$ n\left(t\right)$是白噪声意味着 $E{\left[n\left(t\right)\right]}=0$,否则其功率谱密度将在 $f=0$ 处有冲激。所以
$$E{\left[\xi_{1}\right]}=E{\left[\int_{0}^{T}n\left(t\right)\varphi_{1}\left(t\right)dt\right]}=\int_{0}^{T}E{\left[n\left(t\right)\right]}\varphi_{1}\left(t\right)dt=0$$
$$E\left[\xi_2\right]=E\left[\int_0^Tn\left(t\right)\varphi_2\left(t\right)dt\right]=\int_0^TE\left[n\left(t\right)\right]\varphi_2\left(t\right)dt=0$$
又因为 n(t) 是高斯过程，所以$\xi_1,\xi_2$是服从联合高斯分布的随机变量，故欲$\xi_1$和$\xi_2$统计独立，需$E\begin{bmatrix}\xi_1\xi_2\end{bmatrix}=0$。
$$\begin{aligned}E\left[\xi_1\xi_2\right]&=E\left[\int_0^Tn\left(t\right)\varphi_1\left(t\right)dt\int_0^Tn\left(t\right)\varphi_2\left(t\right)dt\right]\\&=E\left[\int_0^T\int_0^Tn\left(t\right)n\left(t^{\prime}\right)\varphi_1\left(t\right)\varphi_2\left(t^{\prime}\right)dtdt^{\prime}\right]\\&=\int_0^T\int_0^TE{\left[n\left(t\right)n\left(t^{\prime}\right)\right]}\varphi_1\left(t\right)\varphi_2\left(t^{\prime}\right)dtdt^{\prime}\\&=\int_0^T\int_0^T\frac{N_0}2\delta\left(t^{\prime}-t\right)\varphi_1\left(t\right)\varphi_2\left(t^{\prime}\right)dtdt^{\prime}\\&=\frac{N_0}2\int_0^T\varphi_1\left(t\right)\varphi_2\left(t\right)dt\end{aligned}$$
所以 1 $\xi_1$和$\xi_2$统计独立的条件为
$$\intop_0^T\varphi_1\left(t\right)\varphi_2\left(t\right)dt=0$$
即 $\varphi_1(t),\varphi_2(t)$正交。

注：本题的结果表明，白高斯噪声在任意两个正交信号上的投影是独立的两个高斯随机变量。进而还能证明，如果$\varphi_{1}\left(t\right)$和$\varphi_{2}\left(t\right)$的能量为 1，则投影的方差是$\frac{N_{0}}{2}$。更一般化的结论是：白高斯噪声在多维正交信号空间中各维上的投影是独立同分布的高斯随机变量，其均值是 0，方差是$\frac{N_{0}}{2}$（参见课本第六章）

---------------------------------------

3.5. 设$X(t)=X_c\left(t\right)\cos2\pi f_ct-X_s\left(t\right)\sin2\pi f_ct$为窄带高斯平稳随机过程，其均值为0，方差为$\sigma_x^2$。信号$A\cos2\pi f_ct+X(t)$ 经过下图所示的相乘低通电路后成为$Y(t)=u(t)+\nu(t)$,其中$u(t)$是$A\cos2\pi f_ct$ 对应的输出，$\nu(t)$ 是$X(t)$对应的输出。假设$X_{c}(t)$及$X_{s}(t)$的带宽等于低通滤波器通频带。

(1)若$\theta$为常数，求$u(t)$和$\nu(t)$的平均功率之比；

(2)若$\theta$是与$X(t)$独立的 0 均值高斯随机变量，其方差为$\sigma^2$,求$u(t)$和$\nu(t)$的平均功率之比。
\begin{figure}[H]
  \centering
  \includegraphics[scale=0.7]{images/image9.png}
  \caption*{题 3.5 图}
\end{figure}

++++++++++++

解：
$$u\left(t\right)=\left[A\cos2\pi f_ct\times\cos\left(2\pi f_ct+\theta\right)\right]_{LPF}=\frac A2\cos\theta $$
$$\begin{aligned}\nu\left(t\right)&=\left[\left(X_c\left(t\right)\cos2\pi f_ct-X_s\left(t\right)\sin2\pi f_ct\right)\times\cos\left(2\pi f_ct+\theta\right)\right]_{LPF}\\&=\frac12X_c\left(t\right)\cos\theta+\frac12X_s\left(t\right)\sin\theta\end{aligned}$$
给定 $\theta$ 时$u\left(t\right)$的功率为
$$P_u=\frac{A^2\cos^2\theta}4$$
$\nu(t)$的平均功率为
$$P_\nu=\frac{\sigma_X^2}4\cos^2\theta+\frac{\sigma_X^2}4\sin^2\theta=\frac{\sigma_X^2}4$$
故在(1)的条件下
$$\frac{P_u}{P_\nu}=\frac{A^2}{\sigma_X^2}\cos^2\theta\text{ 。}$$
在(2)的条件下，$\nu(t)$的平均功率仍然是$P_\nu=\frac{\sigma_X^2}4$,但此时$u\left(t\right)$的平均功率是
$$P_u=E_\theta\left[\frac{A^2\cos^2\theta}4\right]=\frac{A^2}4\int_{-\infty}^\infty\frac{\cos^2\theta}{\sqrt{2\pi\sigma^2}}e^{-\frac{\theta^2}{2\sigma^2}}d\theta $$
所以
$$\begin{aligned}\frac{P_u}{P_\nu}&=\frac{A^2}{\sigma_X^2}E\left[\cos^2\theta\right]\\&=\frac{A^2}{2\sigma_X^2}E\left[1+\cos2\theta\right]\\&=\frac{A^2}{2\sigma_X^2}\left[1+\int_{-\infty}^\infty\frac1{\sqrt{2\pi\sigma^2}}e^{-\frac{\theta^2}{2\sigma^2}}\cos2\theta d\theta\right]\\&=\frac{A^2}{2\sigma_X^2}\biggl(1+e^{-2\sigma^2}\biggr)\end{aligned}$$

---------------------------------------

6.设$n(t)$是均值为 0、双边功率谱密度为$\frac {N_0}{2}=10^{-6}$W/Hz 的白噪声，$y(t)=\frac {dn(t)}{dt}$,将
$y(t)$通过一个截止频率为 $B$=10Hz 的理想低通滤波器得到 $y_o(t)$,求

(1)$y(t)$的双边功率谱密度；

(2)$y_o\left(t\right)$的平均功率。

++++++++++++

解：(1)$P_y\left ( f\right ) = \frac {N_0}2| j2\pi f| ^2= 2\pi ^2N_0f^2= 3. 95\times 10^{- 5}f^2$W/ Hz

(2)$P_{y_o}=\int_{-B}^BP_y\left(f\right)df=2\int_0^B2\pi^2N_0f^2df=\frac{4\pi^2N_0B^3}3=0.0263$W

---------------------------------------

7. 设 $\xi(t)$是高斯白噪声通过截止频率为$f_H$的理想低通滤波器后的输出，今以2$f_H$的速率
对$\xi(t)$抽样，$\xi_1,\xi_2,\cdotp\cdotp\cdotp,\xi_n$是其中的$n$个抽样值，求这$n$个抽样值的联合概率密度。

++++++++++++

解：首先，因为$\xi(t)$是白高$^{|}$斯过程通过线性系统的输出，故$\xi(t)$是 0 均值的高斯过程。设白噪声的功率谱密度为$\frac{N_0}2$,则$\xi(t)$的功率是$N_0f_H$,所以$\xi(t)$的一维概率密度是
$$p_\xi(x)=\frac1{\sqrt{2\pi N_0f_H}}e^{-\frac{x^2}{2N_0f_H}}$$
$\xi(t)$的功率谱密度为
$$P_\xi\left(f\right)=\begin{cases}\dfrac{N_0}{2}&\left|f\right|<f_H\\0&\text{esle}\end{cases}$$
其自相关函数为
$$R_\xi\left(\tau\right)=\int_{-\infty}^\infty P_\xi\left(f\right)e^{j2\pi f\tau}df=N_0f_H\mathrm{~sinc}\left(2f_H\tau\right)$$
对任意整数$k$有
$$\operatorname{sinc}(k)=\begin{cases}1&k=0\\0&k\neq0\end{cases}$$
说明对于$\xi(t)$的任何两个不相同的样值，如果其时间间隔是$\frac1{2f_H}$的整倍数，则这两个样值
是不相关的，因而是独立的。于是$\xi_1,\xi_2,\cdotp\cdotp\cdotp,\xi_n$的联合概率密度为
$$\begin{aligned}p_{{\xi_{1},\xi_{2},\cdots,\xi_{n}}}\left(x_{1},x_{2},\cdots,x_{n}\right)&=p_{{\xi_{1}}}\left(x_{1}\right)p_{{\xi_{2}}}\left(x_{2}\right)\cdotp\cdotp\cdotp p_{{\xi_{n}}}\left(x_{n}\right)\\&=\left(2\pi N_{0}f_{H}\right)^{{-\frac{n}{2}}}e^{{-\frac{x_{1}^{2}+x_{2}^{2}+\cdots+x_{n}^{2}}{2N_{0}f_{H}}}}\end{aligned}$$

---------------------------------------

3.8. 已知$b(t)$的波形如下图示，$b(t)$所受的干扰是功率谱密度为$\frac{N_0}2$ 的白高斯噪声
(1)画出对$b(t)$匹配的匹配滤波器的冲激响应波形；
(2)求匹配滤波器的最大输出信噪比$\gamma_\mathrm{max}$ ;
(3)求输出信噪比最大时刻输出值的概率密度。
\begin{figure}[H]
  \centering
  \includegraphics[scale=0.7]{images/image10.png}
  \caption*{题 3.8 图}
\end{figure}

++++++++++++

解：(1)在白噪声干扰下，对 $b(t)$匹配的匹配滤波器的冲激响应为$h(t)=b(t_0-t)$,考虑到因
果性，取最佳抽样时刻为$t_0=T$ ,于是$h(t)=b(T-t)$,冲激响应波形如下：
\begin{figure}[H]
  \centering
  \includegraphics[scale=0.75]{images/image11.png}
\end{figure}

(2)匹配滤波器的输出的噪声分量是平稳过程，在任何抽样点上，其平均功率都是
$$\begin{aligned}\sigma^{2}&=\int_{-\infty}^\infty\frac{N_0}2\Big|H\left(f\right)\Big|^2df=\frac{N_0}2\int_{-\infty}^\infty\Big|H\left(f\right)\Big|^2df\\&=\frac{N_0}2\int_{-\infty}^\infty h^2\left(t\right)dt=\frac{N_0}2A^2T\end{aligned}$$
信号分量在最佳取样点$t=T$时刻最大，为
$$\int_{-\infty}^\infty b\left(T-\tau\right)h\left(\tau\right)d\tau=\int_0^Tb\left(T-\tau\right)b\left(T-\tau\right)d\tau=A^2T$$
故最大输出信噪比为
$$\gamma_{\max}=\frac{\left(A^2T\right)^2}{\sigma^2}=\frac{2A^2T}{N_0}\text{。}$$
(3)最佳取样时刻的输出值是
$$y=A^2T_b+\xi $$
其中$\xi$是噪声分量，它是均值为 0 方差为$\sigma^2=\frac{N_0}2A^2T$ 的高斯随机变量$^{\mathcal{C}}$因此 y 的概率密度函数为
$$p(y)=\frac1{\sqrt{\pi N_0A^2T}}e^{-\frac{\left(y-A^2T\right)^2}{N_0A^2T}}$$
---------------------------------------

9. 已知 $X_1$为瑞利分布的随机变量，其概率密度函数为
$$p_{X_1}\left(x_1\right)=\frac{x_1}{\sigma^2}e^{-\frac{x_1^2}{2\sigma^2}},x_1\geq0$$
$X_2$为莱斯分布的随机变量，其概率密度函数为
$$p_{X_2}\left(x_2\right)=\frac{x_2}{\sigma^2}e^{-\frac{x_2^2+A^2}{2\sigma^2}}I_0\left(\frac{Ax_2}{\sigma^2}\right),x_2\geq0$$
另外已知$X_1,X_2$统计独立。求$X_1>X_2$的概率。

++++++++++++

解：
$$\begin{aligned}P\left(X_1>X_2\right)&=\iint_{x_1>x_2}p_{X_1,X_2}\left(x_1,x_2\right)dx_1dx_2\\&=\iint_{x_1>x_2}p_{x_1}\left(x_1\right)p_{X_2}\left(x_2\right)dx_1dx_2\\&=\int_0^\infty p_{X_2}\left(x_2\right)\biggl[\int_{x_2}^\infty p_{X_1}\left(x_1\right)dx_1\biggr]dx_2\\&=\int_0^\infty p_{X_2}\left(x_2\right)\biggl[\int_{x_2}^\infty\frac{x_1^2}{\sigma^2}e^{-\frac{x_1^2}{2\sigma^2}}dx_1\biggr]dx_2\\&=\int_0^\infty\frac{x_2}{\sigma^2}e^{-\frac{x_2^2+A^2}{2\sigma^2}}I_0\biggl(\frac{Ax_2}{\sigma^2}\biggr)\times e^{-\frac{x_2^2}{2\sigma^2}}dx_2\\&=\int_0^\infty\frac{x_2}{\sigma^2}e^{-\frac{2x_2^2+A^2}{2\sigma^2}}I_0\biggl(\frac{Ax_2}{\sigma^2}\biggr)dx_2\end{aligned}$$
令$x= \sqrt 2x_2$ ，$A^{\prime }= \frac A{\sqrt 2}$，则
$$P\left(X_1>X_2\right)=\frac{e^{-\frac{A^{\prime2}}{2\sigma^2}}}2\int_0^\infty\frac x{\sigma^2}e^{-\frac{x^2+A^{\prime2}}{2\sigma^2}}I_0\left(\frac{A^{\prime}x}{\sigma^2}\right)dx$$
由于对于任意给定的$A>0,\sigma^2>0$ ，总有
$$\int_0^\infty p_{X_2}(x)dx=1=\int_0^\infty\frac x{\sigma^2}e^{-\frac{x^2+A^2}{2\sigma^2}}I_0\left(\frac{Ax}{\sigma^2}\right)dx$$
故
$$P\left(X_1>X_2\right)=\frac12e^{-\frac{A^{\prime2}}{2\sigma^2}}=\frac12e^{-\frac{A^2}{4\sigma^2}}$$

---------------------------------------

3.10. 下图中，功率谱密度为$\frac{N_0}2$ 的高斯白噪声 $n(t)$通过线性网络 $H_1(f)$和$H_2(f)$后的输出分别为$\xi_1(t)$和$\xi_2(t)$,问何种$H_1(f)$和$H_2(f)$可保证$\xi_1(t)$和$\xi_2(t)$这两个随机过程独立？
\begin{figure}[H]
  \centering
  \includegraphics[scale=0.8]{images/image12.png}
  \caption*{题 3.10 图 }
\end{figure}

++++++++++++

解：由于$n(t)$是高斯白噪声，所以$\xi_1(t)$和$\xi_2(t)$都是 0 均值的平稳高斯过程，欲此二过程
独立，需$\forall\tau,E\left[\xi_1^*\left(t\right)\xi_2\left(t+\tau\right)\right]=0$
$$\begin{aligned}E\left[\xi_1^*\left(t\right)\xi_2\left(t+\tau\right)\right]&=E\left[\int_{-\infty}^\infty h_1^*\left(u\right)n\left(t-u\right)du\int_{-\infty}^\infty h_2\left(\nu\right)n\left(t+\tau-\nu\right)d\nu\right]\\&=E\left[\int_{-\infty}^\infty\int_{-\infty}^\infty h_1^*\left(u\right)h_2\left(\nu\right)n\left(t+\tau-\nu\right)n\left(t-u\right)dud\nu\right]\\&=\int_{-\infty}^\infty\int_{-\infty}^\infty h_1^*\left(u\right)h_2\left(\nu\right)E\left[n\left(t+\tau-\nu\right)n\left(t-u\right)\right]dud\nu\\&=\int_{-\infty}^\infty\int_{-\infty}^\infty h_1^*\left(u\right)h_2\left(\nu\right)\frac{N_0}2\delta\left(\tau+u-\nu\right)dud\nu\\&=\frac{N_0}2\int_{-\infty}^\infty h_1^*\left(u\right)h_2\left(u+\tau\right)du\end{aligned}$$
由 Paserval 定理
$$\begin{aligned}E\left[\xi_1^*\left(t\right)\xi_2\left(t+\tau\right)\right]&=\frac{N_0}2\int_{-\infty}^\infty h_1^*\left(u\right)h_2\left(u+\tau\right)du\\&=\frac{N_0}2\int_{-\infty}^\infty H_1^*\left(f\right)H_2\left(f\right)e^{j2\pi f\tau}df\\&=0\end{aligned}$$
$\int_{-\infty}^{\infty}[H_1^*(f)H_2(f)]e^{j2\pi f\tau}df=0$即是说0的傅氏变换是$H_1^*(f)H_2(f)$,故所需的条件是
$$\forall f,H_1^*\left(f\right)H_2\left(f\right)=0$$
也即
$$\forall f,H_1\left(f\right)H_2\left(f\right)=0。$$

注：一般我们说“$X(t)$和$Y(t)$不相关”是指两个过程不相关，即对于任意的时间$t_1,t_2$,$X(t_1),X(t_2)$都不相关。如果我们想说对于给定的时间 $t$,两个随机变量$X(t)$和$Y(t)$不相关时，一般的说法是“$X(t)$和$Y(t)$在同一时间不相关”。

---------------------------------------

11. 设有复随机信号 $\xi(t)=e^{j(2\pi f_ct+\theta)}$,其中$\theta$ 等概取值于$\left\{0,\frac\pi3,-\frac\pi3\right\}$这 3 种相位。求 $\xi(t)$的均值$E\begin{bmatrix}\xi(t)\end{bmatrix}$和自相关函数$E\begin{bmatrix} \xi ^* ( t) \xi ( t+ \tau ) \end{bmatrix}$, $\xi ( t)$是不是广义平稳过程？

++++++++++++

解：$\xi(t)=e^j\theta e^{j2\pi f_ct}$,其中$e^j\theta$是随机的。
$$E\left[\xi\left(t\right)\right]=E\left[e^{j\theta}\right]e^{j2\pi f_ct}=\frac13\left(1+e^{j\frac\pi3}+e^{-j\frac\pi3}\right)e^{j2\pi f_ct}=\frac23e^{j2\pi f_ct}$$
$$E\left[\xi^*\left(t\right)\xi\left(t+\tau\right)\right]=E\left[\left(e^{-j\theta}e^{-j2\pi f_ct}\right)\left(e^{j\theta}e^{j2\pi f_c(t+\tau)}\right)\right]=e^{j2\pi f_c\tau}$$
虽然 $\xi(t)$ 的自相关函数与$t$ 无关，但因为数学期望是$t$ 的函数，所以它不是平稳过程。

---------------------------------------

12.设
$$x\left(t\right)=\frac1{\sqrt{N}}\sum_{i=1}^N\cos\left(2\pi it+\theta_i\right)$$
其中$\theta_i,i=1,2,\cdots,N$ 是一组独立同分布的随机变量，$\theta_i$均匀分布在$\left[0,2\pi\right]$内。
(1)求 $x(t)$的数学期望 $m_x(t)=E\begin{bmatrix}x(t)\end{bmatrix}$
(2)求 $x(t)$的自相关函数 $R_x\left(t,t+\tau\right)=E\begin{bmatrix}x(t)x(t+\tau)\end{bmatrix}$。
(3)若$N\to\infty$,求$x(t)$的一维分布。

++++++++++++

解：(1)
$$\begin{aligned}m_{x}\left(t\right)&=E\left[x\left(t\right)\right]\\&=\frac1{\sqrt{N}}\sum_{i=1}^NE\left[\cos\left(2\pi it+\theta_i\right)\right]\\&=\frac1{\sqrt{N}}\sum_{i=1}^N\int_0^{2\pi}\frac1{2\pi}\cos\left(2\pi it+\varphi\right)d\varphi\\&=0\end{aligned}$$
(2)
$$\begin{aligned}R_{x}\left(t,t+\tau\right)&=E\Big[x\left(t\right)x\left(t+\tau\right)\Big]\\&=\frac1NE\left[\left(\sum_{i=1}^N\cos\left(2\pi it+\theta_i\right)\right)\left(\sum_{j=1}^N\cos\left(2\pi jt+2\pi j\tau+\theta_j\right)\right)\right]\\&=\frac1NE\left[\sum_{i=1}^N\sum_{j=1}^N\cos\left(2\pi it+\theta_i\right)\cos\left(2\pi jt+2\pi j\tau+\theta_j\right)\right]\\&=\frac1N\sum_{i=1}^N\sum_{j=1}^NE\Big[\cos\left(2\pi it+\theta_i\right)\cos\left(2\pi jt+2\pi j\tau+\theta_j\right)\Big]\end{aligned}$$
由于$\theta_i$与$\theta_j$独立，所以对于$i\neq j$
$$\begin{aligned}&E\left[\cos\left(2\pi it+\theta_i\right)\cos\left(2\pi jt+2\pi j\tau+\theta_j\right)\right]\\&=E\left[\cos\left(2\pi it+\theta_i\right)\right]E\left[\cos\left(2\pi jt+2\pi j\tau+\theta_j\right)\right]=0\end{aligned}$$
对于$i=j$
$$\begin{aligned}E\left[\cos\left(2\pi it+\theta_i\right)\cos\left(2\pi it+2\pi i\tau+\theta_i\right)\right]&=\frac12E\left[\cos2\pi i\tau+\cos\left(4\pi it+2\pi i\tau+\theta_i\right)\right]\\&=\frac12E\left[\cos2\pi i\tau\right]+\frac12E\left[\cos\left(4\pi it+2\pi i\tau+\theta_i\right)\right]\\&=\frac{\cos2\pi i\tau}2\end{aligned}$$
因此
$$R_x\left(t,t+\tau\right)=\frac1{2N}\sum_{i=1}^N\cos2\pi i\tau=R_x\left(\tau\right)$$
上述结果表明 $x(t)$是广义平稳过程。

(3)由中心极限定律知，$N\to\infty$时，$x(t)$将趋向于高斯分布。其均值为$m_x=0$,方差为
$R_x(0)=\frac12$,故一维分布是
$$p(x)=\frac1{\sqrt{\pi}}e^{-x^2}$$
---------------------------------------

13. 设$D(t)=\sum_{n=-\infty}^{\infty}a_ng\left(t-nT-t_0\right)$,其中码元$a_n$ 等概取值于±1,$a_n,a_m(n\neq m)$互相独立；$t_0$ 是在$(0,T)$内均匀分布的随机变量，且与$a_n$ 独立。$g(t)=\begin{cases}1&0\leq t\leq T\\0&else\end{cases}$为码元波形。求$D(t)$的自相关函数和功率谱密度。

++++++++++++

解：
$$\begin{aligned}R_{D}\left(t,t+\tau\right)&=E\left[D\left(t\right)D\left(t+\tau\right)\right]\\&=E\left[\sum_{n=-\infty}^\infty a_ng\left(t-nT-t_0\right)\sum_{m=-\infty}^\infty a_mg\left(t+\tau-mT-t_0\right)\right]\\&=E\left[\sum_{n=-\infty}^\infty\sum_{m=-\infty}^\infty a_na_mg\left(t-nT-t_0\right)g\left(t+\tau-mT-t_0\right)\right]\\&=\sum_{n=-\infty}^\infty\sum_{m=-\infty}^\infty E\left[a_na_m\right]E\left[g\left(t-nT-t_0\right)g\left(t+\tau-mT-t_0\right)\right]\\&=\sum_{n=-\infty}^\infty E\left[g\left(t-nT-t_0\right)g\left(t+\tau-nT-t_0\right)\right]\end{aligned}$$
这个结果显然是$t$的函数，所以$D(t)$不是平稳过程。因$R_D\left(t,t+\tau\right)=R_D\left(t+T,t+T+\tau\right)$,所以$D(t)$是循环平稳过程，其平均自相关函数为
$$\begin{aligned}\overline{R_{D}}\left(\tau\right)&=\frac1T\int_0^TR_D\left(t,t+\tau\right)dt\\&=E\left[\frac1T\int_0^T\sum_{n=-\infty}^\infty g\left(t-nT-t_0\right)g\left(t+\tau-nT-t_0\right)dt\right]\end{aligned}$$
由于$\sum_{\infty}^{\infty}g\left(t-nT-t_0\right)g\left(t+\tau-nT-t_0\right)$是$t$的周期函数，所以
$$\begin{aligned}\overline{R_D}\left(\tau\right)&=E\left[\frac1T\int_0^T\sum_{n=-\infty}^\infty g\left(t-nT-t_0\right)g\left(t+\tau-nT-t_0\right)dt\right]\\&=E\left[\frac1T\int_{t_0}^{t_0+T}\sum_{n=-\infty}^\infty g\left(t-nT-t_0\right)g\left(t+\tau-nT-t_0\right)dt\right]\\&=E\left[\frac1T\int_0^T\sum_{n=-\infty}^\infty g\left(x-nT\right)g\left(x+\tau-nT\right)dx\right]\\&=\frac1T\int_0^T\sum_{n=-\infty}^\infty g\left(x-nT\right)g\left(x+\tau-nT\right)dx\\&=\frac1T\int_0^Tg\left(x\right)g\left(x+\tau\right)dx\\&=\begin{cases}1-\dfrac{\left|\tau\right|}{T}&\left|\tau\right|\leq T\\\quad0&else\end{cases}\end{aligned}$$
求$\overline{R_D}(\tau)$的傅氏变换，得到
$$P_D\left(f\right)=T\mathrm{sinc}^2\left(fT\right)。$$

---------------------------------------

14. 已知平稳遍历白高斯噪声的功率谱密度为$\frac{N_0}2$。此噪声经过一个带宽为$B$ 的滤波器后成为$n(t)$。对$n(t)$作大量的采样测量，问超过正 2 伏的采样数所占的比率是多少？

++++++++++++

解：$n(t)$是 0 均值的高斯平稳遍历随机过程，其方差为$N_0B$,其一维概率密度函数为
$$p_n(x)=\frac1{\sqrt{2\pi N_0B}}\exp\left(-\frac{x^2}{2N_0B}\right)$$
因此
$$P\left(n>2\right)=\frac12\mathrm{erfc}\left(\frac2{\sqrt{2N_0B}}\right)=\frac12\mathrm{erfc}\left(\sqrt{\frac2{N_0B}}\right)$$

erfc 函数

若随机变量$z\sim N(0,1/2),x>0$,则互补误差函数定义为
$$\mathrm{erfc}\left(x\right)\triangleq P\left(\left|z\right|>x\right)$$
误差函数定义为
$$\mathrm{erf}\left(x\right)\triangleq P\left(\left|z\right|\leq x\right)=1-erfc\left(x\right)$$
matlab 中就有 erfc 函数、erf 函数以及 erf 的反函数。
\begin{figure}[H]
  \centering
  \includegraphics[scale=0.6]{images/image13.png}
\end{figure}

由这个定义显然有
$$P\left(z>x\right)=\frac12\mathrm{erfc}\left(x\right)$$
许多求解误码率的问题可归结为求概率$P(u>a)$或$P(u<-a)$的问题，其中
$a>0$ ,随机变量$u\sim N(0,\sigma^2)$,则
$$P\big(u>a\big)=P\bigg(\frac u{\sqrt{2}\sigma}>\frac a{\sqrt{2}\sigma}\bigg)$$
令$z=\frac u{\sqrt2\sigma}$,则$z\sim N\left(0,\frac12\right)$,因此
$$P\left(u>a\right)=P\left(z>\frac a{\sqrt{2}\sigma}\right)=\frac12\operatorname{erfc}\left(\frac a{\sqrt{2}\sigma}\right)$$

---------------------------------------

15. 已知复平稳遍历高斯噪声$n(t)=x(t)+jy(t)$的实部和虚部是平稳遍历的高斯过程， $x(t)$、$y(t)$的均值为 0,方差为 1。对 $n(t)$作大量的采样测量，问采样结果中顺势幅度超过 2 且相位在$\pm\frac\pi8$之内的采样数所占的比率是多少？

++++++++++++

解：$n(t)$的模 $A(t)=\left|n(t)\right|$服从 Rayleigh 分布，其概率密度函数为
$$p_A\left(A\right)=Ae^{-\frac{A^2}2}$$
$n(t)$的相位$\theta(t)$服从均匀分布，其概率密度函数为
$$p_\theta\left(\theta\right)=\frac1{2\pi}\text{,}\theta\in\left[0,2\pi\right]$$
$A$与$\theta$独立，再由于遍历性可知，所求比率为
$$\begin{aligned}P\Bigg(A>2,\Big|\theta\Big|\leq\frac{\pi}{8}\Bigg)&=P\Big(A>2\Big)P\Bigg(\Big|\theta\Big|\leq\frac{\pi}{8}\Bigg)\\&=\int_2^\infty Ae^{-\frac{A^2}2}dA\times\frac{2\times\pi/8}{2\pi}\\&=\frac1{8e^2}\end{aligned}$$

---------------------------------------

16.平稳随机过程 $X(t)$ 的自相关函数为 $R(\tau)$。已知对任意 $t$,$X(t)$和 $X(t+\tau)$当$\tau\to\infty$ 时不相关。问 $X(t)$ 的平均功率，直流功率，交流功率各为多少？

++++++++++++

解：$X(t)$的平均功率是$E\left [ X^2\left ( t\right ) \right ] = R\left ( 0\right )$, 令 $m_X= E\left [ X\left ( t\right ) \right ] = E\left [ X\left ( t+ \tau \right ) \right ]$,则
$$R\left(\infty\right)=\lim_{\tau\to\infty}E\left[X\left(t\right)X\left(t+\tau\right)\right]=\lim_{\tau\to\infty}E\left[X\left(t\right)\right]E\left[X\left(t+\tau\right)\right]=m_X^2$$
即直流功率是$R(\infty)$。交流功率是总功率扣除直流功率，即$R(0)\boldsymbol{-}R\begin{pmatrix}\infty\end{pmatrix}$。

---------------------------------------

17. 设有随机序列$\left\{a_n\right\}=\left\{a_1,a_2,\cdots,a_n,\cdots\right\}$,其中的元素两两不相关且$E[a_n]=0$
$E\left[a_n^2\right]=1$ 。另有序列$\left\{b_n\right\}=\left\{b_0,b_1,b_2,\cdots,b_n,...\right\}$,其中$b_0=1$,$b_n=\alpha_nb_{n-1}$。证明$\left\{b_n\right\}$
中的元素两两不相关且$E[b_n]=0$、$E[b_n^2]=1$。

++++++++++++

证：对于任意的$n\geq1$,
$$b_n=a_nb_{n-1}=a_na_{n-1}b_{n-2}=\cdots=a_na_{n-1}\cdots a_1b_0=\prod_{i=1}^na_n$$
因此
$$E[b_n]=E\left[\prod_{i=1}^na_i\right]=\prod_{i=1}^nE[a_i]=0$$
$$E\left[b_n^2\right]=E\left[\left(\prod_{i=1}^na_i\right)^2\right]=E\left[\prod_{i=1}^na_i^2\right]=\prod_{i=1}^nE\left[a_i^2\right]=1$$
对于任意的$m\geq1$、$n>m$ ,
$$\begin{aligned}E\left[b_nb_{n-m}\right]&=E\left[\prod_{i=1}^na_i\times\prod_{i=1}^{n-m}a_i\right]\\&=E\left[a_n\times\prod_{i=1}^{n-1}a_i\times\prod_{i=1}^{n-m}a_i\right]\\&=E[a_n]\times E\left[\prod_{i=1}^{n-1}a_i\times\prod_{i=1}^{n-m}a_i\right]\\&=0\end{aligned}$$
得证。

---------------------------------------

18. 已知平稳白高斯噪声的功率谱密度为$\frac{N_0}2$。此噪声经过一个冲激响应为$h(t)$的线性系
统成为$x(t)$。若已知$h(t)$的能量为$E_h$,求$x(t)$的功率。

++++++++++++

解：设$h(t)$的傅氏变换为$H(f)$,则有
$$P_x=\int_{-\infty}^\infty\frac{N_0}2\left|H\left(f\right)\right|^2df=\frac{N_0}2\int_{-\infty}^\infty\left|H\left(f\right)\right|^2df=\frac{N_0}2E_h\text{。}$$

---------------------------------------

\begin{center}
\Large\bfseries{第四章\ 模拟通信
系统 }
\end{center}
---------------------------------------

4.1将模拟信号$m(t)=\sin2\pi f_mt$与载波$c(t)=A_c\sin2\pi f_ct$相乘得到双边带抑制载波调幅
(DSB-SC)信号，设$f_c=6f_m$ 

(1)画出 DSB-SC 的信号波形图； 

(2)写出 DSB-SC 信号的傅式频谱，画出它的振幅频谱图；

(3)画出解调框图。

++++++++++++

解：(1) $s(t)=m(t)c(t)=A_c\sin(2\pi f_mt)\sin(2\pi f_ct)=A_c\sin(2\pi f_mt)\sin(12\pi f_mt)$
\begin{figure}[H]
  \centering
  \includegraphics[scale=0.7]{images/image14.png}
\end{figure}
$(2)s\left(t\right)$可化为
$$s(t)=\frac{A_c}2\Big(\cos10\pi f_mt-\cos14\pi f_mt\Big)$$
于是
$$S(f)=\frac{A_c}4{\left[\delta\left(f-5f_m\right)+\delta\left(f+5f_m\right)-\delta\left(f-7f_m\right)-\delta\left(f+7f_m\right)\right]}$$
\begin{figure}[H]
  \centering
  \includegraphics[scale=0.7]{images/image15.png}
\end{figure}
(3)
\begin{figure}[H]
  \centering
  \includegraphics[scale=0.7]{images/image16.png}
\end{figure}

---------------------------------------

4.2 已知 $s(t)=\cos\left(2\pi\times10^4t\right)+4\cos\left(2.2\pi\times10^4t\right)+\cos\left(2.4\pi\times10^4t\right)$是某个 AM 已调信号的展开式

(1)写出该信号的傅式频谱，画出它的振幅频谱图；

(2)写出 $s(t)$的复包络$s_L(t);$

(3)求出调幅系数和调制信号频率；

(4)画出该信号的解调框图。

++++++++++++

解：(1)s(t)的频谱为
$$\begin{aligned}&S\left(f\right)=\frac12\Big[\delta\Big(f+10\times10^3\Big)+\delta\Big(f-10\times10^3\Big)\Big]+2\Big[\delta\Big(f+11\times10^3\Big)+\delta\Big(f-11\times10^3\Big)\Big]\\&+\frac12\Big[\delta\Big(f+12\times10^3\Big)+\delta\Big(f-12\times10^3\Big)\Big]\end{aligned}$$
振幅频谱图如下
\begin{figure}[H]
  \centering
  \includegraphics[scale=0.7]{images/image17.png}
\end{figure}

由此可见，此调幅信号的中心频率为11KHz，两个边带的频率是 10KHz 和 12KHz,因此

调制信号的频率是 1KHz。载频分量是$4\cos\left(22\pi\times10^3t\right)$。

$(2)s(t)$的复包络 $s_L(t)$ 的频谱为
$$S_L\left(f\right)=\delta\left(f+1000\right)+4\delta\left(f\right)+\delta\left(f-1000\right)$$
$$s_L\left(t\right)=4+2\cos2000\pi t$$
(3)由复包络可以写出调制信号为
$$s\begin{pmatrix}t\end{pmatrix}=\mathrm{Re}\left\{s_L\begin{pmatrix}t\end{pmatrix}e^{j2\pi f_ct}\right\}=4\begin{pmatrix}1+0.5\cos2000\pi t\end{pmatrix}\cos2\pi f_ct$$
因此调幅系数为 0.5,调制信号的频率是 1KHz。

(4) 
\begin{figure}[H]
  \centering
  \includegraphics[scale=0.7]{images/image18.png}
\end{figure}

---------------------------------------

4.3 现有一振幅调制信号$s(t)=(1+A\cos2\pi f_mt)\cos2\pi f_ct$,其中调制信号的频率

$f_{\mathrm{m}}$=5KHz,载频$f_\mathrm{c}$=100KHz,常数A=15。

(1)请问此已调信号能否用包络检波器解调，说明其理由；

(2)请画出它的解调框图；

(3)请画出从该接收信号提取载波分量的框图。

++++++++++++

解：(1)此信号无法用包络检波解调，因为能包络检波的条件是$1+A\cos2\pi f_mt\geq0$,而这里
的A=15 使得这个条件不能成立，用包络检波将造成解调波形失真。

(2)只能用相干解调，解调框图如下
\begin{figure}[H]
  \centering
  \includegraphics[scale=0.7]{images/image19.png}
\end{figure}
(3) 由于接收信号的频谱中有离散的载频分量，所以可以用窄带滤波器滤出载频分量，或利用锁相环作为窄带滤波器。
\begin{figure}[H]
  \centering
  \includegraphics[scale=0.7]{images/image20.png}
\end{figure}

---------------------------------------

4.4已知已调信号波形为 s$(t)=2\cos(2\pi f_mt)\cos(2\pi f_ct)-2\sin(2\pi f_mt)\sin(2\pi f_ct)$,其中调制信号为$m( t) = 2\cos \left ( 2\pi f_mt\right )$, $f_c$是载波频率。

(1)求该已调信号的傅式频谱，并画出振幅谱；

(2)写出该已调信号的调制方式；

(3)画出解调框图；

(4)若调制信号$m(t)$的形式对接收端是未知的，请指出收发端解决载波同步的方法。

++++++++++++

解：(1)由于
$$s(t)=2\cos\left[2\pi\left(f_c+f_m\right)t\right]$$
因此
$$S\left(f\right)=\delta\left(f+f_c+f_m\right)+\delta\left(f-f_c-f_m\right)$$
\begin{figure}[H]
  \centering
  \includegraphics[scale=0.7]{images/image21.png}
\end{figure}
(2)该调制方式为上边带幅度调制。

(3)
\begin{figure}[H]
  \centering
  \includegraphics[scale=0.7]{images/image22.png}
\end{figure}
(4)用$m(t)$对载波$\cos\left(2\pi f_ct+\varphi\right)$作单边带调制(以上边带为例)的结果是
$$s_1\left(t\right)=\mathrm{Re}\left\{\left[m\left(t\right)+j\hat{m}\left(t\right)\right]e^{j\left(2\pi f_ct+\varphi\right)}\right\}=\mathrm{Re}\left\{\left[m\left(t\right)+j\hat{m}\left(t\right)\right]e^{j\varphi}e^{j2\pi f_ct}\right\}$$
令
$$m^{\prime}\left(t\right)=\mathrm{Re}\left\{\left[m\left(t\right)+j\hat{m}\left(t\right)\right]e^{j\varphi}\right\}=m\left(t\right)\cos\varphi-\hat{m}\left(t\right)\sin\varphi $$
则
$$\begin{aligned}m'\left(t\right)+j\hat{m}'\left(t\right)&=\left[m\left(t\right)\cos\varphi-\hat{m}\left(t\right)\sin\varphi\right]+j\left[\hat{m}\left(t\right)\cos\varphi+m\left(t\right)\sin\varphi\right]\\&=m\left(t\right)\left(\cos\varphi+j\sin\varphi\right)+j\hat{m}\left(t\right)\left(\cos\varphi+j\sin\varphi\right)\\&=\left[m\left(t\right)+j\hat{m}\left(t\right)\right]e^{j\varphi}\end{aligned}$$
用 $m^\prime(t)$对载波$\cos2\pi f_ct$ 作单边带调制的结果是
$$s_2\left(t\right)=\mathrm{Re}\left\{\left[m^{\prime}\left(t\right)+j\hat{m}^{\prime}\left(t\right)\right]e^{j2\pi f_ct}\right\}=\mathrm{Re}\left\{\left[m\left(t\right)+j\hat{m}\left(t\right)\right]e^{j\varphi}e^{j2\pi f_ct}\right\}$$
说明这两种情形下，即用$m^{\prime}(t)$对载波$\cos2\pi f_ct$ 作单边带调制和用$m(t)$ 对载波$\cos\left(2\pi f_ct+\varphi\right)$作单边带调制得到的结果是完全一样的。因此若接收端未知调制信号，则绝无可能仅根据接收到的波形识别出发送载波是什么。此时，为了使前面框图中的“载波提取”成为可能，一种常用方法是发送端加导频分量，接收端用窄带滤波器(或锁相环)滤出离散的载频分量。

---------------------------------------

4.5 某单边带调幅信号的载波幅度$A_c=100$,载波频率$f_c=800$KHz,调制信号为
$m(t)=\cos2000\pi t+2\sin2000\pi t$

(1)写出$m(t)$ 的 Hilbert 变换$\hat{m}(t)$表达式； 

(2)写出下单边带调制信号的时域表达式； 

(3)画出下单边带调制信号的振幅频谱。

++++++++++++

解：(1) $\cos200\pi t$ 的 Hilbert 变换是$\sin200\pi t$ ,$\sin2000\pi t$ 的 Hilbert 变换是$-\cos200\pi t$ ,因此$m(t)=\cos2000\pi t+2\sin2000\pi t$ 的 Hilbert 变换是：
$$\hat{m}(t)=\sin2000\pi t-2\cos2000\pi t$$
(2)下单边带调制信号为
$$\begin{aligned}s_{\text{下}}(t)&=\frac{A_c}2m\left(t\right)\cos\left(1600\pi\times10^3t\right)+\frac{A_c}2\hat{m}\left(t\right)\sin\left(1600\pi\times10^3t\right)\\&=50\left(\cos2000\pi t+2\sin2000\pi t\right)\cos\left(1600\pi\times10^3t\right)\\&+50\left(\sin2000\pi t-2\cos2000\pi t\right)\sin\left(1600\pi\times10^3t\right)\\&=50\left[\cos2000\pi t\cos\left(1600\pi\times10^3t\right)+\sin2000\pi t\sin\left(1600\pi\times10^3t\right)\right]\\&+100\left[\sin2000\pi t\cos\left(1600\pi\times10^3t\right)-\cos2000\pi t\sin\left(1600\pi\times10^3t\right)\right]\\&=50\cos\left(2\pi\times799\times10^3t\right)-100\sin\left(2\pi\times799\times10^3t\right)\end{aligned}$$
(3) 由$\cos x=Re\left \{ e^{jx}\right \} $, $\sin x=Re\left\{-je^{jx}\right\}$以及$\operatorname{Re}\left\{z\right\}=\frac {z+z^*}{2}$,得
$$\begin{aligned}s\left(t\right)&=50\cos\left(2\pi\times799\times10^3t\right)-100\sin\left(2\pi\times799\times10^3t\right)\\&=50\operatorname{Re}\left\{e^{j\left(2\pi\times799\times10^3\right)t}\right\}+100\operatorname{Re}\left\{2je^{j\left(2\pi\times799\times10^3\right)t}\right\}\\&=50\operatorname{Re}\left\{\left(1+2j\right)e^{j\left(2\pi\times799\times10^3\right)t}\right\}\\&=25\left\{\left(1+2j\right)e^{j\left(2\pi\times799\times10^3\right)t}+\left(1-2j\right)e^{-j\left(2\pi\times799\times10^3\right)t}\right\}\end{aligned}$$
故
$$S\left(f\right)=25\left[\left(1+2j\right)\delta\left(f-799\times10^3\right)+\left(1-2j\right)\delta\left(f+799\times10^3\right)\right]$$
$$\left|S(f)\right|=25\sqrt{5}\left[\delta\left(f+799\times10^3\right)+\delta\left(f-799\times10^3\right)\right]$$
振幅频谱图如下：
\begin{figure}[H]
  \centering
  \includegraphics[scale=0.7]{images/image23.png}
\end{figure}
---------------------------------------

4.6 下图是一种SSB-AM的解调器，其中载频$f_c$=455KHz。 
(1)若图中 A 点的输入信号是上边带信号，请写出图中各点表示式； 
(2)若图中 A 点的输入信号是下边带信号，请写出图中各点表示式，并问图中解调器应
做何修改方能正确解调出调制信号？
\begin{figure}[H]
  \centering
  \includegraphics[scale=0.7]{images/image24.png}
  \caption*{题 6 图}
\end{figure}

++++++++++++

解：记$m(t)$为基带调制信号，$\hat{m}(t)$为其 Hilbert 变换，不妨设载波幅度为2(2$A_{c}=2$)。

(1)A: $s(t)=m(t)\cos2\pi f_ct-\hat{m}(t)\sin2\pi f_ct=m(t)\cos\left(91\times10^4\pi t\right)-\hat{m}(t)\sin\left(91\times10^4\pi t\right)$ 

B:
$$\begin{aligned} s_{B}(t)&=s\left(t\right)\cos\left(91\times10^4\pi t\right)\\&=m(t)\cos^2\left(91\times10^4\pi t\right)-\hat{m}(t)\sin\left(91\times10^4\pi t\right)\cos\left(91\times10^4\pi t\right)\pi t\\&=\frac{m\left(t\right)}2{\left[1+\cos\left(182\times10^4\pi t\right)\right]}-\frac{\hat{m}\left(t\right)}2\sin\left(182\times10^4\pi t\right)\end{aligned}$$
C: $s_C(t)=\frac12m(t)$

D: 
$$\begin{aligned}s_{D}(t)&=s\left(t\right)\sin\left(91{\times}10^4\pi t\right)\\&=m(t)\sin\left(91{\times}10^4\pi t\right)\cos\left(91{\times}10^4\pi t\right)-\hat{m}(t)\sin^2\left(91{\times}10^4\pi t\right)\\&=\frac{m\left(t\right)}2\sin\left(182{\times}10^4\pi t\right)-\frac{\hat{m}\left(t\right)}2\left[1-\cos\left(182{\times}10^4\pi t\right)\right]\end{aligned}$$

E: 
$s_E(t)=-\frac12\hat{m}(t)$

F: 
$ s_F(t)=-\frac12\hat{\hat{m}}(t)=\frac{m\left(t\right)}2$

---------------------------------------









% 33-64页

G:
$s_{_{G}}(t)=m(t)$\\
(2)当A点输入是下单边带信号时，各点信号如下：\\
A：
$$
\begin{aligned}
S(t) &= m(t)\cos2\pi f_{c}t+\hat{m}(t)\sin2\pi f_{c}t \\
&= m(t)\cos\left(91\times10^{4}\,\pi t\right)+{\hat{m}}(t)\sin\left(91\times10^{4}\,\pi t\right)
\end{aligned}
$$
B:
$$
\begin{aligned}
s_{B}(t) &=s\left(t\right)\mathrm{cos}\left(91\times10^{4}\pi t\right) \\
&=m(t)\cos^{2}\left(91\times10^{4}\pi t\right)+\hat{m}(t)\sin\left(91\times10^{4}\pi t\right)\cos\left(91\times10^{4}\pi t\right)\\
&={\frac{m\left(t\right)}{2}}\Big[1+\cos\big(182\times10^{4}\pi t\Big)\Big]+{\frac{\dot{m}\left(t\right)}{2}}\sin\left(182\times10^{4}\pi t\right)
\end{aligned}
$$
C:
$s_{C}(t)={\frac{1}{2}}m(t)$\\
D:
$$
\begin{aligned}
s_{D}(t) &=s\left(t\right)\sin\left(91\times10^{4}\pi t\right) \\
&=m(t)\sin\left(91\times10^{4}\pi t\right)\cos\left(91\times10^{4}\pi t\right)+{\hat{m}}(t)\sin^{2}\left(91\times10^{4}\pi t\right)\\
&={\frac{m\left(t\right)}{2}}\sin\left(182\times10^{4}\pi t\right)+{\frac{\hat{m}\left(t\right)}{2}}\Bigl[1-\cos\Bigl(182\times10^{4}\pi t\Bigr)\Bigr]
\end{aligned}
$$
E:
$s_{E}(t)={\frac{1}{2}}{\hat{m}}(t)$\\
F:
$s_{F}(t)=\frac{1}{2}\hat{m}(t)=-\frac{m\left(t\right)}{2}$\\
G:
$s_{\scriptscriptstyle G}(t)=0$\\
如欲G点输出m(t)，需将最末端的相加改为相减即可，如下图示：
\begin{figure}[H]
  \centering
  \includegraphics{zwz_images/1.png}
\end{figure}
---------------------------------------\\
4.7 一VSB调幅信号的产生框图如图7(a)所示，其中BPF的传递函数H(f)如图7(b)所示。请画出相应的相干解调框图，并说明解调输出不会失真的理由。\\
\begin{figure}[H]
  \centering
  \includegraphics{zwz_images/2.png}
  \caption*{题7图(a)}
  \nonumber
\end{figure}
\begin{figure}[H]
  \centering
  \includegraphics{zwz_images/3.png}
  \caption*{题7图(b)}
\end{figure}
+++++++++++++++++++++++++++++++++++++\\
解：解调框图如下，其中理想低通的截止频率是 W。
\begin{figure}[H]
  \centering
  \includegraphics{zwz_images/4.png}
\end{figure}
接收信号 r(t) 为带通信号，可写成
$r\left(t\right)=\operatorname{Re}\left\{r_{L}\left(t\right)e^{i2\pi f/\prime}\right\}$
其中$r_{L}\left(t\right)$是$r(t)$的复包络。\\
用载波$2\cos2\pi f_{c}t$对$r(t)$进行相干解调后，得到的输出是$y{\bigl(}t{\bigr)}=\operatorname{Re}{\Bigl\{}r_{L}\left(t\right){\Bigr\}}$。\\
不考虑噪声，那么$r(t)$是 DSB 信号$s\left(t\right)=A m\left(t\right)\cos2\pi f_{c}t$通过带通系统$H(f)$后的输出。DSB 信号$s(t)$的复包络是$s_{\mathit{L}}\left(t\right)={\mathit{A}}m\left(t\right)$，其频谱为$S_{L}\left(f\right)=A M\left(f\right)$。$H(f)$的等效低通是 
$$
H_{L}(f) = \left\{
\begin{array}{cc}
H\left(f+f_{c}\right) & \text{if } -\frac{W}{2} < f < \frac{W}{2} \\
0 & \text{else}
\end{array}
\right.
$$
因此$r(t)$的复包络$r_{L}\left(t\right)$的频谱为
$$R_{L}(f)\!=\!H_{L}(f)S_{L}(f)$$
由于
$$y{\bigl(}t{\bigr)}=\operatorname{Re}\left\{r_{L}{\bigl(}t{\bigr)}\right\}={\frac{r_{L}{\bigl(}t{\bigr)}+r_{L}^{*}\left(t\right)}{2}}$$

\begin{align*}
Y(f) &= \frac{1}{2} \left[ \mathcal{R}_{L}(f) + \mathcal{R}_{L}^*(-f) \right] \\
&= \frac{1}{2} \left[ AM(f)H_L(f) + AM(f) \mathcal{H}_L^*(-f) \right] \\
&= \frac{AM(f)}{2} \left[ H_L(f) + H_L^*( -f) \right] 
\end{align*}

本题条件中$H(f)$是实函数，再由下图可知，$H_{L}\left(f\right)+H_{L}^{\ast}\left(-f_{\mathrm{c}}\right)=1$
\begin{figure}[H]
  \centering
  \includegraphics{zwz_images/5.png}
\end{figure}
因此，$Y{\bigl(}f{\bigr)}={\frac{A}{2}}M\left(f\right)$，即$y(t)={\frac{A}{2}}m\left(t\right)$，说明解调输出没有失真。\\
---------------------------------------\\
4.8 某调频信号$s(t)=10\cos\left(2\pi\times10^{6}t+4\cos200\pi t\right)$，求其平均功率、调制指数、最大频偏以及近似带宽。\\
+++++++++++++++++++++++++++++++++++++\\
解：平均功率为

\begin{align*}
P_s &= \frac{(10)^2}{2} = 50 \ \mathrm{W}
\end{align*}

最大频偏为
\begin{align*}
\Delta f_{\mathrm{max}} &= \max \left\{ \frac{1}{2 \pi} \frac{d}{dt} \left( 4 \cos 200 \pi t \right) \right\} = 400 \ \mathrm{Hz}
\end{align*}

记$m(t)$为调制信号，$K_{F}$为调频灵敏度，则由于$4\cos200\pi t=2\pi K_{F}{\int}_{-m}^{t}m\left(\tau\right)d\tau$，故$m{\bigl(}t{\bigr)}=-{\frac{400}{K_{F}}}\mathrm{sin}\,200\pi t$，所以$m(t)$的频率是$f_{m}=100\mathrm{Hz}$，因此调频指数为

\begin{align*}
\beta_f &= \frac{\Delta f_{\mathrm{max}}}{f_m} = 4 
\end{align*}

近似带宽为
\begin{align*}
B &= 2(1 + \beta f)W \\
&= 2 \times 5 \times 100 \ \mathrm{Hz} \\
&= 1000 \ \mathrm{Hz}
\end{align*}



%4.9
---------------------------------------\\
4.9 调角信号$S(t)=100\cos\bigl(2\pi f_{c}t+4\sin2\pi f_{m}t\bigr)$，其中载频$f_{\mathrm{c}}\!=\!10\mathrm{MHz}$，调制信号的频率是$f_{m}=1000\,\mathrm{Hz}$。\\
(1)假设$s(t)$是$\mathrm{FM}$调制，求其调制指数及发送信号带宽；\\
(2)若调频器的调频灵敏度不变，调制信号的幅度不变，但频率$f_{m}$加倍，重复(1)题；\\
(3)假设$s(t)$是$\mathsf{P M}$调制，求其调制指数及发送信号带宽；\\
(4)若调相器的调相灵敏度不变，调制信号的幅度不变，但频率$f_{m}$加倍，重复(3)题。\\
\\
+++++++++++++++++++++++++++++++++++++\\
解：(1) $4\operatorname{Sin}2\pi f_{m}t=2\pi K_{F}{\int_{-\infty}^{t}m\left(\tau\right)}d\tau$，$m\bigl(t\bigr)=A_{m}\cos2\pi f_{m}t=\frac{4000}{K_{F}}\cos2000\pi t$。最大频偏$\Delta\!f_{\mathrm{max}}=K_{F}\left|m\big(t\big)\right|_{\mathrm{max}}=K_{F}\times{\frac{4000}{K_{F}}}=4000\mathrm{Hz}$。由此可得调制指数为
$$\beta_{f}=\frac{\Delta f_{\mathrm{max}}}{f_{m}}=\frac{4000}{1000}=4$$\\
信号带宽近似$$B=2\left(\Delta\!\!\slash_{\mathrm{max}}+f_{m}\right)=2\left(4000+1000\right)=10000\mathrm{Hz}$$

(2)此时最大频偏不变，但因$\ f_{m}$成为 2000Hz ，所以调制指数成为${\frac{4000}{2000}}=2$，近似带宽成为$2\left(4000+2000\right)=12000\mathrm{HZ}$\\
\\
(3)$4\sin2\pi f_{m}t=K_{p}m\left(t\right)$，$m(t)=\frac{4}{K_{p}}\sin2000\pi t$。\\
调相指数为$\beta_{p}=\operatorname*{max}\left\{4\sin2\pi f_{m}t\right\}=4$\\
最大频偏为$\Delta\!f_{\mathrm{max}}=\mathrm{max}\left\{\frac{1}{2\pi}\frac{d}{d t}\big(4\sin2\pi f_{m}t\big)\right\}=4f_{m}=400\mathrm{Hz}$\\
信号带宽近似为：$B=2\left(\Delta{\cal S}_{\mathrm{max}}+f_{m}\right)=10000\mathrm{Hz}$\\
\\
(4) 此时$s(t)=100\cos\biggl(2\pi f_{c}t+K_{p}\times\frac{4}{K_{p}}\sin4\pi f_{m}t\biggr)=100\cos\bigl(2\pi f_{c}t+4\sin4\pi f_{m}t\bigr)$，调相指数仍然是$\beta_{f}=4$。最大频偏为\\$\Delta\!f_{\mathrm{max}}=\mathrm{max}\left\{\frac{1}{2\pi}\frac{d}{d t}\big(4\sin4\pi f_{m}t\big)\right\}=8f_{m}=800\mathrm{Hz}$\\
信号带宽近似为：\\
$B=2\left(\Delta{\cal f}_{\mathrm{max}}+2{\cal f}_{_m}\right)=2\left(8000+2000\right)=20000\mathrm{Hz}$。\\
---------------------------------------\\
%4.10
4.10 下图中的模拟基带信号$m\left(t\right)$不含直流分量，且其频谱在零频率附近为 0，K 为充分大
的常数，载频$\ f_{c}$远远大于$m(t)$的带宽。写出$s_{1}\left(t\right)$，$s_{2}\left(t\right)$，$s_{3}\left(t\right)$，$s_{4}\left(t\right)$，$s_{5}\left(t\right)$的表达式，并指出它们的调制类型。 
\begin{figure}[H]
  \centering
  \includegraphics{zwz_images/6.png}
  \caption*{题4.10图(a)}
\end{figure}
+++++++++++++++++++++++++++++++++++++\\
解：
$$s_{1}\left(t\right)=m\left(t\right)\cos2\pi f_{c}t+\hat{m}\left(t\right)\sin2\pi f_{c}t$$
它是下边带的 SSB 调制。
$$s_{2}\left(t\right)=m\left(t\right)\cos2\pi f_{c}t+K\cos2\pi f_{c}t=K\left[1+{\frac{m\left(t\right)}{K}}\right]\cos2\pi f_{c}t$$
这是标准调幅（AM）。
\begin{align*}
s_3(t) &= K \cos 2\pi f_c t - \left[ \sum_{-\pi}^{t} m(\tau) d\tau \right] \sin 2\pi f_c t \\
&\qquad = K \left[ \cos 2\pi f_c t - \frac{\int_{-\pi}^{t} m(\tau) d\tau}{K} \sin 2\pi f_c t \right] \\
&\approx K \left[ \cos 2\pi f_c t \cos \left( \frac{1}{K} \int_{-\infty}^{t} m(\tau) d\tau \right) \right. \\
&\quad \qquad \left. - \sin \left( \frac{1}{K} \int_{-\infty}^{t} m(\tau) d\tau \right) \sin 2\pi f_c t \right] \\
&= K \cos \left[ 2\pi f_c t + \frac{1}{K} \int_{-\infty}^{t} m(\tau) d\tau \right]
\end{align*}

这是窄带调频信号。
$$S_{4}\left(t\right)=\cos\biggl[2\pi f_{c}t+K_{\,p}\biggr]_{-\infty}^{t}m\bigl(\tau\bigr)d\tau\biggr]$$
这是频率调制(FM)。$$s_{5}\left(t\right)=\cos\left[2\pi f_{c}t+2\pi K_{f}\int_{-\infty}^{t}\frac{d m\left(\tau\right)}{d\tau}d\tau\right]=\cos\Bigl[2\pi f_{c}t+2\pi K_{f}m\left(t\right)\Bigr]$$
这是相位调制(PM)。 


%4.11
---------------------------------------\\
4.11 双边带调幅信号$s\left(t\right)=m\left(t\right)\cos2\pi f_{c}t$的功率谱密度$P_{s}\left(\left.J\right.\right)$如题 4.11 图(a)示，$s(t)$在传输中受到功率谱密度为$\frac{N_{0}}{2}$的加性白高斯噪声$n_{_W}\left(t\right)$的干扰，接收端用题 4.11 图(b)所示的解调框图进行解调。求解调输出的信噪比。
\begin{figure}[H]
  \centering
  \includegraphics{zwz_images/7.png}
  \caption*{题4.11图(a)}
\end{figure}
\begin{figure}[H]
  \centering
  \includegraphics{zwz_images/8.png}
  \caption*{题4.11图(b)}
\end{figure}
+++++++++++++++++++++++++++++++++++++\\
解：由题 4.11 图(a)可知，$m(t)$的带宽是$\mathcal{W}$，$s(t)$的功率是$$P_{s}=\int_{-\infty}^{\infty}P_{s}\left(f\right)d f=2P_{0}W$$
记$m(t)$的功率为$\ P_{\scriptscriptstyle M}$，由$s\left(t\right)=m\left(t\right)\mathrm{cos}\,2\pi f_{c}t$可知
$$P_{s}={\frac{1}{2}}P_{M}$$
因此
$$P_{M}=4N_{0}{\mathcal W}$$
LPF 输出中的信号分量为$\frac{m\left(t\right)}{2}$，因此输出的信号功率为
$${\frac{P_{M}}{4}}=P_{0}W$$
BPF 输出端的噪声是窄带高斯噪声：
$$n\left(t\right)=n_{c}\left(t\right)\mathrm{cos}\,2\pi f_{c}t-n_{s}\left(t\right)\mathrm{sin}\,2\pi f_{c}t$$
其中${n}_{c}(t)$，${n}_{s}(t)$，${n}(t)$的功率都是
$$P_{n}=N_{0}\times2W=2N_{0}W$$
LPF 输出的噪声分量是$\frac{n_{c}\left(t\right)}{2}$，所以输出的噪声功率是
$${\frac{P_{n}}{4}}={\frac{N_{0}W}{2}}$$
因此输出信噪比是
$$\left(\frac{S}{N}\right)_{o}=\frac{P_{0}W}{\frac{1}{2}N_{o}M}=\frac{2P_{o}}{N_{o}}$$



%4.12
---------------------------------------\\
4.12 某模拟广播系统中基带信号$m(t)$的带宽为$W=10\mathrm{KHz}$，峰均功率比（定义为$\frac{|m(t)\vert_{\mathrm{max}}^{2}}{P_{M}}$，其中$\ P_{M}$是$m(t)$的平均功率）是 5。此广播系统的平均发射功率为 40KW，发射信号经过 80dB 信道衰减后到达接收端，并在接收端叠加了双边功率谱密度为$\frac{N_{0}}{2}=1\,0^{-10}\,\mathrm{W/Hz}$的白高斯噪声。\\
(1)若此系统采用 SSB 调制，求接收机可达到的输出信噪比；\\
(2)若此系统采用调幅系数为 0.85 的标准幅度调制(AM)，求接收机可达到的输出信噪比。\\
\\
+++++++++++++++++++++++++++++++++++++\\
解：记到达接收机的已调信号为$s(t)$，记其功率为$\ P_{\boldsymbol{R}}$，则
$$P_{R}=40\times10^{3}\times10^{-8}=4\times10^{-4}\,\mathrm{W}$$
(1)采用 SSB 时，不妨以下单边带为例，接收机输入端的有用信号为
$$s\left(t\right)=A m\left(t\right)\cos2\pi f_{c}t+A\hat{m}\left(t\right)\sin2\pi f_{c}t$$
其中$f_{c}$是载波频率，${\hat{m}}\left(\,t\,\right)$是$m(t)$的 Hilbert 变换。按照 Hilbert 变换的性质，$m(t)$、${\hat{m}}(t)$的功率都是$\ P_{M}$。故而
$$P_{R}=\frac{A^{2}P_{M}}{2}+\frac{A^{2}P_{M}}{2}=A^{2}P_{M}$$
接收机前端可采用频带为$\left[f_{c}-W,f_{c}\right]$的理想 BPF 限制噪声，于是输入到解调器输入端的噪声为
$$n\left(t\right)=n_{c}\left(t\right)\cos2\pi f_{c}t-n_{s}\left(t\right)\sin2\pi f_{c}t$$
其中${n}_{c}(t)$，${n}_{s}(t)$，${n}(t)$的功率都是
$$P_{n}=N_{0}W=2\times10^{-10}\times10^{4}=2\times10^{-6}\,\mathrm{W}$$
采用理想相干解调时的解调输出为$A m\left(t\right)+n_{c}\left(t\right)$，因此输出信噪比是
$$\left({\frac{S}{N}}\right)_{S B}={\frac{A^{2}P_{M}}{P_{n}}}={\frac{4\times10^{-4}}{2\times10^{-6}}}=200，\text{，即23dB}$$
(2) 采用 AM 调制时，解调器输入的有用信号为：￥$$s\left(t\right)=A\left[1+a m_{n}\left(t\right)\right]\cos2\pi f_{c}t$$
其中调制指数 a = 0.85，$m_{n}\left(t\right)=\frac{m\left(t\right)}{\left|m\left(t\right)\right|_{\mathrm{max}}}$，$m_{n}\left(t\right)$的功率是
$${\cal P}_{M_{n}}=\frac{{\cal P}_{M}}{\left|m\left(t\right)\right|_{\mathrm{max}}^{2}}=\frac{1}{5}$$
接收机前端可采用频带为$\left[f_{c}-W,f_{c}+W\right]$的理想 BPF 限制噪声，于是输入到解调器输入端的噪声为
$$n\left(t\right)=n_{c}\left(t\right)\cos2\pi f_{c}t-n_{s}\left(t\right)\sin2\pi f_{c}t$$
其中${n_{c}}(t)$，${n_{s}}(t)$，${{n}}(t)$的功率都是$$P_{n}=2N_{0}W=2\times2\times10^{-10}\times10^{4}=4\times10^{-6}\text{。}$$
采用理想的包络检波器得到的输出为${A}\Bigl[1+a m_{n}\bigl(t\bigr)\rbrack+n_{c}\left(t\right)$，因此输出信噪比是$\left(\frac{S}{N}\right)_{A M}=\frac{A^{2}a^{2}P_{M_{n}}}{P_{n}}$\\
由此得
$$\left(\frac{S}{N}\right)_{A M}=\frac{2P_{_R}}{1+a^{2}P_{M_{s}}}\times\frac{a^{2}P_{M_{s}}}{P_{n}}=\frac{2\times4\times10^{-4}}{\left(1+\frac{1}{0.85^{2}\times0.2}\right)\times4\times10^{-6}}\approx25.25\text{，即14.02dB}$$


%4.13
---------------------------------------\\
4.13 有 12 路话音信号$m_{1}\left(t\right)$,$m_{2}\left(t\right)$,...,$m_{12}\left(t\right)$，它们的带宽都限制在(0,4000) Hz 范围内。将这 12 路信号以SSB/FDM方式复用为m(t)，再将m(t)通过FM方式传输，如图(a)所示。其中SSB/FDM的频谱安排如图(b)所示。已知调频器的载频为$\ f_{\mathrm{{c}}}$，最大频偏为 480KHz。
\begin{figure}[H]
  \centering
  \includegraphics{zwz_images/9.png}
  \caption*{题4.13图(a)}
\end{figure}
\begin{figure}[H]
  \centering
  \includegraphics{zwz_images/10.png}
  \caption*{题4.13图(b)}
\end{figure}
(1)求 FM 信号的带宽；\\ 
(2)画出解调框图；\\ 
(3)假设 FM 信号在信道传输中受到加性白高斯噪声干扰，求鉴频器输出的第 1 路噪声平均功率与第 12 路噪声平均功率之比。\\
\\
+++++++++++++++++++++++++++++++++++++\\
解：(1) m(t) 的带宽为 48KHz，所以 FM 信号的带宽近似为$2\left(480+48\right)=1056\,\mathrm{KHZ}$\\
(2)解调框图如下。其中$f_{i}=4\left(i-1\right)\mathrm{KHZ}\;,\quad i=1,2,\cdots,12$。第 i 个 BPF 的通带为$\left(f_{i},f_{i}+4\mathrm{KHZ}\right)$。LPF 的截止频率是 4KHz。
\begin{figure}[H]
  \centering
  \includegraphics{zwz_images/11.png}
\end{figure}
(3)鉴频器输出的噪声功率谱密度与$f^{2}$成正比，即
$$P_{n_{o}}\left(f\right)=K f^{2}$$
其中 K 是常系数。落在第 1 路频带范围内的输出噪声功率为
$$P_{1}=\int_{0}^{4000}K f^{2}d f=\frac{\left(4000\right)^{3}K}{3}=\frac{4^{3}\times10^{9}}{3}\,K$$
落在第 12 路频带范围内的输出噪声功率为
$$P_{12}=\int_{4400}^{4800}K f^{2}d f=\frac{\left(48000^{3}-44000^{3}\right)K}{3}=\frac{\left(48^{3}-44^{3}\right)\times10^{9}}{3}\,K$$
所以
$${\frac{P_{12}}{P_{1}}}={\frac{48^{3}-44^{3}}{4^{3}}}=12^{3}-1\;1^{3}=397\text{，相当于26dB}$$


%4.14
---------------------------------------\\
4.14 一射频方向性天线的等效噪声温度Ta为 55K，天线输出经过一个损耗为 2dB的馈线连接
至接收机前置射频放大器输入端，该射频放大器的功率增益$K_{p}$为 20dB，等效带宽为 10MHz，噪声系数为 2dB，设室温为T=300K，求：\\
(1)前置射频放大器的等效噪声温度；\\ 
(2)馈线与前置射频放大器级联的噪声系数及等效噪声温度； \\
(3)求馈线输入点的单边噪声功率谱密度N0值。\\
\\
+++++++++++++++++++++++++++++++++++++\\
解：(1)记$\textstyle T_{\scriptscriptstyle A}$为放大器的等效噪声温度。放大器自身的噪声折合到它的输入端的功率为
$$K T\left(F_{_A}-1\right)B=K T_{_A}B$$
因此
$$T_{_A}=T\left(F_{_A}-1\right)=300\times\left(10^{0.2}-1\right)=175.5\mathrm{K}$$
(2) 设总噪声系数是$F_{_{A F}}$，可用课本上的公式(4.8.28)求出$F_{_{A F}}$：
$$F_{_{A F}}=F_{_1}+{\frac{F_{2}-1}{K_{p a_{1}}}}+{\frac{F_{3}-1}{K_{p a_{1}}K_{p a_{2}}}}+\cdots$$
本题只有两级，第一级馈线的衰减量是${L}=10^{0.2}\left(2\mathrm{dB}\right)$，噪声系数是$F_{1}=L$，增益是$K_{p_{a_{1}}}={\frac{1}{L}}$；第二级前置射频放大器的噪声系数是$F_{2}=F_{1}=L$。因此，
$$F_{_{A F}}=F_{1}+{\frac{F_{2}-1}{K_{_{p a_{1}}}}}=L+{\frac{L-1}{1/L}}=L^{2}=10^{0.4}\text{，即4dB}$$
设馈线和放大器级联的总等效噪声温度是$T_{\mathit{A F}}$，则由式(4.8.33)可得
$$T_{A F}=\left(F_{AF}-1\right)T=\left(10^{0.4}-1\right)\times300=453.6\,\mathrm{K}$$
本题中总噪声系数的求解也可以不直接套用式(4.8.28)。馈线输入的可获噪声功率是$K T B$，输出的噪声功率也是$K T B$，而馈线的损耗是$\boldsymbol{\mathit{L}}$，因此我们可以理解成：馈线输出的噪声中$\frac{K T B}{L}$是输入噪声所致，其余的$P_{L}=K T B-\frac{K T B}{L}=K T B\left(1-\frac{1}{L}\right)$是馈线自身的噪声。馈线自身的噪声$P_{L}$折合到它的输入端是$L P_{L}=K T B\left(L-1\right)$。而放大器自身的噪声折合到放大器输入端是$P_{A}=K T B\left(F_{2}-1\right)=K T B\left(L-1\right)$（本题条件中，$F_{2}=L$），$\ P_{\mathcal{A}}$再折合到馈线输入端是$L P_{A}=K T B\left(L^{2}-L\right)$。所有噪声（馈线和放大器产生的噪声、馈线输入噪声）折合到馈线输入端是$L P_{L}+L P_{4}+K T B=K T B L^{2}$。这说明总噪声系数是$F_{_{A F}}=L^{2}$。\\
(3)如果馈线输入是常温电阻，则前面的推导结果表明馈线输入端的总噪声功率是$P_{0}=K T B L^{2}=K\left(T L^{2}\right)B$，即总噪声温度是$T L_{}^{2}$。扣除常温电阻，那么馈线+放大器的等效噪声温度是$T_{A F}=\left(L^{2}-1\right)T$。今馈线输入的噪声源的温度是$T_{a}=55 K$，于是馈线输入点总的噪声功率是$K\left[T_{a}+T_{A F}\right]B$，因而在 B=1Hz 带宽上的测量到的功率（此即单边噪声功率谱密度，测量时的 1Hz 包含正负频率两部分）为
$$N_{0}=K\left(T_{a}+T_{A F}\right)=1.38\times10^{-23}\times\left[55+300\left(10^{0.4}-1\right)\right]=7.02\times10^{-21}\operatorname{W/Hz}$$



%4.15
---------------------------------------\\
4.15 已知某模拟基带系统中调制信号$m(t)$的带宽是$W{=}5{\mathrm{kHz}}$。发送端发送的已调信号功率是$\ P_{t}$，接收功率比发送功率低 60dB。信道中加性白高斯噪声的单边功率谱密度为$N_{0}=10^{-13}\mathrm{W}/\mathrm{H}\mathrm{Z}$\\
(1)如果采用 DSB-SC，请\\
\indent\text{(a)推导出输出信噪比$\left({\frac{S}{N}}\right)_{o}$和输入信噪比$\left({\frac{S}{N}}\right)_{i}$的关系；}\\
\indent\text{(b)若要求输出信噪比不低于 30dB，发送功率至少应该是多少？}\\
(2)如果采用 SSB，重做第(1)问。\\
+++++++++++++++++++++++++++++++++++++\\
解：(1)(a)解调输入信号可写为
\begin{align*}
r(t) &= A m(t) \cos 2\pi f_c t + n(t) \\
&= A m(t) \cos 2\pi f_c t + n_c(t) \cos 2\pi f_c t - n_s(t) \sin 2\pi f_c t
\end{align*}

输入信噪比为
$$\left(\frac{S}{N}\right)_{i}=\frac{\frac{A^2P_{m}}2}{2N_{0}M}=\frac{A^{2}P_{m}}{4N_{0}M}$$
解调输出为${\mathit{A}}m\left(t\right)+n_{c}\left(t\right)$，输出信噪比为
$$\left(\frac{S}{N}\right)_{o}=\frac{A^{2}\overline{{{m^{2}\left(t\right)}}}}{\overline{{{n}}}_{c}^{2}\left(t\right)}=\frac{A^{2}P_{m}}{2N_{0}M}$$
因此
$${\frac{\left(S/N\right)_{o}}{\left(S/N\right)_{i}}}=2$$
(b)输入信噪比
$$\left(\frac{S}{N}\right)_{i}=\frac{P_{r}}{2N_{0}M}=\frac{P_{t}\times10^{-6}}{2\times10^{-13}\times5\times10^{3}}=10^{3}P_{t}$$
输出信噪比
$$\left(\frac{S}{N}\right)_{o}=2\Biggl(\frac{S_{i}}{N_{i}}\Biggr)=2000P_{t}=10^{3}$$
故$P_{t}=0.5$瓦。\\
(2)(a)解调输入信号可写为
\begin{align*}
 r(t) &= A m(t) \cos 2\pi f_c t + A \hat{m}(t) \sin 2\pi f_c t + n(t)\\
 &= A m(t) \cos 2\pi f_c t + A \hat{m}(t) \sin 2\pi f_c t + n_c(t) \cos 2\pi f_c t - n_s(t) \sin 2\pi f_c t 
 \end{align*} 

输入信噪比为
$$\left(\frac{S}{N}\right)_{i}=\frac{A^{2}\overline{{{m^{2}\left(t\right)}}}}{N_{0}M}$$
解调输出为$A m\left(t\right)+n_{c}\left(t\right)$，输出信噪比为$$\left(\frac{S}{N}\right)_{o}=\frac{A^{2}\overline{{{m^{2}\left(t\right)}}}}{\overline{{{n_{c}^{2}\left(t\right)}}}}=\frac{A^{2}P_{m}}{N_{0}W}$$
因此$\frac{\left(S/N\right)_{o}}{\left(S/N\right)_{i}}=1$\\
(b)输入信噪比
$$\left(\frac{S}{N}\right)_{i}=\frac{P_{r}}{N_{0}W}{=}\frac{P_{t}\times10^{-6}}{10^{-13}\times5\times10^{3}}=2000P_{t}$$
输出信噪比
$$\left(\frac{S}{N}\right)_{o}=\left(\frac{S_{i}}{N_{i}}\right)=2000P_{t}=10^{3}$$
故$P_{t}=0.5$瓦。



%4.16
---------------------------------------\\
4.16 已知某调频系统中，调频指数是$\beta_{f}$，到达接收端的 FM 信号功率是$P_{\boldsymbol{R}}$，信道噪声的单
边功率谱密度是$\textstyle{N_{0}}$，基带调制信号$m\left(t\right)$的带宽是$\textstyle W$，解调输出的信噪比和输入信噪比之比为$3\beta_{f}^{2}\left(\beta_{f}+1\right)$。 \\
\indent\text{(1)求解调输出信噪比；}\\
\indent\text{(2)如果发送端将基带调制信号$m\left(t\right)$变成$2m\left(t\right)$，接收端按照此条件}\\
\indent\text{设计解调器，请问输出信噪比将大约增大多少分贝？}\\
+++++++++++++++++++++++++++++++++++++\\
解：(1)输入信噪比是
$$\left(\frac{S}{N}\right)_{i}=\frac{P_{\!R}}{2N_{\!0}W\!\left(\beta_{\!f}+1\right)}$$
输出信噪比是
$$\left(\frac{S}{N}\right)_{o}=3\beta_{f}^{2}\big(\beta_{f}+1\bigg)\bigg(\frac{S}{N}\bigg)=\frac{3\beta_{f}^{2}P_{R}}{2N_{0}M}$$
(2)此时${\boldsymbol{\beta}}_{f}$变成了${\mathcal{2}}{\boldsymbol{\beta}}_{f}$，输出信噪比是原来的 4 倍，即增加了 6dB。



%4.17
---------------------------------------\\
4.17 两个不包含直流分量的模拟基带信号$m_{1}(t)$、$m_{2}(t)$被同一射频信号同时发送，发送信号
为 $s\left(t\right)=m_{1}\left(t\right)\cos\omega_{c}t+m_{2}\left(t\right)\mathrm{Sin}\;\omega_{c}t+K\cos\omega_{c}t$，其中载频$f_{c}{\mathrm{=}}10\,\mathrm{MHZ}$，K是常数。已
知$m_{1}(t)$与$m_{2}(t)$的傅氏频谱分别为$M_{1}(f)$及$M_{2}(f)$，它们的带宽分别为 5KHz与
10KHz。 \\
(1)请计算 s(t)的带宽；\\
(2)请写出 s(t)的傅氏频谱表示式； \\
(3)画出从s(t)得到m1(t)及m2(t)的解调框图。\\
\\
+++++++++++++++++++++++++++++++++++++\\
解：(1)${s}(t)$由两个 DSB 及一个单频组成，这两个 DSB 的中心频率相同，带宽分别是 10KHz
和 20KHz，因此总带宽是 20KHz；\\
(2) 
\begin{align*}
S(f) &=\frac{1}{2}\left[M_1(f+f_c)+M_1(f-f_c)\right] + \frac{j}{2}\left[M_2(f+f_c)-M_2(f-f_c)\right] \\
&\qquad + \frac{K}{2}\left[\delta(f+f_c)+\delta(f-f_c)\right]
\end{align*}

(3) 
\begin{figure}[H]
  \centering
  \includegraphics{zwz_images/12.png}
\end{figure}


%4.18
---------------------------------------\\
4.18 设$m{\bigl(}t{\bigr)}={\frac{1}{\sqrt{N}}}\sum_{i=1}^{N}\cos2\pi i t$\\
(1)求 m(t) 的平均功率$P_{M}$；\\
(2)若用$m\left(t\right)$进行 AM 幅度调制，求调制效率$~~~~~\eta~~~~~~~~~~~~~~~~~~~~~~~~~~~~~~~~~~~~~~~~~~~~~$作为调制指数${\boldsymbol{a}}$的函数$\eta(a)$，并就$N=3$的情形画出$\eta(a)$。\\
+++++++++++++++++++++++++++++++++++++\\
解：(1)$\cos2\pi i t,i=1,2,\cdots,N$是频率分别为 1、2、…、NHz 的 N 个余弦信号，其功率均为$\frac{1}{2}$。因此$\sum_{i=1}^{N}\cos2\pi it$的功率是$\frac{N}{2}$，$m\left(t\right)=\frac{1}{\sqrt{N}}\sum_{i=1}^{N}\cos2\pi i t$的平均功率是$P_{M}={\frac{1}{2}}$。 \\
(2)$\left|m\left(t\right)\right|_{\mathrm{max}}={\frac{1}{\sqrt{N}}}\sum_{i=1}^{N}1={\sqrt{N}}$，将$m\left(t\right)$的幅度归一化，得
$$m_{n}\left(t\right)=\frac{m\left(t\right)}{\left|m\left(t\right)\right|_{\mathrm{max}}}=\frac{1}{N}\sum_{i=1}^{N}\cos2\pi i t$$
其功率为
$$P_{M_{n}}=\frac{1}{2N}$$
调制指数是 a，因此 AM 信号是
$$s\left(t\right)=A_{c}\left[1+a m_{n}\left(t\right)\right]\cos2\pi f_{c}t$$
其效率为
$$\eta(a)=\frac{a^{2}P_{M_{a}}}{1+a^{2}P_{M_{a}}}=\frac{a^{2}}{\frac{2N}{{1+\frac{a^{2}}{2N}}}}=\frac{a^{2}}{2N+a^{2}}$$
当 N＝3 时
$$\eta(a)={\frac{a^{2}}{6+a^{2}}}$$
a 的取值范围是$\scriptstyle[0,\,1]$，$\eta(a)$的图形如下
\begin{figure}[H]
  \centering
  \includegraphics{zwz_images/13.png}
\end{figure}



%4.19
---------------------------------------\\
4.19 若 SSB 信号发送端的载波初始相位是 0，解调器的本地载波的初始相位是$\varphi$。求解调输出中有用信号的功率占总输出信号功率的比率。\\
+++++++++++++++++++++++++++++++++++++\\
解：问题也相当于发送的载波相位是$-\varphi$，本地载波相位是 0。于是发送信号是
\begin{align*}
s(t) &= A_c m(t) \cos(2\pi f_c t - \varphi) \mp A_c \hat{m}(t) \sin(2\pi f_c t - \varphi) \\
&= \Re \left\{ A_c [m(t) \pm j \hat{m}(t)] e^{j(2\pi f_c t - \varphi)} \right\} \\
&= \Re \left\{ A_c [m(t) \pm j \hat{m}(t)] e^{-j\varphi} e^{j2\pi f_c t} \right\}
\end{align*}

其复包络是$$s_{L}\left(t\right)=A_{c}\left[m\left(t\right)\pm j\hat{m}\left(t\right)\right]e^{-j\varphi}$$
相干解调器的参考载波是$\mathrm{cos}\,2\pi f_{c}t$，所以解调输出是复包络的实部（忽略系数）
$$m_{o}\left(t\right)=\mathrm{Re}\left\{s_{L}\left(t\right)\right\}=A_{c}\left[m\left(t\right)\mathrm{cos}\,\varphi\pm\hat{m}\left(t\right)\mathrm{sin}\,\varphi\right]$$
其中有用信号是$d_{c}m(t)\cos\varphi$，其功率是
$$P_{1}=A_{c}^{2}P_{M}\cos^{2}\varphi$$
无用信号是$\pm A_{c}{\hat{m}}\left(t\right)\sin\varphi$，其功率是
$$P_{2}=A_{c}^{2}P_{M}\sin^{2}\varphi$$
因此有用信号所占的功率比例是
$$\frac{P_{1}}{P_{1}+P_{2}}=\frac{\cos^{2}\varphi}{\cos^{2}\varphi+\sin^{2}\varphi}=\cos^{2}\varphi$$



%4.20
---------------------------------------\\
4.20 设$m_{1}\left(t\right)$、$m_{2}\left(t\right)$是频谱分布在 300~3400Hz 内的两个话音信号，功率都是 1。今发送$s\left(t\right)=m_{1}\left(t\right)\mathrm{cos}\,2\pi f_{c}t-m_{2}\left(t\right)\mathrm{sin}\,2\pi f_{c}t$，问\\
(1)接收端如何解出$m_{1}\left(t\right)$、$m_{2}\left(t\right)$?
(2)若$s(t)$叠加了一个双边功率谱密度为$N_{0}/2$的白高斯噪声，求解调输出为$m_{1}\left(t\right)$这一支路的输出信噪比。\\
+++++++++++++++++++++++++++++++++++++\\
解：(1)
\begin{figure}[H]
    \centering
    \includegraphics{zwz_images/14.png}
\end{figure}
其中 BPF 的中心频率是${f_{c}}$，带宽是 2×3400Hz 。\\
(2) BPF 输出的噪声是
\begin{align*}
n(t) &= \Re \left\{ [n_c(t) + j n_s(t)] e^{j 2 \pi f_c t} \right\} \\
&= n_c(t) \cos 2 \pi f_c t - n_s(t) \sin 2 \pi f_c t
\end{align*}

其平均功率是
$$P_{n}=\overline{{{n^{2}\left(t\right)}}}=\overline{{{n_{c}^{2}\left(t\right)}}}=\overline{{{n_{s}^{2}\left(t\right)}}}=6800N_{0}$$
$m_{1}\left(t\right)$这一路解调输出的有用信号是$m_{1}\left(t\right)$，其功率为 1；输出噪声是$\,n_{c}\left(t\right)$，故输出信噪
比是$\frac{1}{6800N_{\mathrm{o}}}$。 \\



%4.21
---------------------------------------\\
4.21 题 4.21 图所示系统中调制信号$m(t)$的均值为 0，功率谱密度为
$$
P_{M}(f) = \left\{
\begin{array}{cc}
N_{m}\left(1 - \frac{|f|}{f_{m}}\right) & \text{if } |f| \leq f_{m} \\
0 & \text{if } |f| > f_{m}
\end{array}
\right.
$$
白高斯噪声$n_{w}\left(t\right)$的双边功率谱密度为$N_{0}/2$。
\begin{figure}[H]
  \centering
  \includegraphics{zwz_images/15.png}
  \caption*{题4.21图}
\end{figure}
(1)写出${{s}}(t)$的复包络$s_{L}\left(t\right)$，画出其功率谱密度；\\
\indent\text{(2)图中的 BPF 应该如何设计？} \\
\indent\text{(3)求相干解调器的输出信噪比。}\\
+++++++++++++++++++++++++++++++++++++\\
解：(1)由图可见，${{s}}(t)$的表达式为
$$s\left(t\right)=m\left(t\right)A_{c}\cos2\pi f_{c}t+\hat{m}\left(t\right)A_{c}\sin2\pi f_{c}t$$
因此${{s}}(t)$的复包络为
$$
s_{L}(t) = A_{c} \left[m(t) - j \tilde{m}(t)\right]
$$
$m(t) + i\hat{m}(t)$ 是解析信号，因此 $z(t)=m(t) - i\hat{m}(t)=[m(t) + i\hat{m}(t)]^{*}$ 的傅氏谱是
$$
Z(f) = 
\begin{cases}
  2M(f) & \text{if } f < 0 \\
  0     & \text{if } f > 0
\end{cases}
$$
因此$s_{L}\left(t\right)$的功率谱密度是
\begin{align*}
P_{s_L}(f) &= \left\{
\begin{array}{ll}
4A_c^2 P_M(f) & f < 0 \\
0 & f > 0
\end{array}
\right. \\
&=  4N_m A_c^2 \left(1 + \frac{f}{f_m} \right) & f < 0 \\
& & \\
0 & f > 0 
\end{align*}

\begin{figure}[H]
  \centering
  \includegraphics{zwz_images/16.png}
\end{figure}
(2)由于${{s}}(t)$是下单边带信号，所以 BPF 的通频带应该是$\left[f_{c}-f_{m},f_{c}\right]$、、
(3)为方便计算，假设 LPF 的幅度增益是$\frac{2}{A_{c}}$，这个假设对输出信噪比没有影响。此时，相干解调器的输出是
$$y{\ (t){\mid}}=m\left(t\right)+{\frac{1}{A_{c}}}n_{c}\left(t\right)$$
其中$n_{c}\left(t\right)$是 BPF 输出的窄带噪声的同相分量，其功率为$f_{m}N_{0}$。由于$m\left(t\right)$的功率是
$$P_{M}=\int_{-\infty}^{\infty}P_{M}\left(f\right)d f=2\int_{0}^{f_{m}}N_{m}\left(1-\frac{f}{f_{m}}\right)d f=f_{m}N_{m}$$
所以输出信噪比是
$$\left(\frac{\bf S}{\cal N}\right)_{o}=\frac{N_{m}f_{m}}{\frac{1}{{\cal A}_{c}^{2}}\,{\cal N}_{o}f_{m}}=\frac{{\cal A}_{c}^{2}N_{m}}{{\cal N}_{0}}$$


%4.22
---------------------------------------\\
4.22 下图是对 DSB 信号进行相干解调的框图。图中${n}_{w}\left(t\right)$是均值为 0、双边功率谱密度为$\frac{N_{\mathrm{o}}}{2}$的加性白高斯噪声，本地恢复的载波和发送载波有固定的相位差${\boldsymbol{\theta}}$。请导出该系统的输出信噪比表达式。
\begin{figure}[H]
  \centering
  \includegraphics{zwz_images/17.png}
\end{figure}

+++++++++++++++++++++++++++++++++++++\\
解一：LPF 输出的信号分量是
$$m\bigl(t\bigr)\cos2\pi f_{c}t\times2\cos\bigl(2\pi f_{c}t+\theta\bigr)=^{\mathrm{LPF}}m\bigl(t\bigr)\cos\theta$$
其功率为${\overline{{m^{2}\left(t\right)}}}\cos^{2}\theta$。\\
BPF 输出的窄带噪声是
$$n\left(t\right)=n_{c}\left(t\right)\cos2\pi f_{c}t-n_{s}\left(t\right)\sin2\pi f_{c}t$$
其中$n_{c}\left(t\right)$，$n_{s}\left(t\right)$的功率都是$2f_{m}N_{0}$。LPF 输出的噪声分量是 
$$n_{o}\left(t\right)=n_{c}\left(t\right)\mathrm{cos}\,\theta+n_{s}\left(t\right)\mathrm{sin}\,\theta$$
由于$n_{c}\left(t\right)$，$n_{s}\left(t\right)$均值为 0，且在同一时刻不相关，所以输出噪声功率为 
\begin{align*}
E\left[ n_{o}^2(t) \right] &= E\left[ n_c^2(t) \cos^2 \theta \right] + E\left[ n_s^2(t) \sin^2 \theta \right] \\
&= 2N_0 f_m \Bigl[ \cos^2 \theta + \sin^2 \theta \Bigr] = 2N_0 f_m 
\end{align*}

从而输出信噪比为
$$\left(\frac{S}{N}\right)_{o}=\frac{m^{2}\left(t\right)\mathrm{Cos}^{2}\,\theta}{2N_{0}f_{m}}$$
解二：令$\,t=t-\tau$，其中$\tau$满足$2\pi f_{c}\tau=\theta$，再考虑到对$m\left(t\right)$和${n}_{w}\left(t\right)$做任何时延都不影响问题，于是原题演变为下图 
\begin{figure}[H]
  \centering
  \includegraphics{zwz_images/18.png}
\end{figure}
图中${{s}}(t)$的复包络是 
$$s_{L}\left(t\right)=m\left(t\right)e^{-j\theta}$$
BPF 输出端的窄带噪声的复包络为
$$n_{L}\left(t\right)=n_{c}\left(t\right)+j n_{s}\left(t\right)$$
其中$n_{c}\left(t\right)$，$n_{s}\left(t\right)$的功率都是$2f_{m}N_{0}$。BPF 输出的总的信号的复包络为
\begin{align*}
r_{L}(t) &= s_{L}(t) + n_{L}(t) \\
&= m(t) e^{-j\theta} + n_{c}(t) + j n_{s}(t)
\end{align*}

解调输出是
$$y{\big(}t{\big)}=\operatorname{Re}\left\{r_{L}{\big(}t{\big)}\right\}=m{\big(}t{\big)}\cos\theta+n_{c}{\big(}t{\big)}$$
因此输出信噪比是
$$\left(\frac{S}{N}\right)_{o}=\frac{m^{2}\left(t\right)\mathrm{Cos}^{2}\,\theta}{2N_{0}f_{m}}$$



%4.23
---------------------------------------\\
4.23 立体声调频发送端方框图如下所示，其中$m_{_{L}}\left(t\right)$和 $m_{_{R}}\left(t\right)$分别表示来自左右声道的电信号，它们的最高频率都不超过 15kHz。
\begin{figure}[H]
  \centering
  \includegraphics{zwz_images/19.png}
  \caption*{\text{题4.23图}}
\end{figure}
(1)请画出图中$m\left(t\right)$的频谱示意图； \\
\indent\text{(2)请画出立体声调频信号的接收框图。}\\
+++++++++++++++++++++++++++++++++++++\\
\indent\text{解：(1)}
\begin{figure}[H]
  \centering
  \includegraphics{zwz_images/20.png}
\end{figure}
(2)
\begin{figure}[H]
  \centering
  \includegraphics{zwz_images/21.png}
\end{figure}


%4.24
---------------------------------------\\
4.24 12 路话音信号，每路的带宽是 0.3-3.4kHz。今用两级 SSB 频分复用系统实现复用。也即，先把这 12 路话音分 4 组经过第一级 SSB 复用，再将 4 个一级复用的结果经过第二级 SSB 复用。已知第一级复用时采用上边带 SSB，3 个载波频率分别为 0KHz, 4KHz 和8KHz。第二级复用采用下单边带 SSB，4 个载波频率分别为 84KHz, 96KHz，108KHz, 120KHz。试画出此系统发送端框图，并画出第一级和第二级复用器输出的频谱（标出频率值，第一级只画一个频谱）。 \\
+++++++++++++++++++++++++++++++++++++\\
解：
\begin{figure}[H]
  \centering
  \includegraphics{zwz_images/22.png}
\end{figure}

第一级复用器输出的频谱为：
\begin{figure}[H]
  \centering
  \includegraphics{zwz_images/23.png}
\end{figure}

第二级复用器输出的频谱为：
\begin{figure}[H]
  \centering
  \includegraphics{zwz_images/24.png}
\end{figure}


%4.25
---------------------------------------\\
4.25 单边带调幅信号$s_{\mathrm{SSB}}\left(t\right)=A_{c}\left[\,m\left(t\right)\mathrm{cos}\,2\pi f_{c}t\mp\hat{m}\left(t\right)\mathrm{sin}\,2\pi f_{c}t\right]$，其中$m\left(t\right)$是 0 均值平稳随机过程，其自相关函数是$R_{m}\left(\tau\right)$。请：\\
\indent\text{(1)写出$s_{\mathrm{SSB}}\left(t\right)$的复包络$S_{L}\left(t\right)$，求出复包络的功率谱密度；}\\
\indent\text{(2)写出$s_{\mathrm{SSB}}\left(t\right)$的解析信号${\tilde{s}}(t)$，求出解析信号的功率谱密度；}\\
\indent\text{(3)求出$s_{\mathrm{SSB}}\left(t\right)$的功率谱密度。}\\
+++++++++++++++++++++++++++++++++++++\\
解：(1)$s_{\mathrm{SSB}}\left(t\right)$的复包络为
$$s_{L}\left(t\right)=A_{c}\left[m\left(t\right)\pm\hat{m}\left(t\right)\right]$$
其中${\hat{m}}\left(t\right)$是$m\left(t\right)$的希尔伯特变换，式中的±号分别对应上单边带和下单边带。$s_{_{L}}\left(t\right)$的自
相关函数为
\begin{align*}
R_{L}(t, \tau) &= A_c^2 E\biggl[ s_L^*(t) S_L(t+\tau) \biggr] \\
&= A_c^2 E\left\{\left[m\left(t\right) \pm j\hat{m}\left(t\right)\right]^* \left[m\left(t+\tau\right) \pm j\hat{m}\left(t+\tau\right)\right]\right\} \\
&= A_c^2 E\left\{\left[m^*(t) \mp j\hat{m}^*(t)\right] \left[m\left(t+\tau\right) \pm j\hat{m}\left(t+\tau\right)\right]\right\} \\
&= A_c^2 E\Bigl[ m^*(t) m(t+\tau) + \hat{m}^*(t)\hat{m}(t+\tau) \pm j m^*(t)\hat{m}(t+\tau) \mp j \hat{m}^*(t) m(t+\tau) \Bigr] \\
&= A_c^2 \Bigl[ R_m(\tau) + R_{\hat{m}}(\tau) \pm j R_{m\hat{m}}(\tau) \mp j R_{\hat{m}m}(\tau) \Bigr] 
\end{align*}

由希尔伯特变换的性质，有 $$R_{\hat{m}}\left(\tau\right)=R_{m}\left(\tau\right)$$
$$R_{m\hat{m}}\left(\tau\right)=\hat{R}_{m}\left(\tau\right)=-R_{\hat{m}m}\left(\tau\right)$$
其中${\hat{R}}_{m}\left(\tau\right)$是函数$R_{m}\left(\tau\right)$的希尔伯特变换。于是 
$$R_{L}\left(\tau\right)=2A_{c}^{2}\left[R_{m}\left(\tau\right)\pm j\hat{R}_{m}\left(\tau\right)\right]$$
对上式做傅氏变换，设$R_{m}\left(\tau\right)$的傅氏变换是$P_{m}\left(f\right)$。先考虑上单边带的情形，注意到$R_{m}\left(\tau\right)+j\hat{R}_{m}\left(\tau\right)$是$R_{m}\left(\tau\right)$的解析信号，其傅氏变换只有正频率分量，于是
$P_{s_{L},\mathrm{USR}}(f)=\left\{
\begin{array}{cc} 
4\mathcal{A}_{c}^{2}P_{m}(f) & \text{if } f > 0 \\ 
0 & \text{if } f < 0 
\end{array}
\right.$
对于下单边带，$R_{m}\left(\tau\right)-j\hat{R}_{m}\left(\tau\right)$是$R_{m}\left(\tau\right)$的解析信号的共轭，其傅氏变换只有负频率分量，于是
$$ P_{s_{L},\mathrm{LSB}}(f) = 
\begin{cases}
0 & \text{if } f > 0 \\
4A_{c}^{2}P_{m}(f) & \text{if } f < 0
\end{cases} $$
(2)$s_{\mathrm{SSB}}\left(t\right)$的解析信号为
\begin{align*}
\tilde{S}(t) &= S_{\mathrm{SSB}}(t) + j\hat{S}_{\mathrm{SSB}}(t) \\
&= A_c \left[ m(t)\cos 2\pi f_c t \mp \hat{m}(t)\sin 2\pi f_c t \right] + j A_c \left[ m(t)\sin 2\pi f_c t \pm \hat{m}(t)\cos 2\pi f_c t \right] \\
&= A_c m(t) \left(\cos 2\pi f_c t + j \sin 2\pi f_c t \right) \pm j A_c \hat{m}(t) \left(\cos 2\pi f_c t + j \sin 2\pi f_c t \right) \\
&= A_c \left[ m(t) \pm j \hat{m}(t) \right] e^{j 2\pi f_c t} \\
&= s_L(t) e^{j 2\pi f_c t}
\end{align*}

它的自相关函数为
$$R_{\bar{s}}\left(t,\tau\right)=E\Bigl[\tilde{s}^{*}\left(t\right)\tilde{s}\left(t+\tau\right)\Bigr]$$
做傅氏反变换得
$$P_{\bar{s}}\left(J\right)=P_{s_{L},\mathrm{USB}}\left(J-f_{c}\right)$$
对于上单边带调制
$$
P_{\bar{s}}(f) = 
\begin{cases}
4\mathcal{A}_{c}^{2}P_{m}(f - f_{c}) & \text{if } f > f_{c} \\
0 & \text{if } f < f_{c}
\end{cases}
$$
对于下单边带调制
$$
P_{\bar{s}}(f) = 
\begin{cases}
0 & \text{if } f > f_{c} \\
\mathcal{A}_{c}^{2}P_{m}(f - f_{c}) & \text{if } f < f_{c}
\end{cases}
$$
(3)$s_{\mathrm{SSB}}\left(t\right)$可以表示为$$S_{\mathrm{SSB}}\left(t\right)=\mathrm{Re}\left\{\widetilde{S}\left(t\right)\right\}=\frac{1}{2}\biggl[\widetilde{S}\left(t\right)+\widetilde{S}^{*}\left(t\right)\biggr]$$
其自相关函数为
\begin{align*}
R_{\mathrm{SSB}}(\tau) &= \frac{1}{4} E\left\{ \left[\bar{\tilde{S}}(t)+\tilde{\tilde{S}}^{*}(t)\right]^* \left[\tilde{\tilde{S}}(t+\tau)+\tilde{\tilde{S}}^{*}(t+\tau)\right]\right\} \\
&= \frac{1}{4} E\left\{ \left[\tilde{S}^*(t)+s(t)\right] \left[\tilde{S}(t+\tau)+\tilde{S}^*(t+\tau)\right]\right\} \\
&= \frac{1}{4} E\left\{ \tilde{S}^*(t+\tau) + \left[ \tilde{S}^*(t) \tilde{S}(t+\tau) \right]^* + \tilde{S}(t) \tilde{S}(t+\tau) + \left[ \tilde{S}(t) \tilde{S}(t+\tau) \right]^* \right\} \\
&= \frac{1}{4} \left\{ R_{\tilde{s}}(\tau) + R_{\tilde{s}}^*(\tau) + E\left[\tilde{s}(t)\tilde{s}(t+\tau)\right] + E^*\left[\tilde{s}(t)\tilde{s}(t+\tau)\right] \right\}
\end{align*}

其中
\begin{align*}
E\left[\tilde{S}(t)\tilde{S}(t+\tau)\right] &= E\left[S_{L}(t)e^{j2\pi f_{c}t} \times S_{L}(t+\tau)e^{j2\pi f_{c}(t+\tau)}\right] \\
&= E\left[ S_{L}(t)S_{L}(t+\tau) \right] e^{j2\pi f_{c}\tau} e^{j4\pi f_{c}t} \\
&= R_{_{S_{L}^{*}S_{L}}}(\tau) e^{j2\pi f_{c}\tau} e^{j4\pi f_{c}t} 
\end{align*}

其中
\begin{align*}
{\cal{R}}_{s_{L}^{*}s_{L}}(t,t+\tau) &= E\left[s_{L}(t)s_{L}(t+\tau)\right] \\
&= E\left\{\left[m(t) \pm j\hat{m}(t)\right]\left[m(t+\tau) \pm j\hat{m}(t+\tau)\right]\right\} \\
&= E\left[m(t)m(t+\tau) - \hat{m}(t)\hat{m}(t+\tau) \pm \hat{m}(t)m(t+\tau) \pm m(t)\hat{m}(t+\tau)\right] \\
&= R_m(\tau) - R_{\hat{m}}(\tau) \pm j R_{\hat{m}m}(\tau) \pm j R_{m\hat{m}}(\tau) = 0
\end{align*}

因此
$$R_{\mathrm{SSB}}\left(t,\tau\right)=\frac{1}{4}\bigg\{R_{\bar{s}}\left(\tau\right)+R_{\bar{s}}^{\ast}\left(\tau\right)\bigg\}$$
做傅氏变换，注意$R_{\bar{s}}^{*}\left(\tau\right)$的傅氏变换是$P_{s}^{*}\left(-f\right)=P_{s}\left(-f\right)$，因此
$$P_{\mathrm{{SSB}}}\left(f\right)={\frac{1}{4}}{\Bigl[}P_{\bar{s}}\left(f\right)+P_{\bar{s}}\left(-f\right){\Bigr]}={\frac{1}{4}}{\Bigl[}P_{s_{L}}\left(f-f_{c}\right)+P_{s_{L}}\left(-f-f_{c}\right){\Bigr]}$$
对于上单边带 
$$
P_{\mathrm{USB}}(f) = 
\begin{cases}
A_{c}^{2}P_{m}(f - f_{c}) & \text{if } f > f_{c} \\
A_{c}^{2}P_{m}(f + f_{c}) & \text{if } f < -f_{c} \\
0 & \text{otherwise}
\end{cases}
$$
对于下单边带，







%第四章开始
---------------------------------------\\
\begin{center}
\Large\bfseries{第四章\ 模拟通信系统 }
\end{center}


%4.1
---------------------------------------\\
4.1 将模拟信号$m(t) = \sin(2 \pi f_m t)$与载波$c(t)=A_{c}\sin2\pi f_{c}t$相乘得到双边带抑制载波调幅（DSB-SC）信号，设$f_{c}=6f_{m}$\\
(1)画出 DSB-SC 的信号波形图；\\
(2)写出 DSB-SC 信号的傅式频谱，画出它的振幅频谱图；\\
(3)画出解调框图。\\
+++++++++++++++++++++++++++++++++++++\\
解：(1)$s(t)=m(t)c(t)=A_{c}\sin\left(2\pi f_{m}t\right)\sin\left(2\pi f_{c}t\right)=A_{c}\sin\left(2\pi f_{m}t\right)\sin\left(12\pi f_{m}t\right)$
\begin{figure}[H]
  \centering
  \includegraphics{zwz_images/25.png}
\end{figure}
(2)s(t) 可化为
$$s(t)=\frac{A_{c}}{2}\Big(\cos10\pi f_{m}t-\cos14\pi f_{m}t\Big)$$
于是
$$S(f)=\frac{A_{c}}{4}\Big[\delta\big(f-5f_{m}\big)+\delta\big(f+5f_{m}\big)-\delta\big(f-7f_{m}\big)-\delta\big(f+7f_{m}\big)\Big]$$
\begin{figure}[H]
  \centering
  \includegraphics{zwz_images/26.png}
\end{figure}
(3)
\begin{figure}[H]
  \centering
  \includegraphics{zwz_images/27.png}
\end{figure}


%4.2
---------------------------------------\\
4.2 已知$s(t)=\cos\left(2\pi\times10^{4}t\right)+4\cos\left(2.2\pi\times10^{4}t\right)+\cos\left(2.4\pi\times10^{4}t\right)$是某个 AM 已调信号的展开式\\
(1)写出该信号的傅式频谱，画出它的振幅频谱图；\\
(2)写出$s(t)$的复包络$s_{L}(t)$；\\
(3)求出调幅系数和调制信号频率； \\
(4)画出该信号的解调框图。\\
+++++++++++++++++++++++++++++++++++++\\
解：(1)$s(t)$的频谱为
\begin{align*}
S(f) &= \frac{1}{2} \left[\delta(f + 10\times 10^3) + \delta(f - 10\times 10^3)\right] \\
&\quad + 2 \left[\delta(f + 11\times 10^3) + \delta(f - 11\times 10^3)\right] \\
&\quad + \frac{1}{2} \left[\delta(f + 12\times 10^3) + \delta(f - 12\times 10^3)\right]
\end{align*} 

振幅频谱图如下
\begin{figure}[H]
    \centering
    \includegraphics{zwz_images/28.png}
\end{figure}
由此可见，此调幅信号的中心频率为11 ，两个边带的频率是 10KHz 和 12KHz，因此调制信号的频率是 1KHz。载频分量是$4\cos\!\left(22\pi\!\times\!10^{3}t\right)$。\\
\\
(2)s(t) 的复包络$s_{L}\left(t\right)$的频谱为
$${\cal S}_{L}\left(f\right)=\delta\left(f+1000\right)+4\delta\left(f\right)+\delta\left(f-1000\right)$$
$$s_{_L}(t)=4+2\cos2000\pi t$$
(3)由复包络可以写出调制信号为$$s\left(t\right)=\mathrm{Re}\left\{s_{L}\left(t\right)e^{j2\pi f_{c}t}\right\}=4\left(1+0.5\cos2000\pi t\right)\cos2\pi f_{c}t$$
因此调幅系数为 0.5，调制信号的频率是 1KHz。\\
(4)
\begin{figure}[H]
    \centering
    \includegraphics{zwz_images/29.png}
\end{figure}


%4.3
---------------------------------------\\
4.3 现有一振幅调制信号$s(t)=(1+A\cos2\pi f_{m}t)\cos2\pi f_{c}t$，其中调制信号的频率
$f_{\mathrm{m}}{=}{5}\mathrm{K}\mathrm{H}\mathrm{z}$，载频$f_{\mathrm{c}}{=}{100}\mathrm{K}\mathrm{H}\mathrm{z}$，常数A=15。\\
(1)请问此已调信号能否用包络检波器解调，说明其理由； \\
(2)请画出它的解调框图； \\
(3)请画出从该接收信号提取载波分量的框图。 \\
+++++++++++++++++++++++++++++++++++++\\
解：(1)此信号无法用包络检波解调，因为能包络检波的条件是$1+A\cos2\pi f_{m}t\geq0$，而这里的 A=15 使得这个条件不能成立，用包络检波将造成解调波形失真。\\
(2) 只能用相干解调，解调框图如下 
\begin{figure}[H]
    \centering
    \includegraphics{zwz_images/30.png}
\end{figure}
(3) 由于接收信号的频谱中有离散的载频分量，所以可以用窄带滤波器滤出载频分量，或利用锁相环作为窄带滤波器。
\begin{figure}[H]
    \centering
    \includegraphics{zwz_images/31.png}
\end{figure}


%4.4
---------------------------------------\\
4.4 已知已调信号波形为$s(t)=2\cos\bigl(2\pi f_{m}t\bigr)\cos\bigl(2\pi f_{c}t\bigr)-2\sin\bigl(2\pi f_{m}t\bigr)s\mathrm{in}\bigl(2\pi f_{c}t\bigr)$，其中调制信号为$m(t)=2\cos\bigl(2\pi f_{m}t\bigr)$，$f_{c}$是载波频率。\\ 
(1)求该已调信号的傅式频谱，并画出振幅谱；\\ 
(2)写出该已调信号的调制方式； \\
(3)画出解调框图；\\
(4)若调制信号$m(t)$的形式对接收端是未知的，请指出收发端解决载波同步的方法。\\
+++++++++++++++++++++++++++++++++++++\\
解：(1)由于
$$s(t)=2\cos\Bigl[2\pi\big(f_{c}+f_{m}\big)t\Bigr]$$
因此
$$S\left(f\right)=\delta\left(f+f_{c}+f_{m}\right)+\delta\left(f-f_{c}-f_{m}\right)$$
\begin{figure}[H]
    \centering
    \includegraphics{zwz_images/32.png}
\end{figure}
(2)该调制方式为上边带幅度调制。\\
\indent{\text{(3)}}
\begin{figure}[H]
    \centering
    \includegraphics{zwz_images/33.png}
\end{figure}
(4)用$m(t)$对载波$\cos\left(2\pi f_{c}t+\varphi\right)$作单边带调制（以上边带为例）的结果是
$$s_{1}\left(t\right)=\mathrm{Re}\left\{\left[m\left(t\right)+j\hat{m}\left(t\right)\right]e^{j\left(2\pi f_{\ell}+\varphi\right)}\right\}=\mathrm{Re}\left\{\left[m\left(t\right)+j\hat{m}\left(t\right)\right]e^{j\varphi}e^{j2\pi f_{\mathrm{c}}t}\right\}$$
令$$m^{\prime}(t)=\mathrm{Re}\left\{\left[m\big(t\big)+j\hat{m}\big(t\big)\right]e^{j\varphi}\right\}=m\big(t\big)\cos\varphi-\hat{m}\big(t\big)\sin\varphi$$
则
\begin{align*}
m'(t) + j\hat{m}'(t) &= \left[m(t)\cos\varphi - \hat{m}(t)\sin\varphi\right] + j \left[\hat{m}(t)\cos\varphi + m(t)\sin\varphi\right] \\
&= [m(t) - j\hat{m}(t)](\cos\varphi + j\sin\varphi) \\
&= [m(t) + j\hat{m}(t)] e^{j\varphi}
\end{align*}

用${\mathfrak{m}}^{'}(t)$对载波$\cos2\pi f_{c}t$作单边带调制的结果是
$$s_{2}\left(t\right)=\mathrm{Re}\left\{\left[m^{\prime}(t)+j\hat{m}^{\prime}(t)\right]e^{j2\pi f_{c}t}\right\}=\mathrm{Re}\left\{\left[m\left(t\right)+j\hat{m}(t)\right]e^{j\varphi}e^{j2\pi f_{c}t}\right\}$$
说明这两种 情形下，即用${\mathfrak{m}}^{\prime}(t)$对载波$\cos2\pi f_{c}t$作单边带调制 和用$m(t)$对载波$\cos\left(2\pi f_{c}t+\varphi\right)$作单边带调制得到的结果是完全一样的。因此若接收端未知调制信号，则绝无可能仅根据接收到的波形识别出发送载波是什么。此时，为了使前面框图中的“载波提取”成为可能，一种常用方法是发送端加导频分量，接收端用窄带滤波器（或锁相环）滤出离散的载频分量。 \\



%4.5
---------------------------------------\\
4.5 某单边带调幅信号的载波幅度$A_{c}=100$，载波频率$f_{c}=800\mathrm{KHZ}$，调制信号为$m(t)=\cos2000\pi t+2\sin2000\pi t$\\
(1)写出$m(t)$的 Hilbert 变换${\hat{m}}(t)$表达式； \\
 (2)写出下单边带调制信号的时域表达式； \\
 (3)画出下单边带调制信号的振幅频谱。 \\
 \\
 +++++++++++++++++++++++++++++++++++++\\
解：(1)$\cos200\pi t$的 Hilbert 变换是$\sin200\pi t$，$\sin2000\pi t$的 Hilbert 变换是$-\cos200\pi t$，
因此$m(t)=\cos2000\pi t+2\sin2000\pi l$的 Hilbert 变换是：
$${\hat{m}}\left(t\right)=\sin2000\pi t-2\cos20000\pi t$$
(2)下单边带调制信号为
\begin{align*}
s_{B}(t) &= \frac{A_{c}}{2} m(t) \cos(1600 \pi \times 10^{3} t) + \frac{A_{c}}{2} \hat{m}(t) \sin(1600 \pi \times 10^{3} t) \\
&= 50 \left( \cos(2000 \pi t) + 2 \sin(2000 \pi t) \right) \cos \left( 1600 \pi \times 10^{3} t \right) \\
&\quad + 50 \left( \sin(2000 \pi t) - 2 \cos(2000 \pi t) \right) \sin \left( 1600 \pi \times 10^{3} t \right) \\
&= 50 \left[ \cos(2000 \pi t) \cos \left( 1600 \pi \times 10^{3} t \right) + \sin(2000 \pi t) \sin \left( 1600 \pi \times 10^{3} t \right) \right] \\
&\quad + 100 \left[ \sin(2000 \pi t) \cos \left( 1600 \pi \times 10^{3} t \right) - \cos(2000 \pi t) \sin \left( 1600 \pi \times 10^{3} t \right) \right] \\
&= 50 \cos \left( 2 \pi \times 799 \times 10^{3} t \right) - 100 \sin \left( 2 \pi \times 799 \times 10^{3} t \right)
\end{align*}

(3) 由$\mathrm{cos}\,x=\mathrm{Re}\left\{e^{jx}\right\}$，$\sin x=\operatorname{Re}{\big\{}-j e^{j x}{\big\}}$以及$\mathrm{Re}\left\{z\right\}={\frac{z+z^{\star}}{2}}$，得
\begin{align*}
s(t) &= 50 \cos\left(2\pi \times 799 \times 10^{3} t\right) - 100 \sin\left(2\pi \times 799 \times 10^{3} t\right) \\
&= 50 \, \mathrm{Re}\left\{ e^{j(2\pi \times 799 \times 10^{3} t)} \right\} + 100 \, \mathrm{Re}\left\{ 2j e^{j(2\pi \times 799 \times 10^{3} t)} \right\} \\
&= \mathrm{Re}\left\lbrace 25 \left( (1+2j) e^{j(2\pi \times 799 \times 10^{3} t)} + (1-2j) e^{-j(2\pi \times 799 \times 10^{3} t)} \right) \right\rbrace
\end{align*}

故
$$S{\bigl(}f{\bigr)}=2S{\Bigl[}{\bigl(}1+2j{\bigr)}\delta{\bigl(}f-799\times10^{3}{\bigr)}+{\bigl(}1-2j{\bigr)}\delta{\bigl(}f+799\times10^{3}{\bigr)}{\bigr]}$$
$$\left|S\left(f\right)\right|=25\sqrt{5}\left[\delta\left(f+799\times10^{3}\right)+\delta\left(f-799\times10^{3}\right)\right]$$
振幅频谱图如下：
\begin{figure}[H]
    \centering
    \includegraphics{zwz_images/34.png}
\end{figure}



%4.6
---------------------------------------\\
4.6 下图是一种SSB-AM的解调器，其中载频$f_{\mathrm{c}}\!=\!455{\mathrm{KHz}}$。\\
(1)若图中 A 点的输入信号是上边带信号，请写出图中各点表示式； \\
(2)若图中 A 点的输入信号是下边带信号，请写出图中各点表示式，并问图中解调器应做何修改方能正确解调出调制信号？
\begin{figure}[H]
    \centering
    \includegraphics{zwz_images/35.png}
    \caption*{题6图}
\end{figure}
+++++++++++++++++++++++++++++++++++++\\
解：记 m(t) 为基带调制信号，${\hat{m}}(t)$为其 Hilbert 变换，不妨设载波幅度为 2（$2A_{c}=2$）。\\
(1)$\mathrm{A}\colon{\boldsymbol{s}}(t)=m(t)\cos2\pi f_{c}t-{\hat{m}}(t)\sin2\pi f_{c}t=m(t)\cos\left(91\times10^{4}\,{\boldsymbol{\pi}}t\right)-{\hat{m}}(t)\sin\left(91\times10^{4}\,{\boldsymbol{\pi}}t\right)$
B：
\begin{align*}
s_{B}(t) &= s(t) \cos(91 \times 10^{4} \pi t) \\
&= m(t) \cos^{2}(91 \times 10^{4} \pi t) - \hat{m}(t) \sin(91 \times 10^{4} \pi t) \cos(91 \times 10^{4} \pi t) \\
&= \frac{m(t)}{2} \left[ 1 + \cos(182 \times 10^{4} \pi t) \right] - \frac{\hat{m}(t)}{2} \sin(182 \times 10^{4} \pi t)
\end{align*}

C:$ S_{C}(t)\!=\!{\frac{1}{2}}m(t)$\\
D：
\begin{align*}
s_{D}(t) &= s(t) \sin(91 \times 10^{4} \pi t) \\
&= m(t) \sin(91 \times 10^{4} \pi t) \cos(91 \times 10^{4} \pi t) - \hat{m}(t) \sin^{2}(91 \times 10^{4} \pi t) \\
&= \frac{m(t)}{2} \sin(182 \times 10^{4} \pi t) - \frac{\hat{m}(t)}{2} \left[ 1 - \cos(182 \times 10^{4} \pi t) \right]
\end{align*}

$\mathrm{E}\colon\;s_{E}(t)=-\frac{1}{2}\hat{m}(t)$\\
$\mathrm{F}\colon\;\;s_{F}(t)=-\frac12\,\hat{m}(t)=\frac{m\left(t\right)}2$\\
${\mathrm{G}}\colon\;\;s_{G}(t)=m(t)$\\
(2)当 A 点输入是下单边带信号时，各点信号如下：\\
A：
\begin{align*}
s(t) &= m(t) \cos(2 \pi f_{c} t) + \hat{m}(t) \sin(2 \pi f_{c} t) \\
     &= m(t) \cos(91 \times 10^{4} \pi t) + \hat{m}(t) \sin(91 \times 10^{4} \pi t)
\end{align*}

B：
\begin{align*}
s_{B}(t) &= s(t) \cos(91 \times 10^{4} \pi t) \\
&= m(t) \cos^{2}(91 \times 10^{4} \pi t) + \hat{m}(t) \sin(91 \times 10^{4} \pi t) \cos(91 \times 10^{4} \pi t) \\
&= \frac{m(t)}{2} \left[ 1 + \cos(182 \times 10^{4} \pi t) \right] + \frac{\hat{m}(t)}{2} \sin(182 \times 10^{4} \pi t)
\end{align*}

${\mathsf{C}}\!:\ s_{C}(t)\!=\!{\frac{1}{2}}m(t)$\\
D：
\begin{align*}
s_{D}(t) &= s(t) \sin(91 \times 10^{4} \pi t) \\
&= m(t) \sin(91 \times 10^{4} \pi t) \cos(91 \times 10^{4} \pi t) + \hat{m}(t) \sin^2(91 \times 10^{4} \pi t) \\
&= \frac{m(t)}{2} \sin(182 \times 10^{4} \pi t) + \frac{\hat{m}(t)}{2} \left[1 - \cos(182 \times 10^{4} \pi t)\right]
\end{align*}

${\mathrm{E}}\colon\;S_{E}(t)={\frac{1}{2}}{\hat{m}}(t)$\\
$\mathrm{F}\colon\;\;s_{F}(t)=\frac{1}{2}\hat{m}(t)=-\frac{m\left(t\right)}{2}$\\
${\mathrm{G}}\!:\quad S_{G}(t)=0$\\
如欲 G 点输出 m(t) ，需将最末端的相加改为相减即可，如下图示：
\begin{figure}[H]
    \centering
    \includegraphics{zwz_images/36.png}
\end{figure}


%4.7
---------------------------------------\\
4.7 一 VSB 调幅信号的产生框图如图 7(a)所示，其中 BPF 的传递函数H(f)如图 7(b)所示。请画出相应的相干解调框图，并说明解调输出不会失真的理由。 
\begin{figure}[H]
    \centering
    \includegraphics{zwz_images/37.png}
    \caption*{题7图(a)}
\end{figure}
\begin{figure}[H]
    \centering
    \includegraphics{zwz_images/38.png}
    \caption*{题7图(b)}
\end{figure}
+++++++++++++++++++++++++++++++++++++\\
解：解调框图如下，其中理想低通的截止频率是 W。
\begin{figure}[H]
    \centering
    \includegraphics{zwz_images/39.png}
\end{figure}
接收信号${\boldsymbol{r}}(t)$为带通信号，可写成$r\left(t\right)=\mathrm{Re}\left\{r_{L}\left(t\right)e^{j2\pi f_{c}t}\right\}$，其中$r_{L}\left(t\right)$是${\boldsymbol{r}}(t)$的复包络。用载波$2\cos2\pi f_{c}t$对$r(t)$进行相干解调后，得到的输出是$y{\bigl(}t{\bigr)}=\operatorname{Re}{\left\{r_{L}\left(t\right)\right\}}$。\\
不考虑噪声，那么$r(t)$是 DSB 信号$s\left(t\right)=A m\bigl(t\bigr)\mathrm{cos}\,2\pi\ f_{c}t$通过带通系统$H(f)$后的输出。DSB 信号$s(t)$的复包络是${s}_{L}\left(t\right)={A}m\left(t\right)$，其频谱为$S_{L}\left(J\right)=A M\left(J\right)$。$H(f)$的等效低通是
$$
H_{L}(f) = 
\begin{cases}
H(f+f_{c}) & -\frac{W}{2} < f < W \\
0 & \text{otherwise}
\end{cases}
$$
因此${\boldsymbol{r}}(t)$的复包络$r_{L}\left(t\right)$的频谱为
$$R_{L}\left(f\right)=H_{L}\left(f\right)S_{L}\left(f\right)$$
由于
$$y{\big(}t{\big)}=\operatorname{Re}\left\{r_{L}{\big(}t{\big)}\right\}={\frac{r_{L}{\big(}t{\big)}+r_{L}^{*}{\big(}t{\big)}}{2}}$$
\begin{align*}
Y(f) &= \frac{1}{2} \left[ R_{L}(f) + R_{L}^{*}(-f) \right] \\
&= \frac{1}{2} \left[ A M(f) H_{L}(f) + A M(f) H_{L}^{*}(-f) \right] \\
&= \frac{A M(f)}{2} \left[ H_{L}(f) + H_{L}^{*}(-f) \right]
\end{align*}

本题条件中$H(f)$是实函数，再由下图可知，$H_{_{L}}\left(f\right)+H_{_{L}}^{\ast}\left(-f_{c}\right)=1$
\begin{figure}[H]
    \centering
    \includegraphics{zwz_images/40.png}
\end{figure}
因此，$Y(f)=\frac{A}{2}M\left(f\right)$，即$y{\bigl(}t{\bigr)}={\frac{A}{2}}m(t)$，说明解调输出没有失真。


%4.8
---------------------------------------\\
4.8 某调频信号$s(t)=10\cos\bigl(2\pi\times10^{6}t+4\cos200\pi t\bigr)$，求其平均功率、调制指数、最大频偏以及近似带宽。\\
\\
+++++++++++++++++++++++++++++++++++++\\
解：平均功率为
$$P_{s}={\frac{\left(10\right)^{2}}{2}}=50\,\mathrm{W}$$
最大频偏为$$\Delta\!f_{\mathrm{max}}=\mathrm{max}\left\{\frac{1}{2\pi}\frac{d}{d t}\big(4\mathrm{cos}\,200\pi t\big)\right\}=400\mathrm{Hz}$$
记$m(t)$为调制信号，$K_{F}$为调频灵敏度，则由于$4\cos200\pi t=2\pi K_{F}\int_{-\infty}^{t}m\bigl(\tau\bigr)d\tau$，故$m\bigl(t\bigr)=-\frac{400}{K_{F}}\mathrm{sin}\,200\pi t$，所以$m(t)$的频率是$f_{m}=100\mathrm{Hz}$，因此调频指数为
$$\beta_{f}={\frac{\Delta\beta_{\mathrm{max}}}{f_{m}}}=4$$
近似带宽为
$$B=2(1+\beta f)W=2\times5\times100\,\mathrm{HZ}=1000\,\mathrm{HZ}$$



%4.9
---------------------------------------\\
4.9 调角信号$s(t)=100\cos\bigl(2\pi f_{c}t+4\sin2\pi f_{m}t\bigr)$，其中载频$f_{\mathrm{c}}\!=\!10\mathrm{MHz}$，调制信号的频率是$f_{m}=1\,000\,\mathrm{HZ}$。 \\
(1)假设 $s(t)$ 是 FM 调制，求其调制指数及发送信号带宽； \\
(2)若调频器的调频灵敏度不变，调制信号的幅度不变，但频率$f_{m}$加倍，重复(1)题； \\
(3)假设$s(t)$ 是 PM 调制，求其调制指数及发送信号带宽； \\
(4)若调相器的调相灵敏度不变，调制信号的幅度不变，但频率$f_{m}$加倍，重复(3)题。\\
+++++++++++++++++++++++++++++++++++++\\
解：(1)$4\sin2\pi f_{m}t=2\pi K_{F}\int_{-\infty}^{t}m\big(\tau\big)d\tau\;,\;\;\;m\big(t\big)=A_{m}\cos2\pi f_{m}t=\frac{4000}{K_{F}}\cos2000\pi t$。最大频偏$\Delta f_{\mathrm{max}} = K_{F} \left| m(t) \right|_{\mathrm{max}} = K_{F} \times \frac{4000}{K_{F}} = 4000 \, \mathrm{Hz}$。由此可得调制指数为
$$\beta_{f}={\frac{\Delta f_{\mathrm{max}}}{f_{m}}}={\frac{4000}{1000}}=4$$
信号带宽近似为 
$$B=2\left(\Delta{\cal{f}}_{\mathrm{max}}+f_{m}\right)=2\left(4000+1000\right)=1\,000\,\mathrm{MZ}$$
(2)此时最大频偏不变，但因$\textstyle f_{m}$成为 2000Hz ，所以调制指数成为$\frac{400}{2000}=2$，近似带宽成为$2\left(4000+2000\right)=12000\mathrm{Hz}$\\
(3) $4\sin2\pi f_{m}t=K_{p}m(t)$，$m\bigl(t\bigr)=\frac{4}{K_{p}}\mathrm{sin}\,2000\pi l$。\\
调相指数为$\beta_{p}=\operatorname*{max}\left\{4\sin2\pi f_{m}t\right\}=4$\\
最大频偏为$\Delta{f_{\mathrm{max}}}=\mathrm{max}\left\{\frac{1}{2\pi}\frac{d}{d t}\big(4\sin2\pi f_{m}t\big)\right\}=4f_{m}=4000\mathrm{Hz}$\\
信号带宽近似为：$B=2\bigl(\Delta{\cal f}_{\mathrm{max}}+f_{m}\bigr)=1\,0000\,\mathrm{HZ}$\\
\\
(4) 此时$s(t)\!=\!100\mathrm{cos}\!\left(2\pi f_{c}t+K_{p}\!\times\!\frac{4}{K_{p}}\mathrm{sin}\,4\pi f_{m}\!t\right)\!=\!100\mathrm{cos}\!\left(2\pi f_{c}t+4\mathrm{sin}\,4\pi f_{m}t\right)$，调相指数仍然是$\beta_{f}=4$。最大频偏为
$$\Delta f_{\mathrm{max}}=\mathrm{max}\left\{\frac{1}{2\pi}\frac{d}{d t}\big(4\sin4\pi f_{m}t\big)\right\}=8f_{m}=8000\mathrm{Hz}$$
信号带宽近似为： 
$$B=2\left(\Delta{\cal f}_{\mathrm{max}}+2{\cal f}_{m}\right)=2\left(8000+2000\right)=20000\mathrm{HZ}$$


%4.10
---------------------------------------\\
4.10 下图中的模拟基带信号$m(t)$不含直流分量，且其频谱在零频率附近为 0，K 为充分大的常数，载频$f_{c}$远远大于$m(t)$的带宽。写出$s_{1}(t)$，$s_{2}(t)$，$s_{3}(t)$，$s_{4}(t)$，$s_{5}(t)$的表达式，并指出它们的调制类型。
\begin{figure}[H]
    \centering
    \includegraphics{zwz_images/41.png}
    \caption*{题4.10图(a)}
\end{figure}
+++++++++++++++++++++++++++++++++++++\\
解：
$$s_{1}\left(t\right)=m\bigl(t\bigr)\mathrm{cos}\,2\pi f_{c}t+\hat{m}\bigl(t\bigr)\mathrm{sin}\,2\pi f_{c}l$$
它是下边带的 SSB 调制。
$$s_{2}\left(t\right)=m\left(t\right)\cos2\pi f_{c}t+K\cos2\pi f_{c}t=K\left[1+{\frac{m\left(t\right)}{K}}\right]\cos2\pi f_{c}t$$
这是标准调幅（AM）。
\begin{align*}
s_{3}(t) &= K\cos(2\pi f_{c}t) - \left[\int_{-\infty}^{t} m(\tau) \, d\tau\right]\sin(2\pi f_{c}t) \\
&= K \left[\cos(2\pi f_{c}t) - \frac{1}{K} \int_{-\infty}^{t} m(\tau) \, d\tau \sin(2\pi f_{c}t)\right] \\
&\approx K \left[\cos(2\pi f_{c}t)\cos\left(\frac{1}{K} \int_{-\infty}^{t} m(\tau) \, d\tau\right) - \sin\left(\frac{1}{K} \int_{-\infty}^{t} m(\tau) \, d\tau\right)\sin(2\pi f_{c}t)\right] \\
&= K\cos\left[2\pi f_{c}t + \frac{1}{K} \int_{-\infty}^{t} m(\tau) \, d\tau\right]
\end{align*}

这是窄带调频信号。
$$S_{4}\left(t\right)=\cos\biggl[2\pi f_{c}t+K_{p}\int_{-\infty}^{t}m\left(\tau\right)d\tau\biggr]$$
这是频率调制(FM)。
$$s_{s}\left(t\right)=\cos\left[2\pi f_{c}t+2\pi K_{f}\int_{-\infty}^{t}\frac{d m\left(\tau\right)}{d\tau}d\tau\right]=\cos\Bigl[2\pi f_{c}t+2\pi K_{f}m\left(t\right)\Bigr]$$
这是相位调制(PM)。 


%4.11
---------------------------------------\\
4.11 双边带调幅信号$s\left(t\right)=m\left(t\right)\cos2\pi f_{c}t$的功率谱密度$P_{s}(f)$如题 4.11 图(a)示，$s(t)$在传输中受到功率谱密度为$\frac{N_{0}}{2}$的加性白高斯噪声$n_{_W}\left(t\right)$的干扰，接收端用题 4.11 图(b)所示的解调框图进行解调。求解调输出的信噪比。
\begin{figure}[H]
    \centering
    \includegraphics{zwz_images/42.png}
    \caption*{题4.11图(a)}
\end{figure}
\begin{figure}[H]
    \centering
    \includegraphics{zwz_images/43.png}
    \caption*{题4.11图(b)}
\end{figure}
+++++++++++++++++++++++++++++++++++++\\
解：由题 4.11 图(a)可知，$m(t)$的带宽是 W，$s(t)$的功率是 
$$P_{s}=\int_{-\infty}^{\infty}P_{s}\left(f\right)d f=2P_{0}W$$
记 $m(t)$的功率为$\ P_{M}$，由$s\left(t\right)=m\left(t\right)\cos2\pi f_{c}t$可知
$$P_{s}={\frac{1}{2}}P_{o}$$
因此
$$P_{M}=4N_{0}W$$
LPF 输出中的信号分量为$\frac{m(t)}{2}$，因此输出的信号功率为
$${\frac{P_{M}}{4}}=P_{0}W$$
BPF 输出端的噪声是窄带高斯噪声：
$$n\left(t\right)=n_{c}\left(t\right)\cos2\pi f_{c}t-n_{s}\left(t\right)\sin2\pi f_{c}t$$
其中 ${{n}}_{\mathrm{c}}(t)$,${{n}}_{\mathrm{s}}(t)$,${{n}}(t)$ 的功率都是
$$ P_{n} = N_{0} \times 2M = 2N_{0}M $$
LPF 输出的噪声分量是$\frac{n_{c}\left(t\right)}{2}$，所以输出的噪声功率是
$${\frac{P_{n}}{4}}={\frac{N_{0}W}{2}}$$
因此输出信噪比是
$$\left(\frac{S}{N}\right)_{o}=\frac{P_{0}W}{\frac{1}{2}N_{o}M}=\frac{2P_{0}}{N_{0}}$$


%4.12
---------------------------------------\\
4.12 某模拟广播系统中基带信号$m(t)$的带宽为$W=10\mathrm{KHz}$，峰均功率比（定义为$\frac{|m(t)|_{\operatorname*{max}}^{2}}{P_{M}}$，其中$\ P_{M}$是$m(t)$的平均功率）是 5。此广播系统的平均发射功率为 40KW，发射信号经过 80dB 信道衰减后到达接收端，并在接收端叠加了双边功率谱密度为$\frac{N_{0}}{2}=10^{-10}\mathrm{W/Hz}$的白高斯噪声。 \\
(1)若此系统采用 SSB 调制，求接收机可达到的输出信噪比； \\
(2)若此系统采用调幅系数为 0.85 的标准幅度调制(AM)，求接收机可达到的输出信噪比。 \\
\\
+++++++++++++++++++++++++++++++++++++\\
解：记到达接收机的已调信号为$s(t)$，记其功率为$\ P_{R}$，则
$$P_{R}=40\times10^{3}\times10^{-8}=4\times10^{-4}\mathrm{{W}}$$
(1)采用 SSB 时，不妨以下单边带为例，接收机输入端的有用信号为
$$s\left(t\right)=A m\left(t\right)\cos2\pi f_{c}t+A\hat{m}\left(t\right)\sin2\pi f_{c}t$$
其中$f_{c}$是载波频率，${\hat{m}}(t)$是$m(t)$的 Hilbert 变换。按照 Hilbert 变换的性质，$m(t)$、${\hat{m}}(t)$的功率都是$\ P_{M}$。故而
$$P_{R}=\frac{A^{2}P_{M}}{2}+\frac{A^{2}P_{M}}{2}=A^{2}P_{M}$$
接收机前端可采用频带为$\left[f_{c}-W,f_{c}\right]$的理想 BPF 限制噪声，于是输入到解调器输入端的噪声为
$$n\left(t\right)=n_{c}\left(t\right)\cos2\pi f_{c}t-n_{s}\left(t\right)\sin2\pi f_{c}t$$
其中 ${n_{c}(t)}$，${n_{s}(t)}$，${n(t)}$的功率都是
$$P_{n}=N_{0}W=2\times10^{-10}\times10^{4}=2\times10^{-6}\mathrm{W}$$
采用理想相干解调时的解调输出为$A m\left(t\right)+n_{c}\left(t\right)$，因此输出信噪比是
$$\left({\frac{S}{N}}\right)_{S S B}={\frac{A^{2}P_{M}}{P_{n}}}={\frac{4\times10^{-4}}{2\times10^{-6}}}=200\text{，即23dB}$$
(2) 采用 AM 调制时，解调器输入的有用信号为： 
$$s\left(t\right)=A\left[1+a m_{n}\left(t\right)\right]\cos2\pi f_{c}t$$
其中调制指数$a=0.85$，$m_{n}\left(t\right)=\frac{m\left(t\right)}{\left|m\left(t\right)\right|_{\mathrm{max}}}$，$m_{n}\left(t\right)$的功率是
$${ P}_{M_{n}}=\frac{{P}_{M}}{\left|m\left(t\right)\right|_{\mathrm{max}}^{2}}=\frac{1}{5}$$
接收机前端可采用频带为$\left[f_{c}-W,f_{c}+W\right]$的理想 BPF 限制噪声，于是输入到解调器输入端的噪声为 
$$n\left(t\right)=n_{c}\left(t\right)\cos2\pi f_{c}t-n_{s}\left(t\right)\sin2\pi f_{c}t$$
其中${{n}}_{\mathrm{c}}(t)$，${{n}}_{\mathrm{s}}(t)$，${{n}}(t)$的功率都是
$$P_{n}=2N_{0}W=2\times2\times10^{-10}\times10^{4}=4\times10^{-6}$$
采用理想的包络检波器得到的输出为${\mathbf{}}A{\big[}1+a m_{n}\left(t\right){\big]}+n_{c}\left(t\right)$，因此输出信噪比是
$$\left(\frac{S}{N}\right)_{d M}=\frac{A^{2}a^{2}P_{M_{n}}}{P_{n}}$$
由于
$$P_{R}=\frac{A^{2}}{2}\Big[1+a^{2}P_{M_{n}}\Big]=4\times10^{-4}$$
由此得
$$\left(\frac{S}{N}\right)_{A M}=\frac{2P_{_R}}{1+a^{2}P_{M_{s}}}{\times}\frac{a^{2}P_{M_{s}}}{P_{n}}=\frac{2{\times}4{\times}10^{-4}}{\left(1+\frac{1}{0.85^{2}{\times}0.2}\right){\times}4{\times}10^{-6}}{\approx}25.25\text{，即14.02dB}$$


%4.13
---------------------------------------\\
4.13 有 12 路话音信号$m_{1}\left(t\right),m_{2}\left(t\right),\cdots,m_{12}\left(t\right)$，它们的带宽都限制在(0,4000) Hz 范围内。将这 12 路信号以SSB/FDM方式复用为$m(t)$，再将$m(t)$通过FM方式传输，如图(a)所示。其中SSB/FDM的频谱安排如图(b)所示。已知调频器的载频为$f_{\mathrm{c}},$，最大频偏为 480KHz。
\begin{figure}[H]
    \centering
    \includegraphics{zwz_images/44.png}
    \caption*{题4.13图(a)}
\end{figure}
+++++++++++++++++++++++++++++++++++++\\





% 65-94页
\begin{figure}[h]
\centering
\includegraphics[width=\linewidth]{output/image_1_1.png}
\end{figure}
(1)求FM信号的带宽；\\
(2)画出解调框图；\\
(3)假设 FM 信号在信道传输中受到加性白高斯噪声干扰，求鉴频器输出的第 1 路噪声
平均功率与第 12 路噪声平均功率之比。\\
++++++++++++++++++++++++++++++++++++++++++++\\
解：(1)$m(t)$的带宽为 48KHz,所以 FM 信号的带宽近似为 2$\left(480+48\right)=1056$KHz
(2)解调框图如下。其中$f_i= 4( i- 1)$KHz ,$i=1,2,\cdots,12$。第$i$个 BPF 的通带为
$\begin{pmatrix}f_i,f_i+4\text{KHz}\end{pmatrix}$。LPF 的截止频率是 4KHz。
\begin{figure}[h]
\centering
\includegraphics[width=\linewidth]{output/image_1_2.png}
\end{figure}
(3)鉴频器输出的噪声功率谱密度与 $f^2$成正比，即
$P_{n_o}\left(f\right)=Kf^2$

其中$K$是常系数。落在第 1 路频带范围内的输出噪声功率为
$$P_1=\int_0^{4000}Kf^2df=\frac{\left(4000\right)^3K}{3}=\frac{4^3\times10^9}{3}K$$
落在第 12 路频带范围内的输出噪声功率为
$$P_{12}=\int_{44000}^{48000}Kf^{2}df=\frac{\left(48000^{3}-44000^{3}\right)K}{3}=\frac{\left(48^{3}-44^{3}\right)\times10^{9}}{3}K$$
所以

$\frac {P_{12}}P= \frac {48^{3}- 44^{3}}{4^{3}}= 12^{3}- 11^{3}= 397$ ,相当于 26dE
---------------------------------------\\
4.14 一射频方向性天线的等效噪声温度$T_{\alpha}$为 55K ,天线输出经过一个损耗为 2dB的馈线连接至接收机前置射频放大器输入端 ,该射频放大器的功率增益$K_p$为 20dB ,等效带宽为 10MHz , 噪声系数为 2dB,设室温为$T=300$K,求： (1)前置射频放大器的等效噪声温度； (2)馈线与前置射频放大器级联的噪声系数及等效噪声温度；
(3)求馈线输入点的单边噪声功率谱密度$N_0$值。
++++++++++++++++++++++++++++++++++++++++++++\\
13/25

13

---------------------------------------

《通信原理习题解答》

解：(1)记$T_{_A}$为放大器的等效噪声温度。放大器自身的噪声折合到它的输入端的功率为
$$KT\begin{pmatrix}F_A-1\end{pmatrix}B=KT_AB$$
因此
$$T_A=T\begin{pmatrix}F_A-1\end{pmatrix}=300\times\begin{pmatrix}10^{0.2}-1\end{pmatrix}=175.5\text{K}$$
(2) 设总噪声系数是$F_{_{AF}}$,可用课本上的公式(4.8.28)求出$F_{_{AF}}:$
$$F_{AF}=F_{1}+\frac{F_{2}-1}{K_{pa_{1}}}+\frac{F_{3}-1}{K_{pa_{1}}K_{pa_{2}}}+\cdots $$
本题只有两级，第一级馈线的衰减量是$L=10^{0.2}$(2dB),嗓声系数是$F_1=L$,增益是
$K_{p\alpha_{1}}=\frac{1}{L}$;第二级前置射频放大器的噪声系数是$F_2=F_1=L$。因此
$$F_{AF}=F_{1}+\frac{F_{2}-1}{K_{pa_{1}}}=L+\frac{L-1}{1/L}=L^{2}=10^{0.4}\:,\:\text{即4dE}$$
设馈线和放大器级联的总等效噪声温度是$T_{_{AF}}$ ,则由式(4.8.33)可得
$$T_{_{AF}}=\begin{pmatrix}F_{_{AF}}-1\end{pmatrix}T=\begin{pmatrix}10^{0.4}-1\end{pmatrix}\times300=453.6\:\mathrm{K}$$
本题中总噪声系数的求解也可以不直接套用式(4.8.28)。馈线输入的可获噪声功率是$KTB$ ,输出的噪声功率也是$KTB$ ,而馈线的损耗是$L$ ,因此我们可以理解成：馈线输出的操声中$\frac{KTB}L$是输入噪声所致，其余的$P_L=KTB-\frac{KTB}L=KTB\left(1-\frac1L\right)$是馈线自身的時声。馈线自身的噪声$P_L$折合到它的输入端是$LP_L=KTB(L-1)$。而放大器自身的噪声折合到放大器输入端是$P_{_A}=KTB\left(F_{2}-1\right)=KTB\left(L-1\right)$(本题条件中，$F_{_2}=L$ ),$P_{_A}$再折合到馈线输入端是$LP_{_A}=KTB\left(L^2-L\right)$。所有噪声(馈线和放大器产生的噪声、馈线输入噪声 ) 折合到馈线输入端是$LP_L+LP_A+KTB=KTBL^2$。这说明总噪声系数是$F_{AF}=L^2$。
(3)如果馈线输入是常温电阻，则前面的推导结果表明馈线输入端的总噪声功率是$P_0=KTBL^2=K\left(TL^2\right)B$ ,即总噪声温度是$TL^2$。扣除常温电阻，那么馈线+放大器的等效噪声温度是$T_{AF}= \left ( L^2- 1\right ) T$。今 馈 线 输 入 的 噪 声 源 的 温 度 是 $T_\alpha = 55$ K,于是馈线输入点总的噪声功率是$K[T_\alpha+T_{AF}]B$ ,因而在 B=1Hz 带宽上的测量到的功率( 功率谱密度，测量时的 1Hz 包含正负频率两部分)为
$$N_{0}=K\left(T_{a}+T_{AF}\right)=1.38\times10^{-23}\times\left[55+300\left(10^{0.4}-1\right)\right]=7.02\times10^{-21}\:\mathrm{W/Hz}$$
---------------------------------------\\
4.15 已知某模拟基带系统中调制信号$m(t)$的带宽是 W=5kHz。发送端发送的已调信号功率
是$P_i$,接收功率比发送功率低 60dB。信道中加性白高斯噪声的单边功率谱密度为
$N_0=10^{-13}$W/Hz 。

(1)如果采用 DSB-SC,请
a)推导出输出信噪比$\left(\frac SN\right)_o$和输入信嗓比$\left(\frac SN\right)_i$的关系

14/25

14

---------------------------------------

《通信原理习题解答》

(b)若要求输出信噪比不低于 30dB,发送功率至少应该是多少？
(2)如果采用 SSB,重做第(1)问。
++++++++++++++++++++++++++++++++++++++++++++\\
解：(1)(a)解调输入信号可写为
$$\begin{aligned}r\left(t\right)&=Am\left(t\right)\cos2\pi f_{c}t+n\left(t\right)\\&=Am\left(t\right)\cos2\pi f_{c}t+n_{c}\left(t\right)\cos2\pi f_{c}t-n_{s}\left(t\right)\sin2\pi f_{c}t\end{aligned}$$
输入信噪比为
$$\left(\frac{S}{N}\right)_i=\frac{\frac{A^2P_m}{2}}{2N_0W}=\frac{A^2P_m}{4N_0W}$$
解调输出为$Am(t)+n_c(t)$,输出信噪比为
$$\left(\frac{S}{N}\right)_{o}=\frac{A^{2}\overline{m^{2}\left(t\right)}}{\overline{n_{c}^{2}\left(t\right)}}=\frac{A^{2}P_{m}}{2N_{0}W}$$
$$\frac{\left(S/N\right)_o}{\left(S/N\right)_i}=2$$
因此


(b)输入信噪比

输出信噪比
$$\left(\dfrac{S}{N}\right)_i=\dfrac{P_r}{2N_0W}=\dfrac{P_t\times10^{-6}}{2\times10^{-13}\times5\times10^3}=10^3P_t$$
$$\left(\dfrac{S}{N}\right)_o=2\Bigg(\dfrac{S_i}{N_i}\Bigg)=2000P_i=10^3$$
故$P_{\iota}=0.5$瓦。

(2)(a)解调输入信号可写为
$$\begin{aligned}&\text{解调输入信号可与为}\\&r\left(t\right)=Am\left(t\right)\cos2\pi f_{c}t+A\hat{m}\left(t\right)\sin2\pi f_{c}t+n\left(t\right)\\&=Am\left(t\right)\cos2\pi f_{c}t+A\hat{m}\left(t\right)\sin2\pi f_{c}t+n_{c}\left(t\right)\cos2\pi f_{c}t-n_{s}\left(t\right)\sin2\pi f_{c}t\end{aligned}$$
输入信噪比为
$$\left(\frac{S}{N}\right)_{i}=\frac{A^{2}\overline{m^{2}\left(t\right)}}{N_{0}W}$$
解调输出为$Am(t)+n_c(t)$,输出信噪比为
$$\left(\frac{S}{N}\right)_{o}=\frac{A^{2}\overline{m^{2}\left(t\right)}}{\overline{n_{c}^{2}\left(t\right)}}=\frac{A^{2}P_{m}}{N_{0}W}$$
因此$\frac{\left(S/N\right)_o}{\left(S/N\right)_i}=1$

(b)输入信噪比

15/25

---------------------------------------

《通信原理习题解答》

输出信噪比
$$\left(\dfrac{S}{N}\right)_i=\dfrac{P_r}{N_0W}=\dfrac{P_t\times10^{-6}}{10^{-13}\times5\times10^3}=2000P_t$$
$$\left(\dfrac{S}{N}\right)_o=\left(\dfrac{S_i}{N_i}\right)=2000P_i=10^3$$
故$P_{\iota}=0.5$瓦。
---------------------------------------\\
4.16 已知某调频系统中，调频指数是$\beta_f$,到达接收端的 FM 信号功率是$P_R$,信道噪声的单
边功率谱密度是$N_0$,基带调制信号$m(t)$的带宽是 $W$,解调输出的信噪比和输入信噪
比之比为3$\beta_f^{2}\left(\beta_f+1\right)$。(1)求解调输出信噪比； (2)如果发送端将基带调制信号$m(t)$变成 2$m(t)$,接收端按照此条件设计解调
器，请问输出信噪比将大约增大多少分贝？
++++++++++++++++++++++++++++++++++++++++++++\\
解：(1)输入信噪比是

输出信噪比是
$$\left(\dfrac{S}{N}\right)_i=\dfrac{P_R}{2N_0W\left(\beta_f+1\right)}$$
$$\left(\frac{S}{N}\right)_{o}=3\beta_{f}^{2}\Big(\beta_{f}+1\Big)\Bigg(\frac{S}{N}\Bigg)_{i}=\frac{3\beta_{f}^{2}P_{R}}{2N_{0}W}$$
$(2)此时\beta_{j}$变成了2$\beta_f$,输出信噪比是原来的4倍，即增加了6dB。
---------------------------------------\\
4.17 两个不包含直流分量的模拟基带信号$m_1(t)$、$m_2(t)$被同一射频信号同时发送，发送信号
为 $s\left(t\right)=m_{1}\left(t\right)\cos\omega_{c}t+m_{2}\left(t\right)\sin\omega_{c}t+K\cos\omega_{c}t$,其中载频$f_{\mathrm{c}}=10$MHz,$K$是常数。已知$m_1(t)$与$m_2(t)$的傅氏频谱分别为$M_1(f)$及$M_{_2}(f)$,它们的带宽分别为 5KHz与10KHz。
(1)请计算 $s(t)$的带宽；
(2)请写出 $s(t)$的傅氏频谱表示式；
(3)画出从s(t)得到$m_1(t)$及$m_2(t)$的解调框图。
++++++++++++++++++++++++++++++++++++++++++++\\
解：(1)s$\left(t\right)$由两个 DSB 及一个单频组成 ,这两个 DSB 的中心频率相同 ,带宽分别是 10KHz
和 20KHz,因此总带宽是 20KHz ;

(2)
$$\begin{aligned}S\left(f\right)&=\frac{1}{2}\Big[M_{1}\Big(f+f_{c}\Big)+M_{1}\Big(f-f_{c}\Big)\Big]+\frac{j}{2}\Big[M_{2}\Big(f+f_{c}\Big)-M_{2}\Big(f-f_{c}\Big)\Big]\\&+\frac{K}{2}\Big[\delta\Big(f+f_{c}\Big)+\delta\Big(f-f_{c}\Big)\Big]\end{aligned}$$
(3)

16/25

16

---------------------------------------

《通信原理习题解答》

\begin{figure}[h]
\centering
\includegraphics[width=\linewidth]{output/image_5_1.png}
\end{figure}
---------------------------------------\\
4.18 $设 m\left ( t\right ) = \frac 1{\sqrt {N}}\sum _{i= 1}^{N}\cos 2\pi it$
(1)求$m(t)$的平均功率$P_M:$

(2)若用$m(t)$进行 AM 幅度调制，求调制效率$\eta$作为调制指数$\alpha$的函数$\eta(a)$,并就
$N=3$ 的情形画出$\eta(a)$。
++++++++++++++++++++++++++++++++++++++++++++\\
解$:(1)\cos2\pi it,i=1,2,\cdots,N$是频率分别为 1、2、...、NHz 的$N$个余弦信号，其功率均为
$\frac12$。因此$\sum_i=1^N\cos2\pi it$ 的功率是$\frac N2$,$m(t)=\frac1{\sqrt{N}}\sum_{i=1}^N\cos2\pi it$ 的平均功率是$P_M=\frac12$

$(2)\left|m\left(t\right)\right|_{\mathrm{max}}=\frac{1}{\sqrt{N}}\sum_{i=1}^{N}1=\sqrt{N}$ ,将$m(t)$的幅度归一化，得

$m_{n}\left(t\right)=\frac{m\left(t\right)}{\left|m\left(t\right)\right|}=\frac{1}{N}\sum_{i=1}^{N}\cos2\pi it$

其功率为

$P_{M_n}=\frac1{2N}$

调制指数是$a$,因此 AM 信号是

$s\left(t\right)=A_{c}\left[1+am_{n}(t)\right]\cos2\pi f_{c}t$

其效率为
$$\eta(a)=\frac{a^2P_{M_n}}{1+a^2P_{M_n}}=\frac{\frac{a^2}{2N}}{1+\frac{a^2}{2N}}=\frac{a^2}{2N+a^2}$$
$$\eta(a)=\frac{a^2}{6+a^2}$$
$a$ 的取值范围是$[0,1],\eta(a)$的图形如下

当$N=3$时


17/25

---------------------------------------

《通信原理习题解答》

\begin{figure}[h]
	\centering
	\includegraphics[width=\linewidth]{output/image_6_1.png}
\end{figure}
---------------------------------------\\
4.19 若 SSB 信号发送端的载波初始相位是 0,解调器的本地载波的初始相位是$\varphi$。求解调
输出中有用信号的功率占总输出信号功率的比率。
++++++++++++++++++++++++++++++++++++++++++++\\
解：问题也相当于发送的载波相位是$-\varphi$,本地载波相位是 0。于是发送信号是
$$\begin{aligned}s\left(t\right)&=A_{c}m\left(t\right)\cos\left(2\pi f_{c}t-\varphi\right)\mp A_{c}\hat{m}\left(t\right)\sin\left(2\pi\right)\\&=\mathrm{Re}\left\{A_{c}\left[m\left(t\right)\pm j\hat{m}\left(t\right)\right]e^{j\left(2\pi f_{c}t-\varphi\right)}\right\}\\&=\mathrm{Re}\left\{A_{c}\left[m\left(t\right)\pm j\hat{m}\left(t\right)\right]e^{-j\varphi}e^{j2\pi f_{c}t}\right\}\end{aligned}$$
$$s_L\left(t\right)=A_c\left[m\left(t\right)\pm j\hat{m}\left(t\right)\right]e^{-j\varphi}$$
其复包络是


相干解调器的参考载波是$\cos2\pi f_ct$,所以解调输出是复包络的实部(忽略系数)
$$m_o\left(t\right)=\mathrm{Re}\left\{s_L\left(t\right)\right\}=A_c\left[m\left(t\right)\cos\varphi\pm\hat{m}\left(t\right)\sin\varphi\right]$$
其中有用信号是$A_cm(t)\cos\varphi$,其功率是
$$P_1=A_c^2P_M\cos^2\varphi $$
无用信号是$\pm A_c\hat{m}(t)\sin\varphi$,其功率是
$$P_2=A_c^2P_M\sin^2\varphi $$
因此有用信号所占的功率比例是
$$\frac{P_{1}}{P_{1}+P_{2}}=\frac{\cos^{2}\varphi}{\cos^{2}\varphi+\sin^{2}\varphi}=\cos^{2}\varphi $$
---------------------------------------\\
4.20 设$m_1(t)、m_2(t)$是频谱分布在 300~3400Hz 内的两个话音信号，功率都是 1。今发送
$s(t)=m_1(t)\cos2\pi f_ct-m_2(t)\sin2\pi f_ct$,问
(1)接收端如何解出$m_1(t),m_2(t)?$
(2)若$s(t)$叠加了一个双边功率谱密度为$N_0/2$的白高斯噪声，求解调输出为
$m_1(t)$这一支路的输出信噪比。

18/25

18

---------------------------------------

《通信原理习题解答》
++++++++++++++++++++++++++++++++++++++++++++\\
解：(1)
\begin{figure}[h]
	\centering
	\includegraphics[width=\linewidth]{output/image_7_1.png}
\end{figure}
其中 BPF 的中心频率是$f_c$ ,带宽是 2$\times3400$Hz 。

(2) BPF 输出的噪声是
$$\begin{aligned}n\left(t\right)&=\mathrm{Re}\left\{\left[n_{c}\left(t\right)+jn_{s}\left(t\right)\right]e^{j2\pi f_{c}t}\right\}\\&=n_{c}\left(t\right)\cos2\pi f_{c}t-n_{s}\left(t\right)\sin2\pi f_{c}t\end{aligned}$$
其平均功率是
$$P_n=\overline{n^2\left(t\right)}=\overline{n_c^2\left(t\right)}=\overline{n_s^2\left(t\right)}=6800N_0$$
$m_1(t)$这一路解调输出的有用信号是$m_1(t)$,其功率为 1;输出噪声是$n_c(t)$,故输出信噪
比是$\frac1{6800N_{_0}}$。
---------------------------------------\\

4.21 题 4.21 图所示系统中调制信号 $m(t)$的均值为 0,功率谱密度为
$$P_{M}\left(f\right)=\begin{cases}N_{m}\left(1-\frac{\left|f\right|}{f_{m}}\right)&\left|f\right|\leq f_{m}\\0&\left|f\right|>f_{m}\end{cases}$$
白高斯噪声$n_w(t)$的双边功率谱密度为$N_0/2$。
\begin{figure}[h]
	\centering
	\includegraphics[width=\linewidth]{output/image_7_2.png}
\end{figure}
(1)写出$s(t)$的复包络$s_L(t)$,画出其功率谱密度；
(2)图中的 BPF 应该如何设计 ? (3)求相干解调器的输出信噪比。
++++++++++++++++++++++++++++++++++++++++++++\\
解：(1)由图可见，$s(t)$的表达式为
$$s\begin{pmatrix}t\end{pmatrix}=m\begin{pmatrix}t\end{pmatrix}A_c\cos2\pi f_ct+\hat{m}\begin{pmatrix}t\end{pmatrix}A_c\sin2\pi f_ct$$
因此 $s(t)$的复包络为
$$s_L\begin{pmatrix}t\end{pmatrix}=A_c\begin{bmatrix}m\begin{pmatrix}t\end{pmatrix}-j\hat{m}\begin{pmatrix}t\end{pmatrix}\end{bmatrix}$$
$m(t)+j\hat{m}(t)$是解析信号，因此 $z(t)=m\left(t\right)-j\hat{m}\left(t\right)=\left[m\left(t\right)+j\hat{m}\left(t\right)\right]^{*}$的傅氏谱是

19/25

19

---------------------------------------

《通信原理习题解答》
$$Z\left(f\right)=\begin{cases}2M\left(f\right)&f<0\\0&f>0\end{cases}$$
因此 $s_L\left(t\right)$的功率谱密度是
$$\begin{aligned}P_{s_{L}}\left(f\right)&=\begin{cases}4A_{c}^{2}P_{M}\left(f\right)&f<0\\0&f>0\end{cases}\\&=\begin{cases}4N_mA_c^2\left(1+\frac{f}{f_m}\right)&f<0\\\\0&f>0\end{cases}\end{aligned}$$
$P_{u_1}(f)$4$N_mA_c^2$
\begin{figure}[h]
	\centering
	\includegraphics[width=\linewidth]{output/image_8_1.png}
\end{figure}
(2)由于$s\left(t\right)$是下单边带信号，所以 BPF 的通频带应该是$\left[f_{c}-f_{m},f_{c}\right]$
(3)为方便计算，假设 LPF 的幅度增益是$\frac2{A_c}$,这个假设对输出信噪比没有影响。此时，相干调调器的输出是
$$y\begin{pmatrix}t\end{pmatrix}=m\begin{pmatrix}t\end{pmatrix}+\frac{1}{A_c}n_c\begin{pmatrix}t\end{pmatrix}$$
其中$n_c\left(t\right)$是 BPF 输出的窄带噪声的同相分量，其功率为 $f_mN_0$。由于 $m(t)$ 的功率是
$$P_{M}=\int_{-\infty}^{\infty}P_{M}\left(f\right)df=2\int_{0}^{f_{m}}N_{m}\left(1-\frac{f}{f_{m}}\right)df=f_{m}N_{m}$$
所以输出信噪比是
$$\left(\frac{S}{N}\right)_{o}=\frac{N_{m}f_{m}}{\frac{1}{A_{c}^{2}}N_{0}f_{m}}=\frac{A_{c}^{2}N_{m}}{N_{0}}$$
---------------------------------------\\
4.22 下图是对 DSB 信号进行相干解调的框图。图中$n_w(t)$是均值为 0、双边功率谱密度为$\frac{N_0}2$的加性白高斯噪声，本地恢复的载波和发送载波有固定的相位差$\theta$。请导出该系统的一个人口密度的统的输出信噪比表达式。
++++++++++++++++++++++++++++++++++++++++++++\\
\begin{figure}[h]
	\centering
	\includegraphics[width=\linewidth]{output/image_8_2.png}
\end{figure}
解一：LPF 输出的信号分量是
$$m\begin{pmatrix}t\end{pmatrix}\cos2\pi f_ct\times2\cos\begin{pmatrix}2\pi f_ct+\theta\end{pmatrix}\overset{\mathrm{LPF}}{=}m\begin{pmatrix}t\end{pmatrix}\cos\theta $$
其功率为$\overline{m^2(t)}\cos^2\theta.$

BPF 输出的窄带噪声是
$$n\begin{pmatrix}t\end{pmatrix}=n_c\begin{pmatrix}t\end{pmatrix}\cos2\pi f_ct-n_s\begin{pmatrix}t\end{pmatrix}\sin2\pi f_ct$$
其中$n_c\left(t\right),n_s\left(t\right)$的功率都是2$f_mN_0$。LPF 输出的噪声分量是
$$n_o\begin{pmatrix}t\end{pmatrix}=n_c\begin{pmatrix}t\end{pmatrix}\cos\theta+n_s\begin{pmatrix}t\end{pmatrix}\sin\theta $$
由于$n_c\left(t\right),n_s\left(t\right)$均值为 0,且在同一时刻不相关，所以输出噪声功率为
$$\begin{aligned}E\bigg[n_{o}^{2}\left(t\right)\bigg]&=E\bigg[n_{c}^{2}\left(t\right)\cos^{2}\theta\bigg]+E\bigg[n_{s}^{2}\left(t\right)\sin^{2}\theta\bigg]\\&=2N_{0}f_{m}\bigg[\cos^{2}\theta+\sin^{2}\theta\bigg]=2N_{0}f_{m}\end{aligned}$$
从而输出信噪比为
$$\left(\dfrac{S}{N}\right)_o=\dfrac{\overline{m^2\left(t\right)}\cos^2\theta}{2N_0f_m}$$
解二：令$t=t-\tau$,其中$\tau$满足$2\pi f_c\tau=\theta$,再考虑到对$m(t)$和$n_w\left(t\right)$做任何时延都不影响
问题，于是原题演变为下图
\begin{figure}[h]
	\centering
	\includegraphics[width=\linewidth]{output/image_9_1.png}
\end{figure}
图中 $s(t)$的复包络是
$$s_L\begin{pmatrix}t\end{pmatrix}=m\begin{pmatrix}t\end{pmatrix}e^{-j\theta}$$
BPF 输出端的窄带噪声的复包络为
$$n_L\begin{pmatrix}t\end{pmatrix}=n_c\begin{pmatrix}t\end{pmatrix}+jn_s\begin{pmatrix}t\end{pmatrix}$$
其中$n_c\left(t\right),n_s\left(t\right)$的功率都是2$f_mN_0$。BPF 输出的总的信号的复包络为
$$\begin{aligned}r_{L}\left(t\right)&=s_{L}\left(t\right)+n_{L}\left(t\right)\\&=m\left(t\right)e^{-j\theta}+n_{c}\left(t\right)+jn_{s}\left(t\right)\end{aligned}$$
解调输出是
$$y\begin{pmatrix}t\end{pmatrix}=\mathrm{Re}\begin{pmatrix}r_L\begin{pmatrix}t\end{pmatrix}\end{pmatrix}=m\begin{pmatrix}t\end{pmatrix}\cos\theta+n_c\begin{pmatrix}t\end{pmatrix}$$
因此输出信噪比是
$$\left(\dfrac{S}{N}\right)_o=\dfrac{\overline{m^2\left(t\right)}\cos^2\theta}{2N_0f_m}$$
---------------------------------------\\
4.23 立体声调频发送端方框图如下所示，其中$m_L(t)$和$m_R(t)$分别表示来自左右声道的电
信号，它们的最高频率都不超过 15kHz。
21/25


21

---------------------------------------

《通信原理习题解答》

\begin{figure}[h]
	\centering
	\includegraphics[width=\linewidth]{output/image_10_1.png}
\end{figure}
(1)请画出图中$m(t)$的频谱示意图；
(2)请画出立体声调频信号的接收框图。
++++++++++++++++++++++++++++++++++++++++++++\\
\begin{figure}[h]
	\centering
	\includegraphics[width=\linewidth]{output/image_10_2.png}
\end{figure}
\begin{figure}[h]
	\centering
	\includegraphics[width=\linewidth]{output/image_10_3.png}
\end{figure}
---------------------------------------\\
4.24 12 路话音信号，每路的带宽是 0.3-3.4kHz。今用两级 SSB 频分复用系统实现复用。也
即，先把这 12 路话音分 4 组经过第一级 SSB 复用 ,再将 4 个一级复用的结果经过第二级 SSB 复用。已知第一级复用时采用上边带 SSB ,3 个载波频率分别为 0KHz,4KHz 和8KHz。第二级复用采用下单边带 SSB,4个载波频率分别为 84KHz,96KHz,108KHz 120KHz。试画出此系统发送端框图，并画出第一级和第二级复用器输出的频谱(标出频率值，第一级只画一个频谱$_{\mathrm{b}}$
++++++++++++++++++++++++++++++++++++++++++++\\
解：

22/25

---------------------------------------

《通信原理习题解答》

\begin{figure}[h]
	\centering
	\includegraphics[width=\linewidth]{output/image_11_1.png}
\end{figure}
\begin{figure}[h]
	\centering
	\includegraphics[width=\linewidth]{output/image_11_2.png}
\end{figure}
\begin{figure}[h]
	\centering
	\includegraphics[width=\linewidth]{output/image_11_3.png}
\end{figure}
第一级复用器输出的频谱为：
\begin{figure}[h]
	\centering
	\includegraphics[width=\linewidth]{output/image_11_4.png}
\end{figure}
第二级复用器输出的频谱为：
\begin{figure}[h]
	\centering
	\includegraphics[width=\linewidth]{output/image_11_5.png}
\end{figure}
---------------------------------------\\
4.25 单边带调幅信号 $s_\mathrm{SSB}\left(t\right)=A_c\left[m\left(t\right)\cos2\pi f_{\epsilon}t\mp\hat{m}\left(t\right)\sin2\pi f_{\epsilon}t\right]$ ,其中$m\left(t\right)$是 0 均值
平稳随机过程，其自相关函数是$R_{_n}(\tau)$。请：
(1)写出$s_\mathrm{SSB}(t)$的复包络$s_L(t)$,求出复包络的功率谱密度； (2)写出$s_\mathrm{SSB}\left(t\right)$的解析信号$\tilde{s}\left(t\right)$,求出解析信号的功率谱密度； $(3)求出s_{\mathrm{SSB}}(t)$的功率谱密度。

23/25

---------------------------------------

《通信原理习题解答》
++++++++++++++++++++++++++++++++++++++++++++\\
解：(1)$s_{SSB}\left ( t\right )$的 复 包 络 为 
	$$s_L\begin{pmatrix}t\end{pmatrix}=A_c\begin{bmatrix}m(t)\pm\hat{m}(t)\end{bmatrix}$$
	其中$\hat{m}(t)$是$m(t)$的希尔伯特变换，式中的± 号分别对应上单边带和下单边带。$s_L(t)$的自
	$$\begin{aligned}R_{L}\left(t,\tau\right)&=A_{c}^{2}E\left[s_{L}^{*}\left(t\right)s_{L}\left(t+\tau\right)\right]\\&=A_{c}^{2}E\left\{\left[m\left(t\right)\pm j\hat{m}\left(t\right)\right]^{*}\left[m\left(t+\tau\right)\pm j\hat{m}\left(t+\tau\right)\right]\right\}\\&=A_{c}^{2}E\left\{\left[m^{*}\left(t\right)\mp j\hat{m}^{*}\left(t\right)\right]\right]\left[m\left(t+\tau\right)\pm j\hat{m}\left(t+\tau\right)\right]\}\\&=A_{c}^{2}E\left[m^{*}\left(t\right)m\left(t+\tau\right)+\hat{m}^{*}\left(t\right)\hat{m}\left(t+\tau\right)\pm jm^{*}\left(t\right)\hat{m}\left(t+\tau\right)\mp j\hat{m}^{*}\left(t\right)m\left(t+\tau\right)\right]\\&=A_{c}^{2}\left[R_{m}\left(\tau\right)+R_{\hat{m}}\left(\tau\right)\pm jR_{m\hat{m}}\left(\tau\right)\mp jR_{m\hat{m}}\left(\tau\right)\right]\\&\text{中}\text{首尔伯特恋格的性质}\end{aligned}$$
	$$R_{\hat{m}}\left(\tau\right)=R_{m}\left(\tau\right)\\R_{m\hat{m}}\left(\tau\right)=\hat{R}_{m}\left(\tau\right)=-R_{\hat{m}m}\left(\tau\right)$$
	其中$\hat{R}_m\left(\tau\right)$是函数$R_m\left(\tau\right)$的希尔伯特变换。于是
	$$R_{L}\left(\tau\right)=2A_{c}^{2}\left[R_{m}\left(\tau\right)\pm j\hat{R}_{m}\left(\tau\right)\right]$$
	对上式做傅氏变换，设$R_m(\tau)$的傅氏变换是$P_m(f)$。先考虑上单边带的情形，注意到
	
	$R_m\left(\tau\right)+j\hat{R}_m\left(\tau\right)$是$R_m\left(\tau\right)$的解析信号，其傅氏变换只有正频率分量，于是
	$$P_{s_L,\text{USB}}\left(f\right)=\begin{cases}4A_c^2P_m\left(f\right)&f>0\\0&f<0\end{cases}$$
	对于下单边带，$R_m\left(\tau\right)-j\hat{R}_m\left(\tau\right)$是$R_m\left(\tau\right)$的解析信号的共轭，其傅氏变换只有负频率分量，
	于是
	$$P_{s_L,\mathrm{LSB}}\left(f\right)=\begin{cases}\quad0&f>0\\4A_c^2P_m\left(f\right)&f<0\end{cases}$$
	(2) $s_{\mathrm{SSB}}\left ( t\right )$的解析信号为
	$$\begin{aligned}&\tilde{s}\left(t\right)=s_{\mathrm{ssss}}\left(t\right)+j\hat{s}_{\mathrm{sss}}\left(t\right)\\&=A_{c}\left[m\left(t\right)\cos2\pi f_{c}t\mp\hat{m}\left(t\right)\sin2\pi f_{c}t\right]+jA_{c}\left[m\left(t\right)\sin2\pi f_{c}t\pm\hat{m}\left(t\right)\cos2\pi f_{c}t\right]\\&=A_{c}m\left(t\right)\left(\cos2\pi f_{c}t+j\sin2\pi f_{c}t\right)\pm jA_{c}\hat{m}\left(t\right)\left(\cos2\pi f_{c}t+j\sin2\pi f_{c}t\right)\\&=A_{c}\left[m\left(t\right)\pm j\hat{m}\left(t\right)\right]e^{j2\pi f_{c}t}\\&=s_{L}\left(t\right)e^{j2\pi f_{c}t}\end{aligned}$$
	它的自相关函数为
	$$\begin{aligned}R_{\tilde{s}}\left(t,\tau\right)&=E\bigg[\:\tilde{s}^{*}\left(t\right)\tilde{s}\left(t+\tau\right)\bigg]\\&=E\bigg[s_{L}^{*}\left(t\right)e^{-j2\pi f_{c}t}\times s_{L}\left(t+\tau\right)e^{j2\pi f_{c}\left(t+\tau\right)}\bigg]\\&=R_{s_{L}}\left(\tau\right)e^{j2\pi f_{c}\tau}\end{aligned}$$
	做傅氏反变换得
	$$P_{\tilde{s}}\left(f\right)=P_{s_{L},\mathrm{USB}}\left(f-f_{c}\right)$$
	对于上单边带调制
	
	24/25
	
	24
	
	---------------------------------------
	
	《通信原理习题解答》
	$$P_{\tilde{s}}\left(f\right)=\begin{cases}4A_c^2P_m\left(f-f_c\right)&f>f_c\\\quad0&f<f_c\end{cases}$$
	$$P_{\bar{s}}\left(f\right)=\left\{\begin{matrix}0&f>f_c\\4A_c^2P_m\left(f-f_c\right)&f<f_c\end{matrix}\right.$$
	对于下单边带调制
	
	$(3)s_{\mathrm{SSB}}(t)$可以表示为
	$$\begin{aligned}R_{\mathrm{SSB}}\left(t,\tau\right)&=\frac{1}{4}E\left\{\left[\tilde{s}\left(t\right)+\tilde{s}^{*}\left(t\right)\right]^{*}\times\left[\tilde{s}\left(t+\tau\right)+\tilde{s}^{*}\left(t+\tau\right)\right]\right\}\\&=\frac{1}{4}E\left\{\left[\tilde{s}^{*}\left(t\right)+s\left(t\right)\right]\times\left[\tilde{s}\left(t+\tau\right)+\tilde{s}^{*}\left(t+\tau\right)\right]\right\}\\&=\frac{1}{4}E\left\{\tilde{s}^{*}\left(t\right)\tilde{s}\left(t+\tau\right)+\left[\tilde{s}^{*}\left(t\right)\tilde{s}\left(t+\tau\right)\right]^{*}+\tilde{s}\left(t\right)\tilde{s}\left(t+\tau\right)+\left[\tilde{s}\left(t\right)\tilde{s}\left(t+\tau\right)\right]^{*}\right\}\\&=\frac{1}{4}\left\{R_{s}\left(\tau\right)+R_{s}^{*}\left(\tau\right)+E\left[\tilde{s}\left(t\right)\tilde{s}\left(t+\tau\right)\right]+E^{*}\left[\tilde{s}\left(t\right)\tilde{s}\left(t+\tau\right)\right]\right\}\\\text{其中}\end{aligned}$$
	$$s_{\mathrm{SSB}}\left(t\right)=\mathrm{Re}\left\{\tilde{s}\left(t\right)\right\}=\frac{1}{2}\Big[\tilde{s}\left(t\right)+\tilde{s}^{*}\left(t\right)\Big]$$
	$$\begin{aligned}E\Big[\tilde{s}\left(t\right)\tilde{s}\left(t+\tau\right)\Big]&=E\bigg[s_{L}\left(t\right)e^{j2\pi f_{c}t}\times s_{L}\left(t+\tau\right)e^{j2\pi f_{c}\left(t+\tau\right)}\bigg]\\&=E\bigg[s_{L}\left(t\right)s_{L}\left(t+\tau\right)\bigg]e^{j2\pi f_{c}\tau}e^{j4\pi f_{c}t}\\&=R_{s_{L}s_{L}}\left(t,t+\tau\right)e^{j2\pi f_{c}\tau}e^{j4\pi f_{c}t}\end{aligned}$$
	其中
	$$\begin{aligned}R_{s_{L}s_{L}}\left(t,t+\tau\right)&=E\left[s_{L}\left(t\right)s_{L}\left(t+\tau\right)\right]\\&=E\left\{\left[m\left(t\right)\pm j\hat{m}\left(t\right)\right]\times\left[m\left(t+\tau\right)\pm j\hat{m}\left(t+\tau\right)\right]\right\}\\&=E\left[m\left(t\right)m\left(t+\tau\right)-\hat{m}\left(t\right)m\left(t+\tau\right)\pm\hat{m}\left(t\right)m\left(t+\tau\right)\pm m\left(t\right)\hat{m}\left(t+\tau\right)\right]\\&=R_{m}\left(\tau\right)-R_{\hat{m}}\left(\tau\right)\pm jR_{\hat{m}m}\left(\tau\right)\pm jR_{m\hat{m}}\left(\tau\right)=0\end{aligned}$$
	$$R_{\mathrm{SSB}}\left(t,\tau\right)=\frac{1}{4}\left\{R_{\bar{s}}\left(\tau\right)+R_{\bar{s}}^{*}\left(\tau\right)\right\}$$
	做傅氏变换，注意$R_{\bar{s}}^*(\tau)$的傅氏变换是$P_{\bar{s}}^*\left(-f\right)=P_{\bar{s}}\left(-f\right)$,因此
	$$P_{\mathrm{SSB}}\left(f\right)=\frac{1}{4}\Big[P_{s}\left(f\right)+P_{s}\left(-f\right)\Big]=\frac{1}{4}\Big[P_{s_{L}}\left(f-f_{c}\right)+P_{s_{L}}\left(-f-f_{c}\right)\Big]$$
	对于上单边带
	$$P_{\mathrm{USB}}\left(f\right)=\begin{cases}A_c^2P_m\left(f-f_c\right)&f>f_c\\A_c^2P_m\left(f+f_c\right)&f<-f_c\\0&else\end{cases}$$
	对于下单边带，
	$$P_{\mathrm{LSB}}\left(f\right)=\begin{cases}A_c^2P_m\left(f-f_c\right)&0<f<f_c\\A_c^2P_m\left(f+f_c\right)&-f_c<f<0\\0&else\end{cases}$$
	25/25
	
  ---------------------------------------
	
	 第五章 数字信号的基带传输
   ---------------------------------------\\
	5.1 设一数字传输系统传送八进制码元 ,速率为 2400 波特 ,则这时的系统信息速率为多少 ?
	++++++++++++++++++++++++++++++++++++++++\\
	$解：R_b=R_s\log_2M=2400\times\log_28=7200$bps
  ---------------------------------------\\
	5.2 已知：信息代码 1 1 1 0 0 1 0 1
	(1)写出相对码$:\underline1$
	(2)画出相对码的波开图(单极性矩形不归零码)
	++++++++++++++++++++++++++++++++++++++++\\
	 解：(1)写出相对码$:\underline101000110$
	
	
	
	(2)画出相对码的波开图(单极性矩形不归零码)
\begin{figure}[h]
	\centering
	\includegraphics[width=\linewidth]{output/image_14_1.png}
\end{figure}
---------------------------------------\\
5.3 独立随机二进制序列的" $0^{\prime\prime}$、“1”分别由题 3 图中的波形$s_1(t)$及$s_2(t)$表示，已知” O",
“1”等概出现，比特间隔为$T_b$。
(1)若$s_1(t)$如图(a)所示，$s_2(t)=-s_1(t)$,求此数字信号的功率谱密度，并画出图形；
(2)若$s_1(t)$如图(b)所示，$s_2(t)=0$,求此数字信号的功率谱密度，并画出图形。
\begin{figure}[h]
	\centering
	\includegraphics[width=\linewidth]{output/image_14_2.png}
\end{figure}
++++++++++++++++++++++++++++++++++++++++++++\\
解：(1)此时这个数字信号可表示为 PAM 信号
$$s\begin{pmatrix}t\end{pmatrix}=\sum\limits_{n=-\infty}^{\infty}a_{n}g\begin{pmatrix}t-nT_{b}\end{pmatrix}$$
其中序列$\left\{a_n\right\}$以独立等概方式取值于$\pm1,m_\alpha=E\left[a_n\right]=0,\sigma_\alpha^2=E\left[a^2\right]=1;$
$g(t)=s_1(t)$,其傅氏变换是


1/28

---------------------------------------

$$G\bigl(f\bigr)=T_{b}\operatorname{sinc}\bigl(fT_{b}\bigr)$$
所以 $s(t)$的功率谱密度为
$$P_{s}\left(f\right)=\frac{\sigma_{a}^{2}}{T_{b}}\Big|G\Big(f\Big)\Big|^{2}=T_{b}\:\mathrm{sinc}^{2}\Big(fT_{b}\Big)。$$
\begin{figure}[h]
	\centering
	\includegraphics[width=\linewidth]{output/image_15_1.png}
\end{figure}
$(2)此时这个数字信号可表示为$
$$s\begin{pmatrix}t\end{pmatrix}=\sum_{n=-\infty}^{\infty}a_{n}g\begin{pmatrix}t-nT_{b}\end{pmatrix}$$
其中序列$\left\{a_n\right\}$以独立等概方式取值于$\left(0,1\right);g\left(t\right)=s_1\left(t\right)$,其傅氏变换是
$$G\bigl(f\bigr)=\frac{T_{b}}{2}\mathrm{sinc}\biggl(f\frac{T_{b}}{2}\biggr)$$
由于$a_{n}=\frac{1}{2}b_{n}+\frac{1}{2}$
$$s\left(t\right)=\frac{1}{2}\sum_{n=-\infty}^{\infty}b_{n}g\left(t-nT_{b}\right)+\frac{1}{2}\sum_{n=-\infty}^{\infty}g\left(t-nT_{b}\right)$$
$u\left(t\right)=\frac12\sum_{n=-\infty}^{\infty}b_ng\left(t-nT_b\right)-$项的功率谱密度是
$$P_u\left(f\right)=\frac{\left|G(f)\right|^2}{4T_b}=\frac{T_b}{16}\operatorname{sinc}^2\left(f\:\frac{T_b}{2}\right)$$
$\nu\left(t\right)=\frac{1}{2}\sum_{n=-\infty}^{\infty}g\left(t-nT_b\right)$是周期信号，可展成傅氏级数：
$$v\left(t\right)=\frac{1}{2}\sum_{n=-\infty}^{\infty}g\left(t-nT_{b}\right)=\sum_{m=-\infty}^{\infty}c_{m}e^{j2\pi\frac{m}{T_{b}}t}$$
其中
$$\begin{aligned}c_{m}&=\frac{1}{T_{b}}\int_{-T_{b}/2}^{T_{b}/2}\frac{1}{2}\sum_{n=-\infty}^{\infty}g\left(t-nT_{b}\right)e^{-j2\pi\frac{m}{T_{b}}t}dt=\frac{1}{2T_{b}}\int_{-T_{b}/2}^{T_{b}/2}g\left(t\right)e^{-j2\pi\frac{m}{T_{b}}t}dt\\&=\frac{1}{2T_{b}}G\left(\frac{m}{T_{b}}\right)=\frac{\sin\left(\pi\frac{m}{T_{b}}\times\frac{T_{b}}{2}\right)}{4\left(\pi\frac{m}{T_{b}}\times\frac{T_{b}}{2}\right)}=\begin{cases}\pm\frac{1}{2m\pi}&m=\pm1,\pm3,\cdots\\\\\frac{1}{4}&m=0\\\\0&\mathrm{other~}m\end{cases}\end{aligned}$$
2/28

---------------------------------------

所以 $\nu(t)=\frac12\sum_{m=-\infty}^{\infty}g\left(t-nT_{b}\right)$的功率谱密度是
$$P_{\nu}\left(f\right)=\sum_{n=-\infty}^{\infty}\left|c_{n}\right|^{2}\delta\left(f-\frac{n}{T_{b}}\right)=\frac{1}{16}\delta\left(f\right)+\frac{1}{4\pi^{2}}\sum_{k=-\infty}^{\infty}\frac{1}{\left(2k-1\right)^{2}}\delta\left(f-\frac{2k-1}{T_{b}}\right)$$
于是 $s(t)$的功率谱密度为：
$$P_{s}\left(f\right)=\frac{T_{b}}{16}\operatorname{sinc}^{2}\left(f\:\frac{T_{b}}{2}\right)+\frac{1}{16}\:\delta\left(f\right)+\frac{1}{4\pi^{2}}\sum_{k=-\infty}^{\infty}\frac{1}{\left(2k-1\right)^{2}}\:\delta\left(f-\frac{2k-1}{T_{b}}\right)$$
\begin{figure}[h]
	\centering
	\includegraphics[width=\linewidth]{output/image_16_1.png}
\end{figure}
解二 本题还可套用公式(5.2.12):
$$P_s\left(f\right)=\frac{1}{T_b}P_a\left(f\right)\left|G_r\left(f\right)\right|^2$$
其中$P_a\left(f\right)$是序列$R_a\left(m\right)$的离散时间傅氏变换(时域离散，但频域连续，注意和 DFT 的差别 ),$G_T\left(f\right)$是发送脉冲$g_T\left(t\right)$的连续时间傅氏变换。对于独立平稳序列$\left\{a_n\right\}$,我们有式(5.2.14):
$$R_a\left(m\right)=m_a^2+\sigma_a^2\delta\left(m\right)$$
其中$m_a= E[ a_n]$ , $\sigma _\alpha ^2= E\left [ \left ( a_n- m_\alpha \right ) ^2\right ] = E\left [ a_\alpha ^2\right ] - m_\alpha ^2$ , $\delta \left ( m\right )$是离散时间的冲激函数

( Dirac delta function)。对其作傅氏变换，注意常值序列实际是直流的理想取样，所以常值
序列$\left\{m_\alpha^2\right\}$的傅氏变换是$\frac1{T_b}\sum_m\delta\left(f-\frac m{T_b}\right),\delta(m)$的傅氏变换是1。所以有式(5.2.18)
$$P_a\left(f\right)=\sigma_a^2+\frac{m_a^2}{T_b}\sum_m\delta\Bigg(f-\frac{m}{T_b}\Bigg)$$
在第(1)小题中，$\left\{a_n\right\}$是取值于$\left\{\pm1\right\}$的独立等概序列，因此$m_\alpha=0$,$\sigma_\alpha^2=1,P_\alpha\left(f\right)=1$,
$G_{T}\left(f\right)=T_{b}\operatorname{sinc}(fT_{b})$,所以
$$P_{s}\left(f\right)=\frac{1}{T_{b}}P_{a}\left(f\right)\Big|G_{T}\left(f\right)\Big|^{2}=T_{b}\:\mathrm{sinc}^{2}\left(fT_{b}\right)$$
在第(2)小题中，$\{ a_{n}\}$是 取 值 于 $\{ 0, 1\}$的 独 立 等 概 序 列 $, m_{\alpha }= \frac 12, \sigma _{\alpha }^{2}= \frac 14.$
$$P_{a}\left(f\right)=\frac{1}{4}+\frac{1}{4T_{b}}\sum_{m}\delta\Bigg(f-\frac{m}{T_{b}}\Bigg)\:,\:G_{r}\left(f\right)=\frac{T_{b}}{2}\operatorname{sinc}\Bigg(f\:\frac{T_{b}}{2}\Bigg)\:,\:\text{所以}$$
3/28

---------------------------------------

$$\begin{aligned}P_{s}\left(f\right)&=\frac{1}{T_{b}}P_{a}\left(f\right)\left|G_{T}\left(f\right)\right|^{2}\\&=\frac{1}{4T_{b}}\biggl[1+\frac{1}{T_{b}}\sum_{m}\delta\biggl(\:f-\frac{m}{T_{b}}\biggr)\biggr]\times\frac{T_{b}^{2}}{4}\:\mathrm{sinc}^{2}\biggl(\frac{fT_{b}}{2}\biggr)\\&=\frac{1}{16}\biggl\{T_{b}\:\mathrm{sinc}^{2}\biggl(\frac{T_{b}}{2}\biggr)+\sum_{m}\mathrm{sinc}^{2}\biggl(\frac{m}{2}\biggr)\delta\biggl(\:f-\frac{m}{T_{b}}\biggr)\biggr\}\\&=\frac{T_{b}\:\mathrm{sinc}^{2}\biggl(\frac{fT_{b}}{2}\biggr)}{16}+\frac{\delta\left(f\right)}{16}+\frac{\delta\left(f+\frac{1}{T_{b}}\right)}{4\pi^{2}}+\frac{\delta\left(f\pm\frac{3}{T_{b}}\right)}{36\pi^{2}}+\frac{\delta\left(f\pm\frac{5}{T_{b}}\right)}{100\pi^{2}}\cdots\end{aligned}$$
---------------------------------------\\
5.4 假设信息比特 1、0 以独立等概方式出现 ,求数字分相码的功率谱密度。
++++++++++++++++++++++++++++++++++++++++++++\\
解一：数字分相码可以表示成二进制 PAM 信号的形式
$$\begin{aligned}&\text{一进制 PAM 信号的形式}\\&s\left(t\right)=\sum_{n=-\infty}^{\infty}a_{n}g\left(t-nT_{b}\right)\\&\text{双值于}\pm1\:,\:m_{a}=E\Big[a_{n}\Big]=0\:,\:\sigma_{a}^{2}=E\Big[\:a^{2}\:\Big]=1\:;\:g\left(t\right)\text{如下}\end{aligned}$$
$\bar{\text{列}\{a_n\}}$l
图所示
\begin{figure}[h]
	\centering
	\includegraphics[width=\linewidth]{output/image_17_1.png}
\end{figure}
其傅氏变换是
$$\begin{aligned}G\left(f\right)&=-\frac{AT_{b}}{2}\operatorname{sinc}\biggl(f\:\frac{T_{b}}{2}\biggr)e^{j2\pi f\frac{T_{b}}{4}}+\frac{AT_{b}}{2}\operatorname{sinc}\biggl(f\:\frac{T_{b}}{2}\biggr)e^{-j2\pi f\frac{T_{b}}{4}}\\&=-jAT_{b}\sin\frac{\pi fT_{b}}{2}\operatorname{sinc}\biggl(\frac{fT_{b}}{2}\biggr)\end{aligned}$$
$$P_{s}\left(f\right)=\frac{\sigma_{a}^{2}}{T_{b}}\Big|G\Big(f\Big)\Big|^{2}=A^{2}T_{b}\sin^{2}\Bigg(\frac{\pi fT_{b}}{2}\Bigg)\:\mathrm{sinc}^{2}\Bigg(\frac{fT_{b}}{2}\Bigg)。$$
所以


解二：假设二进制 0 映射为-1,1 映射为-1。记信息码序列为$\left\{a_n\right\},a_n\in\left\{\pm1\right\}$,编码结果为 $\{ b_k\}$ $, b_k\in \{ \pm 1\}$ $, \{ a_n\}$ 中 的 第 $n$ 个 $a_n$ 对 应 $\{ b_k\}$ 中 的 $b_{2n}, b_{2n+ 1}$。则 数 字 分 相 码 波 形 可 以 写 成
$$s\begin{pmatrix}t\end{pmatrix}=\sum_{-\infty}^{\infty}b_{k}g\begin{pmatrix}t-kT_{s}\end{pmatrix}$$
4/28

---------------------------------------

其中$T_s= \frac {T_{b}}2$ , $g\left ( t\right ) = \begin{cases} A& {0\leq t< T_{s}}\\ 0& {else} \end{cases}$。


按照数字分相码的编码规则，有下面的关系
$$\begin{cases}b_{2n}=a_n\\b_{2n+1}=-a_n\end{cases}$$
因而
$$\begin{aligned}s\left(t\right)&=\sum_{k=-\infty}^{\infty}b_{k}g\left(t-kT_{s}\right)\\&=\sum_{k\overset{k=-\infty}{\mathrm{is~even}}}^{\infty}b_{k}g\left(t-kT_{s}\right)+\sum_{\overset{k=-\infty}{\mathrm{is~odd}}}^{\infty}b_{k}g\left(t-kT_{s}\right)\\&=\sum_{n=-\infty}^{\infty}a_{n}g\left(t-2nT_{s}\right)-\sum_{n=-\infty}^{\infty}a_{n}g\left(t-2nT_{s}-T_{s}\right)\end{aligned}$$
$令u\left(t\right)=\sum_{n=-\infty}^{\infty}a_{n}g\left(t-2nT_{s}\right),则$
$$s\begin{pmatrix}t\end{pmatrix}=u\begin{pmatrix}t\end{pmatrix}-u\begin{pmatrix}t-T_s\end{pmatrix}$$
所以
$$\begin{aligned}P_{s}\left(f\right)&=P_{u}\left(f\right)\left|1-e^{j2\pi fT_{s}}\right|^{2}\\&=P_{u}\left(f\right)\left|e^{j\pi fT_{s}}\left(e^{-j\pi fT_{s}}-e^{j\pi fT_{s}}\right)\right|^{2}\\&=4P_{u}\left(f\right)\sin^{2}\left(\frac{\pi fT_{b}}{2}\right)\end{aligned}$$
而$u\left(t\right)$是速率为$\frac1{2T_{s}}=\frac1{T_{b}}$的双极性 RZ 码，所以
$$P_{u}\left(f\right)=\frac{1}{T_{b}}\Big|G\big(f\big)\Big|^{2}=\frac{A^{2}T_{b}}{4}\mathrm{sinc}^{2}\Bigg(\frac{fT_{b}}{2}\Bigg)$$
$$P_s\left(f\right)=A^2T_b\operatorname{sinc}^2\Bigg(\frac{fT_b}{2}\Bigg)\:\sin^2\Bigg(\frac{\pi fT_b}{2}\Bigg)$$
故

解三：假设二进制 0 映射为-1,1 映射为-1。记信息码序列为$\left\{a_n\right\},a_n\in\left\{\pm1\right\}$,编码结果为$\left\{b_{k}\right\},b_{k}\in\left\{\pm1\right\},\left\{a_{n}\right\}$中的第 $n$ 个$a_n$ 对应$\left\{b_k\right\}$中的$b_2n,b_{2n+1}$。则数字分相码波形可以写成
$$s\begin{pmatrix}t\end{pmatrix}=\sum_{-\infty}^{\infty}b_{k}g\begin{pmatrix}t-kT_{s}\end{pmatrix}$$
其中$T_s= \frac {T_b}2$ , $g\left ( t\right ) = \begin{cases} A& 0< t< T_{s}\\ 0& else \end{cases}$。按照数字分相码的编码规则，有下面的关系
$$\begin{cases}b_{2n}=a_n\\b_{2n+1}=-a_n\end{cases}$$
$\{b_k\}$的自相关函数为
$$R_b\left(i,i+m\right)=E\left[b_ib_{i+m}\right]$$
5/28

---------------------------------------

若$i$ 为偶数，即若$i=2n$,则有
$$R_{b}\left(2n,2n+m\right)=E\left[a_{n}b_{2n+m}\right]=\begin{cases}1&m=0\\-1&m=1\\0&else\end{cases}$$
若$i$ 为奇数，即若$i=2n+1$,则有
$$R_{b}\left(2n+1,2n+1+m\right)=E\Big[-a_{n}b_{2n+1+m}\Big]=\begin{cases}1&m=0\\-1&m=-1\\0&else\end{cases}$$
可见序列$\left\{b_k\right\}$是非平稳的，是周期为2 的循环平稳。求其平均自相关函数得
$$\begin{aligned}\overline{R_{b}}\left(m\right)&=\frac{1}{2}\Big[R_{b}\left(2n,2n+m\right)+R_{b}\left(2n+1,2n+1+m\right)\Big]\\&=\begin{cases}1&m=0\\-\frac{1}{2}&m=\pm1\\0&else\end{cases}\\\end{aligned}$$
求其离散时间傅氏变换
$$\begin{aligned}P_{b}\left(f\right)&=\sum_{m=-\infty}^{\infty}R_{b}\left(m\right)e^{-j2\pi fmT_{s}}\\&=1-\frac{1}{2}e^{j\pi fT_{b}}-\frac{1}{2}e^{-j\pi fT_{b}}\\&=1-\cos\pi fT_{b}\\&=2\sin^{2}\left(\frac{\pi fT_{b}}{2}\right)\end{aligned}$$
于是
$$\begin{aligned}P_{s}\left(f\right)&=\frac{1}{T_{s}}P_{b}\left(f\right)\left|G\left(f\right)\right|^{2}\\&=\frac{1}{T_{b}/2}\times2\sin^{2}\left(\frac{\pi fT_{b}}{2}\right)\times\left|\frac{AT_{b}}{2}\operatorname{sinc}\left(\frac{fT_{b}}{2}\right)\right|^{2}\\&=A^{2}T_{b}\sin^{2}\left(\frac{\pi fT_{b}}{2}\right)\mathrm{sinc}^{2}\left(\frac{fT_{b}}{2}\right)\end{aligned}$$
---------------------------------------\\
5.5 已知信息代码为：1000000000111001000010,请就 AMI 码、HDB3 码，
Manchester 码三种情形
(1)给出编码结果；
(2)画出编码后的波形； (3)画出提取时钟的框图。
++++++++++++++++++++++++++++++++++++++++++++\\
$解：(1)$
100000000111001000010
信息代码：
+1 0 0 0 0 0 0 0 0 0 0 0 0 - 1 + 1 1 0 0 0 0 0 0 1 0
AMI 码：
+1 0 0 0 $0+ V- B$ 0 $0- V+ 1- 1+ 1$ 0 $0- 1+ B$ 0 $0+ V- 1$ 0
HDB3 码：
Manchester 码：1001010101010101011010100101100101011001
(2)波形如下：


6/28

---------------------------------------

\begin{figure}[h]
	\centering
	\includegraphics[width=\linewidth]{output/image_20_5.png}
\end{figure}
\begin{figure}[h]
	\centering
	\includegraphics[width=\linewidth]{output/image_20_1.png}
\end{figure}
(3)AMI 码和 HDB3 码可以看成是一种双极性的 RZ 信号，经过全波整流后成为单极性 RZ
信号，它包含时钟的线谱分量，故此可直接提取。提取时钟的框图如下：
\begin{figure}[h]
	\centering
	\includegraphics[width=\linewidth]{output/image_20_2.png}
\end{figure}
Macnshester 码经过全波整流后是直流，不能用上述办法。需要先微分，使之成为一种双极性 RZ 信号，然后再用上述办法，不过注意这里提出的是二倍频时钟，需要二分频。提取时钟的框图如下：
\begin{figure}[h]
	\centering
	\includegraphics[width=\linewidth]{output/image_20_3.png}
\end{figure}
---------------------------------------\\
5.6 已知二进制序列的“1”和”0”分别由波形$s_{_1}\left(t\right)=\begin{cases}A&0\leq t\leq T_{b}\\0&else\end{cases}$及$s_{_2}\left(t\right)=0$表示'1" 与”0"等概出现。此信号在信道传输中受到功率谱密度为$\frac{N_0}2$的加性白高斯噪声$n(t)=0$ 的干扰 接收端用下图所示的框图进行接收。图中低通滤波器的带宽是$B\not B$足够大使得 $s_i(t)$ 经过滤波器时近似无失真。
\begin{figure}[h]
	\centering
	\includegraphics[width=\linewidth]{output/image_20_4.png}
\end{figure}
7/28

---------------------------------------

(1)若发送$s_1(t)$,请写出 y(t)表示式，求出抽样值 y 的条件均值$E[y|s_1]$及条件方
差$D\left[y\mid s_{1}\right]$,写出此时 y 的条件概率密度函数 $p_1(y)=p(y\mid s_{1});$
(2) 若发送$s_{2}(t)$,请写出 y$(t)$表示式，求出抽样值 y的条件均值$E[y|s_{2}]$及条件方
差$D\left[y\mid s_{2}\right]$,写出此时 y 的条件概率密度函数 $p_2( y) = p( y\mid s_{2})$ ;
(3)画出$p_1(y)$及$p_2(y)$的图形；
(4)求最佳判决门限$V_{7}$值； (5)推导出平均误比特率。
++++++++++++++++++++++++++++++++++++++++++++\\
解：(1)此时在$0\leq t\leq T_{b}$时间范围内，$y(t)=A+\xi(t)$,$\xi(t)$是白高斯噪声$n(t)$通过低通滤波器后的输出，显然$\xi(t)$是 0 均值的高斯平稳过程 ,其方差为$\sigma^2=N_0B$ 。于是抽样值 y 的条件均值是A，条件方差是$N_0B$ ,条件概率密度函数是 $p_1(y)=\frac1{\sqrt{2\pi N_{\circ}B}}e^{\frac{(y-A)^{2}}{2N_{\circ}B}}$
(2)此时在$0\leq t\leq T_b$时间范围内，$y(t)=\xi(t)$,$\xi(t)$是白高斯噪声$n(t)$通过低通滤波器后输出的 0 均值高斯平稳过程 ,其方差为$\sigma^2=N_0B$ 。因此 ,抽样值 y 的条件均值是 0 ,条${\text{件 方 差 是 }N_{0}B}$ ,其条件概率密度函数是$p_2\left(y\right)=\frac1{.\sqrt{2\pi NB}}e^{-\frac y{2N_{0}B}}$

(3)

\begin{figure}[h]
	\centering
	\includegraphics[width=\linewidth]{output/image_21_1.png}
\end{figure}
(4)最佳门限$V_T$是后验概率相等的分界点，即为方程$P(s_1|y)=P(s_2|y)$的解。由于与“ 0”等概出现，所以$V_T$是$p_1(y)=p_2(y)$的解，即$p_1(y)$和$p_2(y)$的交点。由上图可知$V_T=\frac A2$。
(5)
$$P_b=P\big(s_1\big)P\big(e\:|\:s_1\big)+P\big(s_2\big)P\big(e\:|\:s_2\big)=\frac{1}{2}\Big[\:P\big(e\:|\:s_1\big)+P\big(e\:|\:s_2\big)\Big]$$
8/28

---------------------------------------

$$\begin{aligned}P\left(e\mid s_{1}\right)&=P\left(y<V_{T}\mid s_{1}\right)=P\left(A+\xi<\frac{A}{2}\right)=P\left(\xi<-\frac{A}{2}\right)\\&=\frac{1}{2}\operatorname{erfc}\Bigg(\frac{A/2}{\sqrt{2N_{0}B}}\Bigg)=\frac{1}{2}\operatorname{erfc}\Bigg(\sqrt{\frac{A^{2}}{8N_{0}B}}\Bigg)\\P\left(e\mid s_{2}\right)&=P\left(y>V_{T}\mid s_{2}\right)=P\left(\xi>\frac{A}{2}\right)=P\left(\xi<-\frac{A}{2}\right)=\frac{1}{2}\operatorname{erfc}\left(\sqrt{\frac{A^{2}}{8N_{0}B}}\right)\\& P_{b}=\frac{1}{2}\operatorname{erfc}\left(\sqrt{\frac{A^{2}}{8N_{0}B}}\right)\end{aligned}$$
---------------------------------------\\
5.7 题7图中$s_1(t)$与$s_2(t)$等概出现，白高斯噪声$n_w(t)$的功率谱密度是$\frac{N_0}2$
\begin{figure}[h]
	\centering
	\includegraphics[width=\linewidth]{output/image_22_2.png}
\end{figure}
(1)画出匹配滤波器的冲激响应$h(t);$

(2)求发$s_1(t)$条件下求抽样瞬时值 y 中信号分量的幅度及功率、噪声分量的平均功率 ;
(3)求发$s_2(t)$条件下抽样值 y 的均值、方差和概率密度函数$p_2(y);$
(4)求平均误比特率。
++++++++++++++++++++++++++++++++++++++++++++\\
解：(1)$h(t)=s_{1}\left(T_{b}-t\right)=s_{2}\left(t\right),波形如下$

\begin{figure}[h]
	\centering
	\includegraphics[width=\linewidth]{output/image_22_1.png}
\end{figure}
(2)抽样值 y 中的信号分量是
$$\int_{-\infty}^{\infty}s_{1}\left(T-\tau\right)h\left(\tau\right)d\tau=\int_{-\infty}^{\infty}s_{2}\left(\tau\right)s_{2}\left(\tau\right)d\tau=A^{2}T_{b}=E_{b}$$
此处$E_b=E_1=E_2$是每个比特的能量。抽样瞬时值中信号分量的功率是$E_b^2$。抽样值中的噪
声分量的功率是

9/28

---------------------------------------

$$\sigma^{2}=\int_{-\infty}^{\infty}\frac{N_{0}}{2}\Big|H\left(f\right)\Big|^{2}df=\frac{N_{0}}{2}\int_{-\infty}^{\infty}\Big|S_{2}\left(f\right)\Big|^{2}df=\frac{N_{0}E_{b}}{2}$$
(3)此时 y 中的信号分量是
$$\int_{-\infty}^{\infty}s_{2}\left(T-\tau\right)h\left(\tau\right)d\tau=\int_{-\infty}^{\infty}s_{1}\left(\tau\right)s_{2}\left(\tau\right)d\tau=-A^{2}T_{b}=-E_{b}\\y=-E_{b}+\xi $$
操声分量$\xi$是 0 均值高斯随机变量，其方差是$\frac{N_0E_b}2$。因此，在发送$s_2(t)$的条件下，y 最匀值为$-E_b=-A^{2}T_{b}$、方差为$\frac N{0}E_{b}2$的高斯随机变量，所以$p_2\left(y\right)=\frac1{\sqrt{\pi N_{0}E_{b}}}e^{-\frac{\left(y+E_{b}\right)^{2}}{N_{0}E_{b}}}$ (4)类似地，也可以得到发送$s_1(t)$条件下 y 的概率密度函数为
$$p_{1}\left(y\right)=p_{y|s_{1}}\left(y\:|\:s_{1}\right)=\frac{1}{\sqrt{\pi N_{0}E_{b}}}e^{-\frac{\left(y-E_{b}\right)^{2}}{N_{0}E_{b}}}$$
最佳门限$V_T$是$p_1(y)=p_2(y)$的解，即$V_T=0$。平均错误率为
$$P_b=P\left(s_1\right)P\left(e\mid s_1\right)+P\left(s_2\right)P\left(e\mid s_2\right)=\frac{1}{2}\Big[P\left(e\mid s_1\right)+P\left(e\mid s_2\right)\Big]$$
其中
$$P\left(e\mid s_{1}\right)=P\left(y<0\mid s_{1}\right)=P\left(E_{b}+\xi<0\right)=P\left(\xi<-E_{b}\right)=P\left(\xi>E_{b}\right)\\P\left(e\mid s_{2}\right)=P\left(y>0\mid s_{2}\right)=P\left(-E_{b}+\xi>0\right)=P\left(\xi>E_{b}\right)\\P\left(\xi>E_{b}\right)=\frac{1}{2}\operatorname{erfc}\left(\frac{E_{b}}{\sqrt{2\sigma^{2}}}\right)=\frac{1}{2}\operatorname{erfc}\left(\sqrt{\frac{E_{b}}{N_{0}}}\right)\\ P_{b}=\frac{1}{2}\operatorname{erfc}\left(\sqrt{\frac{E_{b}}{N_{0}}}\right)$$
---------------------------------------\\
5.8 设基带传输系统的发送滤波器、信道和接收滤波器的总传输特性$H(f)$如题 8 图所示，其中$f_1= 1$MHz$, f_2= 3$MHz。试确定该系统无码间干扰传输时的最高码元速率和频带利用率。
\begin{figure}[h]
	\centering
	\includegraphics[width=\linewidth]{output/image_23_1.png}
\end{figure}
++++++++++++++++++++++++++++++++++++++++++++\\
解：由下图可见，当 $f_s=f_1+f_2$ 时$\sum_n=-\infty^\infty H\left(f-nf_s\right)=1$,且不存在更大的$f_s$ 假$\sum_{n=-\infty}^\infty H\left(f-nf_s\right)$等于常数，因此该系统无码间干扰传输时的最高码元速率是

10/28

---------------------------------------

$R_{s}=f_{1}+f_{2}=4$Mbaud,对应的频带利用率是$\frac R{s}{f_{s}}=\frac{4}{3}$Baud/Hz
\begin{figure}[h]
	\centering
	\includegraphics[width=\linewidth]{output/image_24_3.png}
\end{figure}
---------------------------------------\\
5.9 题 9 图中示出了一些基带传输系统的总体传输特性$H\left(f\right)$,若要以 2000 波特的码元速
率传输，请问哪个$H(f)$满足抽样点无码间干扰的条件？
\begin{figure}[h]
	\centering
	\includegraphics[width=\linewidth]{output/image_24_1.png}
\end{figure}
++++++++++++++++++++++++++++++++++++++++++++\\
解：(a)由下图可见，$\sum_n=-\infty^{\infty}H\left(f-2000n\right)$不是常数，因此该系统以 2000 波特的码元速率
传输时不满足抽样点无码间干扰的条件。
\begin{figure}[h]
	\centering
	\includegraphics[width=\linewidth]{output/image_24_2.png}
\end{figure}
11/28

---------------------------------------
\begin{figure}[h]
	\centering
	\includegraphics[width=\linewidth]{output/image_25_3.png}
\end{figure}
(c)由图可知，该系统以 2000 波特的码元速率传输时满足抽样点无码间干扰的条件。
\begin{figure}[h]
	\centering
	\includegraphics[width=\linewidth]{output/image_25_1.png}
\end{figure}
---------------------------------------\\
5.10 设滚降系数为$\alpha=1$ 的升余弦滚降无码间干扰基带传输系统的输入是十六进制码元 其码元速率是$R_{s}=1200$Baud ,求此基带传输系统的截止频率、该系统的频带利用率以及该系统的信息传输速率。
++++++++++++++++++++++++++++++++++++++++++++\\

解：此系统的截止频率为$B=(1+\alpha)\frac{R_{s}}{2}=1200$Hz,频带利用率为
$\frac{R_{s}}{R}=1$Baud/Hz=4bit/s/Hz ,信息传输速率为$R_b=R_s\log_2M=4800$bit/s
---------------------------------------\\
5.11 将二进制双极性不归零脉冲序列通过 16 进制脉孔均衡器$H(f)$及$\alpha=1$ 的根号升余弦滤波器，再将此数字基带信号送至基带信道中传输， 如题 11 图所示。其中网孔均衡的作用是使其输出具有平坦频谱(近似冲激序列 )

\begin{figure}[h]
	\centering
	\includegraphics[width=\linewidth]{output/image_25_2.png}
\end{figure}
设二进制信息速率为$R_b=4$Mb/s,$\{b_n\}$以独立等概方式取值于$\pm1$,16PAM 的幅度取值于
集合$\left\{\pm1,\pm3,\pm5,\cdots,\pm15\right\},g\left(t\right)$是幅度为1，宽度为$T_s$的矩形脉冲。求：
(1)A 点符号速率$R_s$及功率谱密度表达式，并画出功率谱密度图；
(2)B 点功率谱密度表达式，并画出功率谱密度图；
(3)写出$\left|H\left(f\right)\right|^2$表达式。
++++++++++++++++++++++++++++++++++++++++++++\\
$$\text{解:}\left(1\right)R_{s}=\frac{R_{b}}{\log_{2}16}=1\text{MBaud。}$$
因为$\{b_n\}$独立等概，所以序列$\{a_m\}$也是独立等概的。从而有$E[a_m]=0$, $E\left[a_{m}^{2}\right]=\frac{1+3^{2}+\cdots+15^{2}}{8}=85$。由于$g\left(t\right)$的傅氏变换是$G\left(f\right)=T_s\operatorname{sinc}\left(fT_s\right)$,故此点的功率谱密度是



12/28

---------------------------------------

$$P_A\left(f\right)=\frac{E\left[a_m^2\right]}{T_s}\Big|G\left(f\right)\Big|^2=85T_s\text{sinc}^2\left(fT_s\right)$$
\begin{figure}[h]
	\centering
	\includegraphics[width=\linewidth]{output/image_26_2.png}
\end{figure}
(2)B 点信号可写成 $s_B\left(t\right)=\sum_{m=-\infty}^{\infty}a_mg_B\left(t-mT_s\right)$,则依题意脉冲$g_B\left(t\right)$的傅氏变换是
$G_{_B}\left ( f\right ) = \sqrt {H_{\mathrm{rcos}}\left ( f\right ) }$ ,其中
$$H_{\mathrm{rcos}}\left(f\right)=\begin{cases}\frac{T_{s}}{2}\Big(1+\cos\pi fT_{s}\Big)=T_{s}\cos^{2}\Big(\frac{\pi fT_{s}}{2}\Big)&\Big|f\Big|\leq\frac{1}{T_{s}}\\0&\mathrm{else}\end{cases}$$
故此 B 点的功率谱密度是
$$P_{A}\left(f\right)=\frac{E\left[a_{m}^{2}\right]}{T_{s}}\Big|G_{B}\left(f\right)\Big|^{2}=\begin{cases}85\cos^{2}\left(\frac{fT_{s}}{2}\right)&\left|f\right|\leq\frac{1}{T_{s}}\\\\0&0\end{cases}$$
\begin{figure}[h]
	\centering
	\includegraphics[width=\linewidth]{output/image_26_1.png}
\end{figure}
(3)网孔均衡器输出有平坦的功率谱，故需$\left|H(f)G(f)\right|$为常数。由于

$G\left(f\right)=T_{s}$sinc$\left(fT_s\right)$,所以$\left|H\left(f\right)\right|^2=\frac K{\mathrm{sinc}^2\left(fT_s\right)}$,$K$是任意常数
---------------------------------------\\
5.12 某基带传输系统接收端的抽样值为 $y=a+n+i_m$ ,其中$a$ 代表发送信息，它等概取值于$\pm1$。$n$ 代表高斯噪声分量，其均值为 0 ,方差为$\sigma^2$。$i_m$代表码间干扰，$i_m$有3个可能的取 值 : $- \frac 12, 0, \frac 12$, 它 们 的 出 现 概 率 分 别 为 $\frac 14, \frac 12, \frac 14$。若 接 收 端 的 判 决 门 限 是 0, 求 该 系 统 的 平 均 误 比 特 率 。
++++++++++++++++++++++++++++++++++++++++++++\\
解：

13/28

---------------------------------------

$$\begin{aligned}&P_{a}=P\left(e\mid a=1\right)P\left(a=1\right)+P\left(e\mid a=-1\right)P\left(a=-1\right)\\&=\frac{P\left(e\mid a=1\right)+P\left(e\mid a=-1\right)}{2}\\P\left(e\mid a=1\right)&=P\left(1+n+i_{m}<0\right)=P\left(n<-1-i_{m}\right)\\&=P\left(n<-\frac{1}{2}\right)P\left(i_{m}=-\frac{1}{2}\right)+P\left(n<-1\right)P\left(i_{m}=0\right)+P\left(n<-\frac{3}{2}\right)P\left(i_{m}=\frac{1}{2}\right)\\&=\frac{1}{8}\operatorname{erfc}\left(\frac{1/2}{\sqrt{2\sigma^{2}}}\right)+\frac{1}{4}\operatorname{erfc}\left(\frac{1}{\sqrt{2\sigma^{2}}}\right)+\frac{1}{8}\operatorname{erfc}\left(\frac{3/2}{\sqrt{2\sigma^{2}}}\right)\\&=\frac{1}{8}\operatorname{erfc}\left(\frac{1/2}{\sqrt{8\sigma^{2}}}\right)+\frac{1}{4}\operatorname{erfc}\left(\frac{1}{\sqrt{2\sigma^{2}}}\right)+\frac{1}{8}\operatorname{erfc}\left(\frac{3}{\sqrt{8\sigma^{2}}}\right)\\P\left(a\mid a=-1\right)&=P\left(1+n+i_{m}>0\right)=P\left(n>1-i_{m}\right)\\&=P\left(n>\frac{1}{2}\right)P\left(i_{m}=\frac{1}{2}\right)+P\left(n>1\right)P\left(i_{m}=0\right)+P\left(n>\frac{3}{2}\right)P\left(i_{m}=\frac{-1}{2}\right)\\&=\frac{1}{8}\operatorname{erfc}\left(\frac{1/2}{\sqrt{2\sigma^{2}}}\right)+\frac{1}{4}\operatorname{erfc}\left(\frac{1}{\sqrt{2\sigma^{2}}}\right)+\frac{1}{8}\operatorname{erfc}\left(\frac{3/2}{\sqrt{2\sigma^{2}}}\right)\\&=\frac{1}{8}\operatorname{erfc}\left(\frac{1}{\sqrt{8\sigma^{2}}}\right)+\frac{1}{4}\operatorname{erfc}\left(\frac{1}{\sqrt{2\sigma^{2}}}\right)+\frac{1}{8}\operatorname{erfc}\left(\frac{3}{\sqrt{8\sigma^{2}}}\right)\\P\left(a\mid a=-1\right)&=P\left(1+n+i_{m}>0\right)=P\left(n>1-i_{m}\right)\\&=P\left(n>\frac{1}{2}\right)P\left(i_{m}=\frac{1}{2}\right)+P\left(n>1\right)P\left(i_{m}=0\right)+P\left(n>\frac{3}{2}\right)P\left(i_{m}=\frac{-1}{2}\right)\\&=\frac{1}{8}\operatorname{erfc}\left(\frac{1/2}{\sqrt{2\sigma^{2}}}\right)+\frac{1}{4}\operatorname{erfc}\left(\frac{1}{\sqrt{2\sigma^{2}}}\right)+\frac{1}{8}\operatorname{erfc}\left(\frac{3/2}{\sqrt{2\sigma^{2}}}\right)\\\end{aligned}$$
$$P_{e}=\frac{1}{8}\operatorname{erfc}\left(\frac{1}{\sqrt{8\sigma^{2}}}\right)+\frac{1}{4}\operatorname{erfc}\left(\frac{1}{\sqrt{2\sigma^{2}}}\right)+\frac{1}{8}\operatorname{erfc}\left(\frac{3}{\sqrt{8\sigma^{2}}}\right)$$
---------------------------------------\\
5.13 下图中，速率为$R_b=9600$bit/s 的二进制序列经串并变换、D/A 变换后输出 8PAM 信号，然后经过频谱成形电路使输出脉冲具有滚降因子$\alpha=0.5$的根号升余弦频谱，再发送至基带信道。求发送信号的符号速率$R$、要求的基带信道的带宽以及该系统的频带利用率。
\begin{figure}[h]
	\centering
	\includegraphics[width=\linewidth]{output/image_27_1.png}
\end{figure}
++++++++++++++++++++++++++++++++++++++++++++\\
解：发送信号的符号速率为$R_{s}=\frac{R_{b}}{3}=3200$Baud,要求的基带信号的带宽为
$V=\left(1+\alpha\right)\frac{R_{s}}{2}=2400$Hz,频带利用率为$\frac{R_{b}}{W}=\frac{9600}{2400}=4$ bits/s/H
---------------------------------------\\
5.14 二进制信息序列经 MAPM 调制及升余弦滚降频谱成形后通过基带信道传输，已知基
带信道的带宽是 W=3000Hz。若滚降系数$\alpha$分别为 0、0.5、1。请
(1)分别求出该系统无码间干扰传输时的符号速率；
(2) 若 MPAM 的进制数 M 是 16,请写出其相应的二进制信息速率。

14/28

---------------------------------------
++++++++++++++++++++++++++++++++++++++++++++\\
解：(1)$ W=\left(1+\alpha\right)\frac{R_{s}}{2} R_{s}=\frac{2W}{\alpha+1}。因此，\alpha=0时，R_{s}=6000$Baud$;\alpha=0.5时$
$R_s=4000$Baud$;\alpha=1$时$,R_s=3000$Baud.
(2) $R_{b}= R_{s}\times \log _{2}M$ , $\alpha = 0$ 时 , $R_{b}= 6000$baud$\times4=24$Kbit/s ; $\alpha = 0. 5$时 ,
$R_b=16$Kbit/s$;\alpha=1$时$,R_b=12$Kbit/s
---------------------------------------\\
5.15 设计一 M 进制 PAM 系统 其输入的比特速率为14.4 Kbit/s 基带信道带宽 W=2400Hz ,
画出此最佳基带传输系统的框图。
++++++++++++++++++++++++++++++++++++++++++++\\
解：假设系统的滚降系数为$\alpha$,则
$$\begin{aligned}&W=\left(1+\alpha\right)\frac{R_{s}}{2}=\left(1+\alpha\right)\frac{R_{b}}{2\log_{2}M}\\&\frac{1+\alpha}{\log_{2}M}=\frac{2W}{R_{b}}=\frac{2\times2400}{14400}=\frac{1}{3}\\&\alpha=\frac{5}{18}\log_{2}M-1\end{aligned}$$
按$\alpha>0$设计，最小的M是16，相应的滚降系数是$\frac13$。系统框图如下
发送端框图：
\begin{figure}[h]
	\centering
	\includegraphics[width=\linewidth]{output/image_28_3.png}
\end{figure}
接收端框图：
\begin{figure}[h]
	\centering
	\includegraphics[width=\linewidth]{output/image_28_1.png}
\end{figure}
注：在通信的相关文献中，出于语言简洁的目的，我们经常会省略掉一些说明。此
时需注意对问题的理解必须是合理的理解。例如对于下面这个框图
\begin{figure}[h]
	\centering
	\includegraphics[width=\linewidth]{output/image_28_2.png}
\end{figure}
如果上下文中没有其他文字说明的话，合理的理解可以是：

15/28

---------------------------------------

(3)无论输入脉冲是什么形状，经过“脉冲成形”这个方框后，脉冲被改造成了另
外一种形状，改造后的脉冲的傅氏变换是根升余弦函数；
(4)若方框的输入是矩形脉冲，则方框的"矩形脉冲响应”的傅氏变换是根升余弦函数。矩形脉冲响应的定义可类比冲激响应的定义。本课 6.5 节也用到了这个概念。
---------------------------------------\\
5.16 题 16 图是一个数字基带传输系统，图中信道$C(f)$是理想低通，发送滤波$G_{_T}(f)$和接收滤波$G_{_R}(f)$满足$G_T(f)=G_R(f)=\sqrt{H_{\cos}(f)},H_{r\cos}(f)$是升余弦滚降特性。$n_w(t)$是功率谱密度为$\frac{N_0}2$的白高斯噪声，$\{a_n\}$独立等概取值于$\pm1$,A 点平均比特能量光$E_{b}$。

\begin{figure}[h]
	\centering
	\includegraphics[width=\linewidth]{output/image_29_1.png}
\end{figure}
(1)求最佳抽样时刻抽样值中的信号幅度值、瞬时信号功率值及噪声平均功率、抽样时刻
的信噪比；
(2)求系统的平均误比特率。

++++++++++++++++++++++++++++++++++++++++++++\\

解：(1)记$g_T(t)$、$g_R(t)$、$g(t)$分别为$G_T(f)$、$G_R(f),G_T(f)G_R(f)=H_{\mathrm{ros}}(f)$的

傅氏反变换，则 A 点接收信号可写为
$$r\left(t\right)=\sum_{n}a_{n}g_{T}\left(t-nT_{b}\right)+n_{w}\left(t\right)=s\left(t\right)+n_{w}\left(t\right)$$
A 点有用信号 $s(t)$的功率谱密度为
$$P_s\left(f\right)=\frac{\left|G_T\left(f\right)C\left(f\right)\right|^2}{T_b}=\frac{H_{\mathrm{rcos}}\left(f\right)}{T_b}$$
平均每比特的能量是
$$E_{b}=T_{b}\int_{-\infty}^{\infty}P_{s}\left(f\right)df=\int_{-\infty}^{\infty}H_{rcos}\left(f\right)df\:。$$
不妨考虑第 0 个取样及发送$a_0=1$的情况，由于
$$g\begin{pmatrix}0\end{pmatrix}=\int_{-\infty}^{\infty}H_{\mathrm{reos}}\begin{pmatrix}f\end{pmatrix}e^{j2\pi f\times0}df=E_{b}$$
$$y_0=a_0g\begin{pmatrix}0\end{pmatrix}+n=E_b+n$$
所以取样结果为

$n$ 是样值中的噪声分量。因此样值中有用信号的幅度是$E_{_b}$,瞬时信号功率是$E_b^2$,嗓声平均

功率是
$$E\bigg[n^{2}\bigg]=E\bigg[n^{2}\big(t\big)\bigg]=\int_{-\infty}^{\infty}\frac{N_{0}}{2}\bigg|G_{R}\big(f\big)\bigg|^{2}df=\frac{N_{0}}{2}\int_{-\infty}^{\infty}H_{\mathrm{reos}}\big(f\big)df=\frac{N_{0}E_{b}}{2}$$
抽样时刻信噪比为$\frac{E_{b}^{2}}{N_{0}E_{b}/2}=\frac{2E_{b}}{N_{0}}$


(2)由对称性可知最佳判决门限$V_T=0$,因而

16/28

---------------------------------------

$$P\left(e\mid a_0=1\right)=P\left(n<-E_b\right)=\frac{1}{2}\operatorname{erfc}\left(\frac{E_b}{\sqrt{2\times\frac{N_0E_b}{2}}}\right)=\frac{1}{2}\operatorname{erfc}\left(\sqrt{\frac{E_b}{N_0}}\right)$$
同理可得
$$P\bigl(e\:|\:a_0=-1\bigr)=P\bigl(e\:|\:a_0=1\bigr)=\frac{1}{2}\operatorname{erfc}\biggl(\sqrt{\frac{E_b}{N_0}}\biggr)$$
于是
$$P_b=P\bigl(\begin{matrix}a_0=1\end{matrix}\bigr)P\bigl(\begin{matrix}e\:|\:a_0=1\end{matrix}\bigr)+P\bigl(\begin{matrix}a_0=-1\end{matrix}\bigr)P\bigl(\begin{matrix}e\:|\:a_0=-1\end{matrix}\bigr)=\frac{1}{2}\operatorname{erfc}\biggl(\begin{matrix}\sqrt{\frac{E_b}{N_0}}\\\end{matrix}\biggr)$$
注：对于任意的二元星座点，高斯干扰，对称分布的情形，最佳门限是两个星座点
之中点。若星座点之间的欧氏距离为$d$ ,则错误率为
$\frac{1}{2}\operatorname{erfc}\biggl(\begin{matrix}\sqrt{\frac{d^{2}}{8\sigma^{2}}}\end{matrix}$\\
---------------------------------------\\
5.17 一个双极性矩形不归零 2PAM 信号受到功率谱的干扰，若发送正极性脉冲$s_1(t)$和负极性脉冲$s_2(t)$的概率分别是$P(s_1)=1/3$, $P(s_2)=2/3。求$
(1)最佳接收系统的最佳判决门限；
(2)平均误比特率。
++++++++++++++++++++++++++++++++++++++++++++\\
解：(1)设比特能量为$E_{_b}$,脉冲宽度为$T_{_b}$,则
$$\begin{aligned}&\text{,则}\\&s_{1}\left(t\right)=\begin{cases}\sqrt{\frac{E_{b}}{T_{b}}}&0\leq t<T_{b}\\0&\mathrm{else}\end{cases}\\&s_{2}\left(t\right)=-s_{1}\left(t\right)\end{aligned}$$
假设接收机的匹配滤波器对$s_1(t)$匹配，则其冲激响应为$h(t)=Ks_1(T_b-t)=Ks_1(t)$,最佳
取样时刻是$t=T_b$。取$K=1$,则$h(t)=s_1(t)$。

发$s_1$时，抽样值为
$$\begin{aligned}\text{y}&=\int_{0}^{T_{b}}s_{1}\left(\tau\right)\left[s_{1}\left(T_{b}-\tau\right)+n_{w}\left(T_{b}-\tau\right)\right]d\tau\\&=\int_{0}^{T_{b}}\sqrt{\frac{E_{b}}{T_{b}}}\left[\sqrt{\frac{E_{b}}{T_{b}}}+n_{w}\left(T_{b}-\tau\right)\right]d\tau\\&=E_{b}+\underbrace{\sqrt{\frac{E_{b}}{T_{b}}}\int_{0}^{T_{b}}n_{w}\left(T_{b}-\tau\right)d\tau}_{Z}=E_{b}+Z\end{aligned}$$
Z 为高斯随机变量，
$$\begin{array}{c}{\text{文量,}}\\{E\left[Z\right]=E\left[\int_{0}^{T_{b}}n_{w}\left(T_{b}-\tau\right)d\tau\right]=\int_{0}^{T_{b}}E\left[n_{w}\left(T_{b}-\tau\right)\right]d\tau=0\:,}\end{array}$$
17/28

---------------------------------------

$$\begin{aligned}E\left[Z^{2}\right]&=\frac{E_{b}}{T_{b}}E\left[\int_{0}^{T_{b}}n_{w}\left(T_{b}-\tau\right)d\tau\times\int_{0}^{T_{b}}n_{w}\left(T_{b}-\tau^{\prime}\right)d\tau^{\prime}\right]\\&=\int_{0}^{T_{b}}\int_{0}^{T_{b}}E\Big[n_{w}\left(T_{b}-\tau\right)n_{w}\left(T_{b}-\tau^{\prime}\right)\Big]d\tau d\tau^{\prime}\\&=\int_{0}^{T_{b}}\int_{0}^{T_{b}}\frac{N_{0}}{2}\:\delta\left(\tau^{\prime}-\tau\right)d\tau d\tau^{\prime}=\frac{N_{0}E_{b}}{2}\\&=\int_{0}^{T_{b}}\int_{0}^{T_{b}}\frac{N_{0}}{2}\:\delta\left(\tau^{\prime}-\tau\right)d\tau d\tau^{\prime}=\frac{N_{0}E_{b}}{2}\\\end{aligned}$$
所以发送$s_{_1}$条件下 y 的概率密度函数为 $p_1\left(y\right)=\frac1{\sqrt{\pi N_{0}E_{h}}}e^{\frac{\left(y-E_{b}\right)^{*}}{N_{0}E_{b}}}$

司理可得，发送$s_2$条件下 y 的概率密度函数为 $p_2\left(y\right)=\frac1{\sqrt{\pi E_{c}N}}e^{\frac{\left(y+E_{b}\right)}{N_{0}E_{b}}}$
最佳判决门限$V_T$是后验概率$P(s_1|y)$、$P(s_2|y)$相等的分界点，即
$P\Big(s_{1}|V_{T}\Big)=P\Big(s_{2}|V_{T}\Big)$,由贝叶斯准则，有$P\Big(s_1\Big)p\Big(V_{T}|s_{1}\Big)=P\Big(s_{2}\Big)p\Big(V_{T}|s_{2}\Big)$,即

$\frac{1}{3}\frac{1}{\sqrt{\pi N_{0}E_{h}}}e^{\frac{\left(V_{T}-E_{b}\right)^{2}}{N_{0}E_{b}}}=\frac{2}{3}\frac{1}{\sqrt{\pi N_{0}E_{h}}}e^{-\frac{\left(V_{T}+E_{b}\right)^{2}}{N_{0}E_{b}}}$
$\Rightarrow-\frac{\left(V_{T}-E_{b}\right)^{2}}{N_{0}E_{b}}=\ln2-\frac{\left(V_{T}+E_{b}\right)^{2}}{N_{0}E_{b}}$
$\Rightarrow2V_{T}E_{b}=N_{0}E_{b}\ln2-2V_{T}E_{b}$
$\Rightarrow V_T=\frac{N_0\ln2}4$

(2)


平均误码率为：

$P_{b}=P\Big(s_{1}\Big)P\Big(e\Big|s_{1}\Big)+P\Big(s_{2}\Big)P\Big(e\Big|s_{2}\Big)$
$=\frac{1}{6}$erfc$\left(\sqrt{\frac{E_{b}}{N_{0}}}\left(1-\frac{N_{0}\ln2}{4E_{b}}\right)\right)+\frac{1}{3}$erfc$\left(\sqrt{\frac{E_{b}}{N_{0}}}\left(1+\frac{N_{0}\ln2}{4E_{b}}\right)\right)$

18/28

---------------------------------------

注：匹配滤波器可以有任意的常系数$K$ ,不同$K$时所得到的最佳门限值也相差一个系数$K$ ,但最佳的误码率与$K$无关。$K$的最自然选择是$K=1$。另外一个常见的选择是通过$\kappa$的适当选择使得匹配滤波器的冲激响应的能量为 1 ,这种选择更符合信号空间描述中的习惯(见第6章)
---------------------------------------\\
5.18 某数字基带传输系统在抽样时刻的抽样值存在码间干扰，该系统的冲激响应$x(t)$的离散抽样值为$x_{-1}=0.3$,$x_0=0.9$,$x_1=0.3$。若该系统与三抽头迫零均衡器相级联，请求出此迫零均衡器的抽头器的抽头系数$\omega_{-1}$,$\omega_0$,$\omega_1$值，并计算出均衡前后的峰值畸变值。
++++++++++++++++++++++++++++++++++++++++++++\\
解：(1)包含均衡器后总的系统响应是
$\{h_k\}=\{x_k\}\otimes\{\omega_k\}$


$\otimes$代表卷积，即
$$h_k=\sum_{n=-\infty}^\infty\omega_{k-n}x_n=\omega_{k-1}x_1+\omega_kx_0+\omega_{k+1}x_{-1}=0.3\omega_{k-1}+0.9\omega_k+0.3\omega_{k+1}$$
三抽头迫零均衡将使
$$h_k=\begin{cases}1&k=0\\0&k=\pm1\end{cases}$$
由此得线性方程组
$$\begin{cases}0.9\omega_{-1}+0.3\omega_0=0\\0.3\omega_{-1}+0.9\omega_0+0.3\omega_1=1\\0.3\omega_0+0.9\omega_1=0\end{cases}$$
得解为
$$\begin{cases}\omega_{-1}=\omega_1=-\dfrac{10}{21}\\\\\omega_0=-3\omega_1=\dfrac{10}{7}\end{cases}$$
(2)计算均衡前后的峰值畸变
$\textcircled{1}$均衡前的峰值畸变
$$D=\dfrac{1}{x_0}\sum\limits_{k=-1\atop k\neq0}^1\bigl|x_k\bigr|=\dfrac{0.3+0.3}{0.9}=\dfrac{2}{3}$$
$\textcircled{2}$均衡后的总体响应为
$$h_{k}=0.3\omega_{k-1}+0.9\omega_{k}+0.3\omega_{k+1}=\begin{cases}1&k=0\\0&k=\pm1\\-\dfrac{1}{7}&k=\pm2\\0&\text{else}\end{cases}$$
$$D=\dfrac{1}{h_0}\sum\limits_{k=-2\atop k\neq0}^2\bigl|h_k\bigr|=\dfrac{1}{7}+\dfrac{1}{7}=\dfrac{2}{7}$$
均衡后的峰值畸变值为

---------------------------------------\\
5.19 二进制序列$\left\{b_n\right\}$通过加有预编码器的第一类部分响应系统，如题 19 图所示。

19/28





% 97-128
\includegraphics[width=0.5\textwidth]{images_ysy/1.png}

请写出以下的编码及相应电平、判决结果：

输入数据 $\{ b_n\}$ 1 0 0 1 0 1 1 0 0 1 $\cdots$

预编码器输出$\left\{d_n\right\}$

二电平序列$\{a_n\}$

抽样序列 $\{c_n\}$

判决输出 $\begin{Bmatrix}\hat{b}_n\end{Bmatrix}$
+++++++++++++++++++++++++++++++++++++++++++++++++++++++++++
解：

输入数据二电平序列
$$\begin{aligned}&\{b_{n}\}\quad1\quad0\quad0\quad1\quad0\quad1\quad1\quad0\quad0\quad1\quad\cdots\cdots\\&\text{出}\left\{d_{n}\right\}\quad0\:1\quad1\quad1\quad0\quad0\quad1\quad0\quad0\quad0\quad1\quad\cdots\cdots\\&\left\{a_{n}\right\}\quad-1+1\:+1\:+1\:-1\:-1\:+1\:-1\:-1\:-1\:-1\:+1\:\cdots\:\cdots\\&\left\{c_{n}\right\}\quad0\quad2\quad2\quad0\quad-2\quad0\quad0\quad-2\quad-2\quad0\quad\cdots\cdots\\&\left\{\hat{b}_{n}\right\}\quad1\quad0\quad0\quad1\quad0\quad1\quad1\quad0\quad0\quad1\quad\cdots\cdots\end{aligned}$$
抽样序列判决输出
---------------------------------------
5.20 设有数字基带信号如题 20 图示，此信号经过一个限带滤波器后成为限带信号，试画出 从此限带信号中提取符号同步信号的原理框图，并说明原理，并指出这种提取时钟的方法对 限带滤波器有什么要求。

\includegraphics[width=0.5\textwidth]{images_ysy/2.png}
+++++++++++++++++++++++++++++++++++++++++++++++++++++++++++
解：原理框图如下：

\includegraphics[width=0.5\textwidth]{images_ysy/3.png}

限带信号可写成

$$u\begin{pmatrix}t\end{pmatrix}=\sum_{n=-\infty}^{\infty}a_nx\begin{pmatrix}t-nT_s-t_0\end{pmatrix}$$
其中$x(t)$是限带后的脉冲波形，若限带滤波器的特性是$H(f)$,则$x(t)$的频谱为
$$X\begin{pmatrix}f\end{pmatrix}=AT_bH\begin{pmatrix}f\end{pmatrix}\text{sinc}(fT_s)$$
平方后的信号是
$$u^2\left(t\right)=\sum_n\sum_ma_na_mx\left(t-nT_s-t_0\right)x\left(t-mT_s-t_0\right)$$
假设$\left\{a_n\right\}$是取值于土1的独立等概序列，则平方器输出信号的数学期望是
$$\nu\big(t\big)=E\bigg[u^{2}\big(t\big)\bigg]=\sum_{n}\sum_{m}E\big[a_{n}a_{m}\big]x\big(t-nT_{s}-t_{0}\big)x\big(t-mT_{s}-t_{0}\big)=\sum_{m}x^{2}\big(t-mT_{s}-t_{0}\big)$$
令$g(t)=x^2\left(t\right)$,则$\nu(t)$ 等价于$\sum_m\delta\left(t-mT_s\right)$经过一个冲激响应为$g\left(t-t_0\right)$的线性系统的
输出。注意到
$$\sum_m\delta(t-mT_s)=\frac{1}{T_s}\sum_ne^{j2\pi\frac{n}{T_s}t}$$
设$G(f)$是$g(t)$ 的傅氏变换，由于$e^{j2\pi\frac nT_st}$通过$G(f)$后的输出是$G\left(\frac n{T_s}\right)e^{j2\pi\frac nT_st}$,所以
$$\nu\big(t\big)=\frac{1}{T_s}\sum_nG\bigg(\frac{n}{T_s}\bigg)e^{j2\pi\frac{n}{T_s}\big(t-t_0\big)}$$
因此，只要$G\left(\frac1{T_s}\right)\neq0,u^{2}\left(t\right)$中就包含有$f=\frac1{T_s}$的离散线谱分量，用窄带滤波器或锁相
环提出这个分量就获得了同步时钟。由$g\left(t\right)=x^2\left(t\right)$及频域卷积关系得
$$G\Bigg(\frac{1}{T_s}\Bigg)=\int_{-\infty}^{\infty}X\left(f\right)X\Bigg(\frac{1}{T_s}-f\Bigg)df$$
因此，为了能够用上面的框图提出时钟分量，需要的条件是积分$\int_{-\infty}^{\infty}X\left(f\right)X\left(\frac1{T_s}-f\right)df$ 不

为 0。使此积分不为零的必要条件是$X(f)$和$X\left(\frac1T-f\right.^{\prime}$必须有重叠，也就是要求$H(f)$ 和$H\left(\frac1{T_s}-f\right)$必须有重叠，也就是要求$H(f)$ 的频带宽度必须超过最小 Nyquist 带宽$(\frac{1}{2T_{s}})$。这虽然是一个必要条件，不过对于实际系统来说，$H(f)$的频带宽度超过最小Nvquist 带宽一般也是能够用上沭框佟提取时钟的充分条件。
---------------------------------------
5.21 设有二进制 PAM 信号 $s(t)=\sum_{n=-\infty}^na_ng\left(t-T_b\right)$,其中$\left\{a_n\right\}$是平稳独立序列，$a_n\in\left\{\pm1\right\}$,
$P\big(a_{n}=1\big)=p$ ,求$s\big(t\big)$的功率谱密度及功率。
+++++++++++++++++++++++++++++++++++++++++++++++++++++++++++
解：(1)s$(t)$可看成是$x(t)=\sum_{n=-\infty}^\infty a_n\delta(t-nT_b)$通过冲激响应为$g(t)$的线性系统的输出。故
先求 $x(t)$ 的功率谱密度。令$u( t) = x( t) - E[ x( t) ]$, $\nu ( t) = E[ x( t) ]$,则
$$\begin{aligned}&\nu\left(t\right)&&=E\Big[x\Big(t\Big)\Big]=\sum_{n=-\infty}^{\infty}E\Big[a_{n}\Big]\delta\Big(t-nT_{s}\Big)\\&&&=\left(2p-1\right)\sum_{n=-\infty}^{\infty}\delta\left(t-nT_{s}\right)\\&u\left(t\right)&&=x\left(t\right)-E\left[x\left(t\right)\right]=\sum_{n=-\infty}^{\infty}b_{n}\delta\left(t-nT_{b}\right)\end{aligned}$$
其中
$$b_n=\begin{cases}2\bigl(1-p\bigr)&\quad a_n=1\\-2p&\quad a_n=-1\end{cases}$$
序列$\left\{b_n\right\}$是独立平稳序列，其均值为 0,方差为
$$\sigma_b^2=E\bigg[b_n^2\bigg]=4\:p\big(1-p\big)$$
故此，$u(t)$的功率谱密度为
$$P_u\left(f\right)=\frac{\sigma_b^2}{T_b}=\frac{4p\left(1-p\right)}{T_b}$$
由于$\sum_{n=-\infty}^{\infty}\delta\left(t-nT_{b}\right)=\frac{1}{T_{b}}\sum_{m=-\infty}^{\infty}e^{j2\pi\frac{m}{T_{b}}t}$,所以周期信号$\nu(t)$的功率谱密度为
$$P_\nu\left(f\right)=\left(\frac{2p-1}{T_b}\right)^2\sum_{m=-\infty}^\infty\delta\left(f-\frac{m}{T_b}\right)$$
因此$s(t)$的功率谱密度为
$$\begin{aligned}P_{s}\left(f\right)&=P_{x}\left(f\right)\left|G\left(f\right)\right|^{2}\\&=\left[P_{u}\left(f\right)+P_{\nu}\left(f\right)\right]\left|G\left(f\right)\right|^{2}\\&=\frac{4\:p\left(1-p\right)}{T_{b}}\Big|G\left(f\right)\Big|^{2}+\left(\frac{2\:p-1}{T_{b}}\right)^{2}\sum_{m=-\infty}^{\infty}\left|G\left(\frac{m}{T_{b}}\right)\right|^{2}\:\delta\Bigg(f-\frac{m}{T_{b}}\Bigg)\end{aligned}$$
其功率为
$$\begin{aligned}\text{S}&=\int_{-\infty}^{+\infty}P_{s}(f)df\\&=\int_{-\infty}^{+\infty}\left[\frac{4\:p\left(1-p\right)}{T_{b}}\Big|G\left(f\right)\Big|^{2}+\left(\frac{2\:p-1}{T_{b}}\right)^{2}\sum_{m=-\infty}^{\infty}\left|G\left(\frac{m}{T_{b}}\right)\right|^{2}\delta\left(f-\frac{m}{T_{b}}\right)\right]df\\&=\frac{4\:p\left(1-p\right)}{T_{b}}E_{g}+\left(\frac{2\:p-1}{T_{b}}\right)^{2}\sum_{m=-\infty}^{\infty}\left|G\left(\frac{m}{T_{b}}\right)\right|^{2}\end{aligned}$$
解二 本题中$\{a_n\}$ 是平稳独立序列，因此也可以通过将式(5.2.18)代入式(5.2.12)求得功率谱。
这也就是套用公式(5.2.19)。在本题条件下，
$$m_a=E\begin{bmatrix}a_n\end{bmatrix}=p\begin{pmatrix}+1\end{pmatrix}+\begin{pmatrix}1-p\end{pmatrix}\begin{pmatrix}-1\end{pmatrix}=2p-1$$

\begin{equation}
\begin{aligned}
\sigma_{a}^{2} &= E\left[ a_{n}^{2} \right] - m_{a}^{2} = 1 - \left(1 - 2p\right)^{2} = 4p\left(1 - p\right) \\
\text{于是} \\
P_{a}(f) &= \sigma_{a}^{2} + \frac{m_{a}^{2}}{T_{b}} \sum_{m=-\infty}^{\infty} \delta\left(f - \frac{m}{T_{b}}\right) \\
&= 4p\left(1 - p\right) + \frac{\left(1 - 2p\right)^{2}}{T_{b}} \sum_{m=-\infty}^{\infty} \delta\left(f - \frac{m}{T_{b}}\right) \\
\text{因此} \\
P_{s}\left(f\right) &= \frac{1}{T_{b}} P_{a}\left(f\right) \left|G_{T}\left(f\right)\right|^{2} \\
&= \left[\frac{4p\left(1 - p\right)}{T_{b}} + \frac{\left(1 - 2p\right)^{2}}{T_{b}^{2}} \sum_{m=-\infty}^{\infty} \delta\left(f - \frac{m}{T_{b}}\right)\right] \left|G_{T}\left(f\right)\right|^{2} \\
&= \frac{4p\left(1 - p\right)}{T_{b}} \left|G_{T}\left(f\right)\right|^{2} + \frac{\left(1 - 2p\right)^{2}}{T_{b}^{2}} \sum_{m=-\infty}^{\infty} \left|G_{T}\left(\frac{m}{T_{b}}\right)\right|^{2} \delta\left(f - \frac{m}{T_{b}}\right) \\
\text{i} \text{其中} G_{T}\left(f\right) \text{是发送滤波器的频率响应,也即本题中} g(t) \text{的傅氏变换。} \\
\end{aligned}
\end{equation}

\includegraphics[width=0.5\textwidth]{images_ysy/4.png}
+++++++++++++++++++++++++++++++++++++++++++++++++++++++++++
解：由对图中的三角脉冲作傅氏变换即为总体传递函数
$$H\left(f\right)=T_{b}\mathrm{sinc}^{2}\left(fT_{b}\right)e^{-j2\pi fT_{b}}$$
记基带发送脉冲的傅氏变换为$G_T(f)$,匹配滤波器的传递函数为$G_R(f)$,则有
$$H\begin{pmatrix}f\end{pmatrix}=G_T\begin{pmatrix}f\end{pmatrix}C\begin{pmatrix}f\end{pmatrix}G_R\begin{pmatrix}f\end{pmatrix}=G_T\begin{pmatrix}f\end{pmatrix}G_R\begin{pmatrix}f\end{pmatrix}$$
由匹配滤波的关系，有
$$G_R\begin{pmatrix}f\end{pmatrix}=KG_T^*\begin{pmatrix}f\end{pmatrix}e^{-j2\pi ft_0}$$
其中$t_0$ 是最佳取样时刻，从图中可见$t_0=T_b$,因此有
$$H\left(f\right)=K\left|G_T\left(f\right)\right|^2e^{-j2\pi fT_b}=T_b\text{sinc}^2\left(fT_b\right)e^{-j2\pi fT_b}$$
所以
$$\left|G_T\left(f\right)\right|^2=\frac{T_b}{K}\text{sinc}^2\left(fT_b\right)$$
考虑到 $K$ 的任意性，取$K=\frac1T$。合理的解必须保证$g_T(t)$、$g_R(t)$都满足因果性，由此得

它是一个矩形脉冲
$$G_T\left(f\right)=\dfrac{1}{K}G_R\left(f\right)=T_b\text{sinc}\left(fT_b\right)e^{-j\pi fT_b}$$
$$g_T\left(t\right)=\begin{cases}1&0\leq t<T_b\\0&\text{else}\end{cases}$$
---------------------------------------
5.23 若二进制通信系统中接收端抽样判决处的判决量为$y=a+n$ ,其中数字信息$\alpha$等概取值于0或$A$, 噪声$n$服从指数分布，其概率密度函数是$f\left(n\right)=e^{-|n|}$。试求最佳判决门限$V_T$ 及平均比特错误概率$P_b\circ$
+++++++++++++++++++++++++++++++++++++++++++++++++++++++++++
解：发送 0 时$y$的条件概率密度为
$$p_1(y)=e^{-|y|}$$
$$p_2(y)=e^{-|y-A|}$$
\# 发送$A$时$y$的条件概率密度为


最佳门限是似然概率相等的分界点，即有
$$p_1\begin{pmatrix}V_T\end{pmatrix}=p_2\begin{pmatrix}V_T\end{pmatrix}$$
从而得
$$V_T=\frac{A}{2}$$
发送 0 而判错的概率是
$$P\big(e\:|\:0\big)=P\bigg(n>\dfrac{A}{2}\bigg)=\int_{A/2}^{\infty}e^{-n}dn=e^{-\dfrac{A}{2}}$$
发送$A$而判错的概率是
$$P\big(e\mid A\big)=P\bigg(n<-\frac{A}{2}\bigg)=\int_{-\infty}^{-A/2}e^ndn=e^{-\frac{A}{2}}$$
平均比特错误率为
$$P_b=\frac{P\left(e\mid0\right)+P\left(e\mid A\right)}{2}=e^{-\frac{A}{2}}$$
---------------------------------------
5.24 设信源是独立等概的二进制序列，求其 AMI 码的功率谱密度。
+++++++++++++++++++++++++++++++++++++++++++++++++++++++++++
解：记$\left\{b_n\right\}$为二进制信源，$b_n\in\{0,1\}$。记$\left\{a_n\right\}$为 AMI 编码结果，$\alpha_n\in\left\{+1,0,-1\right\}$。则 AMI
码波形可写为
$$s_\text{AMI}\left(t\right)=\sum_na_ng\left(t-nT_b\right)$$
其中
$$g\left(t\right)=\begin{cases}A&\left|t\right|\leq\dfrac{T_b}{2}\\0&\textit{else}\end{cases}$$
是半占空的 NZ 脉冲。

(1)求$\left\{a_n\right\}$的自相关函数$R_a\left(m\right)=E\left[a_na_{n+m}\right]$
首先，对于$m=0$

$$\begin{aligned}R_{a}\left(m\right)&=\sum_{a_{n}\in\{+1,-1\}}\sum_{a_{n+m}\in\{-a_{n},a_{n}\}}a_{n}a_{n+m}\Big[P\big(a_{n+m},I_{m-1}=0\:|\:a_{n}\big)+P\big(a_{n+m},I_{m-1}=1\:|\:a_{n}\big)\Big]P\big(a_{n}\big)\\&=\sum_{a_{n}\in\{+1,-1\}}\Big[-P\big(a_{n+m}=-a_{n},I_{m-1}=1\:|\:a_{n}\big)+P\big(a_{n+m}=a_{n},I_{m-1}=0\:|\:a_{n}\big)\Big]P\big(a_{n}\big)\\&P(z-z_{n}I_{n}-1|z_{n}-|1)\quad\text{定}\quad\text{ }I_{n-1}=\text{ 站 }\end{aligned}$$
$P(a_{n+m}=-a_n,I_{m-1}=1|a_n=\pm1)$实 际 就 是 在$b_n=1$的 条 件 下 , $b_{n+m}=1$且$\left\{b_{n+1},b_{n+2},\cdots,b_{n+m-1}\right\}$ 中有偶数个 1 的概率。由于$\left\{b_n\right\}$ 是独立等概序列，所以如下 3 个随机事件：“$b_n=1$”、“$b_{n+m}=1$”以及“$\left\{b_{n+1},b_{n+2},\cdots,b_{n+m-1}\right\}$中有偶数个1"是两两独立的， 因此
$$P\big(a_{n+m}=-a_n,I_{m-1}=1\:|\:a_n=\pm1\big)=P\big(b_{n+m}=1\big)P\big(I_{m-1}=1\big)=\frac{1}{2}P\big(I_{m-1}=1\big)$$
同理
$$P\big(a_{n+m}=a_n,I_{m-1}=0\mid a_n=\pm1\big)=P\big(b_{n+m}=1\big)P\big(I_{m-1}=0\big)=\frac{1}{2}P\big(I_{m-1}=0\big)$$
因此
$$\begin{aligned}R_{a}\left(m\right)&=\sum_{a_{n}\in\{+1,-1\}}\biggl[-\frac{1}{2}P\bigl(I_{m-1}=1\bigr)+\frac{1}{2}P\bigl(I_{m-1}=0\bigr)\biggr]P\bigl(a_{n}\bigr)\\&=\frac{P\left(I_{m-1}=0\right)-P\left(I_{m-1}=1\right)}{2}\biggl[1-P\bigl(a_{n}=0\bigr)\biggr]\\&=\frac{P\left(I_{m-1}=0\right)-P\left(I_{m-1}=1\right)}{4}\end{aligned}$$
若$m$是偶数，则连续$m-1$个独立等概二进制序列中出现偶数个 1 的概率为
$$P\left(I_{m-1}=1\right)=\frac{C_{m-1}^{0}+C_{m-1}^{2}+\cdots+C_{m-1}^{m-2}}{2^{m-1}}$$
出现奇数个 1 的概率是
$$P\big(I_{m-1}=0\big)=\frac{C_{m-1}^1+C_{m-1}^3+\cdots+C_{m-1}^{m-1}}{2^{m-1}}$$
利用关系$C_{m-1}^j=C_{m-1}^{m-1-j}$可得
$$P\big(I_{m-1}=0\big)=\frac{C_{m-1}^{m-2}+C_{m-1}^{m-4}+\cdots+C_{m-1}^{(m-1)-(m-1)}}{2^{m-1}}=P\big(I_{m-1}=1\big)$$

对于$m$ 是奇数的情形同理可得$P(I_{m-1}=0)=P(I_{m-1}=1)$。于是得到$R_a( m) = 0$ $( m> 1) $。

再利用自相关函数是偶函数这一点可得到序列$\left\{a_n\right\}$的自相关函数为
$$R_a\left(m\right)=\begin{cases}\frac{1}{2}&m=0\\-\frac{1}{4}&m=\pm1\\0&else\\\end{cases}$$
(2)求$u\left ( t\right ) = \sum _{n}a_{n}\delta \left ( t- nT_{b}\right ) \textbf{的 功 率 谱 密 度 }$ $u(t)$的自相关函数为

$$\begin{aligned}R_{u}\left(t,t+\tau\right)&=E\bigg[u\big(t\big)u\big(t+\tau\big)\bigg]\\&=E\bigg[\sum_{n}a_{n}\delta\big(t-nT_{b}\big)\times\sum_{j}a_{j}\delta\big(t+\tau-jT_{b}\big)\bigg]\\&=E\bigg[\sum_{n}\sum_{j}a_{n}a_{j}\delta\big(t-nT_{b}\big)\delta\big(t+\tau-jT_{b}\big)\bigg]\\&=\sum_{n}\sum_{j}R_{a}\left(n-j\right)\delta\left(t-nT_{b}\right)\delta\left(t+\tau-jT_{b}\right)\\\end{aligned}$$
代入前面得到的结果
$$R_{u}\left(t,t+\tau\right)=\sum_{n}\left[\frac{\delta\left(t+\tau-nT_{b}\right)}{2}-\frac{\delta\left(t+\tau-nT_{b}-T_{b}\right)}{4}-\frac{\delta\left(t+\tau-nT_{b}+T_{b}\right)}{4}\right]\delta\left(t-nT_{b}\right)$$
这个结果表明$u(t)$是循环平稳过程，求其平均自相关函数：
$$\begin{aligned}\overline{R_{u}}\left(\tau\right)&=\frac{1}{T_{b}}\int_{-T_{b}/2}^{T_{b}/2}R_{u}\left(t,t+\tau\right)dt\\&=\frac{1}{T_{b}}\int_{-T_{b}/2}^{T_{b}/2}\left[\frac{\delta\left(t+\tau\right)}{2}-\frac{\delta\left(t+\tau-T_{b}\right)}{4}-\frac{\delta\left(t+\tau+T_{b}\right)}{4}\right]\delta\left(t\right)dt\\&=\frac{1}{4T_{b}}\Big[2\delta\left(\tau\right)-\delta\left(\tau-T_{b}\right)-\delta\left(\tau+T_{b}\right)\Big]\end{aligned}$$
做傅氏变换得到$u(t)$的平均功率谱密度为
$$P_u\left(f\right)=\frac{1}{4T_b}\Big[2-e^{-j2\pi fT_b}-e^{j2\pi fT_b}\Big]=\frac{1-\cos2\pi fT_b}{2T_b}=\frac{\sin^2\pi fT_b}{T_b}$$
(3)求 AMI 信号的功率谱密度
$s_\text{AMI}{(t)}$可以看成是$u(t)$通过一个冲激响应为$g(t)$的线性系统的输出。由于$g(t)$ 的傅氏变

换是
$$G\big(f\big)=\frac{AT_b}{2}\text{sinc}\bigg(\frac{fT_b}{2}\bigg)$$
所以 AMI 信号的功率谱为
$$P_{_{AMI}}\left(f\right)=\frac{A^2T_b}{4}\sin^2\left(\frac{fT_b}{2}\right)\sin^2\pi fT_b$$

解二 本题在得到 AMI 码序列$\left\{a_{_m}\right\}$的自相关函数$R_{_a}\left(m\right)$后，可求出$\left\{a_{_m}\right\}$的功率谱$P_a\left(f\right)$
它是对$R_a(m)$的离散时间傅氏变换：
$$\begin{aligned}P_{a}\left(f\right)&=\sum_{m=-\infty}^{\infty}R_{a}\left(m\right)e^{-j2\pi fmT_{b}}\\&=\frac{1}{2}-\frac{1}{4}e^{-j2\pi fT_{b}}-\frac{1}{4}e^{-j2\pi fT_{b}}\\&=\frac{1}{2}\big(1-\cos2\pi fT_{b}\big)=\sin^{2}\big(\pi fT_{b}\big)\end{aligned}$$
于是，半占空归零脉冲的 AMI(RZ)码的平均功率谱为

$$\begin{aligned}
P_{s}\left(f\right)& =\frac{1}{T_b}P_a\big(f\big)\big|G_T\big(f\big)\big|^2 \\
&=\frac1{T_b}\sin^2\left(\pi fT_b\right)\Bigg(\frac{AT_b}2\Bigg)^2\sin^2\Bigg(\frac{fT_b}2\Bigg) \\
&=\frac{A^2T_b}4\sin^2\left(\pi fT_b\right)\text{sinc}^2\left(\frac{fT_b}2\right)
\end{aligned}$$
---------------------------------------
第六章 数字信号的频带传输
---------------------------------------
6.1 在题 6.1 图中，A点信号是幅度为 1 的单极性不归零码，二进制序列独立等概，速率为$R_{b}=1Mbit/s_{\text{,B点信号是OOK，载波频率是}}f_{c}=100\mathrm{MHz}_{\text{。请给出A、B两点的功率谱}}$ 密度，并画出双边带功率谱密度图。

\includegraphics[width=0.5\textwidth]{images_ysy/5.png}
+++++++++++++++++++++++++++++++++++++++++++++++++++++++++++
解：A点信号等价于幅度为$\overset{\pm\frac12}{\operatorname*{=}}\frac12_{\text{的双极性不归零信号叠加了一个幅度为}}\overset{\large{\frac12}}{\operatorname*{=}}$的直流，因此该点
的功率谱密度为
$$P_{A}\left(f\right)=\frac{1}{4T_{b}}\Big|T_{b}\operatorname{sinc}\Big(fT_{b}\Big)\Big|^{2}+\frac{1}{4}\delta\Big(f\Big)=\frac{R_{b}}{4}\operatorname{sinc}^{2}\Bigg(\frac{f}{R_{b}}\Bigg)+\frac{\delta\Big(f\Big)}{4}$$

\includegraphics[width=0.5\textwidth]{images_ysy/6.png}

$$\begin{aligned}
P_{B}\left(f\right)& =A^2\times\frac{P_A\left(f-f_c\right)+P_A\left(f+f_c\right)}{4} \\
&=\frac{A^2R_b}{16}\Bigg[\mathrm{sinc}^2\Bigg(\frac{f-f_c}{R_b}\Bigg)+\mathrm{sinc}^2\Bigg(\frac{f+f_c}{R_b}\Bigg)\Bigg]+\frac{A^2}{16}\Bigg[\delta\left(f-f_c\right)+\delta\left(f-f_c\right)\Bigg]
\end{aligned}$$

\includegraphics[width=0.5\textwidth]{images_ysy/7.png}
---------------------------------------
6.2 已知二进制OOK数字通信系统中发送的二元信号是$^{s_1}(t)=A\cos2\pi f_ct$、$^{s_2(t)=0}$,持续时间为$^{0\leq t<T_b}$。OOK信号传输中受到功率密度为$N_0/2\text{ 的加性高斯白噪声}^{n(t)}$的干扰，

接收信号为

$$r(t)=s_i(t)+n(t)\:,\quad i=1,2$$
(1)请分别画出最佳相干接收及非相干接收框图；
(2)请分别推导出它们的平均误比特率(设$^{S_1(t)}$与$^{S_2(t)}$等概出现，$f_c>>1/T_b)$。

+++++++++++++++++++++++++++++++++++++++++++++++++++++++++++
解：(1)最佳相干接收机框图如下

\includegraphics[width=0.5\textwidth]{images_ysy/8.png}

最佳非相干接收机框图如下

\includegraphics[width=0.5\textwidth]{images_ysy/9.png}


(2)最佳相干：

发送$s_1(t)_\text{时抽样值}$
$$y=\int_{0}^{T_{b}}r\left(t\right)\cos2\pi f_{c}tdt=\frac{AT_{b}}{2}+\underbrace{\int_{0}^{T_{b}}n\left(t\right)\cos2\pi f_{c}tdt}_{z}=\frac{AT_{b}}{2}+Z$$
$Z$是 0 均值高斯随机变量，其方差为
$$\begin{aligned}\text{C}&=E\left[\int_{0}^{T_{b}}n\left(t\right)\cos2\pi f_{c}tdt\times\int_{0}^{T_{b}}n\left(t\right)\cos2\pi f_{c}tdt\right]\\&=\int_{0}^{T_{b}}\int_{0}^{T_{b}}E\Big[n\Big(t\Big)n\Big(t'\Big)\Big]\cos2\pi f_{c}t\cos2\pi f_{c}t'dtdt'\\&=\int_{0}^{T_{b}}\int_{0}^{T_{b}}\frac{N_{0}}{2}\delta\Big(t'-t\Big)\cos2\pi f_{c}t\cos2\pi f_{c}t'dtdt'\\&=\int_{0}^{T_{b}}\frac{N_{0}}{2}\cos^{2}2\pi f_{c}tdt=\frac{N_{0}T_{b}}{4}\end{aligned}$$
因此$^Z\sim \mathcal{N} \left ( 0, \sigma ^2\right ) _\text{, 故 }$

发送$s_2(t)_\text{时，}$
$$p_1\left(y\right)=p\left(y\mid s_1\right)=\sqrt{\frac{2}{N_0T_b}}e^{-\frac{2\left(y-\frac{AT_b}{2}\right)^2}{N_0T_b}}$$
$$y=\int_0^{T_b}r\left(t\right)\cos2\pi f_ctdt=\underbrace{\int_0^{T_b}n\left(t\right)\cos2\pi f_ctdt}_{Z}=Z$$
$$p_2\left(y\right)=p\left(y\mid s_2\right)=\sqrt{\frac{2}{N_0T_b}}e^{-\frac{2y^2}{N_0T_b}}$$
同理可得

$$P\Big(e\mid s_{2}\Big)=P\Bigg(Z>\frac{AT_{b}}{4}\Bigg)=P\Bigg(Z<-\frac{AT_{b}}{4}\Bigg)=\frac{1}{2}\mathrm{erfc}\Bigg(\sqrt{\frac{A^{2}T_{b}}{8N_{0}}}\Bigg)=P\Big(e\mid s_{1}\Big)$$
故平均误比特率为

$s_{1}(t)_{\text{与}}s_{2}(t)_{\text{等概出现，故最佳门限}}V_{T\text{ 是}}p_{1}(y)=p_{2}(y)_{\text{的解，可得}}V_{T}=\frac{AT_{b}}{4}$
$P\Big(e\mid s_{1}\Big)=P\Big(\frac{AT_{b}}{2}+Z<\frac{AT_{b}}{4}\Big)=P\Big(Z<-\frac{AT_{b}}{4}\Big)$
$=\frac{1}{2}$erfc$\left(\frac{\frac{AT_b}4}{\sqrt{2\times\frac{N_0T_b}4}}\right)=\frac12$erfc$\left(\sqrt{\frac{A^2T_b}{8N_0}}\right)$
$$P_e=\frac{1}{2}\operatorname{erfc}\Bigg(\sqrt{\frac{A^2T_b}{8N_0}}\Bigg)=\frac{1}{2}\operatorname{erfc}\Bigg(\sqrt{\frac{E_b}{2N_0}}\Bigg)$$
最佳非相干：
$s_1(t)_\text{的复包络是}$
$$s_{1,L}\left(t\right)=\begin{cases}A&0\leq t\leq T_b\\0&else\end{cases}$$
设带通滤波器的冲激响应是$^{h(t)}$,其复包络是$^{h_L(t)}$,则带通滤波器的等效基带冲激响应是
$$\begin{aligned}&\text{小冲做响应足}\\&h_{eq}\left(t\right)=\frac{1}{2}h_{L}\left(t\right)=\frac{1}{2}s_{1,L}\left(T_{b}-t\right)e^{j\theta}=\begin{cases}\dfrac{1}{2}\:Ae^{j\theta}&0\leq t\leq T_{b}\\[2ex]0&else\end{cases}\end{aligned}$$
因此
$$h\begin{pmatrix}t\end{pmatrix}=\mathrm{Re}\left\{h_{L}\begin{pmatrix}t\end{pmatrix}e^{j2\pi f_{c}t}\right\}=\begin{cases}A\cos\bigl(2\pi f_{c}t+\theta\bigr)&0\leq t\leq T_{b}\\0&else\end{cases}$$
其中$\theta_\text{ 是体现非相干的一个任意相移。}s_1(t)_{\text{通过带通滤波器后的复包络是}}$
$$\int_{-\infty}^{\infty}h_{eq}\left(\tau\right)s_{1,L}\left(t-\tau\right)d\tau=\int_{0}^{T_{b}}\frac{Ae^{j\theta}}{2}s_{1,L}\left(t-\tau\right)d\tau $$

在最佳取样时刻的输出是
$$\int_0^{T_b}\frac{Ae^{j\theta}}{2}s_{1,L}\left(T_b-\tau\right)d\tau=\frac{Ae^{j\theta}}{2}\int_0^{T_b}Ad\tau=\frac{A^2T_b}{2}e^{j\theta}=E_1e^{j\theta}=2E_be^{j\theta}$$
白高斯噪声通过带通滤波器的输出是窄带高斯噪声，其复包络为
$$n_L\begin{pmatrix}t\end{pmatrix}=n_c\begin{pmatrix}t\end{pmatrix}+jn_s\begin{pmatrix}t\end{pmatrix}$$
其方差为
$$\sigma^{2}=\int_{-\infty}^{\infty}\frac{N_{0}}{2}\Big|H\left(f\right)\Big|df=\frac{N_{0}}{2}\int_{-\infty}^{\infty}h^{2}\left(t\right)dt=\frac{A^{2}T_{b}N_{0}}{4}=\frac{N_{0}E_{1}}{2}=N_{0}E_{b}$$
包络检波器在采样点的输出是
$$\begin{aligned}&y=\left|2E_{b}e^{j\theta}+n_{c}\left(T_{b}\right)+jn_{s}\left(T_{b}\right)\right|=\left|2E_{b}+\left[n_{c}\left(T_{b}\right)+jn_{s}\left(T_{b}\right)\right]e^{-j\theta}\right|\\&=\left|2E_{b}+n_{c}^{\prime}\left(T_{b}\right)+jn_{s}^{\prime}\left(T_{b}\right)\right|\end{aligned}$$

$$\text{其中}n_c^{\prime}\begin{pmatrix}t\end{pmatrix}_\text{和}n_s^{\prime}\begin{pmatrix}t\end{pmatrix}_\text{的方差是}\sigma^2=\frac{A^2T_bN_0}4=N_0E_b\text{。在大信噪比条件下,}$$
$$y=\left|2E_{b}+n_{c}^{\prime}\left(T_{b}\right)+jn_{s}^{\prime}\left(T_{b}\right)\right|\approx2E_{b}+n_{c}^{\prime}\left(T_{b}\right)\\p_{1}\left(y\right)=p\left(y\mid s_{1}\right)=\frac{1}{\sqrt{2\pi N_{0}E_{b}}}e^{\frac{\left(y-2E_{b}\right)^{2}}{2N_{0}E_{b}}}$$
发送$^{S_2}(t)_{\text{时采样点的输出是}}$
$$y=\begin{vmatrix}n_c\begin{pmatrix}T_b\end{pmatrix}+jn_s\begin{pmatrix}T_b\end{pmatrix}\end{vmatrix}$$
$$p_2\left(y\right)=p\left(y\mid s_2\right)=\frac{y}{\sigma^2}e^{-\frac{y^2}{2\sigma^2}}=\frac{y}{N_0E_b}e^{-\frac{y^2}{2N_0E_b}}$$
$$V_{T}=\frac{2E_{b}}{2}=E_{b}$$
$s_1(t)_\text{与}{s_2}(t)_\text{等概出现，故最佳门限}{V_T\text{近似为}}$
$$P\left(e\mid s_{1}\right)=P\left(n_{c}^{\prime}\left(T_{b}\right)<-E_{b}\right)=\frac{1}{2}\operatorname{erfc}\left(\frac{E_{b}}{\sqrt{2N_{0}E_{b}}}\right)=\frac{1}{2}\operatorname{erfc}\left(\sqrt{\frac{E_{b}}{2N_{0}}}\right)$$
$$P\big(e\:|\:s_{2}\big)=\int_{E_{b}}^{\infty}p_{2}\big(y\big)dy=\int_{E_{b}}^{\infty}\frac{y}{N_{0}E_{b}}e^{-\frac{y^{2}}{2N_{0}E_{b}}}dy=\int_{\frac{E_{b}}{2N_{0}}}^{\infty}e^{-t}dt=e^{-\frac{E_{b}}{2N_{0}}}$$
故平均误比特率为
$$P_{e}=\frac{1}{4}\operatorname{erfc}\Bigg(\sqrt{\frac{E_{b}}{2N_{0}}}\Bigg)+\frac{1}{2}\operatorname{exp}\Bigg(-\frac{E_{b}}{2N_{0}}\Bigg)\approx\frac{1}{2}\operatorname{exp}\Bigg(-\frac{E_{b}}{2N_{0}}\Bigg)$$
---------------------------------------
6.3 已知 2FSK 系统的两个信号波形为
$$s_{i}\left(t\right)=\sin\left(2i\pi t\right)\:,\:i=1,2\:,\:0\leq t<T_{b}$$
其中$T_b= 1s$ , $s_1( t)$与$s_2(t)_\text{等概出现。}$
(1)画出两信号的波形图，求出平均比特能量$^{E_b}$及两信号波形的相关系数$\rho_i$
2)若 2FSK信号在信道传输中受到功率谱密度为$\frac{N_0}2$加性白高斯噪声的干扰 请画出最佳轻
收框图，并推导其误比特率。
+++++++++++++++++++++++++++++++++++++++++++++++++++++++++++
解：(1)两信号波形如下

\includegraphics[width=0.5\textwidth]{images_ysy/10.png}

$s_{\mathrm{l}}(t$

$s_2\left(t\right)_{\text{的能量分别为}}$

$$E_i=\int_0^1\bigl(\sin2\pi it\bigr)^2\:dt=\frac{1}{2}\:,\:i=1,2$$
平均比特能量为$E_b=\frac{E_1+E_2}2=\frac12$。

相关系数为
$$\rho_{12}=\dfrac{1}{E_b}\int_0^{T_b}s_1\big(t\big)s_2\big(t\big)dt=2\int_0^1\big(\sin2\pi t\big)\big(\sin4\pi t\big)dt=0$$
(2)在白高斯噪声干扰下的最佳接收框图如下

\includegraphics[width=0.5\textwidth]{images_ysy/11.png}
图中的判决量 \\
$$\begin{aligned}
&\text{图中的判决量} \\
&\text{1}&& \nu=\int_0^1r\left(t\right)\left(\sin2\pi t-\sin4\pi t\right)dt \\
&&&=\int_0^1\left(\sin2\pi it\right)\left(\sin2\pi t-\sin4\pi t\right)dt+\underbrace{\int_0^1n_w\left(t\right)\left(\sin2\pi t-\sin4\pi t\right)dt}_Z \\
&&&=\frac{a_i}2+Z \\
&a_i=\\\text{其中}&& \begin{cases}1&i=1\\-1&i=2\end{cases},Z\sim\mathcal{N}\left(0,\sigma^2\right)_\text{是均值为 }0\text{ 的高斯随机变量,其方差为} \\
&&&\sigma^2=E\left[\left(\int_0^1n_w\left(t\right)\left(\sin2\pi t-\sin4\pi t\right)dt\right)^2\right] \\
&&&=E\left[\int_0^1\int_0^1n_w\left(t\right)\left(\sin2\pi t-\sin4\pi t\right)n_w\left(t^\prime\right)\left(\sin2\pi t^\prime-\sin4\pi t^\prime\right)dt^\prime dt\right] \\
&&&=\int_0^1\int_0^1E\left[n_w\left(t\right)n_w\left(t^{\prime}\right)\right](\sin2\pi t-\sin4\pi t)(\sin2\pi t^{\prime}-\sin4\pi t^{\prime})dt^{\prime}dt \\
&&&=\int_0^1\int_0^1\frac{N_0}2\delta\left(t-t'\right)\left(\sin2\pi t-\sin4\pi t\right)\left(\sin2\pi t'-\sin4\pi t'\right)dt'dt \\
&&&=\frac{N_0}2\int_0^1\left(\sin2\pi t-\sin4\pi t\right)^2dt=\frac{N_0}2 \\
&&&\begin{pmatrix}t\end{pmatrix}\text{、}s_2\begin{pmatrix}t\end{pmatrix}_\text{时判决量}y\text{的概率密度函数分别为}
\end{aligned}$$

$$\begin{cases}p_1\left(y\right)=\dfrac{1}{\sqrt{\pi N_0}}e^{-\frac{\left(2y-1\right)^2}{4N_0}}\\\\p_2\left(y\right)=\dfrac{1}{\sqrt{\pi N_0}}e^{-\frac{\left(2y+1\right)^2}{4N_0}}\end{cases}$$
最佳判决门限$V_{T\text{满足}}p_{1}(V_{T})=p_{2}(V_{T})$,故有$V_{T}=0$。此时
$$P\big(e\mid s_{1}\big)=P\bigg(Z<-\frac{1}{2}\bigg)=\frac{1}{2}\mathrm{erfc}\bigg(\frac{1/2}{\sqrt{2\sigma^{2}}}\bigg)=\frac{1}{2}\mathrm{erfc}\bigg(\frac{1}{\sqrt{4N_{0}}}\bigg)\\P\big(e\mid s_{2}\big)=P\bigg(Z>\frac{1}{2}\bigg)=P\bigg(Z<-\frac{1}{2}\bigg)=\frac{1}{2}\mathrm{erfc}\bigg(\frac{1}{\sqrt{4N_{0}}}\bigg)$$
因此平均误比特率是
$$P_b=\dfrac{1}{2}\text{erfc}\Bigg(\dfrac{1}{\sqrt{4N_0}}\Bigg)$$
---------------------------------------
6.4 在题 6.4 图中，A点信号是幅度为 1 的双极性不归零信号，二进制序列独立等概，速率为$R_{b}=1$Mbit/s$_\text{,B点信号是BPSK，载波频率是}f_{c}=100$MHz。请给出A、B两点的功率谱密度。并画出功率谱密度图。

\includegraphics[width=0.5\textwidth]{images_ysy/12.png}

+++++++++++++++++++++++++++++++++++++++++++++++++++++++++++
解：A 点信号的功率谱密度为
$$P_{A}\left(f\right)=\frac{1}{T_{b}}\Big|T_{b}\operatorname{sinc}\big(fT_{b}\big)\Big|^{2}=R_{b}\operatorname{sinc}^{2}\Bigg(\frac{f}{R_{b}}\Bigg)$$

\includegraphics[width=0.5\textwidth]{images_ysy/13.png}

$$\begin{aligned}&\text{B 点 BPSK 信号的功率谱为}\\&P_{B}\left(f\right)=A^{2}\times\frac{P_{A}\left(f-f_{c}\right)+P_{A}\left(f+f_{c}\right)}{4}=\frac{A^{2}R_{b}}{4}\left[\mathrm{sinc}^{2}\left(\frac{f-f_{c}}{R_{b}}\right)+\mathrm{sinc}^{2}\left(\frac{f+f_{c}}{R_{b}}\right)\right]\end{aligned}$$

\includegraphics[width=0.5\textwidth]{images_ysy/14.png}
---------------------------------------
6.5 已知BPSK系统的两个信号波形为$^{s_1}(t)=A\cos\left(2\pi f_ct\right)_\text{和}{s_2}(t)=-s_1(t)$,持续时间是

$0\leq t< T_b$ , $s_1( t)$ 与 $s_2( t)$等概出现。

(1)求平均比特能量$E_{b\text{ 及两信号波形的相关系数}}\rho_{i}$

(2)若BPSK信号在信道传输中受到功率谱密度为$N_0/2$ 加性白高斯噪声的干扰，请画出最

佳接收框图，并推导其误比特率。
+++++++++++++++++++++++++++++++++++++++++++++++++++++++++++
解：(1)平均比特能量为$E_b=\frac{A^2T_b}2$,相关系数为
$$=\frac{1}{E_{b}}\int_{0}^{T_{b}}s_{1}\left(t\right)s_{2}\left(t\right)dt=\frac{1}{E_{b}}\int_{0}^{T_{b}}-s_{1}^{2}\left(t\right)dt=-1$$
(2)在白高斯噪声干扰下的最佳接收框图如下

\includegraphics[width=0.5\textwidth]{images_ysy/15.png}


$$\begin{aligned}
&s_{i}\left(t\right)=a_{i}A\cos2\pi f_{c}t_{\text{,其中}}a_{i}& =\begin{cases}1&i=1\\-1&i=2\end{cases} \\
&&y=\int_{0}^{T_{b}}\frac{2}{A}r\left(t\right)\cos2\pi f_{c}tdt=2\int_{0}^{T_{b}}a_{i}\left(\cos2\pi f_{c}t\right)^{2}dt+\underbrace{\frac{2}{A}\int_{0}^{T_{b}}n_{w}\left(t\right)\left(\cos2\pi f_{c}t\right)dt}_{Z} \\
=a_{i}T_{b}+Z \\
Z是均值为0的高&& \text{斯随机变量,其方差为} \\
\sigma^{2}&& =\frac{4}{A^{2}}E\left[\left(\int_{0}^{T_{b}}n_{w}\left(t\right)\cos^{2}\left(2\pi f_{c}t\right)dt\right)^{2}\right] \\
&&=\frac{4}{A^{2}}E\biggl[\int_{0}^{T_{b}}\int_{0}^{T_{b}}n_{w}\bigl(t\bigr)\bigl(\cos2\pi f_{c}t\bigr)n_{w}\bigl(t'\bigr)\bigl(\cos2\pi f_{c}t'\bigr)dt'dt\biggr] \\
&&=\frac{4}{A^{2}}\int_{0}^{T_{b}}\int_{0}^{T_{b}}E\Big[n_{w}\big(t\big)n_{w}\big(t'\big)\Big]\cos2\pi f_{c}t\cos2\pi f_{c}t'dt'dt \\
&&=\frac{4}{A^{2}}\int_{0}^{T_{b}}\int_{0}^{T_{b}}\frac{N_{0}}{2}\delta\left(t-t^{\prime}\right)\cos2\pi f_{c}t\cos2\pi f_{c}t^{\prime}dt^{\prime}dt \\
&&=\frac{2N_{0}}{A^{2}}\int_{0}^{T_{b}}\left(\cos2\pi f_{c}t\right)^{2}dt=\frac{N_{0}T_{b}}{A^{2}}
\end{aligned}$$

$$\begin{gathered}
\text{因此发送}s_1\begin{pmatrix}t\end{pmatrix}\text{、}s_2\begin{pmatrix}t\end{pmatrix}\text{时判决量}y\text{的概率密度函数分别为} \\
\begin{cases}p_1\left(y\right)=\frac{A}{\sqrt{2\pi N_0T_b}}e^{-\frac{A^2\left(y-T_b\right)^2}{2N_0T_b}}\\\\p_2\left(y\right)=\frac{A}{\sqrt{2\pi N_0T_b}}e^{-\frac{A^2\left(y+T_b\right)^2}{2N_0T_b}}\end{cases} \\
_\text{最佳判决门限}V_T\text{ 满足 }p_1\left(V_T\right)=p_2\left(V_T\right)_\text{,故有}V_T=0_\text{。此时} \\
P\big(e | s_1\big)=P\big(Z<-T_b\big)=\frac{1}{2}\operatorname{erfc}\bigg(\frac{T_b}{\sqrt{2\sigma^2}}\bigg)=\frac{1}{2}\operatorname{erfc}\bigg(\sqrt{\frac{A^2T_b}{2N_0}}\bigg)=\frac{1}{2}\operatorname{erfc}\bigg(\sqrt{\frac{E_b}{N_0}}\bigg) \\
P\left(e\mid s_{2}\right)=P\left(Z>T_{b}\right)=P\left(e\mid s_{1}\right) \\
因此平均误比特率是 \\
P_{b}=\frac{1}{2}\mathrm{erfc}\Bigg(\sqrt{\frac{A^{2}T_{b}}{2N_{0}}}\Bigg)=\frac{1}{2}\mathrm{erfc}\Bigg(\sqrt{\frac{E_{b}}{N_{0}}}\Bigg) 
\end{gathered}$$
---------------------------------------
6.6 在题 6.6 图中，BPSK信号是由双极性不归零信号调制正弦载波而成，理想带通滤波器BPF的带宽恰好能使BPSK信号的主瓣通过。恢复载波和发送载波之间有固定的相位偏差$\theta$, LPF用干滤除二倍载频分量。取样时刻是码元的中点。假设BPF对BPSK信号引起的失真可以忽略，求误比特率。

\includegraphics[width=0.5\textwidth]{images_ysy/16.png}
+++++++++++++++++++++++++++++++++++++++++++++++++++++++++++
解：不妨以本地载波为参考，设发送载波为$^{\cos\left(2\pi f_ct+\theta\right)}$,本地恢复载波为$^{\cos2\pi f_ct}$。BPF输出的BPSK信号可写成$\operatorname{Re}\left\{\pm Ae^{j\theta}e^{j2\pi f_ct}\right\}_{\text{,其中}A\text{是BPSK信号的幅度，“土”对应发送}}$ $“1”、“0”$。
BPF输出的噪声$^{n(t)}$可以表示成$^n(t)=\operatorname{Re}\left\{\tilde{n}(t)e^{j2\pi f_ct}\right\}$,其中$\tilde{n}\left(t\right)=n_c\left(t\right)+jn_s\left(t\right)$, $n(t)、n_c(t)、n_s(t)_\text{都是 0 均值的平稳高斯过程，方差都是}\sigma^2=N_0B_\text{,其中}{N_0\text{ 是白噪声}}$ 的单边功率普密度，$B$是BPF的等效噪声带宽。
BPF 的总输出是
$$r\begin{pmatrix}t\end{pmatrix}=\mathrm{Re}\left\{\begin{bmatrix}\pm Ae^{j\theta}+n_c\begin{pmatrix}t\end{pmatrix}+jn_s\begin{pmatrix}t\end{pmatrix}\end{bmatrix}e^{j2\pi f_ct}\right\}$$
用载波$\cos2\pi f_ct$ 解调得到的是
$$\mathrm{Re}\left\{\tilde{r}\left(t\right)\right\}=\mathrm{Re}\left\{\pm Ae^{j\theta}+n_{c}\left(t\right)+jn_{s}\left(t\right)\right\}=\pm A\cos\theta+n_{c}\left(t\right)$$
采样结果是$y= \pm A+ \xi _\text{, }\xi _\text{ 是 对 }{n_c}( t) _\text{的 采 样 结 果 , }\xi \sim \mathcal{N} \left ( 0, N_0B\right ) $。
发$+A$ 而错的概率是

$$P\big(\xi<-A\cos\theta\big)=\dfrac{1}{2}\text{erfc}\bigg(\dfrac{A\cos\theta}{\sqrt{2\sigma^2}}\bigg)=\dfrac{1}{2}\text{erfc}\bigg(\sqrt{\dfrac{A^2\cos^2\theta}{2N_0B}}\bigg)$$
发$-A$ 而错的概率是
$$P\big(\xi>A\cos\theta\big)=\dfrac{1}{2}\text{erfc}\bigg(\sqrt{\dfrac{A^2\cos^2\theta}{2N_0B}}\bigg)$$
故此平均错误率是
$$P_b=\dfrac{1}{2}\operatorname{erfc}\Bigg(\sqrt{\dfrac{A^2\cos^2\theta}{2N_0B}}\Bigg)$$
$$s\left(t\right)=\sum_{n=-\infty}^{\infty}a_{n}g\left(t-nT_{b}\right)\cos\left(2\pi f_{c}t+\theta\right),$$
---------------------------------------
6.7 速率为$^{R_b\text{ 的BPSK信号为}}$
其中$T_b=1/R_b$是比特间隔，
$$a_n\in\begin{Bmatrix}+1,-1\end{Bmatrix}\text{代表信息比特,载频}$$
$$f_c>>\frac{1}{T_b} , g\big(t\big)=\begin{cases} A&0\leq t<T_b\\ 0&else\end{cases}$$
,
$$\left\{a_n\right\}$$
是独立等概序列，$\theta_\text{ 是任意的一个载波相位。}$

立等概序列，$\theta_\text{是任意的一个载波相位。}$

$(1)\text{求}^{s\left(t\right)}_{\text{的功率谱密度，并画图表示；}}$

(2)请画出 BPSK 接收机提取载波的两种方框图，并说明原理。
$$b(t)=\sum_{n=-\infty}^{\infty}a_{n}g\left(t-nT_{b}\right)$$
+++++++++++++++++++++++++++++++++++++++++++++++++++++++++++
解：(1)$s(t)_\text{是数字基带信号}$
$)_{\text{对载波}}\cos\left(2\pi f_{c}t+\theta\right)_{\text{的DSB调制。}}$


$b(t)_{\text{ 的功率谱密度是}}$
$$P_b\left(f\right)=\frac{1}{T_b}\Big|AT_b\text{sinc}\Big(fT_b\Big)\Big|^2=A^2T_b\text{sinc}^2\Big(fT_b\Big)=\frac{A^2}{R_b}\text{sinc}^2\Bigg(\frac{f}{R_b}\Bigg)$$
因此$^{s(t)}$的功率谱密度是
$$P_{s}\left(f\right)=\frac{P_{b}\left(f-f_{c}\right)+P_{b}\left(f+f_{c}\right)}{4}=\frac{A^{2}}{4R_{b}}\left[\mathrm{sinc}^{2}\left(\frac{f-f_{c}}{R_{b}}\right)+\mathrm{sinc}^{2}\left(\frac{f+f_{c}}{R_{b}}\right)\right]$$

\includegraphics[width=0.5\textwidth]{images_ysy/17.png}

(2)A 平方环法
\includegraphics[width=0.5\textwidth]{images_ysy/18.png}

原理：$s\left(t\right)=\sum_{n=-\infty}^{\infty}a_ng\left(t-nT_b\right)\cos\left(2\pi f_ct+\theta\right)$经过平方器后成为
$$\nu\left(t\right)=\left(\sum_{n=-\infty}^{\infty}a_{n}g\left(t-nT_{b}\right)\cos\left(2\pi f_{c}t+\theta\right)\right)^{2}=\cos^{2}\left(2\pi f_{c}t+\theta\right)\left(\sum_{n=-\infty}^{\infty}a_{n}g\left(t-nT_{b}\right)\right)^{2}$$
$=\frac{1+\cos\left(4\pi f_{c}t+2\theta\right)}{2}\sum_{m=-\infty}^{\infty}\sum_{n=-\infty}^{\infty}a_{n}a_{m}g\left(t-mT_{b}\right)g\left(t-nT_{b}\right)$
$=\frac{1+\cos\left(4\pi f_{c}t+2\theta\right)}{2}\sum_{-}^{\infty}a_{n}^{2}g^{2}\left(t-nT_{b}\right)=\frac{A^{2}}{2}+\frac{A^{2}}{2}\cos\left(4\pi f_{c}t+2\theta\right)$
将$^{\nu(t)}$通过一个中心频率为$^{2f_c\text{的滤波器以限制噪声，送入锁相环。锁相环锁定时VCO的}}$ 输出近似是$-\sin\left(4\pi f_ct+2\theta\right)_\text{,经过 2 分频后为}{-\sin\left(2\pi f_ct+\theta\right)}_\text{或者}{\sin\left(2\pi f_ct+\theta\right)}_{,}$ 移相后得到$\cos\left(2\pi f_ct+\theta\right)_\text{或者}-\cos\left(2\pi f_ct+\theta\right)_\text{(相位模糊)。}$


B Costas 环法

\includegraphics[width=0.5\textwidth]{images_ysy/19.png}

原理：假设VCO输出的是$^{-\sin\left(2\pi f_ct+\theta+\varphi\right)}$,则上支路(Q支路)低通滤波的输出将是$-a_n\sin\varphi_\text{,下支路(I支路)低通滤波器的输出是}{\alpha_n\cos\varphi_\text{,环路滤波器的输入是}{-\sin2\varphi_0}}$。设$^{\varphi}$在 0 附近，则当$^{\varphi>0}$时，VCO将降低频率使得$^{\varphi}$减小，当$^{\varphi<0}$时VCO将增加频率使得$\varphi_\text{增加，因此稳态情况下}\varphi_{\text{近似为 }0\text{,使得恢复载波近似为}}\cos(2\pi f_ct+\theta)$。另外还有一个稳定点是$\varphi=\pi$,表示相位模糊。
---------------------------------------
6.8 假设BPSK信号在信道传输中受到双边功率谱密度为$N_0/2=10^{-10}$W/Hz的加性白高斯噪声的干扰，发送信号比特能量$E_b=A^2T_b/2$,其中$^{T_b}$是比特间隔，$A$是信号幅度。最佳接收的误比特率$P_b=10^{-3}$,请求出在输入信息速率$^{R_b\text{分别为 10 kbit/s、100 kbit/s、1 Mbit/s条件下的BPSK}}$

$$\text{于}\frac{1}{2}\operatorname{erfc}\Bigg(\sqrt{\frac{E_{_b}}{N_{_0}}}\Bigg)=1\:0^{-3}\:,\:\text{有}10\lg\frac{E_{_b}}{N_{_0}}=6.8\mathrm{dB}\:)$$
信号幅度$A$值。(注：对于
+++++++++++++++++++++++++++++++++++++++++++++++++++++++++++
解：BPSK最佳接收时的误比特率为$\begin{aligned}P_{b}&=\frac12\mathrm{erfc}\left(\sqrt{\frac{E_{b}}{N_{0}}}\right)\\\text{,今}P_{b}&=10^{-3},\text{由此得}\end{aligned}\frac{E_{b}}{N_{0}}=10^{0.66}$
$$\frac{A^{2}T_{b}}{2}=10^{0.68}\times2\times10^{-10}\quad,\quad A=\sqrt{4\times10^{-9.32}R_{b}}=2\times10^{-4.66}\sqrt{R_{b}}$$
代入不同的速率得知速率分别为 10Kbps、100Kbps、1Mbps 时对应的$A$分别为0.0044、0.0138

及0.0438。
---------------------------------------
6.9 -DPSK数字通信系统，信息速率为$^{R_b}$,输入数据为 110100010110...
(1)写出相对码(设：相对码的第一个比特为 1);
(2)画出 DPSK 发送框图；
(3)写出 DPSK 发送信号的载波相位 (设：第一个比特的 DPSK 信号的载波相位为0);
(4)画出 DPSK 信号的功率谱图(设：输入数据是独立等概序列);
(5)画出 DPSK 的非相干接收框图。
+++++++++++++++++++++++++++++++++++++++++++++++++++++++++++
$解：(1)$

绝对码：110100010110

相对码：1011000011011

(2)DPSK 发送框图

\includegraphics[width=0.5\textwidth]{images_ysy/20.png}

(3)载波相位：$0\pi00\pi\pi\pi\pi00\pi00$

(4)输入数据独立等概时，差分编码的结果也是独立等概的。此时 DPSK 的功率谱和 BPSK

是一样的。

\includegraphics[width=0.5\textwidth]{images_ysy/21.png}

(5)DPSK 的非相干接收框图：

\includegraphics[width=0.5\textwidth]{images_ysy/22.png}
---------------------------------------
6.10 某数字通信系统的信息速率为 3Mbit/s,发送载波为 2GMHz,要求发射频谱限于 (2000 $\pm1)$ MHz 范围内，该系统采用 70MHz 的中频。请给出合理的系统设计，画出发送框图和接收框图。
+++++++++++++++++++++++++++++++++++++++++++++++++++++++++++
解：需要在 $B=2\mathrm{MHz~}_{\text{带宽内发送}R_{\mathrm{b}}=3\mathrm{Mbps,}}$ 因此频带利用率必须达到$\begin{aligned}\frac{R_{b}}{B}=1.5\mathrm{bps/Hz}\\\mathrm{lBaud}\mathrm{Hz-lox~Abnc/Hz}\end{aligned}$ 采用$M$进制二维调制时的最大频带利用率是$^{\text{1Baud}/\mathrm{Hz}=\log_2M\mathrm{bps/Hz}}$。故此必须$\log_{2}M\geq1.5,M\geq2^{1.5}$,取$M=4$,则符号速率是$R_s= \frac {R_{b}}2= 1. 5$MBaud。为了满足限带的要求，必须采用滚降技术。设滚降系数为$\alpha_\text{,则有}B=R_s\left(1+\alpha\right)$,由此得$\alpha=\frac{B}{R_s}-1=\frac{2}{1.5}-1=\frac{1}{3}$。设计采用升余弦滚降。再考虑到收发匹配的问题，则可画出收发

框图如下：

发送框图
\includegraphics[width=0.5\textwidth]{images_ysy/23.png}

接收框图

\includegraphics[width=0.5\textwidth]{images_ysy/24.png}
---------------------------------------
6.12 已知一OQPSK调制器的输入数据是速率为$R_{b}=2\mathrm{Mbps}_\text{的独立等概二进制序列。请}$
(1)画出 OQPSK 调制器的框图 (假设成型滤波器的冲激响应为矩形不归零脉冲);
(2)若输入数据为 11100100...,请画出OQPSK调制器中同相和正交支路的基带信号$^{I(t)}$
$\sum_\text{及}{Q(t)}_\text{波形；}{Q(t)}$
(3)画出 OQPSK 调制信号的双边功率谱密度并标上频率之。

+++++++++++++++++++++++++++++++++++++++++++++++++++++++++++
$解：(1)$

\includegraphics[width=0.5\textwidth]{images_ysy/25.png}

$$\text{(2)串并变换后I路的数据是 1100...,Q路的数据是 1010......,}I\begin{pmatrix}t\end{pmatrix}_\text{及}{Q\begin{pmatrix}t\end{pmatrix}}_\text{的波形如下:}$$
\includegraphics[width=0.5\textwidth]{images_ysy/26.png}

\text{(3)OQPSK 的功率谱密度和 QPSK 一致,其图如下:}
\includegraphics[width=0.5\textwidth]{images_ysy/27.png}
---------------------------------------
6.13 某 4ASK信号的产生框图如题 6.13 图所示，其中$^{\{b_m\}}$是速率为 IMbps的独立等概二进
制序列，$b_m\in \{ 1, 0\} $, $a_n\in \{ - 3, - 1, + 1, + 3\} $, $\{ a_n\}$是独立等概的四进制序列。
$$
g\left(\begin{matrix}t\end{matrix}\right)=\begin{cases} \sqrt{\frac{2}{T_s}} & 0\leq t\leq T_s\\ 0 & \text{else} \end{cases}
$$

求解A点和B点的功率谱密度，以及画图。
\includegraphics[width=0.5\textwidth]{images_ysy/28.png}
+++++++++++++++++++++++++++++++++++++++++++++++++++++++++++
解：A点信号是 4PAM信号。因为$E\left[a_n\right]=0,E\left[a_n^2\right]=\frac{\left(-3\right)^2+\left(-1\right)^2+1^2+3^2}4=5,g(t)$
的傅氏变换为$G(f)=\sqrt{2T_s}\operatorname{sinc}(fT_s)e^{-j\pi JT_s}$,所以A点的功率谱密度为
$$P_A\left(f\right)=\frac{5}{T_s}\Big|\sqrt{2T_s}\:\mathrm{sinc}\Big(fT_s\Big)e^{-j\pi fT_s}\Big|^2=10\:\mathrm{sinc}^2\Big(fT_s\Big)$$

\includegraphics[width=0.5\textwidth]{images_ysy/29.png}

$$\begin{aligned}&\text{B 点信号是 A 点信号的 DSB 调制,因此}\\&P_{B}\left(f\right)=\frac{A^{2}}{4}\Big[P_{A}\left(f-f_{c}\right)+P_{A}\left(f+f_{c}\right)\Big]=2.5A^{2}\left[\mathrm{sinc}^{2}\left(\frac{f-f_{c}}{R_{s}}\right)+\mathrm{sinc}^{2}\left(\frac{f+f_{c}}{R_{s}}\right)\right]\end{aligned}$$

\includegraphics[width=0.5\textwidth]{images_ysy/30.png}
---------------------------------------
6.14 已知 4ASK信号表示式为$s_i(t)=a_if_1(t)_\text{,其中}{a_i\in\{-3,-1,+1,+3\}}_\text{,}{f_1(t)}_\text{是归一化}$
基函数
$$f_{1}\left(t\right)=\sqrt{\frac{2}{T_{s}}}\cos2\pi f_{c}t\quad0\leq t\leq T_{s}$$
(1)请画出 4ASK 信号的信号空间图(星座图);
(2)4ASK信号在信道传输中受到功率谱密度为$N_0/2\text{ 的白高斯噪声}^{n_w}(t)_\text{的干扰，接收信号为}$
$$r\begin{pmatrix}t\end{pmatrix}=s_i\begin{pmatrix}t\end{pmatrix}+n_w\begin{pmatrix}t\end{pmatrix}$$
其最佳接收框图如题 6.14 图所示

\includegraphics[width=0.5\textwidth]{images_ysy/31.png}


(a)若发$^{s_1}(t)$,请写出收端在$^{t=T_s\text{时刻抽样值}r_\text{i 的表示式，并求出均值}E[r_1|s_1]}$、方
差$D\left[r_1\mid s_1\right]_\text{以及条件概率密度函数}p\left(r\mid s_1\right);$
(b)若发$s_2(t)_\text{,请写出}p(r\mid s_2)_\text{的表达式}{:}$ (c)若发$s_3(t)_\text{,请写出}p(r|s_3)_{\text{的表达式}i}$ (d)若发$S_4(t)_\text{,请写出}p(r\mid s_4)_\text{的表达式};$ (e)请画出$p(r|s_1)$、$p(r|s_2)$、$p(r|s_3)$、$p(r|s_4)$图；
$$\begin{aligned}&\text{(f)在}P\big(s_1\big)=P\big(s_2\big)=P\big(s_3\big)=P\big(s_4\big)\\&\text{图表示;}\end{aligned}$$\text{条件下,用最大似然(ML)准则划分判决域,并用}(g)求发送$^{S_{2}}(t)_{\text{而收端错判的概率}}P(e\mid s_{2});$
(h)求出该系统的平均误符号率；
+++++++++++++++++++++++++++++++++++++++++++++++++++++++++++
解：（1）
\includegraphics[width=0.5\textwidth]{images_ysy/32.png}

(2)(a)
$$r_{1}=\int_{0}^{T_{s}}r\left(t\right)f_{1}\left(t\right)dt=\int_{0}^{T_{s}}\left[s_{i}\left(t\right)+n_{w}\left(t\right)\right]f_{1}\left(t\right)dt=a_{i}+\underbrace{\int_{0}^{T_{s}}n_{w}\left(t\right)f_{1}\left(t\right)dt}_{Z}\\Z\sim\mathcal{N}\left(0,\frac{N_{0}}{2}\right)_{\circ}$$
发送$s_1(t)_\text{时，抽样值的表示式为}$
$$r_1=-3+Z$$
$$E[r_1\mid s_1]=-3$$
$$D\left[r_{1}\mid s_{1}\right]=D\left[Z\right]=\frac{N_{0}}{2}\\p\left(r_{1}\mid s_{1}\right)=\frac{1}{\sqrt{\pi N_{0}}}e^{-\frac{\left(r_{1}+3\right)^{2}}{N_{0}}}$$
$$p\left(r_{1}\mid s_{2}\right)=\frac{1}{\sqrt{\pi N_{0}}}e^{-\frac{\left(r_{1}+1\right)^{2}}{N_{0}}}\\p\left(r_{1}\mid s_{3}\right)=\frac{1}{\sqrt{\pi N_{0}}}e^{-\frac{\left(r_{1}-1\right)^{2}}{N_{0}}}\\p\left(r_{1}\mid s_{4}\right)=\frac{1}{\sqrt{\pi N_{0}}}e^{-\frac{\left(r_{1}-3\right)^{2}}{N_{0}}}$$
$( b) ( c) ( d) :$ 同理可得


(e)
\includegraphics[width=0.5\textwidth]{images_ysy/33.png}

(f)ML准则下，最佳判决域的边界是似然概率$p\left(r_1|s_i\right)_\text{相等的点，也就是相邻星座点的距离}$
中点。由此得 4 个判决域的 3 个分界点是-2、0、+2, 如下图示

\includegraphics[width=0.5\textwidth]{images_ysy/34.png}

(g)发送$s_2(t)_\text{时的判决量为}$
$$r_1=a_2+Z=-1+Z$$
判决出错的概率就是$r_{_1\text{落在}}D_{_2\text{之外的概率，也即}}|Z|\geq1$概率，因此
$$P\big(e\:|\:s_{2}\big)=P\big(\big|Z\big|\geq1\big)=\mathrm{erfc}\bigg(\frac{1}{\sqrt{N_{0}}}\bigg)$$
$$P\big(e\:|\:s_{3}\big)=\mathrm{erfc}\bigg(\frac{1}{\sqrt{N_{0}}}\bigg)_{\circ}$$
(h)同理可得

发送$s_1(t)_\text{时的判决量为}$
$$r_1=a_1+Z=-3+Z$$
判决出错的概率就是$^{r_1}$落在$^{D_1}$之外的概率，也即$Z\geq1$概率，因此
$$P\big(e\:|\:s_{1}\big)=P\big(Z\geq1\big)=\frac{1}{2}\operatorname{erfc}\Bigg(\frac{1}{\sqrt{N_{0}}}\Bigg)$$
同理可得$P(e\mid s_4)=\frac12$erfc$\left(\frac1{\sqrt{N_0}}\right)_\text{。因此平均符号错误率为}$
$$P_s=\frac{\frac{1}{2}+1+1+\frac{1}{2}}{4}\mathrm{erfc}\Bigg(\frac{1}{\sqrt{N_0}}\Bigg)=\frac{3}{4}\mathrm{erfc}\Bigg(\frac{1}{\sqrt{N_0}}\Bigg)$$
---------------------------------------
6.15 一数字通信系统在收端抽样时刻的抽样值为
$$y=s_i+n\quad i=1,2,3$$
其中$s_{1}=-2$、$s_{2}=0$、$s_3=2$。$s_1,s_2,s_3$等概出现。$n$是符合标准正态分布的随机变量。如

果按照ML准则进行判决，请写出(1)发 $s_{1\text{时 的 错 判 概 率 }}P( e\mid s_{1})$ ; (2)发$s_2$时 的 错 判 概 率 $P( e\mid s_2) ;$ (3)发 $s_3$时 的 错 判 概 率 $P( e\mid s_3) ;$ (4)系统的平均错判概率$P_e$。
+++++++++++++++++++++++++++++++++++++++++++++++++++++++++++
解：ML 准则下的最佳判决域如下图示

\includegraphics[width=0.5\textwidth]{images_ysy/35.png}

$$\begin{aligned}
&\text{根据判决域可得} \\
&&&P\left(e\mid s_{1}\right)=P\left(y\not\in D_{1}\mid s_{1}\right)=P\left(n>1\right)=\frac{1}{2}\operatorname{erfc}\left(\frac{1}{\sqrt{2}}\right) \\
&&&P\left(e\mid s_{2}\right)=P\left(y\not\in D_{2}\mid s_{2}\right)=P\left(\left|n\right|>1\right)=\mathrm{erfc}\left(\frac{1}{\sqrt{2}}\right) \\
&&&P\left(e\mid s_{3}\right)=P\left(y\not\in D_{3}\mid s_{3}\right)=P\left(n<-1\right)=\frac{1}{2}\mathrm{erfc}\left(\frac{1}{\sqrt{2}}\right) \\
&因此平均错判概率为 \\
&&&P_{e}=\frac{\frac{1}{2}+1+\frac{1}{2}}{3}\operatorname{erfc}\left(\frac{1}{\sqrt{2}}\right)=\frac{2}{3}\operatorname{erfc}\left(\frac{1}{\sqrt{2}}\right) \\
---------------------------------------
&6.16-8PSK及8QAM的星座图如题6.16图所示
\end{aligned}$$

\includegraphics[width=0.5\textwidth]{images_ysy/36.png}

(1)给定$A$,求 8PSK 和 8QAM 星座图中圆的半径$r$、$a$、$b$

(2)给定$A$,假设星座点等概出现，求 8PSK 和 8QAM 各自的平均发送功率。
$$\begin{aligned}A&=2r\sin\frac{\pi}{8}\\A&=2r\sin\frac{\pi}{8}\end{aligned},\text{因此}\begin{aligned}r&=\frac{A}{2\sin\frac{\pi}{8}}=1.3066A\\\end{aligned}。$$
+++++++++++++++++++++++++++++++++++++++++++++++++++++++++++
解：(1)对于8PSK，
对于8QAM，可列出方程

解得
$$\begin{cases}A^2=a^2+a^2\\\\A^2=a^2+b^2-2ab\cos\dfrac{\pi}{4}\end{cases}$$
$$\begin{cases}&a=\dfrac{A}{\sqrt{2}}=0.7071A\\b=\dfrac{\sqrt{3}+1}{2}A=1.366A\end{cases}$$

(2)8PSK 的平均符号能量为
$$E_{8PSK}=r^{2}=\frac{A^{2}}{\left(2\sin\frac{\pi}{8}\right)^{2}}=\frac{A^{2}}{2-\sqrt{2}}=1.7071A^{2}$$
因此平均发送功率为
$$P_{8PSK}=\frac{E_{8PSK}}{T_s}=1.7071\frac{A^2}{T_s}$$
8QAM 的平均符号能量为
$$E_{8QAM}=\frac{a^2+b^2}{2}=\frac{3+\sqrt{3}}{4}A^2=1.1830A^2$$
因此平均发送功率为
$$P_{8QAM}=\frac{E_{8QAM}}{T_s}=1.1830\frac{A^2}{T_s}$$
---------------------------------------
6.17 某 8QAM调制器输入的信息速率为$R_{_b}=90\mathrm{Mbps}_\text{,求符号速率}{R_s}_{。}$
+++++++++++++++++++++++++++++++++++++++++++++++++++++++++++
$R_s=\frac{R_b}{\log_28}=30\mathbf{M}_\text{符号/秒。}$

---------------------------------------
6.18 速率为$^{R_b\text{的二进制数据经过MQAM调制后以}R_s=2400\text{波特的符号速率在 300-3300Hz}}$

的话音频带内传输。请分别就$R_{_b}=9600\mathrm{bps}_\text{和}{R_b}=14400\mathrm{bps}_\text{两种情形画出调制框图，并}$
画出输出的单边功率谱密度。
+++++++++++++++++++++++++++++++++++++++++++++++++++++++++++
解：信道带宽是 $B=3000\mathrm{Hz}_\text{,由于}^{B=\begin{pmatrix}1+\alpha\end{pmatrix}R_{s}}$, 所以滚降系数为$\alpha=\frac B{R_{s}}-1=0.25_{\circ}$ 对于$R_{b}=9600\mathrm{bps}_{\text{的情形，每调制符号需携带4比特信息，因此是16QAM。对于}}$ $R_{b}=14400\mathrm{bps}_{\text{的情形，每调制符号需携带6比特信息，因此是64QAM。相应的调制框图}}$ 如下：

\includegraphics[width=0.5\textwidth]{images_ysy/37.png}

\includegraphics[width=0.5\textwidth]{images_ysy/38.png}

如果发送功率相同，则$R_{b}=9600\mathrm{bps}_\text{和}{R_{b}}=14400\mathrm{bps}_{\text{两种情形下的发送信号功率谱}}$
密度完全相同，单边功率谱密度如下图所示：

\includegraphics[width=0.5\textwidth]{images_ysy/39.png}
---------------------------------------
第七章 信源和信源编码
---------------------------------------
7.1 设一信源由六个不同的独立符号组成：
\includegraphics[width=0.5\textwidth]{images_ysy/40.png}

试求：
(1)信源符号熵 H (X) =?
(2)若信源每秒发送 1000 个符号，求信源每秒传送的信息量应为多少？
(3)若信源各符号等概出现，求信源最大熵$H_{\max}\left(X\right)_{=?}$
+++++++++++++++++++++++++++++++++++++++++++++++++++++++++++
解：(1)
$$\begin{aligned}H\left(X\right)&=-\frac{1}{2}\log_{2}\frac{1}{2}-\frac{1}{4}\log_{2}\frac{1}{4}-\frac{2}{32}\log_{2}\frac{1}{32}-\frac{1}{8}\log_{2}\frac{1}{8}-\frac{1}{16}\log_{2}\frac{1}{16}\\&=\frac{1}{2}\times1+\frac{1}{4}\times2+\frac{1}{8}\times3+\frac{1}{16}\times4+\frac{2}{32}\times5\\&=\frac{31}{16}=1.9375\mathrm{bit/symbol}\\&375\mathrm{~bit/符号}\times1000\text{ 符号/秒 }=1937.5\mathrm{~bit/s}\end{aligned}$$
$(2)R_b=$
$$H_{\mathrm{max}}\left(X\right)=6\times\left(-\frac{1}{6}\log_{2}\frac{1}{6}\right)=\log_{2}6\approx2.585\mathrm{bit/symbol}$$
---------------------------------------
7.2 已知两个二进制随机变量$X$和$Y$服从下列联合分布：
$$P\big(X=Y=0\big)=P\big(X=0,Y=1\big)=P\big(X=Y=1\big)=\frac{1}{3}$$
$_\text{试求：}H(X)_{\text{、}}H(Y)_{\text{、}}H(X|Y)_{\text{、}}H(Y|X)_{\text{和}}H(X,Y)_{\text{。}}$
+++++++++++++++++++++++++++++++++++++++++++++++++++++++++++
解：由联合分布可得到边际分布为
$$P\big(X=0\big)=\sum_YP\big(X=0,Y\big)=\dfrac{2}{3}$$
$$P\bigl(X=1\bigr)=1-P\bigl(X=0\bigr)=\frac{1}{3}$$
$$P\big(Y=1\big)=\sum_XP\big(X,Y=1\big)=\dfrac{2}{3}$$
$$P\big(Y=0\big)=1-P\big(Y=0\big)=\frac{1}{3}$$

$$\text{因此}\\H\left(X\right)=-\frac{2}{3}\log_{2}\frac{2}{3}-\frac{1}{3}\log_{2}\frac{1}{3}=\log_{2}3-\frac{2}{3}=0.9183\mathrm{bit/symbol}\\H\left(Y\right)=-\frac{2}{3}\log_{2}\frac{2}{3}-\frac{1}{3}\log_{2}\frac{1}{3}=H\left(X\right)=0.9183\mathrm{bit/symbol}$$

$$\begin{aligned}
H\left(X\mid Y\right)& =E\bigg[-\log_{2}P\big(X\mid Y\big)\bigg]=-E\bigg[\log_{2}\frac{P\big(X,Y\big)}{P\big(Y\big)}\bigg] \\
&=\sum_{X}\sum_{Y}P\big(X,Y\big)\log_{2}\frac{P\big(Y\big)}{P\big(X,Y\big)} \\
&=\frac13\log_2\frac{1/3}{1/3}+\frac13\log_2\frac{2/3}{1/3}+\frac13\log_2\frac{2/3}{1/3} \\
&=\frac{2}{3}\mathrm{bit/symbol} \\
H\left(X\mid Y\right)& =E\Big[-\log_{2}P\big(Y | X\big)\Big]=-E\Bigg[\log_{2}\frac{P\big(X,Y\big)}{P\big(X\big)}\Bigg] \\
&=\sum_{X}\sum_{Y}P\big(X,Y\big)\log_{2}\frac{P\big(X\big)}{P\big(X,Y\big)} \\
&=\frac13\log_2\frac{2/3}{1/3}+\frac13\log_2\frac{2/3}{1/3}+\frac13\log_2\frac{1/3}{1/3} \\
&=\frac{2}{3}\mathrm{bit/symbol} \\
H(X,& Y\big)=H\left(X\right)+H\left(Y\mid X\right)=\log_{2}3=1.585\mathrm{bit} 
\end{aligned}$$
---------------------------------------
7.3 已知下列联合事件的概率表如下：

<table>
	<tbody>
		<tr>
			<th>$P(A_bB_j)$ $A_i$</th>
			<th>$B_1$</th>
			<th>$B_2$</th>
			<th>$B_3$</th>
		</tr>
		<tr>
			<td>$A_1$</td>
			<td>0.10</td>
			<td>0.08</td>
			<td>0.13</td>
		</tr>
		<tr>
			<td>$A_2$</td>
			<td>0.05</td>
			<td>0.03</td>
			<td>0.09</td>
		</tr>
		<tr>
			<td>$A_3$</td>
			<td>0.05</td>
			<td>0.12</td>
			<td>0.14</td>
		</tr>
		<tr>
			<td>$A_4$</td>
			<td>0.11</td>
			<td>0.04</td>
			<td>0.06</td>
		</tr>
	</tbody>
</table>

试求：
$$\begin{aligned}&(1)P\big(A_{i}\big)=?\:P\big(B_{j}\big)=?\\&(2)H\left(A\right)=?\:H\left(B\right)=?\:H\left(A,B\right)=?\\&(3)I\left(A;B\right)=?\end{aligned}$$
+++++++++++++++++++++++++++++++++++++++++++++++++++++++++++
$解：(1)P(A_{1})=0.10+0.08+0.13=0.31,同理可得$
$\begin{aligned}&P(A_{2})=0.17;\:&\mathrm{P(A_{3})=0.31~;}&&P(A_{4})=0.21;\\&P(B_{1})=0.31\:;\:&P(B_{2})=0.27\:;\:&P(B_{3})=0.42\end{aligned}$
$$H(A)=-\sum_{\mathrm{i}=1}^{4}P(A_{i})\log_{2}P(A_{i})=0.524\times2+0.435+0.473=1.95\mathrm{~bit/symbol}$$
(2)
$$\begin{aligned}
&&H(A)=-\sum_{\mathrm{i=1}}^{4}P(A_{i})\log_{2}P(A_{i})=0.524\times2+0.435+0.473=1.95\mathrm{~bit/symbol} \\
&&H(B)=-\sum_{j=1}^{3}P\Big(B_{j}\Big)\log_{2}P\Big(B_{j}\Big)=0.524+0.51+0.526=1.56\mathrm{~bit/symbol} \\
&&H(A,B)=-\sum_{\mathrm{i}=1}^{4}\sum_{\mathrm{j}=1}^{3}P(A_{i}B_{j})\log_{2}P(A_{i}B_{j})=1.0061+0.6804+0.98+0.7796=3.446 \mathrm{bit} \\
&&I\left(A;B\right)=H\left(A\right)+H\left(B\right)-H\left(A,B\right)=1.95+1.56-3.446=0.064 \mathrm{bit} \\
---------------------------------------
&_\text{7.4 证明:}I\left(X,Y\right)=H\left(X\right)+H\left(Y\right)-H\left(X,Y\right) \\
+++++++++++++++++++++++++++++++++++++++++++++++++++++++++++
&\text{证:} \\
&I(X;Y)& \text{-} && =H(X)-H(X | Y) \\
&&=H\left(X\right)-E\left[-\log P\left(X\mid Y\right)\right] \\
&&=H\left(X\right)+E\left[\log\frac{P\left(X,Y\right)}{P\left(Y\right)}\right] \\
&&=H\left(X\right)+E\Big[-\log P\Big(Y\Big)\Big]-E\Big[-\log P\Big(X,Y\Big)\Big] \\
&&=H\left(X\right)+H\left(Y\right)-H\left(X,Y\right) \\
&\# \\
---------------------------------------
&\text{7.5 已知一信源}
\end{aligned}$$

\includegraphics[width=0.5\textwidth]{images_ysy/41.png}

试求：

(1)信源熵$H(X)=?$

(2)若进行Huffman编码，试问如何编码？并求编码效率$\eta_{=?}$
+++++++++++++++++++++++++++++++++++++++++++++++++++++++++++
解：(1)
$$H(X)=-\sum_{i=1}^3P(x_i)\log_2P(x_i)=0.5184+0.5301+0.46438=1.51\text{bit/symbol}$$
(2)Huffman 编码方式如下：

\includegraphics[width=0.5\textwidth]{images_ysy/42.png}

$$\text{平均码长}\quad\overline{K}=\sum_{i=1}^3P(x_i)k_i=0.45\times1+0.35\times2+0.2\times2=1.55\text{bit}$$








% 129-159页


$$\eta=\frac{H(X)}{\overline{K}}=\frac{1.51}{1.55}=97.4\%$$
编码效率
---------------------------------------
7.6已知一信源
\begin{figure}
    \centering
    \includegraphics[width=1\linewidth]{赵培志/d4b46fb7c60298289b6cc7c9232c846.png}
\end{figure}


试求：

(1)信源熵$H(X)=?$

(2)若进行Huffman编码，试问如何编码？并求编码效率$\eta=?$
++++++++++++
解：(1)

$$H(X)=-\sum_{i=1}^{7}P\bigl(x_{i}\bigr)\log_{2}P\bigl(x_{i}\bigr)=2.63\text{bit/symbol}$$

(2）编码方式如下
\begin{figure}
    \centering
    \includegraphics[width=1\linewidth]{赵培志/dba23daf26d0e17ab4778a6536864cd.png}
    
    
\end{figure}



平均码长$$\overline{K}=\sum_{i=1}^{7}P(x_{i})k_{i}=2.72\mathrm{bit}$$



编码效率
$$\eta=\frac{H(X)}{\overline{K}}=\frac{2.63}{2.72}=96.7\%$$---------------------------------------
7.7试确定能重构信号$x(t)=$sinc$(2000t)$ 所需的最低取样频率$f_{s}$值。
++++++++++++
解： $x(t)$ 的傅氏变换是

$$X\left(f\right)=\begin{cases}\frac{1}{2000}&\left|f\right|\leq1000\\0&else\end{cases}$$

其带宽是1000Hz ，因此所需的最低取样率是$f_{s}=2000\mathrm{Hz}_{o}$
---------------------------------------
$$\begin{aligned}&\text{7.8 已知信号}s\begin{pmatrix}t\end{pmatrix}=10\cos200\pi t\cos2000\pi t\text{,对}s\begin{pmatrix}t\end{pmatrix}_\text{以}f_s\text{ 的速率进行理想采样得到抽样}\\&\text{信号}s_s\begin{pmatrix}t\end{pmatrix}=\sum_{n=-\infty}^\infty s\begin{pmatrix}nT_s\end{pmatrix}\delta\begin{pmatrix}t-nT_s\end{pmatrix}_\text{。将}S_s\begin{pmatrix}t\end{pmatrix}_\text{通过一个}\end{aligned}$$

(1)截止频率为$f_{H}$ 的理想低通滤波器；

(2)中心频率为$f_c$ ，带宽为B的理想带通滤波器。

后输出还是$s(t)$ ，求相应的最小抽样频率和对应的滤波器参数。
++++++++++++
解： $s(t)$ 的最高频率分量是$f_m=1100$Hz ，因此用低通恢复时，最小取样率是$f_s=2200$Hz 。对应低通滤波器的截止频率是1100Hz

$(2)^{s(t)}$ 的带宽是$W=200$Hz ，最高频率是$f_{m}=1100$Hz ，因此最小需要的抽样频率是。其中$n=\left\lfloor\frac{f_m}{W}\right\rfloor=5$ 。所以$f_s=\frac{2200}5=440$Hz 。对应要求理想带通滤波器带宽B是200Hz。( $B=1100-900=200$Hz)
---------------------------------------
7.9已知一个12路载波电话占有频率范围$60\sim108$ HZ ，求出其最低取样频率$f_{\mathrm{smin}}$
++++++++++++
解：

$$f_{s\min}=\frac{2f_H}{n}$$

其中

$$n=\left\lfloor\frac{f_H}{B}\right\rfloor=\left\lfloor\frac{108}{108-60}\right\rfloor=2$$

所以

$$f_{s\min}=\frac{2f_H}{2}=f_H=108\text{KHz}$$
---------------------------------------
7.10将一个幅度为3.25V. 频率为$800Hz$ 的正弦信号输入到取样频率为$8KHz$ 的取样保持器然后再通过一个量化特性如下图所示的8电平均匀量化器。试画出量化器输出的波形

\begin{figure}
    \centering
    \includegraphics[width=1\linewidth]{赵培志/4b6cc16b385f94a526a6aa34bd95494.png}
\end{figure}
++++++++++++
解：抽样频率是正弦信号频率的10倍，每个正弦周期内有10个样点，抽样值及其量化电平如下表所示：



\begin{figure}
    \centering
    \includegraphics[width=1\linewidth]{赵培志/2b37756bcaaf6fbe8fb74f50d6817f4.png}
\end{figure}
量化器输出如下图示：


\begin{figure}
    \centering
    \includegraphics[width=1\linewidth]{赵培志/dfa4dc5d50a151bc69a2dfb1d7505f8.png}
\end{figure}
---------------------------------------
7.11 已知模拟信号取样值x的概率密度函数$p(x)$ 图所示。将x经过一个4电平均匀量化器得到输出是$y\in\left\{y_{1},y_{2},y_{3},y_{4}\right\}$ ，试求:

(1）量化器输出信号的平均功率$S_q=E[y^2]$
(2）量化噪声$e=y-x$ 平均功率$N_q=E\biggl[e^2\biggr]$
(3）量化信噪比$\frac{S_{q}}{N_{q}} $的分贝值。

\begin{figure}
    \centering
    \includegraphics[width=1\linewidth]{赵培志/9d76c0b6614b1120340d327f9418f0f.png}
\end{figure}
++++++++++++
解：依题意，4个量化区间是
$$D_{1}=\left[-1,-\frac{1}{2}\right],\:D_{2}=\left[-\frac{1}{2},0\right],\:D_{3}=\left[0,\frac{1}{2}\right],\:D_{4}=\left[\frac{1}{2},1\right]$$

4个量化电平是

$$y_1=-\frac{3}{4}\:,\quad y_2=-\frac{1}{4}\:,\quad y_3=\frac{1}{4}\:,\quad y_4=\frac{3}{4}$$

如下图示：



\begin{figure}
    \centering
    \includegraphics[width=1\linewidth]{赵培志/799e8c33ef4d94413be3d8a0d5ab656.png}
\end{figure}

取样值$x$ 落入4个区间的概率是

$$P_1=P_4=\int_{0.5}^1p\big(x\big)dx=\frac{1}{8}$$

$$P_2=P_3=\int_0^{0.5}p\left(x\right)dx=\frac{3}{8}$$

因此：

$$\begin{aligned}
&S_{q}&& =E\bigg[y^2\bigg]=\sum_{i=1}^4y_iP_i \\
&&&=2\times\left[\left(\frac{1}{4}\right)^{2}\times\frac{3}{8}+\left(\frac{3}{4}\right)^{2}\times\frac{1}{8}\right]=2\times\left(\frac{1}{4}\right)^{2}\times\frac{3}{8}\times\left(1+3\right)=\frac{3}{16} \\
&N_{q}=E\bigg[\big(y-x\big)^{2}\bigg] \\
&&&=2\times\left[\int_0^{1/2}\left(x-\frac{1}{4}\right)^2\left(1-x\right)dx+\int_{1/2}^1\left(x-\frac{3}{4}\right)^2\left(1-x\right)dx\right]
\end{aligned}$$

在第一个积分中令$ t=x-\frac{1}{4}\quad $,在第二个积分中令$t=x-\frac{3}{4}\quad,\text{得}$
$$\begin{aligned}
N_{q}& =2\times\left[\int_{-1/4}^{1/4}t^{2}\left(\frac{3}{4}-t\right)dt+\int_{-1/4}^{1/4}t^{2}\left(\frac{1}{4}-t\right)dt\right]  \\
&=2\int_{-1/4}^{1/4}t^{2}\left(1-2t\right)dt \\
&=2\int_{-1/4}^{1/4}t^{2}dt \\
&=\frac{1}{48} \\
\end{aligned}$$

或者在第二个积分中令$t=1-x$ ，得
$$\begin{aligned}
N_{q}& =2\times\left[\int_{0}^{1/2}\left(x-\frac{1}{4}\right)^{2}\left(1-x\right)dx+\int_{0}^{1/2}\left(\frac{1}{4}-t\right)^{2}tdx\right]  \\
&=2\int_{0}^{1/2}\Bigg(x-\frac{1}{4}\Bigg)^{2}\:dx \\
&=2\int_{-1/4}^{1/4}t^{2}dt \\
&=\frac{1}{48}
\end{aligned}$$



$$\frac{S_{q}}{N_{q}}=\frac{\frac{3}{16}}{\frac{1}{48}}=\frac{3\times48}{16}=9
$$
$$\left[\frac{S_{q}}{N_{q}}\right]_{\mathrm{dB}}=10\lg9=9.542\mathrm{dB}$$
---------------------------------------
7.12将正弦信号$x(t)=\sin\bigl(1600\pi t\bigr)$ 以 8KHz速率取样后输入于A律13折线PCM编码器此PCM编码器得设计输入电压范围是[-1,+1] 。求在一个正弦信号周期内所有取档$x_n=\sin\frac{n\pi}5$ $n=0,1,\cdots,9$ 的DPCM编综码输出码组
++++++++++++
解：

\begin{figure}
    \centering
    \includegraphics[width=1\linewidth]{赵培志/3718224589e73f6cc2243dfa45ebc92.png}
\end{figure}
---------------------------------------
7.13某A律13折线PCM编码器的设计输入范围是(-5,5)V. 。若抽样脉冲幅度$x=+1.2$V ，请按照CCITTG.711建议进行PCM编码

（1）求编码器的输出码组； 
（2）求解码器输出的量化电平值，并计算量化误差
（3）写出对应于A律13折线PCM码组的均匀量化线性编码的码组（13位码）
++++++++++++
解： $(1)\:x>0$ 所以极性码为1;
$\begin{vmatrix}x\end{vmatrix}=1.2\text{V},$折分$\text{为}\frac{1.2}5\times4096=983.04$量化单位。

\begin{figure}
    \centering
    \includegraphics[width=1\linewidth]{赵培志/6cebfb97f070670afdb24cd0592ab9e.png}
\end{figure}

---------------------------------------

从此图可见段落码是101

此段落的量化间隔是 1 =32 化单位。因此段内码是$\left\lfloor\frac{983.04-512}{\Delta}\right\rfloor=14=\left(1110\right)_{2}$ 因此输出码组是11011110

(2）所在量化区间的起点是960单位，终点是992单位，因此量化电平是976个量化单位， 肌$976\times\frac{5}{4096}=1.1914$V 。量化误差为983.04-976=7.04量化单位,即0.0086伏

$$(3)976=512+256+128+64+16=2^9+2^8+2^7+2^6+2^4$$
因此编码结果是(10011101000），第1位是极性码。
---------------------------------------
7.14在CD播放机中，假设音乐是均匀分布，抽样速率为$44.1kHz$ ，采用每抽样16比特的均匀量化线性编码进行量化、编码。试确定存储50分钟时间段的音乐所需要的比特数和字节数，并求出量化信噪比的分贝值。
++++++++++++
解：(1)$44.1\times10^{3}\times\frac{16}{8}\times50\times60\mathrm{bit}=2.1168\mathrm{Gbit=264.6MByte}$


(2)量化级数是$M=2^{16}$ ，对于均匀量化器，当输入为均匀分布时的量化信噪比为

$$\frac{S}{N_q}=M^2=2^{32}$$

换算成分贝值为

$$10\lg\left(\frac{S}{N_q}\right)=320\lg2\approx96.33\mathrm{dB}$$
---------------------------------------
7.15已知$x(n)$ 是平稳随机序列，其自相关特性为
$$R_{x}\left(m\right)=\begin{cases}1&m=0\\c_{1}&m=\pm1\\c_{2}&m=\pm2\\0&else\end{cases}$$

若DPCM采用了如下的一个二阶线性预测

$$\hat{x}\big(n\big)=\alpha_1x\big(n-1\big)+\alpha_2x\big(n-2\big)$$

(1)请确定使预测误差最小的系数 $\alpha_{1}^{*}$ 和$\alpha_2^{*}$ ;

（2）给出最小的均方误差

++++++++++++
解：预测误差为

$$e\begin{pmatrix}n\end{pmatrix}=x\begin{pmatrix}n\end{pmatrix}-\hat{x}\begin{pmatrix}n\end{pmatrix}=x\begin{pmatrix}n\end{pmatrix}-\alpha_1x\begin{pmatrix}n-1\end{pmatrix}-\alpha_2x\begin{pmatrix}n-2\end{pmatrix}$$

均方误差是



$$\begin{aligned}
\sigma_{e}^{2}& =E\bigg[e^{2}\left(n\right)\bigg]=E\bigg[\bigg|x(n)-\alpha_{1}x\big(n-1\big)-\alpha_{2}x\big(n-2\big)\bigg|^{2}\bigg] \\
&=E\bigg[x^{2}\left(n\right)+\alpha_{1}^{2}x^{2}\left(n-1\right)+\alpha_{2}^{2}x^{2}\left(n-2\right)\bigg] \\
&+E\Big[-2\alpha_{1}x\big(n\big)x\big(n-1\big)-2\alpha_{2}x\big(n\big)x\big(n-2\big)+2\alpha_{1}\alpha_{2}x\big(n-1\big)x\big(n-2\big)\Big] \\
&=R_{x}\left(0\right)\left(1+\alpha_{1}^{2}+\alpha_{2}^{2}\right)+2\left(-\alpha_{1}+\alpha_{1}\alpha_{2}\right)R_{x}\left(1\right)-2\alpha_{2}R_{x}\left(2\right) \\
&=1+\alpha_{1}^{2}+\alpha_{2}^{2}-2\alpha_{1}c_{1}+2\alpha_{1}\alpha_{2}c_{1}-2\alpha_{2}c_{2}
\end{aligned}$$

求极小值，通过令$$\frac{\delta\sigma_{e}^{2}}{\delta\alpha_{1}}=0\:,\:\frac{\delta\sigma_{e}^{2}}{\delta\alpha_{2}}=0\:;$$
得到

$$\begin{cases}2\alpha_1^*-2\Big(1-\alpha_2^*\Big)c_1=0\\2\alpha_2^*-2c_2+2c_1\alpha_1^*=0\end{cases}$$

解得

$$\begin{cases}\alpha_1^*=\frac{c_1-c_1c_2}{1-c_1^2}=\frac{c_1\left(1-c_2\right)}{1-c_1^2}\\\\\alpha_2^*=\frac{c_2-c_1^2}{1-c_1^2}\end{cases}$$

将$\alpha_{1}^{*}$ ， $\alpha_{2}^*$ 代入$\sigma_{c}^{2}$ 式，得到

$$\sigma_{e\min}^{2}=\frac{2\left(1-c_{2}\right)c_{1}^{4}+\left(c_{2}^{2}+2c_{2}-3\right)c_{1}^{2}+1-c_{2}^{2}}{\left(1-c_{1}^{2}\right)^{2}}$$
---------------------------------------
7.16若已知离散信源协方差矩阵

$$\Phi_{_x}=\begin{bmatrix}a&b&b&b\\b&a&b&b\\b&b&a&b\\b&b&b&a\end{bmatrix}$$

试求：

(1)由K-L变换后所得的协方差$\Phi_{y}=?$

(2)用Walsh变换代替K-L变换，求输出协方差$\Phi_{y=?}$
++++++++++++
解： $_{(1)}\Phi_{x}$ 的特征值是$a+3b,a-b,a-b,a-b$ ，根据K-L变换的性质可得到变换后的协方差矩阵为
$$\Phi_y=\begin{bmatrix}a+3b&0&0&0\\0&a-b&0&0\\0&0&a-b&0\\0&0&0&a-b\end{bmatrix}$$
$_{(2)}$$$\begin{gathered}H_{WH}=\begin{bmatrix}+1&+1&+1&+1\\+1&-1&+1&-1\\+1&+1&-1&-1\\+1&-1&-1&+1\end{bmatrix}\\H_{WH}^{T}=\begin{bmatrix}+1&+1&+1&+1\\+1&-1&+1&-1\\+1&+1&-1&-1\\+1&-1&-1&+1\end{bmatrix}\end{gathered}$$$$\Phi_y=H_{WH}\Phi_xH_{_{WH}}^T=\begin{bmatrix}4a+12b&0&0&0\\0&4a-4b&0&0\\0&0&4a-4b&0\\0&0&0&4a-4b\end{bmatrix}$$---------------------------------------
7.17 若已知离散信源协方差矩阵 $$\Phi_x=\begin{bmatrix}a&b&b&0\\b&a&b&b\\b&b&a&b\\0&b&b&a\end{bmatrix}$$试求：经离散傅氏变换后（DFT）的协方差矩阵$\Phi_{y}=_{?}$
++++++++++++
解：$$\begin{gathered}A_{DFT}=\frac{1}{2}\begin{bmatrix}1&1&1&1\\1&-j&-1&j\\1&-1&1&-1\\1&j&-1&-j\end{bmatrix}\\A_{DFT}^{H}=\begin{bmatrix}A_{DFT}^*\end{bmatrix}^T=\frac{1}{2}\begin{bmatrix}1&1&1&1\\1&+j&-1&-j\\1&-1&1&-1\\1&-j&-1&+j\end{bmatrix}\end{gathered}$$


注意记号$()^H$表示矩阵的共轭专置。

$$\Phi_{y}=A_{DFT}\Phi_{x}A_{DFT}^{H}$$

$$\begin{aligned}
&=\frac{1}{2}\begin{bmatrix}1&1&1&1\\1&-j&-1&j\\1&-1&1&-1\\1&j&-1&-j\end{bmatrix}\times\begin{bmatrix}a&b&b&0\\b&a&b&b\\b&b&a&b\\0&b&b&a\end{bmatrix}\times\frac{1}{2}\begin{bmatrix}1&1&1&1\\1&+j&-1&-j\\1&-1&1&-1\\1&-j&-1&+j\end{bmatrix} \\
&=\frac{1}{4}\begin{bmatrix}a+2b&a+3b&a+3b&a+2b\\a-b-bj&\left(b-a\right)j&b-a&\left(a-b\right)j-b\\a&b-a&a-b&-a\\a-b+bj&\left(a-b\right)j&b-a&\left(b-a\right)j-b\end{bmatrix}\times\begin{bmatrix}1&1&1&1\\1&+j&-1&-j\\1&-1&1&-1\\1&-j&-1&+j\end{bmatrix} \\
&=\frac{1}{4}\begin{bmatrix}4a+10b&-b+jb&0&-b-jb\\-b-jb&4\left(a-b\right)&b-jb&-j2b\\0&b+jb&4a-2b&b-jb\\-b+bj&j2b&b+jb&4\left(a-b\right)\end{bmatrix}
\end{aligned}$$

---------------------------------------

第八章信道

8.1已知信道的结构如题8.1图所示，求信道的冲激响应和传递函数，并说明是恒参信道还是随参信道，何种信号经过信道有明显失真？何种信号经过信道的失真可以忽略？

\begin{figure}
    \centering
    \includegraphics[width=1\linewidth]{赵培志/98154c8a0239b2f5374e3cc28b3daf3.png}
\end{figure}
++++++++++++
解：由图可知信道的冲激响应为$h(t)=\delta(t)+\delta(t-\tau)$ ，传递函数为$H\left(f\right)=1+e^{-j2\pi f\tau}$ 冲激函数和传递函数不随时间变化，可知该信道为恒参信道。此信道是一个二径信道，信道的相干带宽为$B_c=1/\tau$ 。因此，若输入信号的带宽远小于$1/\tau$ 时，输出的失真将可以忽略； 若输入信号的带宽与$1/\tau$ 可比或者更大时，信号经过信道会有明显失真
---------------------------------------
8.2已知信道的传输特性如题8.2（a）图所示，其输入为已调信号$s(t)=m(t)$ cos $\omega_ct$ s(t)的频谱密度如题8.2（b）图所示，且$W<\Delta\omega$ 。求：信道的输出信号，并说明有无失真，

\begin{figure}
    \centering
    \includegraphics[width=1\linewidth]{赵培志/e1e41ce47b6904e4fc3b0d79dac6a8e.png}
\end{figure}
++++++++++++
解：信道的传输特性为

---------------------------------------

$$H\left(\omega\right)=\begin{cases}e^{-j\left(\omega-\omega_c\right)t_0}&\left|\omega-\omega_c\right|\leq\Delta\omega/2\\e^{-j\left(\omega+\omega_c\right)t_0}&\left|\omega+\omega_c\right|\leq\Delta\omega/2\\0&else\end{cases}$$

输入信号的傅利叶变换为

$$S(\omega)=\frac{1}{2}M(\omega+\omega_c)+\frac{1}{2}M(\omega-\omega_c)$$

因而输出信号的频谱为
$$Y(\omega)=S(\omega)H(\omega)=\frac{1}{2}M(\omega+\omega_{c})e^{-j(\omega+\omega_{c})t_{0}}+\frac{1}{2}M(\omega-\omega_{c})e^{-j(\omega-\omega_{c})t_{0}}$$

作傅氏反变换得

$$y(t)=\frac{1}{2}m(t-t_{0})e^{-j\omega_{c}t}+\frac{1}{2}m(t-t_{0})e^{j\omega_{c}t}=m(t-t_{0})\cos\omega_{c}t\:,$$

对比输入$s(t)=m(t)\cos\omega_{c}t$ ，信号整体有失真，但包络$m(t)$ 无失真。
---------------------------------------
8.3令恒参信道模型如题8.3图所示，求其时延特性和群时延特性，并说明它们对信号传输的影响。

\begin{figure}
    \centering
    \includegraphics[width=1\linewidth]{赵培志/5e701bf9ea2fb6ba68de437bfd5e087.png}
\end{figure}

++++++++++++

解：信道的传输特性为

$$H\left(f\right)=\frac{\frac{1}{j2\pi fC}}{R+\frac{1}{j2\pi fC}}=\frac{1}{1+j2\pi fRC}$$

幅频特性为

$$\begin{vmatrix}H\left(f\right)\end{vmatrix}=\frac{1}{\sqrt{1+\left(2\pi fRC\right)^{2}}}$$

相频特性为

$$\varphi\bigl(f\bigr)=-\tan^{-1}\bigl(2\pi fRC\bigr)$$

时延特性为

$$\tau(f)=\frac{\varphi(f)}{2\pi f}=-\frac{\tan^{-1}\left(2\pi fRC\right)}{2\pi f}$$

群时延特性为

$$\tau_G\left(f\right)=\frac{d\varphi\left(f\right)}{2\pi df}=-\frac{RC}{1+\left(2\pi fRC\right)^2}$$

由于时延特性和群时延特性都不是常数，信号通过此信道会发生失真。但当$f<<\frac{1}{2\pi RC}$ 时
$$H\begin{pmatrix}f\end{pmatrix}=\begin{vmatrix}H\begin{pmatrix}f\end{pmatrix}\end{vmatrix}e^{j\varphi(f)}\approx e^{-j2\pi fRC}$$



此时该信道近似是一个时延为$RC$ 的无失真信道。
---------------------------------------
8.4设随参信道的最大多径时延扩展等于$5\mu s$ ，若欲传输100kbit/s的信息速率，问下列调制方式下，哪个可以不考虑使用均衡器？

1) BPSK

(2)将数据经过串并变换为1000个并行支路后用1000个不同的载频按BPSK调制方式传输。（假设各载频间互相正交）
++++++++++++
解：采用BPSK传输时的符号间隔是$10\mu s$ ，多径时延和符号间隔相比不可忽略，存在码间干扰，故需要使用均衡器。

采用1000路并行传送时，每路的数据速率100bit/s，符号间隔是$10ms$ ，多径时延近似可忽略。故此可以不需要均衡器。
---------------------------------------
8.5已知在高斯信道理想通信系统传送某一信息所需带宽为$10^6Hz$ ，信噪比为20dB：若将所需信噪比降为$10dB$ ，求所需信道带宽。
++++++++++++
解：根据Shannon公式

$$C=B\log\left(1+\frac{S}{\sigma^2}\right)$$

可求得信噪比为20dB时的信道容量为

$$C=10^6\log_2\left(1+100\right)\mathrm{bps}。$$

当信噪比降为10dB时，为了保持相同的信道容量，要求
$$B\log_2(1+10)=10^6\log_2(1+10^2)$$

由此可得

$$B=\frac{10^6\log_2101}{\log_211}=1.92\times10^6\text{Hz。}$$
---------------------------------------
8.6计算机终端输出的数据经编码调制后通过电话信道传输，电话信道带宽为$3.4kHz$ ，加性高斯噪声信道输出的信噪比S/N等于20dB。该终端每次输出128个可能符号之一，输出序列中的各符号相互统计独立、等概出现。

（1）计算信道容量：

（2）计算理论上容许的终端输出数据的最高符号速率。
++++++++++++
解：(1) $C=3.4\times10^{3}\log_{2}(1+10^{2})=22.638kbit/s\:;$ 
     （2）最高符号速率为$C/\log_{2}128=3.234$k 波特。
---------------------------------------
8.7分别对10路话音信号（每路信号的最高频率分量为$4kHz$ ）以奈奎斯特速率取样，取样后均匀量化、线性编码，然后将此10路PCM信号与两路180kbit/s的数据以时分多路方式复用，时分复用的输出送至BPSK解调器。若要求BPSK信号的功率谱主瓣宽度为（ $100\pm1$ MHz，其中100 MHz 是载波频率，请求出每路PCM信号的最大量化电平数M值。
++++++++++++
解：每路话音的抽样速率为$8kHz$ ，若将抽样值量化为$b$ 个二进制码元，则有$2^b$ 个量化电平码元速率为$8b\times10^3$bps 。10路PCM信号与两路180kbit/s的数据时分复用，输出速率为

$R=8b\times10^{4}+2\times180\times10^{3}=(8b+36)\times10^{4}$ bit/s。进行BPSK调制后功率谱主瓣宽度应为$2R=2(8b+36)\times10^{4}$ Hz ，实际主瓣带宽为$2MHz$ ，因此$(8b+36)\times10^{4}=1\times10^{6}$ ，因此$b=8$ ，最大量化电平数$M=2^{b}=256$ 。

---------------------------------------

8.8某一待传输的图片约含$2.5\times10^6$ 个象素，为了很好地重现图片，需要将每象素量化为16个亮度电平之一，假若所有这些亮度电平等概出现且互不相关，并设加性高斯噪声信道中的信噪比为30dB，试计算用3分钟传送一张这样的图片所需的最小信道带宽（假设不压缩编码）。
++++++++++++
解：该图片的信息量为

$$I=2.5\times10^6\times\left(-\log_2\frac{1}{16}\right)=10^7\text{bit}$$

3分钟传送该图片所需信道容量为

$$C=\frac{10^7}{180}bit/s$$

由仙农公式$C=B\log_{2}\left(1+\frac{S}{\sigma^{2}}\right)$ 得

$$\frac{10^7}{180}=B\times\log_2\begin{pmatrix}1+1000\end{pmatrix}$$

所以

$$B\approx5.6\mathrm{kHz}$$
---------------------------------------
8.9对10路模拟信号分别进行A律13折线PCM，然后进行时分复用，再经过$\alpha=0.5$ 升余弦脉冲成形滤波器无码间干扰传输，该升余弦基带系统的截止频率为$480kHz$ （1）求该系统的最大信息传输速率

（2）求允许每路模拟信号的最高频率分量$f_{H}$ 值。
++++++++++++
解：（1）由公式$\frac{R_s}2(1+\alpha)=f_m$ ，得$\frac{R_{s}}{2}(1+0.5)=480k$ 所以最大信息传送速率为$R_{s}=640kbps$

（2）10路模拟信号数字化以后的总速率不能超过640kbps，因此每路模拟信号数字化后的速率不能超过64kbps。A律13折线编码是对采样值进行8比特编码，因此每路模拟信号的采样率不得高于64/8=8ksps 。按Nyquist速率抽样，则每路模拟信号的最高频率分量不得超过$4kHz$
---------------------------------------
8.10给定信道如下图所示

\begin{figure}
    \centering
    \includegraphics[width=1\linewidth]{赵培志/af01f5ad41635f91be2c3ba59058389.png}
\end{figure}


图中的$X\in\{0,1\}$ 是信道输入， $Y\in\{0,1,E\}$ 是信道输出， $\begin{pmatrix}1-\mu\end{pmatrix}$ 、 $\mu$ 是信道转移概率$P(Y|X)$。

(1)若已知$X$ 的概率分布是$P(X=0)=p$ ， $P\big(X=1\big)=1-p$ 。试求

(a)Y的概率分布；
(b)H(Y);
(c)H(Y|X) ；
(d)I(Y;X); 
(2)求此信道的信道容量。
++++++++++++
\textbf{解}：(1) $$\begin{aligned}
&\mathrm{(a)} P\big(Y=0\big)=\sum_{x}P\big(Y=0 | x\big)P\big(x\big)=p\big(1-\mu\big) \\
&P\big(Y=1\big)=\sum_xP\big(Y=1 | x\big)P\big(x\big)=\big(1-p\big)\big(1-\mu\big) \\
&P\big(Y=E\big)=\sum_xP\big(Y=E | x\big)P\big(x\big)=p \mu+\big(1-p\big)\mu=\mu 
\end{aligned}$$$$\begin{aligned}
& \mathrm{(b)} \\
&H(Y)&& =\sum_vP\big(Y\big)\log P\big(Y\big) \\
&&&=-P(Y=E)\log P(Y=E)-P(Y=0)\log P(Y=0)-P(Y=1)\log P(Y=1) \\
&&&=-\mu\log\mu-p\left(1-\mu\right)\log\left[p\left(1-\mu\right)\right]-(1-p)(1-\mu)\log\left[(1-p)(1-\mu)\right] \\
&&&=-\mu\log\mu-p\left(1-\mu\right)\left[\log p+\log\left(1-\mu\right)\right] \\
&&&-(1-p)(1-\mu)[\log(1-p)+\log(1-\mu)] \\
&&&=-\mu\log\mu-(1-\mu)\log(1-\mu)-p(1-\mu)\log p-(1-p)(1-\mu)\log(1-p) \\
&&&=h_{2}(\mu)+(1-\mu)h_{2}(p)
\end{aligned}$$
其中$h_2\left(x\right)\triangleq-x\log x+\left(1-x\right)\log\left(1-x\right)$


$(\mathfrak{c})\text{给定 }X=0\text{ 条件下,}Y\text{只可能是 }0\text{ 和 }E\text{,于是}$$$\begin{aligned}
H\big(Y | X=0\big)& =E_{Y|X=0}\Big[-\log P\Big(Y | X=0\Big)\Big] \\
&=-\sum_{Y\in\{0,E\}}P\big(Y\mid X=0\big)\log P\big(Y\mid X=0\big) \\
&=h_{2}(\mu)
\end{aligned}$$同理可得$H\left(Y\mid X=1\right)=h_2\left(\mu\right)$

于是$$\begin{aligned}
H\left(Y|X\right)& =E_{x}\left\{E_{Y|X=x}\Big[-\log P\big(Y | X=x\big)\Big]\right\} \\
&=h_{2}(\mu)
\end{aligned}$$

$$\begin{aligned}&(\mathrm{d})I\left(Y;X\right)=H\left(Y\right)-H\left(Y\mid X\right)=\left(1-\mu\right)h_{2}\left(p\right)\\&(2)C=\max_{p}I\left(X;Y\right)=\max_{p}\left\{\left(1-\mu\right)h_{2}\left(p\right)\right\}=\left(1-\mu\right)\max_{p}\left\{h_{2}\left(p\right)\right\}=1-\mu\end{aligned}$$---------------------------------------
8.11 已知 BSC 信道的误码率为μ，若通过此信道发送 7 个比特，求如下事件发生的概率： 
(1)7 个都对；
(2)至多有 1 个错； 
(3)至少有 2 个错； 
++++++++++++
$\begin{aligned}&\text{解:(1)}P_1=\left(1-\mu\right)^7\\&(2)P_2=\left(1-\mu\right)^7+7\mu\left(1-\mu\right)^6=\left(1-\mu\right)^6\left(1+6\mu\right)\\&(3)P_3=1-P_2=1-\left(1-\mu\right)^6\left(1+6\mu\right)\end{aligned}$
---------------------------------------
第九章信道编码

9.1求下列二元码字之间的汉明距离：
（1）0000，0101
（2）01110，11100
（3）010101，101001 
（4）1110111，1101011 

++++++++++++
解：根据汉明距离的定义可得知上述4种情况下的汉明距依次为2、2、4、3。
---------------------------------------
9.2某码字的集合为

0000000      1000111      0101011      0011101 
1101100      1011010      0110110      1110001 
试求：

（1）该码字集合的最小汉明距离：；
（2）根据最小汉明距离确定其检错和纠错能力。
++++++++++++
解一：

（1）通过两两比较（共有$C_8^2=28$ 种组合）这8个码字可得最小汉明距离为4。

（2）由$t+1=4$ ，得该码可以保证检3位错； 
          由$2t+1=4$ ，得该码可以保证纠1位错。

 解二：

(1)就本题的具体情况，可以验证这8个码字构成了线性码。事实上，令$\mathbf{c}_{1}=10000111$ $\mathbf{c}_2=0101011$ 、 $\mathbf{c}_3=0011101$ ，则$\mathbf{c}_1$ 、 $\mathbf{c}_2$ 、 $\mathbf{c}_3$ 线性无关，而$1101100=\mathbf{c}_1+\mathbf{c}_2$ ， $1011010=\mathbf{c}_1+\mathbf{c}_3$ $0110110=\mathbf{c}_2+\mathbf{c}_3,1110001=\mathbf{c}_1+\mathbf{c}_2+\mathbf{c}_3$o 再由线性码的最小码距是非0码的最小码重的性质得知这8个码字之间的最小汉明距离为4。

(2）同解一

注：由最小码距$d_{\mathrm{min}}$ 所确定的纠错或检错能力是保证的纠错或检错能力，不表示其它的错误图样不能被纠正或者检出。特别在检错的情形下，由$d_{\mathrm{min}}$ 所确定的错能力往往过干保守。以本题为例，这个码实际可以检出一切的1、2、3、5、67比特错，在4比特错中，可以检出28种错，即所有4比特错误图案中除去和许用码字相同的那些。
---------------------------------------
9.3假设二进制对称信道的差错率$P=10^{-2}$

（1）（5.1）重复码通过此信道传输，不可纠正的错误出现的概率是多少？ 
（2）（4.3）偶校验码通过此信道传输，不可检出错误的出现概率是多少？
++++++++++++
解：（1)（5.1）重复码中发生3个或者更多错误时不可纠正，因此不可纠正错误的出现概率为
$$P_{1}=C_{5}^{3}P^{3}(1-P)^{2}+C_{5}^{4}P^{4}(1-P)^{1}+C_{5}^{5}P^{5}(1-P)^{0}\approx9.85\times10^{-6}$$

（2）（4.3）偶校验码中发生偶数个错时不可检出，这样的概率是
$$P_{2}=C_{4}^{2}P^{2}(1-P)^{2}+C_{4}^{4}P^{4}(1-P)^{0}\approx5.88\times10^{-4}。$$
---------------------------------------
9.4等重码的所有码字都具有相同的汉明重量请问这样的等重码是否是线性码吗？请说明理由。



++++++++++++
答：因为该码的所有码字都有相同数目的”1”，因此它不包括全0码字，但线性码必然包含全0码字，所以该码不是线性码
---------------------------------------
9.5若已知一个（7，4）码生成矩阵为

$$G=\begin{pmatrix}1&0&0&0&1&1&1\\0&1&0&0&1&0&1\\0&0&1&0&0&1&1\\0&0&0&1&1&1&0\end{pmatrix}$$

请生成下列信息组的码字：（1）（0100）（2）（0101）（3）（1110）（4）（1001）
++++++++++++
解： $\mathbf{c}=\mathbf{u}G$ ，将以上信息带入公式，得到码字分别为：

（1）0100101 （2）0101011 （3） 1110001 （4） 1001001 
---------------------------------------
9.6已知一个系统（7，4）汉明码监督矩阵如下：

$$H=\begin{pmatrix}1&1&1&0&1&0&0\\0&1&1&1&0&1&0\\1&1&0&1&0&0&1\end{pmatrix}$$

试求：（1）生成矩阵$G$ ：

（2）当输入信息序列m= （110101101010）时求输出码序列$\mathrm{c=}$ ？
++++++++++++
解： （1） 由$H$ 阵求出G阵：

$$G=\begin{pmatrix}1&0&0&0&1&0&1\\0&1&0&0&1&1&1\\0&0&1&0&1&1&0\\0&0&0&1&0&1&1\end{pmatrix}$$

（2） 首先将m 分组，四位码一组，不足的用0补，

得$m_{1}=1\:10\:1$ ， $m_2=0110$ ， $m_3=1010$ ，则$C_{i}=m_{i}\cdot G$
$$\therefore C_1=\begin{pmatrix}1&1&0&1\end{pmatrix}\begin{pmatrix}1&0&0&0&1&0&1\\0&1&0&0&1&1&1\\0&0&1&0&1&1&0\\0&0&0&1&0&1&1\end{pmatrix}=\begin{pmatrix}1&1&0&1&0&0&1\end{pmatrix}$$

$$C_{2}=m_{2}\cdot G=(0\quad1\quad1\quad0\quad0\quad0\quad1)\\C_{3}=m_{3}\cdot G=(1\quad0\quad1\quad0\quad0\quad1\quad1)$$
---------------------------------------
9.7试将下列生成矩阵

$$G=\begin{pmatrix}0&0&0&1&0&1&1\\0&0&1&0&1&1&0\\0&1&0&1&1&0&0\\1&0&1&1&0&0&0\end{pmatrix}$$

（1）化为系统码的生成矩阵$G=(I{:}Q)$ 形式；

（2）写出典型监督矩阵。

++++++++++++

解：（1）利用G的初等行变换来完成相应的变换

1.1行$\leftarrow\rightarrow4$ 行，2行$\leftarrow\rightarrow3$ 行，得：

$$G'=\begin{pmatrix}1&0&1&1&0&0&0\\0&1&0&1&1&0&0\\0&0&1&0&1&1&0\\0&0&0&1&0&1&1\end{pmatrix}$$

2. 2行+4行，1行+3行+4行，得：

$$G"=\begin{pmatrix}1&0&0&0&1&0&1\\0&1&0&0&1&1&1\\0&0&1&0&1&1&0\\0&0&0&1&0&1&1\end{pmatrix}$$

故$G^{\prime\prime}$ 就是所求的系统码生成阵。

（2）典型监督矩阵为

$$H=\begin{pmatrix}1&1&1&0&1&0&0\\0&1&1&1&0&1&0\\1&1&0&1&0&0&1\end{pmatrix}$$
---------------------------------------
9.8已知某线性分组码生成矩阵为

$$G=\begin{pmatrix}0&0&1&1&1&0&1\\0&1&0&0&1&1&1\\1&0&0&1&1&1&0\end{pmatrix}$$

试求：

（1）系统码生成矩阵$G=(I{:}Q)$ 表达形式；

（2）写出典型监督矩阵$H$ ； 
（3）若译码器输入$y=$ （0011111），请计算其校正子$S_{0}$ 。
（4）若译码器输入 $y=$（1000101），请计算其校正子$S_{0}$。
++++++++++++
解：（1）由初等行变换得：

$$G''=\begin{pmatrix}1&0&0&1&1&1&0\\0&1&0&0&1&1&1\\0&0&1&1&1&0&1\end{pmatrix}$$

（2）相应的校验矩阵为

$$H''=\begin{pmatrix}1&0&1&1&0&0&0\\1&1&1&0&1&0&0\\1&1&0&0&0&1&0\\0&1&1&0&0&0&1\end{pmatrix}$$

(3) $S=yH^{T}=(\:0010\:)$

(4) $S=yH^{T}=$ (1011)
---------------------------------------
9.9已知某（7，3）码生成矩阵为

---------------------------------------

$$G=\begin{pmatrix}1&0&0&1&0&1&1\\0&1&0&1&1&1&0\\0&0&1&0&1&1&1\end{pmatrix}$$

试求：可纠正差错图案和对应伴随式。
++++++++++++
解：先写出此码的校验矩阵为

$$H=\begin{pmatrix}1&1&0&1&0&0&0\\0&1&1&0&1&0&0\\1&1&1&0&0&1&0\\1&0&1&0&0&0&1\end{pmatrix}$$

$H$ 有4行，且各行线性无关，因此伴随式有16种不同的取值。除全0伴随式外，另外15种非0伴随式表示可以纠正15种错误图样。合理设计中应当选择错误个数最少的错误图案作为可纠正的错误图案。因此15种可纠正错误图案当中有7种单比特错，8种为双比特错或3比特错。

设从左到右的比特位置序号是从7到1。若单比特错误位置在第$i$ 位，则对应的伴随式是$H$ 的第i列。若双比特错误位置在第i $j$ 位（ $i\neq j$ ），则其对应伴随式是$H$ 的第$i$ 列和第$j$ 列之和。同样3比特错的伴随式是$H$ 中对应3个位置的列之和。由此可以根据伴随式写出相应的可纠正错误图样，结果如下

\begin{figure}
    \centering
    \includegraphics[width=1\linewidth]{赵培志/1d379e20827b99f44d26e4991293630.png}
\end{figure}
---------------------------------------
9.10在下表中列出了4种（3.2）码，这4个码组是线性码吗？是循环码吗？请分别说明理由

\begin{figure}
    \centering
    \includegraphics[width=1\linewidth]{赵培志/6960b5d7933539d8eb66f8a669b5520.png}
\end{figure}

++++++++++++

解：

C是线性码，因为它满足加法自封闭性；它也是循环码，因为它的码字任意移位后仍是可用码字。

C是线性码，因为它满足加法自封闭性；但它不是循环码，因为它的码字任意移位后不是可用码字。C不是线性码，因为它不满足加法自封闭性；它的码字任意移位后仍是可用码字，所以它是

非线性的循环码，它不是通常意义下的循环码C不是线性码，因为它不满足加法自封闭性；它也不具备循环封闭性：它的码字任意移位后

不是可用码字。
---------------------------------------
9.11下列中的所有多项式都是系数在GF(2）上的多项式，请计算下列各式：
$$\begin{aligned}
&\text{(1)} \left(x^{4}+x^{3}+x^{2}+1\right)+\left(x^{3}+x^{2}\right) \\
&\text{(2)} \left(x^{3}+x^{2}+1\right)\left(x+1\right) \\
&\text{(3)} \left(x^{4}+x\right)\mathrm{mod}\left(x^{2}+1\right) \\
&\text{(4)} \left(x^{4}+1\right)\mathrm{mod}\left(x+1\right) \\
&\text{(} 5)\left(x^{3}+x^{2}+x+1\right)\mathrm{mod}\left(x+1\right) 
\end{aligned}$$
++++++++++++
解：在GF(2)上的加法就是模2加，特别注意减法即是加法， 1+1=0。
$$\begin{aligned}
&1)\left(x^{4}+x^{3}+x^{2}+1\right)+\left(x^{3}+x^{2}\right) =x^{4}+\left(x^{3}+x^{3}\right)+\left(x^{2}+x^{2}\right)+1 \\
&=x^{4}+1 \\
&2)\:\left(x^{3}+x^{2}+1\right)\left(x+1\right)=x\left(x^{3}+x^{2}+1\right)+\left(x^{3}+x^{2}+1\right) \\
&=\left(x^{4}+x^{3}+x\right)+\left(x^{3}+x^{2}+1\right) \\
&=x^{4}+x^{2}+x+1 \\
&3)\:x^{4}+x=x^{4}+x^{2}+x^{2}+1+1+x=x^{2}\left(x^{2}+1\right)+\left(x^{2}+1\right)+x+1
\end{aligned}$$

所以$\left(x^{4}+x\right)$mod$\left(x^{2}+1\right)=x+1$

(4)由于$1^4+1=0$ ，说明$x=1$ 是多项式$x^4+1$ 的根，所以x4 $x^4$ $x^4+1$ 包含因式$x+1$ ，所以
$$\begin{pmatrix}x^4+1\end{pmatrix}\mathrm{mod}\begin{pmatrix}x+1\end{pmatrix}=0$$

(5)同(4)，将由于$x=1$ 代入$x=1$ 是多项式$x^{3}+x^{2}+x+1$ 得0 ，说明$x=1$ 是多项式$x^3+x^2+x+1$ 的根，所以$\left(x^{3}+x^{2}+x+1\right)$mod$\left(x+1\right)=0$
---------------------------------------
9.12对于系数在GF(2)上的多项式，计算$x^7$ 除以$g(x)=x^{3}+x+1$ 的余式
++++++++++++
\textbf{解}：可用长除法来做，过程如下，结果是1 
\begin{figure}
    \centering
    \includegraphics[width=1\linewidth]{赵培志/d66c59943a7cdaa94c2583c8c29c5d2.png}
\end{figure}
---------------------------------------
9.13 若已知一个(7,3)循环码生成多项式为$g(x)=x^4+x^3+x^2+1$,试求其生成矩阵。
++++++++++++
解：循环码的生成矩阵的多项式形式为$$G\left(x\right)=\left[\begin{array}{c}x^2g(x)\\xg(x)\\g(x)\end{array}\right]=\left[\begin{array}{c}x^6+x^5+x^4+x^2\\x^5+x^4+x^3+x\\x^4+x^3+x^2+1\end{array}\right]$$即$$G=\begin{pmatrix}1&1&1&0&1&0&0\\0&1&1&1&0&1&0\\0&0&1&1&1&0&1\end{pmatrix}$$
---------------------------------------
9.14已知：

$$x^7+1=(x^3+x^2+1)(x^3+x+1)(x+1)$$

(1)写出(7,3)循环码的生成多项式，

(2)写出(7,4)循环码的生成多项式，

(3)写出(7,6)偶校验码的生成多项式，

(4)写出(7,1)重复码的生成多项式。
++++++++++++
解：因为(n,k)循环码的生成多项式的最高次幂为$n-k$,故(7,3)循环码的生成多项式有两个，
分别是$x^4+x^2+x+1$和$x^4+x^3+x^2+1$。
同理，(7,4)循环码的生成多项式也有两个，分别是$x^3+x^2+1$ 和 $x^3+x+1$ 。
因为偶校验码仍为循环码(循环移任意位后仍为可用码),故它的生成多项式为一次多项式
即$x+1$。
因为(7,1)重复码仍为循环码，它的生成多项式的最高次幂为 6,即为
$$x^6+x^5+x^4+x^3+x^2+x+1$$

---------------------------------------
9.15(1)试证明(15,5)循环码的生成多项式为

$$g(x)=x^{10}+x^8+x^5+x^4+x^2+x+1$$

(2)并求对应的$h(x)=?$

(3)若信息码多项式为$u(x)=x^4+x+1$,请写出系统码多项式。

++++++++++++
解：(1)$n=15、k=5$,首先$g(x)$的次数等于$n-k$。其次用$g(x)$去除$x^{15}+1$,结果是

$$\frac{x^{15}+1}{g(x)}=x^5+x^3+x+1$$

即$g(x)$可整除$x^{15}+1$。因此$g(x)$是(15,5)循环码的生成多项式。
$$\begin{aligned}&(2)\:h(x)=\frac{x^{15}+1}{g\left(x\right)}=x^{5}+x^{3}+x+1\\&(3)u(x)x^{n-k}\mathrm{~mod~}g(x)=r(x)=x^{8}+x^{7}+x^{6}+x\:,\\&\text{故系统码}c(x)=u(x)x^{n-k}+r(x)=x^{14}+x^{11}+x^{10}+x^{8}+x^{7}+x^{6}+x\end{aligned}$$
---------------------------------------
9.16已知(15,11)循环码的生成多项式为$g(x)=x^{4}+x+1$ ，试求：
(1)h(x)=?

(2)写出该系统码生成矩阵$G=(I{:}Q)$ 形式。
++++++++++++
解：(1)$$h(x)=\frac{x^{15}+1}{g\left(x\right)}=x^{11}+x^8+x^7+x^5+x^3+x^2+x+1$$
(2) 系统码的生成矩阵的矩阵形式为
$$\left.\text{}\\G=\left(\begin{array}{c}x^{14}+\left(x^{14}\right)\mathrm{mod}\:g\left(x\right)\\x^{13}+\left(x^{13}\right)\mathrm{mod}\:g\left(x\right)\\\vdots\\x^{4}+\left(x^{4}\right)\mathrm{mod}\:g\left(x\right)\end{array}\right.\right)=\left(\begin{array}{c}x^{14}+x^{3}+1\\x^{13}+x^{3}+x^{2}+1\\x^{12}+x^{3}+x^{2}+x+1\\x^{11}+x^{3}+x^{2}+x\\x^{10}+x^{2}+x+1\\x^{9}+x^{2}+x\\x^{8}+x^{2}+1\\x^{7}+x^{3}+x+1\\x^{6}+x^{3}+x^{2}\\x^{3}+x^{2}+x\\x^{4}+x+1\end{array}\right)$$
---------------------------------------
9.17已知（15.7）循环码的生成多项式为
$$g(x)=x^8+x^7+x^6+x^4+1$$
(1)求$h(x)=?$
(2)若信息码组$u=(0011001)$ ，请写出系统码的码组。
++++++++++++
解：(1)$$h(x)=\frac{x^{15}+1}{g\left(x\right)}=x^7+x^6+x^4+1$$
(2)$$u(x)=x^{4}+x^{3}+1, u(x)x^{n-k} \mathrm{mod} g(x)=r(x)=x^{7}+x^{6}+x^{5}+x^{4}+x^{2}+x$$
故系统码$$c(x)=u(x)x^{n-k}+r(x)=x^{12}+x^{11}+x^8+x^7+x^6+x^5+x^4+x^2+x$$
即为0 0 1 1 0 0 1 1 1 1 0 1 1 0
---------------------------------------
解：
9.18 若(7,3)循环码的生成多项式是$g(x)=x^4+x^2+x+1$,请画出系统码编码电路。
\begin{figure}
    \centering
    \includegraphics[width=1\linewidth]{赵培志/40258a260ccc6faf0a432b0e5683e5a.png}
\end{figure}
图中开关在前三时钟周期内接在实线位置，后四时钟周期内接在虚线位置。
---------------------------------------
9.19 用八进制形式表示下列多项式.

$$\begin{aligned}&(1)x^{4}+x^{3}+1\\&(2)x^{5}+x^{2}+1\\&(3)x^{6}+x^{4}+x^{3}+x+1\\&(4)x^{9}+x^{5}+x^{3}+x^{2}+1\end{aligned}$$

++++++++++++
解：(1)表示为二进制是 11001,故八进制表示为 31
(2)表示为二进制是 100101,故八进制表示为 45 (3)表示为二进制是 1011011,故八进制表示为 133 (4)表示为二进制是 1000101101,故八进制表示为 1055
---------------------------------------
 9.20 将下列八进制数写成多项式。(1)23;(2)45;(3)105;(4)721
++++++++++++
解：(1)八进制 23 转化为二进制是 010011,对应得多项式表示为$x^4+x+1;$
(2)八进制 45 转化为二进制是 100101,多项式表示为$x^5+x^2+1$
(3)八进制 105 转化为二进制是 001000101,多项式表示为$x^6+x^2+1$
(4)八进制 721 转化为二进制是 111010001,多项式表示为$x^8+x^7+x^6+x^4+1$
---------------------------------------
9.21 某(3,1,3)卷积码$g_1=100,g_2=101,g_3=111$ ,画出该码的编码器。
++++++++++++
\begin{figure}
    \centering
    \includegraphics[width=1\linewidth]{赵培志/b084d9b1751443e8e45b77e670becd3.png}
\end{figure}
注：本课以及其他许多文献中采用这样的约定：

在分组码中，二进制比特向量的左起第一位是最高有效位（MSB.mos significantbit）右起第一比特是最低有效位（LSB,least significantbit）表达为多项式时，LSB代表0次项。例如：1011对应的多项式是$x^3+x+1$

在卷积码中，二进制比特向量的右起第一位是最高有效位（MSB）左起第一比特是最低有效位（LSB）表达为多项式时，LSB代表0次项。例如：1011对应的多项式是$1+x^{2}+x^{3}$

无论是卷积码还是分组码，发送次序是从LSB开始
---------------------------------------
9.22已知一个(2,1,5)卷积码$g^{1}=(11101)$ $g^2=(10011)$ ，则请

(1)画出编码器框图； (2)写出该码生成多项式$g(x)$ ： (3）写出该码生成矩阵$G$ （4）若输入信息序列为11010001，求输出码序列$c=7$
++++++++++++
解：(1) $n=2,k=1,m=4$

\begin{figure}
    \centering
    \includegraphics[width=1\linewidth]{赵培志/0d9d960a5f6fe76ca23f7e4801cbd26.png}
\end{figure}

$$g_1(x)=1+x+x^2+x^4\:;\:g_2(x)=1+x^3+x^4$$
(3)
$$G=\begin{pmatrix}11&10&10&01&11\\&11&10&10&01&11\\&&11&10&10&01&11\\0&&&...&...&...&...&...\end{pmatrix}$$

(4)

$$c=\begin{pmatrix}1\:10\:10001\end{pmatrix}\:,\:c\begin{pmatrix}x\end{pmatrix}=1+x+x^3+x^7$$

上支路的输出为$c\left(x\right)g_{1}\left(x\right)=1+x^{8}+x^{9}+x^{11}$ ，即10000001101
下支路的输出为$c\big(x\big)g_{2}\big(x\big)=1+x+x^{5}+x^{6}+x^{10}+x^{11}$ ，即110001100011
串并变换后的输出是:11 01 00 00 00 01 01 00 10 10 01 11
---------------------------------------
9.23已知一个（3.1.3）卷积码

$$g_1(x)=1+x+x^2,g_2(x)=1+x+x^2,g_3(x)=1+x^2$$

则试：

(1)画出该码编码器框图； 
(2）画出状态图、树图；
(3）求该码自由距离。
++++++++++++
$$\text{解:(1)}\:n=3,k=1,m=2$$

\begin{figure}
    \centering
    \includegraphics[width=1\linewidth]{赵培志/b813c0c092a0205bcc629f8f081b3d7.png}
\end{figure}
(2)
\begin{figure}
    \centering
    \includegraphics[width=1\linewidth]{赵培志/f9b8ad5a82e7f96cd00389c99f8562e.png}
\end{figure}

(3）观察树图可以看到：非0路径首次离开$a$ 必然经过$b$ ，首次回到$a$ 必然经过$c_{0}$ 因此所有自由路径一定是$ab\cdots ca$ 的形式。路径abca的码重是8，而其他所有形如$ab\cdots ca$ 的路径的码重不可能比abca 更轻。因此自由距为$d_{f}=8$
---------------------------------------
9.24已知一个(2.1.3）卷积码编码器结构如下图所示，试

(1)写出$g^{1}$ 、 $g^2$ 与生成矩阵G； (2)画出状态图、树图、网格图。

\begin{figure}
    \centering
    \includegraphics[width=1\linewidth]{赵培志/683556c9f9b1445364bc0618a8129dc.png}
\end{figure}
++++++++++++
解：(1) $g_{1}(x)=1+x^{2}\:;\:g_{2}(x)=x+x^{2}$

$$G=\begin{pmatrix}10&01&11&&&\mathbf{0}\\&10&01&11\\&&10&01&11\\\mathbf{0}&&&\cdots&\cdots&\cdots\end{pmatrix}$$(2)状态图

\begin{figure}
    \centering
    \includegraphics[width=1\linewidth]{赵培志/d3fe8163967b14b68a64e759a6b2adf.png}
\end{figure}
树图（图中a、b、c、 $d$ 分别代表状态00、10、01、11）


\begin{figure}
    \centering
    \includegraphics[width=1\linewidth]{赵培志/3d7c7348201328785d600eb5cd36dc9.png}
\end{figure}
网格图如下：
\begin{figure}
    \centering
    \includegraphics[width=1\linewidth]{赵培志/07b4e0d64ead48c11780744b75f47b7.png}
\end{figure}
---------------------------------------
9.25已知一卷积码编码器结构如下图所示，试求：
$(1)(n,k,K)=?$

(2)  $g^{1}=?\: g^{2}=?$ ，生成矩阵$G=?$ 
(3)若$x=(10111)$ ，求输出$c=?$

\begin{figure}
    \centering
    \includegraphics[width=1\linewidth]{赵培志/1.png}
\end{figure}
++++++++++++
解：（1）n=2,k=1,K=4 (2) $g^{1}=(1000)\:;\:g^{2}=(1101)$

$$G=\begin{pmatrix}11&01&00&01&&&\mathbf{0}\\&11&01&00&01&&\\&&11&01&00&01&\\\mathbf{0}&&&\cdots&\cdots&\cdots&\cdots\end{pmatrix}$$
(3)考虑编码器状态会0，则上支路的输出是10111000，下支路的输出是11110011，串并变换后的输出序列是（ 1101 11 11 10 00 01 01
---------------------------------------
9.26CDMA移动通信IS-95制式的下行信道采用（2.1.9）卷积码.其生成多项式为
$$g^{1}=(753)_{8}\:;\quad g^{2}=(561)_{8}\:;$$

问题1111111111111111111111111111其中$\left(x\mathrm{x}x\right)_{8}$ 表示八进制。请画出该编码器的框图
++++++++++++
解$:(753)_{8}=(1111101011)_{2}$ ， $(561)_8=(101110001)_{2\sigma}$ ：由此可画出编码器框图如下：

\begin{figure}
    \centering
    \includegraphics[width=1\linewidth]{赵培志/fd4b9ed0b69a8aea2b168d93c3a0deb.png}
\end{figure}
---------------------------------------
9.27移动通信IS-95制式上行信道采用（3，1.9）卷积码，其生成多项式为$g^{1}=(557)_{8}$ g' =(557)s $g^2=(663)_8$ $g^2=(663)_8$ $g^{1}=(557)_{8}\:;\: g^{2}=(663)_{8}\:;\: g^{3}=(711)_{8}$ g²= (711)s $g^{3}=(711)_{8}$ 。请画出该编码器的框图。
++++++++++++
解： $(557)_8=(101101111)_2$ ， $(663)_8=(110110011)_2$ ， $(711)_8=(111001001)_2$。由此可画出编码器框图如下：

\begin{figure}
    \centering
    \includegraphics[width=1\linewidth]{赵培志/7ad1c8c66d3b70ea7bbb47d1ff79089.png}
\end{figure}
---------------------------------------
9.28某信源的信息速率为9.6kbit/s,信源输出通过一个1/2码率的卷积编码器后用4PSK调制信号传输（设载波频率为1MHz），4PSK采用了滚降系数为1的频谱成形，请问

（1)4PSK的符号速率是多少？
  (2)请画出在信道中传送的4PSK信号功率谱（标上频率值）
++++++++++++
解：(1)9.6kbps的信息经过1/2码率的编码器之后速率变为19.2kbps，由于采用4PSK，所以符号速率是19.2/2=9.6kBaud。
       (2)滚降系数为1，所以调制后的信号带宽为$19.2\mathrm{kHZ}_{0}$ 其功率谱图如下
\begin{figure}
    \centering
    \includegraphics[width=1\linewidth]{赵培志/eb4aab02e310a914e3ab9fecb934687.png}
\end{figure}

---------------------------------------
 第十章正交码与伪随机码

10.1利用m 序列的移位相加特性证明双极性m 序列的周期性自相关函数为二值函数，且主副峰之比等于码长（周期）。
++++++++++++
证：m序列的移位相加特性特性是说，单极性m序列和它的移位相加后仍然是m序列。相加的结果在一个周期内1比0多一个。双极性m序列是把0、表示的m序列映射为土1表示，其中0映射为+1，1映射为-1。对于双极性m序列，一个周期内-1比+1多一个。在这种映射下模2加运算变成了乘法运算，如下表所示：

\begin{figure}
    \centering
    \includegraphics[width=1\linewidth]{赵培志/59c1eb7441fe511dd0aabc2a8c1d9e1.png}
\end{figure}

因此m 序列的移位相加特性特性对于双极性m 序列表现为： m 序列和它的移位相乘后仍然是m 序列。

$$R\bigl(\:j\bigr)=\frac{1}{p}\sum_{k=1}^{p}b_{k}b_{k+j}$$
，其中$b$ 的下周期为p的双极性m序列的周期性自相关函数定义为标按模p运算，即$b_{k+p}=b_k$

当/为$p$ 的整倍数时， $b_{k+j}=b_j$ . $b_kb_{k+j}=1$ ，因此$R(j)=R(0)=1$ ，这是自相关函数的主峰值；

对于$0<j<p$ ，令$c_{k}=b_{k}b_{k+j}$ ，则m序列的移位相加特性表明序列$\{c_k\}$ 也是m序列
$$\sum_{k=1}^pb_kb_{k+j}=\sum_{k=1}^pc_k$$
C 表示一个周期内求和，由于-1 比+1多一个，所以c=一1，从而得

$$R\bigl(j\bigr)=\begin{cases}&1&j\bmod p=0\\-\frac{1}{p}&else\\\end{cases}$$

这就证明了周期为$p$ 的m 序列的周期性自相关函数为二值函数，且主副峰之比等于码长
---------------------------------------
10.2已知线性反馈移存器序列的特征多项式为$f(x)=x^{3}+x+1$ ，求此序列的状态转移图， 并说明它是否是m序列。
++++++++++++
解：该序列的发生器逻辑框图为：

\begin{figure}
    \centering
    \includegraphics[width=1\linewidth]{赵培志/efda33d0bbe06ed610a77a32cc2234a.png}
\end{figure}



定义状态为向量$\mathbf{s}=\left(s_{1},s_{2},s_{3}\right)$ ，假设起始状态是100，则状态转移图如下：

\begin{figure}
    \centering
    \includegraphics[width=1\linewidth]{赵培志/c8b024944410378414ee999faf6696f.png}
\end{figure}

由于其周期为$2^{3}-1=7$ ，所以此序列是m序列
---------------------------------------
10.3已知m序列的特征多项式为$f(x)=x^{4}+x+1$ ，写出此序列一个周期中的所有游程。
++++++++++++
解：该m 序列的周期为15，一个周期为100011110101100，共有8个游程：

1             000          1111           0         10         11           00

其中长度为1的游程有4个；长度为2的游程有2个；长度为3的游程有1个；长度为4的游程有1个
---------------------------------------
10.4写出长度为8的所有瑞得麦彻序列和沃尔什序列，并比较说明它们之间的关系。
++++++++++++
解：长度为8的瑞得麦彻序列有4个，如下表所示

\begin{figure}
    \centering
    \includegraphics[width=1\linewidth]{赵培志/19ad095951beabfdc095d66df86495a.png}
\end{figure}

长度为8的沃尔什序列有8个，若以哈达马矩阵的行号为编号，则这8个码如下表所示：

\begin{figure}
    \centering
    \includegraphics[width=1\linewidth]{赵培志/25338c326816bfaa60d1a1c827fc18c.png}
\end{figure}
容易看到Walsh码的集合W包含了Rademacher码的集合$R_{0}$ Rademacher码除去全0码字外，其余三个线性不相关。另外由双极性Walsh码对乘法封闭这个特性可知单极性Walsh 码对模2加法封闭，由此可知W是一个线性空间，R是其一组基。因此每个Walsh码一定可以表示为R的线性组合，即
$$\mathbf{w}_{i}=\mathbf{b}G$$

此处


$$G=\begin{pmatrix}0&1&0&1&0&1&0&1\\0&0&1&1&0&0&1&1\\0&0&0&0&1&1&1&1\end{pmatrix}$$

二进制向量$\mathbf{b}=(b_0b_1b_2)$ 和Walsh码编号i的关系如下表：

\begin{figure}
    \centering
    \includegraphics[width=1\linewidth]{赵培志/18ffc1f1f1bd4e42a5d8137e112906a.png}
\end{figure}

也就是说， $b$ 是$i$ 的自然二进制表示
---------------------------------------
10.5已知优选对$m_{1}$ 、 $m_2$ 的特征多项式分别为$f_1(x)=x^3+x+1$ 和$f_{2}(x)=x^{3}+x^{2}+1$ ，写出由此优选对产生的所有Gold码，并求其中两个的周期互相关函数
++++++++++++
解：特征多项式为$f_1(x)=x^3+x+1$ 的m 序列的一个周期为1110100；特征多项式为$f_{2}(x)=x^{3}+x^{2}+1$ 的m序列的一个周期为1110010。由此生成的Gold码为
问题1111111
$\mathrm{G1:1110100\oplus1110010=0000110}$ $\mathrm{G4:1110100\oplus0101110=1011010}$ $\mathrm{G5:1110100\oplus0010111=1100011}$ 
G6 : 11101001001011 = 0111111 G7 : 11101001100101 = 001000

再加上原有的两个m 序列：

G8 : 1110100 G9 : 1110010

一共有9个。

考虑G1和G2的互相关。将这两个码的双极性形式为
g1 : 1 1 1 1 -1 -1 1

g2 : -1 1 1 -1 -1 1 -1

其互相关函数为：

$$R_{12}\left(k\right)=\sum_{i=0}^{6}g_{1}\left(i\right)g_{2}\left(i+k\right)\:,\:k\in\left\{0,1,2,\cdots,6\right\}$$

其中$i+k$ 按mod7计算。通过具体计算可得

$$R_{12}\left(k\right)=\begin{cases}-1&\quad k\in\left\{0,1,3,5\right\}\\3&\quad k=2\\-5&\quad k=4\end{cases}$$

---------------------------------------

10.6采用m 序列测距，已知时钟频率等于1MHz，最远目标距离为3000km ，求m 序列的长度（一周期的码片数）

解：m序列一个周期的时间长度即为可测量的最大时延值。m序列收发端与最远目标的往返时间为3000km$\times2/(3\times10^5$km/s)=0.02s ，因此m序列的周期应该大于20ms。由于序列发生器的时钟频率为$1MHz$ ，所以m序列长度应大于0.02sx 1MHz$=2\times10^4$

10.7已知某线性反馈移存器序列发生器的特征多项式为$f(x)=x^3+x^2+1$ ，请画出此序列发生器的结构图，写出它的输出序列（至少包括一个周期），并指出其周期是多少
++++++++++++
解：此序列发生器的结构为

\begin{figure}
    \centering
    \includegraphics[width=1\linewidth]{赵培志/d1ee227f75a286af2dfda356c6b6564.png}
\end{figure}
输出序列为：..10111001011100.. 周期为：7

---------------------------------------


\end{document}